\documentclass[11pt,a4paper]{article}
% Report A of batch 60 (with batch 61): seven manuscripts merged; base = manuscript 05.
% Preamble: manuscript 05's, extended by the macros of manuscripts 01-04, 06 and 07.
\usepackage[T1]{fontenc}
\usepackage[utf8]{inputenc}
\usepackage{amsmath,amsthm,mathtools}
\usepackage{newtxtext,newtxmath}
\usepackage[margin=27mm,headheight=15pt]{geometry}
\usepackage{microtype}
\usepackage{booktabs,array,longtable,tabularx}
\usepackage{float}% manuscript 02 places two tables with [H]
\usepackage{xcolor}
\usepackage{enumitem}
\usepackage{listings}
\usepackage{fancyhdr}
\usepackage{titlesec}
\usepackage[most]{tcolorbox}
\usepackage{xurl}
\usepackage[colorlinks=true,linkcolor=blue!45!black,citecolor=blue!45!black,urlcolor=blue!45!black]{hyperref}
\usepackage{aliascnt}
\usepackage[capitalise,noabbrev]{cleveref}
\definecolor{ink}{HTML}{163746}
\definecolor{accent}{HTML}{236B79}
\definecolor{pale}{HTML}{F0F5F7}
\titleformat{\section}{\Large\bfseries\color{ink}}{\thesection}{0.7em}{}
\titleformat{\subsection}{\large\bfseries\color{ink}}{\thesubsection}{0.7em}{}
\pagestyle{fancy}
\fancyhf{}
\fancyhead[L]{\small\scshape Canonical Diophantine Certificates}
\fancyhead[R]{\small Research report (seven manuscripts)}
\fancyfoot[C]{\thepage}
\renewcommand{\headrulewidth}{0.3pt}
\setlength{\parskip}{4pt plus 1pt}
\setlength{\parindent}{1.15em}
\setlist{itemsep=3pt,topsep=5pt}
\setlength{\emergencystretch}{2.2em}
\allowdisplaybreaks[2]
\setcounter{tocdepth}{1}
\newtheorem{theorem}{Theorem}[section]
\newaliascnt{lemma}{theorem}
\newtheorem{lemma}[lemma]{Lemma}
\aliascntresetthe{lemma}
\crefname{lemma}{Lemma}{Lemmas}
\newaliascnt{proposition}{theorem}
\newtheorem{proposition}[proposition]{Proposition}
\aliascntresetthe{proposition}
\crefname{proposition}{Proposition}{Propositions}
\newaliascnt{corollary}{theorem}
\newtheorem{corollary}[corollary]{Corollary}
\aliascntresetthe{corollary}
\crefname{corollary}{Corollary}{Corollaries}
\theoremstyle{definition}
\newaliascnt{definition}{theorem}
\newtheorem{definition}[definition]{Definition}
\aliascntresetthe{definition}
\crefname{definition}{Definition}{Definitions}
\newaliascnt{example}{theorem}
\newtheorem{example}[example]{Example}
\aliascntresetthe{example}
\crefname{example}{Example}{Examples}
\newaliascnt{question}{theorem}
\newtheorem{question}[question]{Research question}
\aliascntresetthe{question}
\crefname{question}{Research question}{Research questions}
\theoremstyle{remark}
\newaliascnt{remark}{theorem}
\newtheorem{remark}[remark]{Remark}
\aliascntresetthe{remark}
\crefname{remark}{Remark}{Remarks}
\crefname{part}{Part}{Parts}
\crefname{table}{Table}{Tables}
% ---- common macros (all manuscripts agree) ----
\newcommand{\N}{\mathbb N}
\newcommand{\Z}{\mathbb Z}
\newcommand{\Q}{\mathbb Q}
\newcommand{\R}{\mathbb R}
\newcommand{\eps}{\varepsilon}
\newcommand{\bits}{\operatorname{bits}}
\newcommand{\supp}{\operatorname{supp}}
\newcommand{\rank}{\operatorname{rank}}
\newcommand{\Run}{\operatorname{Run}}
\newcommand{\NF}{\operatorname{NF}}
\newcommand{\A}{\mathcal A}
\newcommand{\Par}{\operatorname{Par}}
% \code: one definition for all manuscripts (01's); 02's \texttt{\small #1} and
% 05's \nolinkurl{#1} print the same file and declaration names.
\newcommand{\code}[1]{\texttt{\detokenize{#1}}}
% ---- manuscript 05 (base) ----
\newcommand{\Net}{\mathcal N}
\newcommand{\Sol}{\operatorname{Sol}_{\N}}
\newcommand{\pos}[1]{\left(#1\right)_{+}}
\newcommand{\ind}{\mathbf 1}
\newcommand{\parikh}{\operatorname{Par}}
\newcommand{\sharpP}{\mathsf{\#P}}
\newcommand{\ce}{\text{c.e.}}
\newtcolorbox{resultbox}{colback=pale,colframe=accent,boxrule=0.6pt,arc=1mm,left=3mm,right=3mm,top=2mm,bottom=2mm}
\lstset{basicstyle=\ttfamily\small,breaklines=true,columns=fullflexible,keepspaces=true,frame=single,rulecolor=\color{black!20},showstringspaces=false}
% ---- manuscript 01 ----
\newcommand{\dom}{\operatorname{dom}}
\newcommand{\Tr}{\operatorname{Tr}}
\newcommand{\CF}{\operatorname{CF}}
\newcommand{\minuspart}[1]{#1^{-}}
\newcommand{\pluspart}[1]{#1^{+}}
\newcommand{\boxguard}{\ell_a\le x\le u_a}
% ---- manuscript 02 ----
\newcommand{\sgn}{\operatorname{sgn}}
\newcommand{\Exit}{\operatorname{Exit}}
\newcommand{\Nonterm}{\operatorname{Nonterm}}
\newcommand{\PolyPos}{\operatorname{PolyPos}}
\newcommand{\dd}{\mathop{}\!\mathrm{d}}
% ---- manuscript 03 ----
\newcommand{\Mem}{\mathsf{Mem}}
\newcommand{\Valid}{\mathsf{Valid}}
\newcommand{\cmp}{\mathsf{cmp}}
\newcommand{\cE}{\mathcal E}
\newcommand{\cW}{\mathcal W}
\newcommand{\cP}{\mathcal P}
\newcommand{\htg}{\mathsf{halt}}
\newcommand{\app}{\mathsf{app}}
\newcommand{\one}{\mathbf 1}
\newcommand{\repo}[1]{\nolinkurl{#1}}
% ---- manuscript 04 ----
\newcommand{\App}{\operatorname{App}}
\newcommand{\cl}{\operatorname{cl}}
\DeclareMathOperator{\poly}{poly}
% ---- manuscript 06 ----
\newcommand{\B}{\{0,1\}}
\newcommand{\Lex}{\operatorname{Lex}}
\newcommand{\len}{\operatorname{len}}
\newcommand{\sig}{\operatorname{sig}}
\newcommand{\Summ}{\operatorname{res}}
\newcommand{\capTwo}{\operatorname{cap}_2}
\newcommand{\Reach}{\operatorname{Reach}}
\newcommand{\depth}{\operatorname{depth}}
\newcommand{\vars}{\operatorname{vars}}
\newcommand{\ones}{\boldsymbol{1}}
\newcommand{\zeros}{\boldsymbol{0}}
\newcommand{\kw}[1]{\texttt{#1}}
\newcommand{\abs}[1]{\lvert #1\rvert}
\newcommand{\norm}[1]{\lVert #1\rVert_\infty}
% ---- manuscript 07 ----
\newcommand{\val}{\operatorname{val}}
\newcommand{\ff}[2]{(#1)_{\underline{#2}}}
\newcommand{\CT}{\mathcal A}
\DeclareMathOperator{\lc}{lc}
% \word: 07's \mathtt (06 defined \mathsf but never used it)
\newcommand{\word}[1]{\mathtt{#1}}
% ---- merge ----
\newcommand{\wtag}{\textup{[write]}}% marks text written in the batch-60/61 write phase
\newcommand{\srcnote}[1]{\par\smallskip\noindent{\small\itshape #1}\par\smallskip}
\hypersetup{pdftitle={Canonical Diophantine Certificates},pdfauthor={Research report merged from seven manuscripts prepared for Vladimir Reshetnikov},pdfsubject={Witness-faithful Diophantine representations of bounded computation: trace classes, accelerators, polynomial trajectories, memory logs, queues and rewriting}}
\newcommand{\srctag}[1]{\texorpdfstring{\hspace{0.5em}{\normalfont\small\color{accent}[#1]}}{ [#1]}}

\begin{document}
\hypersetup{pageanchor=false}
\begin{titlepage}
\vspace*{2mm}
{\color{accent}\large\scshape ProveIt research report \textperiodcentered{} merged from seven manuscripts\par}
\vspace{5mm}
{\color{ink}\fontsize{26}{31}\selectfont\bfseries Canonical Diophantine Certificates\par}
\vspace{4mm}
{\large Witness-faithful polynomial representations of bounded computation:
trace classes, accelerators, polynomial trajectories, memory logs, queues and rewriting\par}
\vspace{4mm}
{Research report prepared for Vladimir Reshetnikov. Manuscripts dated September 30, 2026; merged in the batch-60/61 write phase.\par}
\vspace{3mm}
\hrule
\vspace{2mm}
\begin{abstract}
\wtag{} An existential Diophantine representation says that a computation exists; it need not preserve which computation occurred, how many computations occurred, or how many arithmetic witnesses describe one of them. This report collects seven manuscripts that construct \emph{witness-faithful} (single-fold or parsimonious) polynomial representations of bounded computation and locate the boundary beyond which such representations are as hard as the open single-fold and finite-fold problems.

For Petri nets and interval-guarded translations, an explicit convex quadratic integer polynomial has its natural zeros in bijection with the commutation classes of enabled executions of exact length~$T$; its nonnegative real zero set is a bounded rational polytope whose lattice points are exactly those zeros (Part~II, manuscript~05). Quartic second routes (manuscripts~04 and~06), a Cartier--Foata variant for bounded causal height (manuscript~01), exact all-marking commutation tests, an exponential automaton-state lower bound and $\sharpP$-completeness of counting accompany it (Parts~I and~II). Single-fold quartics of fixed arity describe multisets, fixed-word powers and nested repetitions of counter actions independently of the repetition counts (Part~III); the relation ``$p(i)\ge0$ for all $i\le T$'' with $O((d+1)^3)$ witnesses independent of~$T$ (Part~IV); random-access memory logs with $24L$ witnesses (Part~V); cyclic-tag, deletion-tag and fixed-read FIFO executions with one coordinate per consumed symbol and one guard slack, together with a proof that faithful polynomial word memory needs two natural coordinates (Part~VI); and FRACTRAN runs and scheduled SKI reductions (Part~VII). Part~VIII collects further substrates.

Part~IX proves, with proofs from four manuscripts and a sketch in a fifth, that a single-fold (finite-fold) polynomial for universal halting exists exactly when every computably enumerable set has a single-fold (finite-fold) Diophantine representation, which is an open problem that this report does not solve.
\end{abstract}
\vspace{2mm}
\begin{resultbox}
\small\textbf{Status.} \wtag{} A merge of seven AI-assisted, unrefereed manuscripts, six delivered in batch~60 and one in batch~61 of ProveIt's incoming reports (\cref{cdc:sec:provenance}). The theorems carry ordinary mathematical proofs; the accompanying Python programs perform finite exact checks and are not proofs. No statement of this report is formalized in Lean or Rocq, although the report sits beside ProveIt's formal Hilbert's-tenth-problem project; the project declarations it cites are named in the README. Global priority is not claimed for any construction, and no resolution of the general single-fold or finite-fold problem is claimed. Every manuscript's own status statement is printed in \cref{cdc:sec:manuscripts}.
\end{resultbox}
\end{titlepage}
\hypersetup{pageanchor=true}
\setcounter{page}{1}

\tableofcontents
\newpage

% ---- manuscript 05, lines 119-153 ----
\section{The research problem and its relation to ProveIt}
\label{cdc:sec:intro}
\subsection{Beyond an existence-only translation}
A conventional application of the Matiyasevich--Robinson--Davis--Putnam theorem turns an effectively enumerable reachability predicate into a polynomial solvability problem. It does not, by itself, identify polynomial solutions with executions. One execution can acquire many auxiliary encodings; a nondeterministic execution may have many schedules; and a scheduling quotient can introduce another layer of multiplicity.

For a compiler intended for proofs, counting, or optimization, the natural stronger specification is a bijection
\begin{equation}\label{cdc:eq:goal}
\{\hbox{computational objects of interest}\}
\ \longleftrightarrow\
\{z\in\N^n:P(z)=0\}.
\end{equation}
The entire right-hand tuple must be determined, not merely its projection to a few distinguished coordinates. Otherwise even the number of solutions of the equation can be unrelated to the number of computations.

Here the main computational objects are executions modulo specified swaps of adjacent transitions. We resolve four concrete design questions. First, which swaps preserve enabledness at \emph{every} marking? Second, how can one enforce one representative per commutation class without expanding an exponentially large automaton? Third, can all arithmetic auxiliaries be unique while keeping the polynomial degree small? Fourth, what does this genuinely bounded construction imply---and not imply---about a single polynomial representing arbitrarily long universal computations?

\subsection{Repository baseline and the extension developed here}
The inspected ProveIt snapshot has commit
\begin{center}
\small\code{e8bb0931d67f80d9fce87a8cddb0f661ff19f956}.
\end{center}
Its Hilbert-tenth-problem development contains theorems
\path{boundedForall_dioph}, \path{exactIter_dioph}, and
\path{existsExactIter_dioph}; the last two provide Diophantine exact iteration and finite reachability under appropriate Diophantine substitutions \cite{repo-trace}. These are the correct existing interfaces to reuse for existence-level consequences.

The present article adds an explicit \emph{semantic fibre} specification. We keep the transition word, its states, the normal-form recognizer, and every auxiliary variable visible, prove unique extension, and only then collect the constraints into one polynomial. The bounded construction is elementary and does not need MRDP. Its relation to the repository is therefore complementary: it supplies a structured front end and a stronger correctness theorem, rather than another proof of existential trace closure.

The repository's status register distinguishes completed results from pending formalization \cite{repo-status}. No repository build was performed for this article. Statements about its interfaces are based on the inspected sources, not on a new independent kernel audit. No result below relies on an unverified numerical optimization in the repository.

\subsection{Theorems, classical ingredients, and novelty}
The exact encoding, its size accounting, its unique-extension proof, the graph realization, the memory lower bound, and the particular counting reduction are all supplied in full. They form a concrete extension suitable for review and implementation. We do not claim that an exhaustive literature search establishes priority for each individual observation.

The trace-monoid forbidden-factor characterization of lexicographic normal forms is classical; an explicit published form appears in Barth\'elemy \cite{barthelemy}. Sum-of-squares conjunction and arithmetic encodings of Boolean constraints are also standard ideas. Register-machine Diophantine representations have a long history, including Jones--Matiyasevich \cite{jm1984}. The contribution pursued here is the \emph{combined guarantee}: exact semantic commutation, canonical representative selection, no free witnesses, bounded real geometry, and explicit separation from fixed-dimensional finite-fold compression.

The distinction between polynomial representability and finite-fold or single-fold representability is essential. Matiyasevich discusses this distinction and the associated conjectures \cite{mat2010}; Cantone, Cuzziol, and Omodeo discuss recent single-fold specification methods and the danger of introducing uncontrolled witnesses \cite{cco2024}. Our boundary theorem concerns these stronger specifications. It is not presented as their solution.


\subsection{How this report is organized}\label{cdc:sec:organization}
\wtag{} This report merges seven manuscripts written independently for the same purpose, witness-faithful Diophantine compilation of bounded computation. Six (numbered 01--06 here) were delivered in batch~60 of ProveIt's incoming reports and pin the repository at \texttt{e8bb0931d}; the seventh (07) was delivered in batch~61, pins \texttt{725d2ebb6}, and names manuscripts~04 and~06 as its companions. Manuscript~05, whose introduction the preceding subsections reproduce, is the base: its trace-class theorem has the weakest hypothesis and the strongest conclusion among the four manuscripts that prove one (\cref{cdc:part:traces}). The seven titles, abstracts and status statements, and every manuscript's own introduction, are printed in \cref{cdc:sec:manuscripts}; \cref{cdc:sec:provenance} records their archives, pins and commits.

After the conventions (\cref{cdc:conv:sec}), the material is arranged by subject in nine Parts:
\begin{center}
\small
\begin{tabularx}{\textwidth}{@{}lXl@{}}
\toprule
Part & Subject & Manuscripts\\
\midrule
I & Exact commutation and resource algebra (\cref{cdc:part:resource}) & 05, 01, 06, 07\\
II & Canonical trace classes (\cref{cdc:part:traces}) & 05; 04, 06 (second routes); 01\\
III & Compressed repetition and accelerators (\cref{cdc:part:accel}) & 01, 06\\
IV & Polynomial trajectories (\cref{cdc:part:poly}) & 02\\
V & Memory logs, RAM and graph evaluation (\cref{cdc:part:memory}) & 03\\
VI & Queues: tag systems and FIFO networks (\cref{cdc:part:queues}) & 07\\
VII & FRACTRAN, SKI and term rewriting (\cref{cdc:part:rewriting}) & 04, 02, 01, 05\\
VIII & Other substrates (\cref{cdc:part:substrates}) & 05, 04, 01, 02, 06\\
IX & The boundary of fixed-arity compression (\cref{cdc:part:boundary}) & all seven\\
\bottomrule
\end{tabularx}
\end{center}
The back matter collects each manuscript's implementation and validation record (\cref{cdc:sec:validation}), its formalization proposals (\cref{cdc:sec:formal}), the merged research questions (\cref{cdc:sec:questions}), the remaining conclusions, the manuscripts' appendices, the provenance record and one merged bibliography.

A result proved by several manuscripts is printed once, at the place named in \cref{cdc:sec:provenance}, and credits every manuscript that proves it; a genuinely different proof or construction of the same result is kept as a marked second route. Two such mergers change the shape of the text: the universal-halting equivalence is one theorem, \cref{cdc:bd:thm:universal}, and the absence of a computable witness-height bound is one proposition, \cref{cdc:bd:prop:nobound}. At the places where the other manuscripts stated them, a note points to the merged statement and the manuscript's own proof follows.

\section{Conventions}\label{cdc:conv:sec}
Throughout, $\N=\{0,1,2,\ldots\}$. Every witness ranges over $\N$; every
polynomial has integer coefficients, and subtraction in a displayed
polynomial is subtraction in $\Z$, never truncated subtraction. A
\emph{residual system} is a finite list of integer polynomials required to
vanish; its single equation is the sum of their squares, of degree at most
four when every residual is quadratic and at most two when every residual
is affine. A node of an arithmetic circuit names a polynomial subexpression
and is not a witness; only an added coordinate is. Parameters are natural
except in Part~IV, whose coefficient parameters are signed integers. In
Part~VI the program, the input and terminal words and the horizon are data
of the construction, not unknowns.

A representation is \emph{single-fold} if each parameter tuple has at most
one witness, \emph{finite-fold} if it has finitely many, and
\emph{parsimonious} for a set of semantic objects if its witnesses are in a
specified bijection with them. ``Canonical certificate'' (Part~V) means a
single-fold representation.

\paragraph{Words used in two senses.}
\begin{itemize}
\item \emph{Certificate.} Here, a natural witness tuple of a polynomial. In
ProveIt's Hilbert's-tenth-problem project a ``straight-line certificate'' of
the 1980 Theorem~5 is a program whose arithmetic operations are counted; the
project's 75-operation universal certificate is of that kind. No result of
this report concerns operation counts.
\item \emph{Height.} Part~II's \emph{causal height} $H$ (from manuscript 01)
is the number of Cartier--Foata layers of a trace. In Part~VI, $H$ is an
upper bound on queue length. Elsewhere a \emph{witness height} bounds the
magnitude of witness coordinates.
\item \emph{Length} $T$ is the number of steps of a bounded execution in
Parts~II, VI, VII and~VIII; in Part~IV $T$ is the horizon of $\{0,\ldots,T\}$.
\item \emph{Memory.} Part~V certifies logs of a random-access memory.
Part~VI's \emph{word memory} is a numerical encoding of a finite word; its
dimension is the number of natural coordinates of the encoding.
\item \emph{Universal.} The Jones articles' \emph{universal pair}
$(\nu,\delta)$ bounds unknowns and degree of an existence representation;
Theorem~\ref{cdc:bd:thm:universal} concerns a universal polynomial that is
moreover single-fold or finite-fold, which is equivalent to an open problem.
\end{itemize}
\wtag{} Three further collisions need the same care. The letter $H$ has two more meanings: in Part~VII it is the total context depth of an SKI schedule (manuscript~04), and in Part~II $H_{i,a}$ (manuscript~05) and $H_{a,i}$ (manuscript~06) are linear forms in the selectors of one step. In Part~II, ``memory'' (\cref{cdc:thm:memory,cdc:wf:thm:memory}) counts the states or bits of a deterministic finite-state recognizer of normal forms, which is neither Part~V's nor Part~VI's notion. Finally, ``canonical'' qualifies three different objects: a canonical \emph{representative} (the lexicographically least word of a trace class, or the Cartier--Foata normal form), a canonical \emph{witness} (the unique auxiliary tuple of a single-fold representation, Part~V's ``canonical certificate''), and canonical \emph{padding} (a uniquely forced way of filling a fixed-length record). The tempting false reading of a ``canonical certificate for a trace'' is that the trace has one witness among all the traces with the same endpoints; the true reading, throughout Part~II, is that each trace class has exactly one witness, so the number of zeros at fixed endpoints is the number of classes.

\paragraph{Letters that change meaning between Parts.}
\wtag{} Each Part keeps the letters of the manuscript it prints. The dictionary below lists the letters that change meaning; the tempting false reading is to carry a letter across a Part boundary.
\begingroup\small
\begin{longtable}{@{}p{0.12\textwidth}p{0.83\textwidth}@{}}
\toprule
Symbol & Meaning by manuscript (and Part)\\
\midrule
\endhead
places, transitions & 05: $d$ places, $m$ transitions (Parts I, II, VIII); 01: $p$ coordinates, $m$ actions, alphabet $A$ (Parts I--III); 04: $p$ places, $r$ transitions, and in FRACTRAN $r$ fractions (Parts II, VII); 06: $p$ resources, $q$ transition labels (Parts I--III); 07: $p$ is the period of a cyclic-tag program, and also the number of resources in the first-failure product and of queues in a FIFO network (Parts I, VI).\\
$m$ & 01: number of actions; 05: number of transitions; 06: an endpoint marking ($m\xrightarrow{w}n$); 04: marking variables $m_{t,i}$; 07: appendant lengths $m_i$.\\
$u,v$ & 01: upper guard $u_a$ and displacement $v_a$; 05: upper guard $u_a$ (interval guards), production $v_a$, and the auxiliaries $u_{i,a}$; 06: least input $u_a$ and its output $v_a$; 07: guard slack $u$; 04: production matrix $V$.\\
$D$ & 01: dependence relation $D=(A\times A)\setminus I$; 05: sequential demand $D_{ab}$; 06: linear forms $D_{a,i}$; 07: total number of consumed symbols $D$ and prefix factors $D_i$.\\
$d$ & 05: number of places; 01: $d=\sum_a|D(a)|$; 02: nominal degree bound; 07: deletion number.\\
$N$ & 05: final marking; 01: witness counts $N_{\mathrm{acc}}$, $N_H$; 03: number of padded records; 06: number of syntax nodes; 07: $N_{B,d}$, the number of halted words.\\
$B$ & 01, 04: witness-height bounds; 04: stack base; 05: coordinate bounds $B_p$; 07: alphabet size.\\
$K$ & 01: universal halting $K(e,x)$; 05: label-additive cost $K(w)$; 07: $K=(B-1)!$.\\
$C$ & 03: number of comparators; 05: linear forms $C_{i,a}$; 07: the largest number of symbols read in one event.\\
$L$ & 01: loss totals $L_j$ (and $L=n_a+2n_b$ in an example); 03: log length; 05: lower box $L_{ab}$; 07: length residual.\\
$G$ & 01: gain totals $G_j$; 07: the first-failure guard.\\
$I$ & a symmetric irreflexive independence relation in every manuscript that uses one; 05 and 06 name maximal ones $I_{\max}$ and $I_*$.\\
\bottomrule
\end{longtable}
\endgroup

\paragraph{Reading the merged text.}
\wtag{} Text printed from a manuscript keeps that manuscript's wording. In it, ``this article'', ``this paper'', ``the present report'' and ``below'' refer to that manuscript; theorem, section and equation numbers are this report's. Where a manuscript names one of its files, the shipped name in \path{code/} or \path{data/} is printed (the README maps every delivered name to its shipped name), except inside a displayed command listing, which keeps the delivered layout because the programs import one another by their delivered names. A manuscript's repository citations are to the pinned commit it names. Renamed macros and symbols: manuscript~02's \verb|\code| (small typewriter) and manuscript~05's \verb|\code| (a non-breaking URL) print as manuscript~01's \verb|\code|; the pin macros of manuscripts 01, 02, 04 and~07 (01's and 07's share the name \verb|\repoSHA| but hold different commits) are printed as the literal commit identifiers; manuscript~06 defined a \verb|\word| macro it never used, so manuscript~07's is the only one. No mathematical symbol was renamed. The rest of this section prints each manuscript's own conventions.

% ---- manuscript 05, lines 154-197 ----
\subsection{Manuscript 05: Conventions and witness-preserving arithmetic}
\label{cdc:sec:conventions}
Throughout, $\N=\{0,1,2,\ldots\}$. All unknowns are nonnegative integers unless a real relaxation is explicitly being considered. Polynomial equalities, including their subtractions, are interpreted in $\Z$ or $\R$, not using truncated subtraction on $\N$. A transition label identifies a transition, not merely its visible action name; distinct transitions with identical effects remain distinct labels.

\begin{definition}[Representations and fibres]
A polynomial $P(x,z)\in\Z[x,z]$ represents a relation $R$ if
\[
R(x)\quad\Longleftrightarrow\quad\exists z\in\N^n\;P(x,z)=0.
\]
It is \emph{finite-fold} if for each fixed $x$ there are only finitely many such $z$, and \emph{single-fold} if there is at most one. A representation of a set of computational objects is \emph{parsimonious} here when an explicitly defined decoding map is a bijection from its entire natural zero set to those objects.
\end{definition}

A finite number $n$ of witness \emph{coordinates} says nothing about the cardinality of a fibre. Even a one-variable equation such as $0z=0$ has infinitely many natural solutions. Conversely, the polynomials in the main theorem generally have many solutions at a fixed pair of endpoints: one for each execution class. Their unique-extension property is conditional on the selected canonical word.

\begin{lemma}[Unique slack and scalarization]\label{cdc:lem:slack}
Suppose $L_1(z),\ldots,L_s(z)$ are integer polynomials. Then
\[
\sum_{j=1}^s L_j(z)^2=0\quad\Longleftrightarrow\quad
L_j(z)=0\text{ for every }j
\]
over $\Z$ and over $\R$. For integer-valued expressions $A,B$, the inequality $A\le B$ is equivalent to the existence of a \emph{unique} $s\in\N$ satisfying $A+s=B$.
\end{lemma}
\begin{proof}
Every square is nonnegative. A sum of nonnegative reals is zero exactly when every summand is zero. The slack is necessarily $B-A$, which is a natural number exactly under the stated inequality.
\end{proof}

\begin{lemma}[Affine Boolean conjunction]\label{cdc:lem:and}
For fixed $p,q\in\{0,1\}$, the natural solutions of
\begin{equation}\label{cdc:eq:and}
p=u+s_1,\qquad q=u+s_2,\qquad
1+u=p+q+s_3
\end{equation}
consist of exactly one tuple $(u,s_1,s_2,s_3)$, with $u=pq$.
For real $p,q\in[0,1]$, the same equations imply
\[
\max(0,p+q-1)\le u\le\min(p,q).
\]
\end{lemma}
\begin{proof}
Nonnegativity gives $u\le p$, $u\le q$, and $u\ge p+q-1$. For Boolean inputs these bounds determine $u$: it is zero when either input is zero and one when both are one. Each slack then has a prescribed value. The displayed real bounds are precisely the same inequalities.
\end{proof}

This small gadget is responsible for the quadratic, rather than quartic, final polynomial. We do not use a nonunique decomposition into sums of squares to enforce nonnegativity: natural-variable domains provide it directly. Replacing all natural variables by arbitrary integer variables would invalidate the later one-hot and slack arguments. Replacing a natural variable by a sum of four squares preserves existence but not the claimed bijection.


% ---- manuscript 01, lines 74-90 ----
\subsection{Manuscript 01: Three different meanings of an exact representation}
Throughout, $\N=\{0,1,2,\ldots\}$. For an integer polynomial $P(x,w)$ and fixed parameter tuple $x$, let
\[
 Z_P(x)=\{w\in\N^k:P(x,w)=0\}.
\]
We distinguish three contracts.
\begin{definition}
An \emph{existence representation} of a relation $R$ satisfies
$R(x)\iff Z_P(x)\ne\varnothing$.
It is \emph{single-fold} if $|Z_P(x)|\le1$ for every $x$, and \emph{finite-fold} if $|Z_P(x)|<\infty$ for every $x$.
Given a set $\mathcal C(x)$ of semantic objects, a \emph{parsimonious representation of $\mathcal C$} is a specified bijection
$Z_P(x)\simeq\mathcal C(x)$.
\end{definition}
A parsimonious history representation need not be single-fold: there can be several genuinely different histories between the same endpoints. The goal is not to identify those histories, but to avoid multiplying each of them by irrelevant choices. Conversely, a single-fold decision representation need not display a recognizable computation at all.

The universal single-fold and finite-fold questions are not solved by ordinary MRDP. Matiyasevich distinguishes ordinary polynomial representability from single-fold \emph{exponential} representability and asks whether the two improvements can be combined \cite{mat2010}. We use this distinction as a precise boundary, rather than claiming that deterministic execution automatically yields a unique polynomial witness.


% ---- manuscript 02, lines 206-340 ----
\subsection{Manuscript 02: Single-fold arithmetic without hidden choices}
\label{cdc:pt:sec:arithmetic}

\begin{definition}[Single-fold representation]
For a relation $R(\boldsymbol a)$, a polynomial representation
\[
 R(\boldsymbol a)\iff
 \exists\boldsymbol w\in\N^m\;P(\boldsymbol a,\boldsymbol w)=0
\]
is \emph{single-fold} if for every fixed input $\boldsymbol a$ there is at
most one $\boldsymbol w$ satisfying the equation. Thus a true instance has
exactly one witness and a false instance has none. It is \emph{finite-fold}
if every input has only finitely many witnesses.
\end{definition}
Our convention is $0\in\N$. Free parameters may range over $\Z$, while
existential coordinates always range over $\N$. To obtain an entirely
natural-input relation, replace each signed free parameter $a$ by
$a^+-a^-$, with $a^+,a^-\in\N$. These are \emph{free inputs}; their
possible redundancy does not create extra existential witnesses for a fixed
input pair. When a signed quantity itself is existentially quantified, its
pair must be normalized as below.

\subsubsection{Canonical signed coordinates}
\begin{lemma}[Signed normal form]
\label{cdc:pt:lem:signed}
Every $v\in\Z$ has a unique representation
\[
 v=v^+-v^-,\qquad v^+,v^-\in\N,\qquad v^+v^-=0.
\]
It is $v^+=\max(v,0)$ and $v^-=\max(-v,0)$.
\end{lemma}
\begin{proof}
The product condition implies at least one coordinate is zero. If $v>0$,
then $v^-$ cannot be positive or $v^+$ zero, so $(v^+,v^-)=(v,0)$.
The case $v<0$ is symmetric, and $v=0$ forces both coordinates to vanish.
\end{proof}
Without the product condition, the pairs $(v^++k,v^-+k)$ would give
infinitely many encodings. This elementary issue is not cosmetic: it is
exactly the sort of witness multiplicity that the construction must avoid.

\begin{lemma}[A total single-fold sign test]
\label{cdc:pt:lem:sign}
For a given integer $v$, the following equations have exactly one solution
$(b_-,b_0,b_+,s)\in\N^4$:
\begin{align}
 b_\epsilon(b_\epsilon-1)&=0\quad(\epsilon\in\{-,0,+\}),\nonumber\\
 b_-+b_0+b_+&=1,\label{cdc:pt:eq:sign-gadget}\\
 v&=(b_+-b_-)(s+1),\nonumber\\
 b_0s&=0.\nonumber
\end{align}
The three bits respectively indicate $v<0$, $v=0$, and $v>0$.
\end{lemma}
\begin{proof}
Exactly one of the three bits is $1$. If $b_+=1$, the magnitude equation
says $v=s+1>0$ and determines $s=v-1$. If $b_-=1$, it says
$v=-(s+1)<0$ and determines $s=-v-1$. If $b_0=1$, it forces $v=0$;
the last equation then forces $s=0$. The indicated assignment exists in
each case, and the cases are mutually exclusive.
\end{proof}
All equations in \eqref{cdc:pt:eq:sign-gadget} are quadratic or linear when $v$ is
a linear expression in existing wire coordinates. In particular, the sign
of a negative integer needs neither a square-root witness nor a choice of
squares representing an inequality.

\subsubsection{Total circuits and the quartic conversion}
\begin{lemma}[Canonical circuit lemma]
\label{cdc:pt:lem:circuit}
Let $\Phi(\boldsymbol a)$ be a fixed quantifier-free Boolean combination of
integer polynomial equalities and inequalities. Given an arithmetic and
Boolean circuit of size $S$ for $\Phi$, one can effectively construct
$O(S)$ natural auxiliary coordinates and quadratic integer residuals
$Q_1,\ldots,Q_N$ such that
\[
 \Phi(\boldsymbol a)
 \iff \exists\boldsymbol u\in\N^{O(S)}\;
 \sum_{j=1}^N Q_j(\boldsymbol a,\boldsymbol u)^2=0.
\]
The auxiliary tuple is unique whenever it exists. The resulting polynomial
has degree at most four in \emph{all} its free and existential variables.
\end{lemma}
\begin{proof}
Evaluate the entire circuit in a fixed topological order. Represent every
signed arithmetic output by its canonical pair from
Lemma~\ref{cdc:pt:lem:signed}. For an addition gate, impose
\[
 z^+-z^-=(x^+-x^-)+(y^+-y^-),\qquad z^+z^-=0.
\]
For multiplication, replace the right-hand side by
$(x^+-x^-)(y^+-y^-)$. These equations have degree at most two, and the
output pair is uniquely determined by the previously determined inputs.
Fixed integer constants and signed free inputs are allowed in the same
formulas. Affine expressions may be retained as expressions rather than
stored as extra wires.

For an atomic comparison, first evaluate its polynomial difference and then
use Lemma~\ref{cdc:pt:lem:sign}. Equality is the zero bit; strict inequality is
one of the other bits. Non-strict inequality is the sum of the appropriate
two bits. The Boolean operations are represented by
\[
 \neg b=1-b,\qquad b\wedge c=bc,\qquad
 b\vee c=b+c-bc.
\]
Each output is uniquely determined and is automatically a bit when its
inputs are bits. Store outputs as additional wires as necessary to keep
every defining residual quadratic. Finally require the truth-output bit to
be $1$.

All internal computations, including those belonging to a disjunct that is
not needed to make $\Phi$ true, are still evaluated. Consequently they
have unique values and cannot introduce unconstrained witnesses.
Induction along the circuit proves that all internal coordinates are
uniquely determined by $\boldsymbol a$ before the final truth assertion is
made. That assertion holds precisely when $\Phi$ is true.

A sum of squares of integers is zero exactly when each summand is zero.
Squaring the quadratic residuals and summing them yields the advertised
single equation and the degree bound. The number of residuals and variables
per gate is bounded by an absolute constant.
\end{proof}

\begin{remark}[What is and is not being eliminated]
The lemma removes Boolean structure and comparisons from a \emph{fixed
finite} formula. It does not eliminate an arbitrary bounded universal
quantifier, nor a variable-length circuit. The next sections provide a
special finite formula for one such quantifier, using polynomial degree.
\end{remark}

\begin{example}[An unsafe disjunction]
The relation $a=0\vee b>0$ can be expressed by
$\exists s\in\N\;a(b-s-1)=0$. But when $a=0$, every $s$ is allowed.
A product representing a disjunction preserves existence but not
uniqueness. A total sign circuit has no such freedom. Canonical padding of
inactive certificate components will play the same role below.
\end{example}


% ---- manuscript 03, lines 169-195 ----
\subsection{Manuscript 03: Conventions}
Throughout, $\N=\{0,1,2,\ldots\}$. Every polynomial has coefficients in $\Z$, but
all input and witness coordinates range over $\N$. Subtraction in a displayed
polynomial is subtraction in $\Z$, \emph{not} truncated natural subtraction.
These conventions are essential to the comparison gadget.

A \emph{residual system} is a finite list $F_1,\ldots,F_R$ of integer
polynomials, required to vanish. Its associated polynomial is
\begin{equation}\label{cdc:mem:eq:sos}
 P=\sum_{i=1}^{R}F_i^2.
\end{equation}
At natural, integer, or real arguments, $P=0$ is equivalent to simultaneous
vanishing of all residuals. In particular, quadratic residuals produce a
polynomial of degree at most four. We keep both representations: the residual
list exposes the semantics, while \eqref{cdc:mem:eq:sos} supplies the single equation.

\begin{definition}[Canonical certificate]
For a relation $R(x,y)$, a certificate $F(x,y,z)=0$ is \emph{canonical} when
$R(x,y)$ holds exactly when there exists a witness $z$, and for each fixed
$(x,y)$ there is at most one such $z$. If $R$ is the graph of a total function,
there is exactly one pair $(y,z)$ for each $x$.
\end{definition}
The word \emph{parsimonious} will mean that projection from polynomial solutions
to the specified semantic executions is a bijection. Distinct labelled
nondeterministic choices count as distinct executions, even when they later
merge into the same machine state.


% ---- manuscript 04, lines 108-148 ----
\subsection{Manuscript 04: Arithmetic interfaces and canonical auxiliary variables}\label{cdc:wf:sec:interfaces}
\subsubsection{Domains, residuals, and witnesses}
Throughout, $\N=\{0,1,2,\ldots\}$. Every witness variable ranges over $\N$, but all polynomial expressions are evaluated in $\Z$. In particular, subtraction is ordinary integer subtraction, not truncated natural subtraction. A \emph{residual system} is a finite list
\begin{equation}\label{cdc:wf:eq:residual}
 F_1(\mathbf x)=0,\ldots,F_E(\mathbf x)=0,
 \qquad F_i\in\Z[\mathbf x].
\end{equation}
The single associated polynomial is
\begin{equation}\label{cdc:wf:eq:sos}
 P(\mathbf x)=\sum_{i=1}^{E}F_i(\mathbf x)^2.
\end{equation}
Since an integer square is nonnegative, $P=0$ if and only if every residual vanishes. If $\deg F_i\leq2$, then $\deg P\leq4$. This conversion introduces no variables and preserves the entire solution set, not just solvability.

A displayed sum of squares is already an explicit polynomial. Expanding it into a list of monomials can enlarge its description. Our exact equation counts count the residuals; they are not claims about a minimum number of operations or about the number of monomials in a fully expanded quartic.

\begin{definition}[Witness-faithfulness]
For fixed external data $d$, a residual system is witness-faithful for a set $\mathcal C_d$ if it comes with mutually inverse maps between $\mathcal C_d$ and all its natural solutions. If the objects are classes under an equivalence relation, the maps must be defined on classes, rather than on arbitrary representatives.
\end{definition}

\subsubsection{Parsimonious quadratization}
The following familiar circuit principle will also prevent several common multiplicity errors.
\begin{lemma}[Unique-extension quadratization]\label{cdc:wf:lem:quadratization}
Given a finite system of integer polynomial equations in natural variables $\mathbf x$, one can effectively construct a system of equations of degree at most two in $(\mathbf x,\mathbf w)$ such that every original solution has exactly one extension $\mathbf w$, and no other solutions exist. Consequently the construction preserves finite-fold and single-fold fibers with respect to any distinguished input coordinates.
\end{lemma}
\begin{proof}
Write each polynomial as $F^+-F^-$, where $F^+$ and $F^-$ have nonnegative integer coefficients. Construct acyclic arithmetic circuits for these two polynomials using natural constants, addition, and multiplication. For each gate introduce a fresh natural variable $w$. An addition gate gives $w-u-v=0$ and a multiplication gate gives $w-uv=0$; a constant gate gives $w-c=0$. Finally equate the two output variables.

For any assignment to $\mathbf x$, evaluating the circuits gives an extension satisfying all gate equations. Conversely, topological induction through the circuit forces each gate variable to equal that value. The final output equality is precisely $F=0$. Applying the construction to all residuals proves the assertion. Combining the quadratic residuals by \eqref{cdc:wf:eq:sos} preserves this unique extension.
\end{proof}

An arbitrary integer variable should not instead be replaced by $u-v$ with unrestricted $u,v\in\N$: the replacements $(u+k,v+k)$ immediately give infinitely many encodings of the same integer. The lemma avoids this problem by computing the positive and negative polynomial parts as \emph{functions} of the original variables.

\begin{lemma}[Canonical branch activation]\label{cdc:wf:lem:activation}
Suppose each labeled branch has a natural-witness certificate with at most one auxiliary assignment for fixed input and output. A finite branch choice can be compiled into quadratic residuals so that the resulting solutions correspond exactly to valid labeled choices, with no multiplicity from inactive branches.
\end{lemma}
\begin{proof}
Use one-hot natural selectors $x_a$ with $x_a(x_a-1)=0$ and $\sum_a x_a=1$. For every branch-local witness $w_{a,j}$ impose $(1-x_a)w_{a,j}=0$. Thus all inactive branch witnesses are zero. Compute the positive and negative parts of each branch residual by the always-active circuits of Lemma~\ref{cdc:wf:lem:quadratization}; their outputs $u,v$ are uniquely determined even when that branch is inactive. Gate the final equality by $x_a(u-v)=0$, a quadratic equation. In the selected branch the original certificate holds and has its assumed unique extension; elsewhere the witnesses are fixed at zero and the circuit wires are forced. This proves the bijection.
\end{proof}

The branch \emph{label} is part of the computational object in this lemma. If two labels describe the same semantic event, a further canonicalization is needed. Also, directly multiplying a selector by a quadratic residual gives a cubic: computing its two polynomial parts first is what retains quadratic residual degree.


% ---- manuscript 06, lines 232-308 ----
\subsection{Manuscript 06: Domains, schedules, and three kinds of uniqueness}
\label{cdc:cs:sec:definitions}

Throughout, $\N=\{0,1,2,\ldots\}$. Every unknown in a Diophantine equation ranges
over $\N$, while coefficients and polynomial subtraction belong to $\Z$.
Thus $x-y$ in an equation is ordinary integer subtraction, never truncated
natural subtraction. Vector inequalities, minima, and products explicitly
called coordinatewise are taken coordinate by coordinate.

\begin{definition}[Token transition]
Fix $p\geq1$. A transition $a$ is a pair
\[
u_a,v_a\in\N^p.
\]
It is enabled at marking $m$ when $m\geq u_a$, and its result is
$m-u_a+v_a$. A finite ordered alphabet $A$ labels the transitions; distinct
labels may have identical numerical behavior. Write $q=|A|\geq1$.
\end{definition}

The model is exactly the firing semantics of a weighted Petri net once labels
are assigned to transitions. It also describes multiset rewriting with a fixed
finite species set. There is no test for the absence of tokens, no reset,
copying operation, or multiplication of the current marking. An explicit
input requirement permits a transition to consume and restore a catalyst.
Such a transition cannot in general be replaced by its net displacement.

For a word $w=a_1\cdots a_L$, write
$m\xrightarrow{w}n$ when every successive transition is enabled and the final
marking is $n$. The empty word $\eps$ is executable at every marking and leaves
it unchanged.

\begin{definition}[Parametrized repetition expression]
\label[definition]{cdc:cs:def:expression}
Fix parameter names $k_1,\ldots,k_r$. An expression has grammar
\[
E::=\eps\mid a\mid(EF)\mid E^{k_j}.
\]
An assignment $\mathbf{k}\in\N^r$ determines a finite word
$\llbracket E\rrbracket_{\mathbf{k}}$ by concatenation and ordinary word
powers. In particular, $E^0$ expands to $\eps$, whether or not the expansion of
$E$ could execute from the given marking. The same parameter may occur at many
nodes. The repetition depth is the maximum number of repetition nodes on a
root-to-leaf path.
\end{definition}

We count each empty expression, letter, concatenation, and repetition as a
syntax node. Unless stated otherwise, the syntax is a tree. The template is
fixed before the free numerical parameters are supplied. This is not a
program that chooses a different syntactic body in response to its current
counter values. Parameter-dependent choices will occur only inside the
\emph{certificate compiler}, and their values will be uniquely constrained.

\begin{definition}[Single-fold auxiliary representation]
A polynomial $Q(x,z)\in\Z[x,z]$ represents a relation $R(x)$ single-fold when
\[
\#\{z\in\N^s:Q(x,z)=0\}=
\begin{cases}1,&R(x),\\0,&\neg R(x).\end{cases}
\]
A representation is finite-fold when this set is finite for every $x$ and
nonempty exactly on $R$.
\end{definition}

The declaration of the free tuple $x$ is part of the theorem. In our compressed
compiler it includes the starting marking, ending marking, and \emph{all}
repetition parameters. If the parameters are subsequently quantified, the
number of solutions can increase. A second uniqueness notion concerns the
word represented by a parameter tuple. A third concerns the word selected
inside a commutation class. These three notions will be kept separate.

\begin{example}[An infinite multiplicity hidden by projection]
For $E=(a^{k_1})^{k_2}$, every tuple with $k_1=0$ or $k_2=0$ expands to the empty
word. Our compiler has exactly one auxiliary solution for each such tuple
and equal endpoints. If $(k_1,k_2)$ become witnesses too, the empty execution
has infinitely many certificates. No local uniqueness gadget can remove this
multiplicity without changing the parameterization.
\end{example}


% ---- manuscript 07, lines 125-144 ----
\subsection{Manuscript 07: Arithmetic and word interfaces}
\subsubsection{Natural witnesses, integer evaluation}
Throughout, $\N=\{0,1,2,\ldots\}$. Witness coordinates lie in $\N$, but polynomial coefficients and evaluation lie in $\Z$. In particular, every subtraction in this article is signed subtraction. Formal queue lengths after an illegal event may be negative; silently replacing them by truncated natural subtraction would invalidate several arguments.

\begin{definition}
A \emph{residual system} is a finite list $R_1,\ldots,R_s\in\Z[X_1,\ldots,X_k]$. Its associated polynomial is
\[
 P=\sum_{j=1}^{s}R_j^2.
\]
Its natural zero set is the set of $\mathbf x\in\N^k$ satisfying every residual equation $R_j(\mathbf x)=0$.
\end{definition}
The equivalence follows because each square is a nonnegative integer. A quadratic residual system therefore gives a polynomial of degree at most four. Introducing an extra coordinate for a circuit gate changes the witness count; simply naming a polynomial subexpression does not.

\subsubsection{Arithmetic circuits are polynomial syntax}
A directed acyclic arithmetic circuit with integer constants and addition, subtraction, and multiplication denotes one ordinary integer polynomial. Reusing a node is a way of writing a polynomial compactly. It is not existential quantification over the node's value. Our JSON artifacts make this distinction explicit: only the listed input variables are witnesses; internal nodes are evaluated deterministically.

Consequently an $O(T)$ circuit need not expand to $O(T)$ monomials. For example, a product of $T$ affine factors may have exponentially many monomials when expanded. Circuit operations on large integers are also not unit-cost bit operations. We report degree, number of witness coordinates, arithmetic size, and witness bit height separately.

The horizon $T$, the input word, and the program are parameters of the \emph{construction}, not additional unknowns. A symbol such as $2^T$ is an integer coefficient computed before solving the equation. Products indexed by $i<T$ are finite syntactic products for that particular polynomial.


\section{The seven manuscripts}\label{cdc:sec:manuscripts}
\wtag{} Each subsection below gives one manuscript's title, author line and date as delivered, its abstract and status statement verbatim, and then its own introduction (except manuscript~05's, which is \cref{cdc:sec:intro}). The introductions describe each manuscript's results in its own terms and numbering-free words; where they name a theorem, the reference points to where this report prints it.

\subsection{Manuscript 05 (base): Canonical Trace Polytopes}\label{cdc:sec:ms05}
\srcnote{\wtag{} Subtitle: ``Witness-Preserving Diophantine Representations of Computational Substrates''. ``Research memorandum prepared for Vladimir Reshetnikov'', September 30, 2026; pinned at \texttt{e8bb0931d}. Batch~60, manuscript~05 (\path{canonical_trace_polytopes.zip}). Printed in \cref{cdc:sec:intro} (introduction), \cref{cdc:conv:sec} (conventions), Parts~I, II, VII, VIII and~IX, and the back matter.}

% ---- manuscript 05, lines 99-113 ----
\paragraph{Abstract.}
An existential Diophantine representation says that a computation exists; it need not preserve which computation occurred, how many computations occurred, or how many arithmetic witnesses describe it. We construct a stronger, bounded-horizon representation. For a Petri net with $d$ places, $m$ distinctly labelled transitions, and horizon $T$, an explicit convex quadratic integer polynomial has its nonnegative integer zeros in bijection with the commutation classes of enabled length-$T$ executions between prescribed endpoints. The construction uses
\[
(2T+1)d+(7T+1)m
\]
unknowns. Its nonnegative real zero set is a bounded rational polytope, and every slack variable is uniquely determined by the canonical execution. A smaller quartic presentation is also given.

We derive an exact all-marking criterion for permissible commutations, realize arbitrary independence graphs by unit-weight nets, and prove an exponential automaton-state lower bound that explains the use of Boolean memory rather than an explicitly expanded automaton. Counting classes remains $\sharpP$-complete even for unit-weight, $1$-safe nets and their maximal static commutation relation. The same witness-preserving arithmetic methods apply to bounded counter programs, Boolean circuits, cellular automata, and explicitly size-bounded rewriting.

Finally, we prove a precise boundary theorem: a single finite-fold polynomial representation of the decidable verifier for canonical first-halting universal computations is equivalent to the general finite-fold Diophantine conjecture; the analogous statement holds for single-fold representations. The bounded constructions do not cross this boundary. Complete mathematical proofs, an exact-arithmetic reference implementation, reproducible finite checks, and a formalization and research agenda are included.


\begin{resultbox}
\small\textbf{Status.} This is a proved theorem package and executable research prototype, not a claimed resolution of the general finite-fold or single-fold conjecture. Classical ingredients are identified in the text. Global priority for the combined results has not been established, and the new proofs have not been checked in Lean.
\end{resultbox}

\srcnote{\wtag{} Manuscript~05's introduction is \cref{cdc:sec:intro}.}

\subsection{Manuscript 01: Canonical Diophantine Certificates for Causal Computation}\label{cdc:sec:ms01}
\srcnote{\wtag{} Subtitle: ``Exact commuting-block acceleration, quartic trace quotients, and the boundary of single-foldness''. ``Research report prepared for the ProveIt project'', 30 September 2026; pinned at \texttt{e8bb0931d}. Batch~60, manuscript~01 (\path{Diophantine_Causal_Computation.zip}). Printed in Parts~I, II, III, VII, VIII and~IX and the back matter.}

% ---- manuscript 01, lines 51-59 ----
\paragraph{Abstract.}
Ordinary Diophantine representability records whether a computation exists; it need not preserve how many computational objects exist. We develop explicit ordinary polynomial representations which control the entire witness fibre for a substantial class of computational substrates. An action is an integer translation restricted to a coordinatewise lower-and-upper guard. This includes weighted Petri transitions and finite-control counter instructions with zero tests.

Our first result gives the exact initial-resource box in which \emph{every} ordering of a prescribed multiset of actions can execute. When the active actions commute as partial maps, the same box characterizes \emph{some} ordering. For Petri transitions it reduces to total resource loss plus the largest retained resource requirement. The resulting accelerator is a single-fold polynomial of degree at most four, with witness arity independent of the multiplicities. A concrete three-action example uses eleven witnesses even for multiplicities $10^{100},10^{80},10^{90}$.

Our second construction combines these exact resource tests with the classical Cartier--Foata normal form. For each fixed causal-height bound, the nonnegative roots of an explicit quartic are in bijection with executable trace classes, not with their interleaved schedules. All intermediate markings, guard slacks, and normal-form auxiliaries are uniquely determined. We give complete proofs, exact variable and equation counts, block-composition extensions, an executable compiler, and differential tests. Finally, we prove that a single fixed-arity single-fold representation for universal halting would imply the full single-fold MRDP principle. The bounded-height family and commuting-block accelerator do not make that unproved leap.

\vspace{0.3em}
\noindent\textbf{Status.} This is an unrefereed mathematical development with reproducible exact-integer tests, not a Lean-checked extension of ProveIt. The particular resource formulas and compiler assemblies are derived and proved here; their priority in the wider literature is not established. Classical arithmetization, trace normal forms, and the finite-fold conjecture are explicitly distinguished from these constructions.

% ---- manuscript 01, lines 64-73 ----
\subsubsection{The representation problem beyond MRDP}\label{cdc:ct:sec:motivation}
\paragraph{What is already available in ProveIt}
The starting point is the ProveIt repository at the pinned revision
\begin{center}\path{e8bb0931d67f80d9fce87a8cddb0f661ff19f956}.\end{center}
Its Hilbert-tenth development already supplies MRDP, finite-witness integer polynomials, and Diophantine closure under exact finite iteration. In particular, the inspected source \code{Diophantine/Common/DiophantineTrace.lean} contains the interfaces \code{boundedForall_dioph}, \code{exactIter_dioph}, and \code{existsExactIter_dioph}. These are existence statements about Diophantine representations; they do not assert uniqueness of auxiliary witnesses \cite{repo-h10,repo-trace}.

This is an important distinction in selecting a research target. Merely observing that finite executions of an effective substrate form a computably enumerable relation, and then applying MRDP, would add little to that development. Nor would another general register-machine-to-Diophantine proof be new: the repository includes the Jones--Matiyasevich 1984 route, and independent mechanized work uses Minsky machines and FRACTRAN \cite{repo-h10,jm1984,lf2022}.

The present target is \emph{fibre control}: eliminate witnesses that encode nothing more than a choice of interleaving, an oversized coding base, an unconstrained inactive branch, or a noncanonical padding length. This is useful both mathematically, when counting solutions, and architecturally, when designing verified compilers that should preserve a specified semantic object rather than only existential truth.


% ---- manuscript 01, lines 91-109 ----
\paragraph{Questions answered in this report}
The concrete questions are as follows. Can one characterize the exact resources needed to serialize a multiset without inspecting all its permutations? Can this become an ordinary low-degree polynomial with one complete witness tuple? Can a polynomial count causal traces rather than scheduling permutations? And which step would be needed to turn these constructions into a universal single-fold theorem?

The answers are given by the envelope theorem (Theorem~\ref{cdc:ct:thm:envelope}), the fixed-arity accelerator (Theorem~\ref{cdc:ct:thm:accelerator}), the causal-height compiler (Theorem~\ref{cdc:ct:thm:history}), and the universal-halting equivalence (Theorem~\ref{cdc:ct:thm:universal}). These are scoped representation results, not a claim to have settled a previously named major open conjecture.

\begin{center}
\begin{tabular}{p{0.25\textwidth}p{0.62\textwidth}}
\toprule
Component & Role and provenance \\\midrule
MRDP and exact iteration & Existing repository infrastructure; used as context, not as the proof of the explicit compilers.\\
Cartier--Foata layers & Classical normal form; a self-contained proof of the version used here is supplied.\\
Resource envelope & Direct extremal-prefix derivation for arbitrary repeated box-guarded translations.\\
Quartic compilers & Explicit constructions with proofs of all auxiliary-witness uniqueness claims.\\
Global compression & Reduced to the single-fold or finite-fold MRDP principle; not claimed to follow from the bounded construction.\\
Formal verification & Proposed integration work. The attached Python tests are not proof-assistant verification.\\
\bottomrule
\end{tabular}
\end{center}


\subsection{Manuscript 02: Canonical Diophantine Certificates for Polynomial Trajectories}\label{cdc:sec:ms02}
\srcnote{\wtag{} Subtitle: ``Single-fold acceleration, finite-difference sign tables, and the boundary of general computation''. ``Research manuscript for the ProveIt project. Prepared with ChatGPT'', September 30, 2026; pinned at \texttt{e8bb0931d}. Batch~60, manuscript~02 (\path{canonical_diophantine_certificates.zip}). Printed in \cref{cdc:conv:sec} (its Section~2) and Parts~IV, VII, VIII and~IX and the back matter; the first two subsections of its introduction open Part~IV (\cref{cdc:pt:sec:problem}).}

% ---- manuscript 02, lines 59-88 ----
\paragraph{Abstract.}
The existence of a Diophantine representation of finite computation does not
by itself control the number of representing witnesses. We develop a
canonical certificate construction for a substantial restricted class of
computations: trajectories polynomial in the iteration counter, with
arbitrary fixed Boolean combinations of polynomial guards.

The main construction eliminates a bounded universal quantifier over one
integer variable without exponentiation and without an existentially chosen
Chinese-remainder modulus. For every fixed degree bound $d$, the relation
$\forall i\in\{0,\ldots,T\}\;p(i)\geq0$ has a single-fold representation by
one integer polynomial of degree at most four, with $O((d+1)^3)$ natural
witnesses, independently of $T$. The geometric part consists of exactly
$(d+1)(2d+1)$ coordinates. These record the unique maximal sign intervals of
$p,\Delta p,\ldots,\Delta^d p$, with canonical padding. Monotonicity between
successive difference levels reduces verification to endpoint tests.

We prove single-fold representations for exact guarded execution and first
exit of polynomially iterable maps, including unitriangular nonlinear
updates, unipotent affine maps, and a certified eventually quasi-polynomial
affine extension. A canonical root bound also gives single-fold certificates
for nontermination from a specified input. We provide a reproducible
integer-only certificate generator, checker, and compiler to explicit
quadratic residuals whose squared sum is the promised quartic.

The contribution is a constructive restricted result, not a solution of the
general finite-fold or single-fold representation problem. The proofs are
given here; historical priority of this particular construction has not been
established, and no proof-assistant verification is claimed.


% ---- manuscript 02, lines 94-96 ----
\subsubsection{The problem addressed and the contribution}
\label{cdc:pt:sec:intro}


% ---- manuscript 02, lines 161-205 ----
\paragraph{What is inherited from ProveIt}
\label{cdc:pt:subsec:repo}
We inspected the ProveIt repository at commit \path{e8bb0931d67f80d9fce87a8cddb0f661ff19f956}. Its
\path{Computability/HilbertTenthProblem} development contains an MRDP
interface and a finite-trace interface. In particular,
\code{DiophantineTrace.lean} provides ordinary Diophantine closure for
bounded universal quantification, exact iteration, and finite reachability
\cite{repo-trace}. These interfaces do not state uniqueness of the
auxiliary witnesses. The project also maintains corrected editions of the
Jones articles, including the Jones--Matiyasevich register-machine treatment
of exponential Diophantine representation \cite{repo-h10}.

We use this as a precise project-level motivation: ordinary finite
reachability is already available, so another invocation of MRDP would not
constitute the intended advance. The new target is a restricted
\emph{multiplicity-preserving acceleration interface}. None of our proofs
requires assuming an unproved repository theorem. Conversely, source
inspection is not a fresh compilation or transitive axiom audit of the
repository, and we make no such claim.

\paragraph{Relation to prior work and novelty status}
There are two relevant bodies of prior work. Single-fold specifications
require care with arithmetic auxiliary variables and disjunction; the
recent treatment by Cantone, Cuzziol, and Omodeo discusses these issues and
the role of exponentiation \cite{cco2024}. Our total-circuit normalization is
an explicit implementation-oriented formulation of elementary ingredients,
not a claim that canonical arithmetic gadgets are new.

Polynomial and triangular loop analysis also has substantial prior
literature. Hark, Frohn, and Giesl study triangular weakly nonlinear loops
with polynomial Boolean guards and reductions for termination
\cite{hfg2024}. Lommen, Meyer, and Giesl use closed forms and stabilization
thresholds in modular complexity analysis \cite{lmg2024}. We therefore do
not claim to introduce polynomial loop acceleration, polynomial closed
forms, or input-dependent termination bounds.

The result developed here is the particular \emph{canonical finite-difference
certificate}, its explicit cubic witness bound, and its use in single-fold
ordinary polynomial representations of execution. The cited papers were
checked for the surrounding landscape, not exhaustively mined to prove
historical priority. Accordingly, this manuscript offers complete proofs of
a specific constructive theorem, while leaving priority and publication
novelty to further expert comparison. Its usefulness does not depend on
labeling a standard consequence as a breakthrough.


\subsection{Manuscript 03: Linear-Size Canonical Diophantine Certificates for Random-Access and Graph Computation}\label{cdc:sec:ms03}
\srcnote{\wtag{} Subtitle: ``Quartic equations, unique witnesses, and exact path weights''. ``Research note prepared for Vladimir Reshetnikov'', September 30, 2026; pinned at \texttt{e8bb0931d}. Batch~60, manuscript~03 (\path{ProveIt_Linear_Size_Diophantine_Certificates.zip}). Printed in \cref{cdc:conv:sec} and Parts~V and~IX and the back matter.}

% ---- manuscript 03, lines 50-84 ----
\paragraph{Abstract.}
MRDP settles the existence of Diophantine descriptions of effective computation,
but does not automatically preserve individual executions or control artificial
witness multiplicity. We give two explicit canonical compilers for finite
random-access memory logs over the natural numbers. Every valid log has exactly
one auxiliary witness, and every invalid log has none.

The stronger scalar-size result uses exactly $24L$ witnesses and $22L+1$
quadratic residuals for a length-$L$ log. Their sum of squares is a sparse
polynomial of degree at most four. A proposed sorted list is certified by
packing complete records in a uniquely fixed radix and comparing two products
at a base large enough to make coefficient carries impossible. The product
identity is deterministic and exact, not a probabilistic fingerprint. All
large powers are computed by quadratic chains. If the input coordinates have
$b$ bits, the largest witness has $O(L^2(b+\log(L+2)))$ bits.

A complementary low-height compiler uses a canonical two-equation comparison
gadget and a sorting network. Its exact counts are $10C+2L+3(N-1)$ witnesses
and $10C+3L+4(N-1)+1$ residuals for $N$ padded records and $C$ comparators.
It gives explicit $O(L\log^2(L+2))$ size, or $O(L\log(L+2))$ using optimal
networks, while keeping every witness to $b+O(\log(L+2))$ bits.

We prove parsimonious composition with RAM instruction certificates, apply it
to counter, cellular, stack, and immutable combinator-graph computation, and
derive exact execution-count and first-halting probability formulas. We also
prove the boundary: fixed-variable single-fold compression of one universal
halting substrate is equivalent to the general single-fold conjecture. Our
constructions are horizon-indexed families, not that compression. Both memory
backends emit actual sparse quartics and pass reproducible exhaustive and
adversarial tests. The proofs are ordinary mathematical proofs, not newly
Lean-verified results; literature-wide priority is not claimed.

\vspace{0.3em}
\noindent\textbf{Keywords.} Diophantine representation; random-access memory;
parsimonious compilation; sorting networks; combinatory logic; finite traces.

% ---- manuscript 03, lines 89-168 ----
\subsubsection{The research target and its relation to ProveIt}
\paragraph{What is already available}
The repository snapshot used for this investigation is
\begin{center}
\texttt{e8bb0931d67f80d9fce87a8cddb0f661ff19f956}.
\end{center}
Its Hilbert-tenth documentation reports MRDP in both directions, a
register-machine encoding, and reusable interfaces for bounded quantification,
exact iteration, and finite reachability. Its combinatory-logic documentation
also carefully distinguishes positive-step simulations from operational
reflection and full abstraction. These are the relevant starting points
\cite{repo-h10,repo-status,repo-cl}. We inspected these documents at the pinned
revision; we did not rebuild the repository or independently audit its entire
formal development.

The classical register-machine route is represented by Jones and Matiyasevich
\cite{jm1984}. Carneiro's Lean development supplies relevant historical context
for machine-checked Diophantine exponentiation \cite{carneiro}. Neither the
classical theorem nor an existential library interface automatically gives a
one-to-one correspondence between program executions and polynomial witnesses.
That missing contract is the target here.

A statement that a Turing-complete substrate is Diophantine, obtained only by
composing its interpreter with MRDP, would not be a new mathematical result.
Instead we ask a quantitative and intensional question:
\begin{quote}
Can arbitrary finite memory access be represented by a sparse quartic equation,
with exactly one auxiliary witness for each valid access log, without expanding
the entire memory address space or introducing a factor proportional to address
bit width in the number of arithmetic constraints?
\end{quote}
The answer is yes. Sections~\ref{cdc:mem:sec:gadgets}--\ref{cdc:mem:sec:linear} give two complete
constructions: a low-height sorting-network compiler and a linear-size
large-base-product compiler, both with exact counts and executable artifacts.

\paragraph{Prior art and the actual contribution}
Sorted memory transcripts and routing-network checks are established techniques.
In particular, the RAM arithmetization in \emph{SNARKs for C} checks local
transitions and a transcript ordered by address and timestamp, using routing
constraints \cite{snark}. We do not claim to invent that architecture or
quasilinear verification. The distinction pursued here is an exact
natural-number polynomial with a fully specified, unique auxiliary fiber,
including the comparison and routing data. A deterministic comparison network
makes this contract transparent. This is a different guarantee from merely
exhibiting some valid routing.

Batcher's sorting construction and the asymptotically optimal sorting-network
theorem are also established results \cite{batcher,aks}. Our proofs identify the
precise arithmetic interface they need and quantify the complete resulting
Diophantine certificate. The explicit bitonic version is proved and implemented
below. Its improved network bound uses the external optimal-network theorem;
no AKS implementation is supplied. The separate linear-size construction in
Section~\ref{cdc:mem:sec:linear} does not use that theorem or any sorting network.
It uses an elementary carry-free integer evaluation argument to certify a
unique sorted record list, paying for fewer coordinates with greater height.

\begin{center}
\begin{tabularx}{\textwidth}{@{}p{0.24\textwidth}X@{}}
\toprule
Status & Content \\
\midrule
Established background & MRDP; sorting-based RAM verification; Batcher and AKS
sorting networks. \\
Proved in this note & Two exact canonical memory compilers; linear scalar size;
witness and residual counts; size--height tradeoffs; composition and weighted-path consequences. \\
Executable artifact & A general symbolic memory-log compiler, sparse quartic
export, canonical witness generation, and deterministic regression tests. \\
Not established here & Literature-wide novelty of every formulation;
fixed-variable single-fold representation of arbitrary halting; a formal
verification of this new compiler; full abstraction of ProveIt's existing
lambda-to-combinator translations. \\
\bottomrule
\end{tabularx}
\end{center}

This status distinction matters. The result is a concrete research extension
with complete proofs, not a declaration that a famous open conjecture has been
solved. Its potential value to ProveIt is a stronger reusable specification:
existence, uniqueness, decoding, and quantitative bounds travel together.


\subsection{Manuscript 04: Witness-Faithful Diophantine Compilation}\label{cdc:sec:ms04}
\srcnote{\wtag{} Subtitle: ``Concurrent traces, priority arithmetic, and context-local rewriting''. ``Research article prepared for Vladimir Reshetnikov'', September 30, 2026; pinned at \texttt{e8bb0931d}. Batch~60, manuscript~04 (\path{Witness_Faithful_Diophantine_Compilation.zip}). Printed in \cref{cdc:conv:sec} and Parts~II, VII, VIII and~IX and the back matter.}

% ---- manuscript 04, lines 45-51 ----
\paragraph{Abstract.}
Existential Diophantine representability is already available for every computably enumerable computational semantics. A more discriminating question is whether arithmetic solutions faithfully represent the intended computational objects, without spurious choices of auxiliaries or repeated encodings of the same concurrent execution. We develop three explicit constructions addressing this question. For a Petri net with $p$ places and $r$ labeled transitions, exact $T$-step executions between fixed endpoints, modulo a specified disjoint-support commutation relation, are in bijection with the natural zeros of a polynomial of degree at most four in $(p+r)(2T+1)$ variables. The construction uses a quadratic, $r$-bit monitor for lexicographic trace normal forms; a matching linear-in-alphabet memory lower bound follows from an explicit family requiring exponentially many deterministic automaton states. For an ordered $r$-fraction FRACTRAN program, an exact run is represented uniquely using $T(7r+2)+1$ variables and $T(8r+3)+2$ quadratic residuals. For SKI rewriting along a fixed schedule, a context-local encoding uses whole-subterm integer codes, with size linear in the schedule and its total context depth, rather than in the number of tree nodes traversed by a full tableau.

All constructions include two-way proofs, exact resource counts, and executable exact-integer certificates. We distinguish these bounded, witness-faithful constructions from fixed-arity single-fold and finite-fold representations of unbounded computation. The latter are not obtained by determinism alone. The article identifies the precise compression frontier, gives a witness-height obstruction, and proposes a research program linking these constructions to the existing computability developments in ProveIt. Global priority for the particular constructions is not established; no resolution of the general finite-fold or single-fold problem is claimed.

\vspace{3pt}
\noindent\textbf{Keywords.} Diophantine representation; arithmetic certificates; Petri nets; trace monoids; FRACTRAN; combinatory logic; single-fold representations.

% ---- manuscript 04, lines 66-107 ----
\subsubsection{The problem beyond existential representability}
\paragraph{What is already settled}
The Matiyasevich--Robinson--Davis--Putnam theorem, abbreviated MRDP, identifies computably enumerable subsets of the natural numbers with Diophantine sets \cite{mat1970}. Thus a computational substrate with an effective finite-execution semantics generally presents no new \emph{existence} problem: its finite reachability relation is computably enumerable and hence Diophantine. The constructive route through register machines has a substantial history; Jones and Matiyasevich give an exponential-Diophantine register-machine construction, and Larchey-Wendling and Forster's full mechanization uses FRACTRAN as an intermediate language \cite{jm1984,lf2019}. Re-proving that FRACTRAN, SKI reduction, or Petri-net reachability is merely Diophantine would therefore miss the interesting issue.

The inspected ProveIt snapshot already contains MRDP interfaces, bounded universal quantification, exact iteration of a Diophantine relation, and finite reachability \cite{repo-h10,repo-trace}. Its combinatory-logic development supplies concrete operational semantics and compiler simulations \cite{repo-cl}. These are appropriate foundations, but ordinary existential equivalence says nothing about how many arithmetic witnesses encode one execution.

We instead ask for an explicit map
\[
 \{\text{intended computational objects}\}
 \longleftrightarrow
 \{\text{natural solutions of an arithmetic certificate}\}
\]
that is a \emph{bijection}. Depending on the substrate, the intended object is a labeled sequential run, a commutation class of concurrent runs, or a reduction along a specified schedule. The distinctions matter. Two different firing words may describe the same concurrent trace, while two distinct traces may have the same endpoint and the same transition counts.

\paragraph{The contributions and their scope}
The principal construction is a direct arithmetic quotient of Petri-net executions by independent swaps. It does not first enumerate a potentially exponential-state normal-form automaton. A bit for each transition records exactly the information needed to prohibit a nonminimal linearization. The resulting polynomial has one solution per concurrent trace class, not one solution per scheduling permutation.

The second construction handles \emph{priority}. FRACTRAN is deterministic only after the least applicable fraction is selected. A polynomial that merely allows any applicable fraction models a different language. We build the least-index choice into a cumulative prefix circuit and force every division and inactive slack to be canonical.

The third construction handles \emph{locality}. A single integer represents a whole SKI subtree. A prescribed contextual contraction need decompose only the path to the redex and the finite pattern at its root. The result is a compact certificate in variable count, but not necessarily in bit length. This is a useful separation of two notions of succinctness, not a claim that integer coding removes all costs of large terms.

\begin{center}
\small
\begin{tabular}{p{.25\linewidth}p{.39\linewidth}p{.27\linewidth}}
\toprule
Construction & Object represented & Natural-solution multiplicity \\
\midrule
Petri net, fixed $T$ & A class of enabled words under the chosen independent swaps & Exactly one per class \\
FRACTRAN, fixed $T$ & The exact priority-respecting run & Zero or one \\
SKI, fixed schedule & The scheduled contextual reduction & Zero or one \\
Unbounded reachability & Existence of some finite execution & MRDP supplies existence, not these multiplicity guarantees \\
\bottomrule
\end{tabular}
\end{center}

The exact compiler theorems, their joint interpretation, and the resource accounting are the results developed in this article. The normal-form criterion, circuit encodings, and general MRDP implications are established background or elementary consequences, and are identified as such. The constructions have not undergone independent peer review or a new proof-assistant check. The accompanying tests are reproducibility evidence, not substitutes for the proofs.

\paragraph{Related work and a necessary distinction}
Lexicographic trace normal forms have a classical forbidden-factor characterization, often attributed to Anisimov and Knuth; a modern use appears in Lohrey, Stober, and Wei\ss\ \cite{lohrey2024}. Counting trace classes is also an active subject. De Colnet, Meel, and Mathur study counting and sampling the classes intersected by a regular language, presented by a finite automaton, and obtain complexity and approximation results \cite{decolnet2026}. Our purpose is different: a succinct operational model is compiled into an exact polynomial with controlled fibers. No new general trace-counting principle is asserted.

A subtlety is that a class intersecting an arbitrary language need not have its lexicographically least member in that language. Our Petri-net endpoint languages are \emph{trace-closed} for the specified independence relation. This closure, proved below, is what makes normal-form filtering preserve the desired set of classes.


\subsection{Manuscript 06: Saturation and Single-Fold Diophantine Compilation of Concurrent Counter Programs}\label{cdc:sec:ms06}
\srcnote{\wtag{} Subtitle: ``Resource summaries, canonical traces, and nested repetition schemes''. ``Research article. Prepared for Vladimir Reshetnikov'', September 30, 2026; pinned at \texttt{e8bb0931d}. Batch~60, manuscript~06 (\path{diophantine_counter_programs.zip}). Printed in \cref{cdc:conv:sec} and Parts~I, II, III, VIII and~IX and the back matter.}

% ---- manuscript 06, lines 79-114 ----
\paragraph{Abstract.}
\noindent We give an explicit, witness-preserving route from structured concurrent
counter executions to ordinary polynomial equations over the nonnegative integers.
A firing word is summarized by the least initial resources it needs and the
resources it leaves at that least input. This classical partial-translation
algebra admits exact composition and arbitrary-power formulas with unique
quadratic witnesses. Independently, we construct a Boolean transfer signature
for lexicographic trace normality and prove that every signature satisfies
$S^2=S^3$. Consequently, in an arbitrarily nested repetition expression,
normality depends on each repetition parameter only through the three cases
$0$, $1$, and at least $2$.

Combining the two summaries gives a polynomial of degree at most four with
$O((p+q)N)$ auxiliary variables for a schedule with $N$ syntax nodes, $p$
resources, and $q$ transition labels. For fixed endpoints and repetition
parameters, there is exactly one auxiliary solution precisely when the
specified execution is valid and, optionally, already canonical. The expanded
execution length does not occur in this variable bound. A separate bounded-word
construction gives a genuine bijection between polynomial solutions and
executable trace classes.

We also locate a precise expressiveness boundary. Flat repetition schemes have
Presburger-definable execution relations, whereas separated increment/decrement
schemes of nesting depth two represent arbitrary systems of quadratic equations.
For one counter, separated depth-$d$ schemes correspond exactly to degree-$d$
polynomial equations. Finite-fold and single-fold Diophantine representability
are equivalent to the corresponding bounded-multiplicity parameterizations by
separated one-counter schemes of depth four. These equivalences identify the
remaining obstruction; they do not resolve the finite-fold or single-fold
conjectures. An executable compiler and 49,722 finite checks accompany the proofs.


\noindent\textbf{Status.} The article contains mathematical proofs and tested
prototype code, not a new kernel-checked formalization. Classical ingredients are
identified explicitly. Historical priority for the combined constructions and
formulations is not established by this study.

% ---- manuscript 06, lines 122-231 ----
\subsubsection{Purpose, context, and the main results}
\label{cdc:cs:sec:intro}

A general existence theorem can hide exactly the information needed by a
compiler. The Matiyasevich--Robinson--Davis--Putnam theorem says that every
computably enumerable relation has a Diophantine representation \cite{mat1970}. It does not,
by itself, say how to preserve a particular execution, how many witnesses that
execution acquires, how large those witnesses are, or whether two witnesses
encode the same concurrent behavior. These distinctions become essential when
Diophantine representation is used as a computational substrate rather than
only as an undecidability theorem.

The motivating repository, \emph{ProveIt}, already documents MRDP in both
directions, finite-trace and bounded-quantifier interfaces, register-machine
encodings, and universal Diophantine equations \cite{repo}. Its relevant status register
also distinguishes completed formal results from remaining obligations
\cite{repo}. Reproving that a computable interpreter has a Diophantine graph
would add little to that program. We instead develop a direct compiler with a
precise contract: a specified structured execution receives \emph{one and only
one} auxiliary certificate, with a degree and size bound visible from its
syntax.

There are two logically different sources of multiplicity. An arithmetic
encoding can create many witnesses for one computation. Independently, a
concurrent computation can have many interleavings, or a repetition expression
can have many parameter tuples that expand to the same word. We eliminate the
first multiplicity. For all words of a fixed length we can also eliminate the
interleaving multiplicity. We do not silently identify either achievement with
eliminating multiplicity from arbitrary existential repetition parameters.

\paragraph{The three principal constructions}

For a finite collection of token-consuming and token-producing transitions, a
word induces a partial translation of $\N^p$. Its domain is a principal upper
orthant. The least input and its output provide an exact resource summary.
Composition uses coordinatewise cancellation; repetition has a closed formula
including the exceptional zero-power case. All local choices can be expressed
by uniquely solvable quadratic systems.

The second construction concerns a different object: the lexicographically
smallest word in a trace class. The familiar normal-form language is regular,
but constructing a full transition table on all subsets of the alphabet would
be wasteful. We describe its block action with $3q+1$ Boolean coordinates,
subject to a disjointness invariant. This description admits linear-size
composition circuits. More importantly, its powers stabilize after the second
power. The stabilization remains valid inside arbitrary nested repetitions.

The third construction combines these branches without conflating them. The
resource summary is invariant under the permitted commutations. The
normal-form signature is intentionally \emph{not} invariant: it distinguishes
the chosen representative from the other words in the same trace. Their
quadratic constraints can nevertheless be joined in a single sum of squares.

\begin{theorem}[Summary of the compilation results]
\label{cdc:cs:thm:overview}
Let $E$ be a fixed parametrized repetition expression with $N$ syntax nodes
over $q$ transitions acting on $p$ nonnegative counters. Let $I$ be a fixed
symmetric, irreflexive relation of globally commuting transition labels.
\begin{enumerate}[label=\textnormal{(\roman*)},leftmargin=2em]
\item Whether the expansion of $E$ is lexicographically least in its $I$-trace
class is unchanged when every repetition parameter $k$ is replaced by
$\min(k,2)$.
\item One can construct an integer polynomial of degree at most four with
$O((p+q)N)$ auxiliary variables whose natural solutions, for fixed endpoints
and repetition parameters, form a singleton exactly when the specified
execution is valid and already lexicographically least. Without the normality
requirement the auxiliary bound is $O(pN)$.
\item For fixed length $T\geq 1$, one can instead construct a quartic
polynomial with
\[
q(2T+1)+p(7T-4)
\]
witness variables whose solutions are in bijection with the executable
$I$-trace classes of length $T$ between the specified endpoints.
\end{enumerate}
\end{theorem}

Part (i) is proved in \cref{cdc:cs:sec:saturation}, part (ii) in
\cref{cdc:cs:sec:compiler}, and part (iii) in \cref{cdc:cs:sec:counting}. A second group of
results, in \cref{cdc:cs:sec:hierarchy,cdc:cs:sec:finitefold}, characterizes exactly where
unbounded parameter search restores the full difficulty of Diophantine
solvability.

\paragraph{What is classical, and what is proved here}

The algebra of one partially blind counter is the bicyclic monoid; products of
such storage components are standard in the study of valence automata
\cite{zetzsche}. Lexicographic trace normal forms go back to Anisimov and Knuth
\cite{anisimov-knuth}. The conversion of arithmetic circuits into quadratic
systems, followed by a sum of squares, is a standard degree-reduction method.
Single-fold representation and its closure issues have an extensive literature
\cite{mat2010,cco2024}. None of these ingredients is
presented as a new discovery.

The contributions established in this article are the explicit synchronized
construction, its local-to-global uniqueness proof, the compact normality
signature and its cap-at-two consequence for nested schemes, the stated
witness and trace-counting bounds, and the solution-preserving substrate
normal forms. We give full proofs rather than relying on novelty labels.
A targeted literature review did not establish historical priority for every
equivalent formulation. In particular, a useful new implementation or
packaging of known algebra should not be confused with a first theorem about
that algebra.

All statements below are independent of any unverified theorem elsewhere in
ProveIt. The repository is the motivation and a prospective integration target.
The source audit used its tree at commit
\texttt{e8bb0931d67f80d9fce87a8cddb0f661ff19f956}; it was not a rebuild or an
independent audit of the repository's proof kernel dependencies.


\subsection{Manuscript 07: Causality Without Tableaux: Compact Diophantine Certificates for Queue Universality}\label{cdc:sec:ms07}
\srcnote{\wtag{} Subtitle: ``Product guards, tag systems, and optimal polynomial word-memory dimension''. ``Research article. Prepared for Vladimir Reshetnikov'', September 30, 2026; pinned at \texttt{725d2ebb6}. Batch~61, its only manuscript (\path{Causal_Diophantine_Queue_Compilation.zip}); numbered 07 in this report. Printed in \cref{cdc:conv:sec} and Parts~I, VI and~IX and the back matter.}

% ---- manuscript 07, lines 61-71 ----
\paragraph{Abstract.}
\noindent We construct ordinary integer polynomials whose natural zeros faithfully encode finite computations of cyclic tag systems and general deletion-tag systems. The main cyclic-tag construction uses exactly $T+1$ variables for a $T$-step execution, has degree at most $\max\{4,2T\}$, and admits a linear-size arithmetic circuit. Its variables are the consumed bits and one uniquely determined slack; no intermediate queue is quantified. The key is a causal reconstruction theorem: one global word identity becomes sufficient when accompanied by a product guard. A first-failure lemma proves that this guard is positive for legal executions and exactly zero for every illegal candidate, even if later formal resource values become negative.

Two alternatives expose the cost tradeoff. A bounded-root construction removes the slack at the expense of a larger circuit. A local radix construction yields a quartic polynomial in $3T-2$ variables. The general deletion-$d$ construction uses $dT+1$ variables, including all consumed symbols, with an explicit degree bound and explicit denominator clearing. Fixed-read FIFO networks admit a corresponding consumed-symbol-plus-one bound.

A separate obstruction explains the radix architecture. Finite words admit a faithful two-coordinate natural-number encoding with affine insertion at either end. No faithful one-coordinate natural-number encoding can make both binary prepends and even one fixed append polynomial. An elementary positive-leading-coefficient centralizer argument proves this dimension lower bound.

The article gives complete mathematical proofs, executable compilers, an independent exact-integer verifier, finite exhaustive checks, and a proposed formalization and research program. These are horizon-indexed, witness-faithful constructions, not a fixed-arity finite-fold representation of unbounded halting. Historical priority for the particular constructions is not established; no general finite-fold conjecture is claimed solved.


\noindent\textbf{Verification status.} The proofs below are mathematical proofs and the supplied Python artifacts have been tested. This package does not contain a new kernel-checked Lean or Rocq formalization. Repository audit reports are distinguished from checks performed for this article.

% ---- manuscript 07, lines 87-124 ----
\subsubsection{The question beyond representability}
\paragraph{Why another Diophantine construction?}
The Matiyasevich--Robinson--Davis--Putnam theorem identifies computably enumerable sets of natural numbers with Diophantine sets. Thus, once a computational substrate has an effective finite-execution semantics, the mere existence of a Diophantine representation of its halting relation is not a new universality theorem. The interesting questions concern the structure of the representing polynomial: its variables, degree, coefficient size, circuit size, and the relationship between its solutions and actual computational histories. The formal developments of Carneiro and of Larchey-Wendling and Forster provide important precedents for making such reductions precise \cite{carneiro,lf2019}.

ProveIt's Hilbert-tenth-problem development already supplies the forward and reverse MRDP interfaces, finite polynomial representations, and trace infrastructure \cite{repo-mrdp,repo-h10}. We therefore do not repackage ``this machine is effective, hence MRDP applies'' as a new result. Instead we ask three concrete questions.

Can a certificate for a variable-length queue avoid storing every intermediate queue? Can a global equality between the words consumed and produced certify an execution without permitting a symbol to create itself after an illegal read? Finally, can a natural-number memory with polynomial operations at both ends really be reduced from two coordinates to one? Theorems~\ref{cdc:qc:thm:stream}, \ref{cdc:qc:thm:first-failure}, and \ref{cdc:qc:thm:one-dimensional} answer these questions in explicit models.

Cyclic tag systems are a useful test case: their binary queue semantics is very small, yet the class is computationally universal. Cook uses them as an intermediate substrate in his universality proof for Rule~110 \cite{cook}. Tag-system simulation also has a quantitative literature; Woods and Neary establish polynomial-time simulation results rather than merely abstract universality \cite{woods-neary}. We use these existing universality results as context. Our proofs concern arithmetic certificates for the precisely defined queue dynamics below.

\paragraph{The results at a glance}
For a fixed cyclic-tag program, input word, terminal word, and positive horizon $T$, the main construction represents the exact $T$-step execution relation. For empty terminal word it represents first halting at exactly $T$. Every successful instance has exactly one natural solution.

\begin{center}
\small
\begin{tabularx}{\textwidth}{@{}Xccc@{}}
\toprule
Construction & Variables & Degree bound & Residuals\\
\midrule
Causal word identity & $T+1$ & $\max\{4,2T\}$ & $T+3$\\
Bounded-root, no slack & $T$ & $2\max\{2,T,H\}$ & $2T+2$\\
Local radix, empty endpoint & $3T-2$ & $4$ & $3T$\\
\bottomrule
\end{tabularx}
\end{center}
Here $H$ is an explicit upper bound on queue length, defined in Section~\ref{cdc:qc:sec:tradeoffs}. The first and third rows have $O(T)$ arithmetic-circuit size for a fixed program; the second has $O(T(1+H))$ size. All three denote ordinary polynomials, not expressions with variable exponents or an unbounded product operator.

For deletion-$d$ tag systems over an alphabet of size $B$, we give $dT+1$ variables and degree at most
\[
 2\max\{B,\ T(B-1),\ d(T-1)(B-1),\ 1\}.
\]
The separate memory theorem is an actual lower bound, but in a carefully specified architecture: faithful word encodings with polynomial insertion maps. No lower bound on the number of variables of arbitrary Diophantine representations is asserted.

\paragraph{Relation to earlier work and scope of novelty}
\wtag{} Manuscripts~04 and~06 of this report treat concurrent resource traces, priority arithmetic in FRACTRAN, context-local SKI rewriting and compressed counter schedules (Parts~II, III and~VII). This article instead treats FIFO word content, temporal availability of produced symbols, and the dimension of polynomial word memory. The companion articles are not proof dependencies.

Radix encodings, sums of squares, finite-table interpolation, and polynomial centralizers are established mathematical tools. Our contribution is the explicit construction and proof of the displayed resource guarantees, the causal obstruction and its repair, and the sharply stated word-memory consequence. The centralizer proof is self-contained; the broader theory has a substantial literature, for example Atela and Hu \cite{atela-hu}. A targeted literature check does not establish historical priority for every formulation here. These results should be evaluated as proved constructions and a proved architectural obstruction, not as a claim that a named longstanding conjecture has been settled.


\part{Exact commutation and resource algebra}\label{cdc:part:resource}

\paragraph{Source and scope.} \wtag{} This Part prints four manuscripts' treatments of one question: when does a sequence of guarded resource updates execute, and when may two adjacent updates be exchanged at every state? Manuscript~05 gives the exact all-marking commutation criterion for Petri transitions, realizes every independence graph, and, in a subsection moved here from its section on other substrates, extends the criterion to interval-guarded translations (\cref{cdc:thm:commute,cdc:thm:graph,cdc:thm:interval}). Manuscript~01 proves the same box criterion for its guarded translations (\cref{cdc:ct:prop:boxes}) and its scalar Petri classification, and then the exact resource envelope of all serializations of a multiset (\cref{cdc:ct:thm:envelope}). Manuscript~06 develops the resource-summary algebra of consume/produce words: principal partial translations, composition by cancellation, exact powers, and its own global commutation test (\cref{cdc:cs:thm:commutation}). Manuscript~07 closes the Part with its first-failure product (\cref{cdc:qc:thm:first-failure}), used by the queue certificates of Part~VI.

\paragraph{Duplicates.} \wtag{} The commutation criterion is proved three times. Manuscript~05's \cref{cdc:thm:commute,cdc:thm:interval} are printed as the main statements, since they name the maximal relation and cover interval guards; manuscript~01's \cref{cdc:ct:prop:boxes} has the same box formulas \eqref{cdc:ct:eq:pair-lower}--\eqref{cdc:ct:eq:pair-upper} as \cref{cdc:thm:interval} and is kept, with its own proof, as a second route; manuscript~06's \cref{cdc:cs:thm:commutation} states the criterion through least inputs, $\min(u_b,v_a)=\min(u_a,v_b)$, which is the counter form of $D_{ab}=D_{ba}$. Manuscript~01's envelope and manuscript~07's first-failure product are different theorems, proved by the same induction on formal prefix states (as is manuscript~01's \cref{cdc:ct:prop:word-power} in Part~III); manuscript~06 exponentiates one word exactly (\cref{cdc:cs:thm:power}), and manuscript~01 does so in \cref{cdc:ct:prop:word-power}.

\paragraph{Setting and hypotheses.} \wtag{} Each manuscript keeps its own model. Manuscript~05: a Petri net with $d$ places and $m$ transitions, consumption $c_a$ and production $v_a$; interval guards $l_a\le x\le u_a$ with translation updates in \cref{cdc:thm:interval}. Manuscript~01: box-guarded translations $F_a(x)=x+v_a$ on $\ell_a\le x\le u_a$ in $\N^p$, with lower and upper guards, $m$ actions. Manuscript~06: token-consuming and token-producing transitions on $p$ counters with $q$ labels, and no upper guards. Manuscript~07: candidate event sequences on $p$ resources whose formal integer histories are allowed to become negative after an illegal step.

% ---- manuscript 05, lines 198-266 ----
\section{An exact commutation test for Petri transitions\srctag{05}}
\label{cdc:sec:commutation}
\subsection{Model and sequential demands}
Let $\Net$ have $d\ge0$ places and a finite ordered alphabet
$\A=\{1,\ldots,m\}$ of transitions, where $m\ge1$. Transition $a$ has consumption and production vectors
\[
c_a,v_a\in\N^d.
\]
At a marking $x\in\N^d$ it is enabled exactly when $x\ge c_a$ coordinatewise, and its result is
\[
a(x)=x-c_a+v_a.
\]
Thus every label acts as a partial function. In a word $ab$, transition $a$ is performed first.

For a vector $y\in\Z^d$, write $y_+$ for its coordinatewise positive part. Define the \emph{sequential demand}
\begin{equation}\label{cdc:eq:demand}
D_{ab}=c_a+(c_b-v_a)_+
      =\max(c_a,c_a+c_b-v_a).
\end{equation}

\begin{theorem}[Exact all-marking commutation criterion]\label{cdc:thm:commute}
The two-step word $ab$ is enabled at $x$ if and only if $x\ge D_{ab}$. Transitions $a,b$ commute as partial functions if and only if
\begin{equation}\label{cdc:eq:imax}
D_{ab}=D_{ba}.
\end{equation}
Consequently
\[
I_{\max}=\{(a,b):a\ne b,\ D_{ab}=D_{ba}\}
\]
is the maximal symmetric irreflexive, context-independent relation whose adjacent swaps preserve enabledness and final markings for all starting markings. It can be computed using $O(dm^2)$ coordinate comparisons and arithmetic operations.
\end{theorem}
\begin{proof}
The two enabling conditions are $x\ge c_a$ and
$x-c_a+v_a\ge c_b$. Their conjunction is exactly
$x\ge\max(c_a,c_a+c_b-v_a)=D_{ab}$. Whenever both orders are enabled they have the same result,
\[
x-c_a+v_a-c_b+v_b.
\]
They are therefore the same partial function exactly when their domains agree. Two principal coordinate orthants $\{x:x\ge D\}$ and $\{x:x\ge E\}$ agree exactly when $D=E$: inclusion in one direction forces the corresponding coordinate inequality by testing the corner of the smaller orthant. This proves necessity as well as sufficiency.

A swap inside a longer word has a common prefix marking. Equality of the two partial maps preserves whether the two-step fragment exists and, if it exists, the marking at the beginning of the unchanged suffix. Hence every pair in the displayed relation is sound in all contexts. Conversely, a relation sound in every context is sound for the empty prefix and suffix at every marking, so every one of its pairs satisfies \eqref{cdc:eq:imax}. Computing both demand vectors for each pair gives the operation bound.
\end{proof}

The operation bound is not a claim of constant-cost arithmetic on unbounded integers; the bit complexity includes the arc-weight bit lengths. Only additions, comparisons, and positive parts of input-size integers are needed.

\begin{example}[Why effect vectors are insufficient]
On one place let $a$ produce a token and let $b$ consume a token:
\[
(c_a,v_a)=(0,1),\qquad(c_b,v_b)=(1,0).
\]
The net changes $+1$ and $-1$ commute algebraically, but $ab$ is enabled at marking zero and $ba$ is not. Indeed $D_{ab}=0$ while $D_{ba}=1$.
\end{example}

\begin{remark}[Sequential commutation is not simultaneous independence]
The relation in \cref{cdc:thm:commute} means equality of sequential partial maps. It is not a claim that the two transitions can fire simultaneously, nor that their token occurrences are causally independent in an unfolding. Two pure consumers sharing a place satisfy $D_{ab}=c_a+c_b=D_{ba}$, although they compete for the same resource. The quotient in this article is explicitly the sequential commutation quotient.
\end{remark}

\subsection{There is no hidden graph restriction}
\begin{theorem}[Realizing arbitrary independence graphs]\label{cdc:thm:graph}
Every symmetric irreflexive relation $I$ on $m$ labels is exactly $I_{\max}$ for a Petri net with unit arc weights and at most $\binom m2$ places.
\end{theorem}
\begin{proof}
For each unordered nonedge $\{a,b\}$, with $a<b$, introduce a private place $p_{ab}$. Transition $a$ produces one token there, transition $b$ consumes one token there, and all other transitions leave it untouched. At this coordinate $D_{ab}=0$ and $D_{ba}=1$, so the nonedge is not in $I_{\max}$.

If $\{a,b\}$ is an edge of $I$, no private place is assigned to this pair. At every existing place, at most one of these two transitions acts. A transition and an identity action have equal demands in either order. Thus all coordinates satisfy $D_{ab}=D_{ba}$, and the pair belongs to $I_{\max}$. The arcs have weights zero or one, and there is one place per nonedge.
\end{proof}

For a complete graph, the construction has zero places and all transitions are identity actions. An unused dummy place can be added when a convention requires a nonempty place set; this has no semantic effect. The theorem concerns commutation at all markings, not boundedness or safety from a particular initial marking.


\subsection{Exact commutation with interval guards}\label{cdc:sec:interval}
\srcnote{\wtag{} This subsection is moved from manuscript~05's section on other substrates (\cref{cdc:sec:substrates}, Part~VIII), where it follows the treatment of counter and register programs with zero tests; its corollary for guarded programs is in Part~II (\cref{cdc:cor:guarded}). Manuscript~01's \cref{cdc:ct:prop:boxes} below proves the same box formulas.}

% ---- manuscript 05, lines 666-693 ----
The Petri demand formula cannot simply be reused after adding zero tests. Domains now contain equalities as well as lower bounds, and a previously commuting update may destroy a guard. The needed stronger test is still explicit for interval-guarded translations.

\begin{theorem}[Exact commutation with interval guards]\label{cdc:thm:interval}
Suppose transition $a$ has domain
\[
D_a=\{x\in\N^d:l_a\le x\le u_a\}
\]
and update $x\mapsto x+\delta_a$, where $l_a\in\N^d$,
$u_a\in(\N\cup\{\infty\})^d$, $l_a\le u_a$, and
$l_a+\delta_a\in\N^d$. Define
\begin{align*}
L_{ab}&=\max(l_a,l_b-\delta_a),\\
U_{ab}&=\min(u_a,u_b-\delta_a),
\end{align*}
using $\infty-t=\infty$. The domain of $ab$ is the box
$[L_{ab},U_{ab}]\cap\N^d$, declared empty if any lower endpoint exceeds its upper endpoint. Two transitions commute as partial maps exactly when either both composed boxes are empty, or both are nonempty and
\[
L_{ab}=L_{ba},\qquad U_{ab}=U_{ba}.
\]
The maximal such relation is computable with $O(dm^2)$ coordinate operations.
\end{theorem}
\begin{proof}
The two domain conditions are $l_a\le x\le u_a$ and
$l_b\le x+\delta_a\le u_b$. Intersecting their coordinate intervals gives exactly the stated box. The final update, when defined, is $x+\delta_a+\delta_b$ in either order. Empty compositions agree as partial functions. For nonempty integer boxes, equality of domains is equivalent to equality of all lower and upper endpoints: each finite extreme is attained, and an unbounded coordinate is distinguished from a bounded one by an arbitrarily large coordinate value. This proves the criterion and the direct pairwise algorithm.
\end{proof}

Lower guards express positivity and consumption; upper guard zero expresses a zero test. Finite control can itself be a coordinate with lower and upper endpoints both equal to the source state, and update equal to target minus source. Thus the interval model includes the usual zero-testing counter programs, not just ordinary Petri nets. Reset and copy updates are not translations and are not covered by this theorem.


% ---- manuscript 01, lines 110-280 ----
\section{Guarded translations and exact commutation\srctag{01}}\label{cdc:ct:sec:actions}
\subsection{A substrate with natural-number states}
Fix a dimension $p\ge1$ and a finite nonempty action alphabet $A$ of size $m$. Each action $a$ has a lower guard $\ell_a\in\N^p$, a displacement $v_a\in\Z^p$, and an upper guard $u_a\in(\N\cup\{\infty\})^p$. We require, coordinatewise,
\begin{equation}\label{cdc:ct:eq:guard-valid}
 \ell_{aj}\ge\max(0,-v_{aj}),\qquad \ell_{aj}\le u_{aj}.
\end{equation}
The action is the partial map
\begin{equation}\label{cdc:ct:eq:action}
 F_a(x)=x+v_a\quad\text{defined precisely when }\ell_a\le x\le u_a.
\end{equation}
Condition~\eqref{cdc:ct:eq:guard-valid} guarantees natural outputs. An infinite upper guard imposes no upper restriction.

For a word $w=a_1\cdots a_t$, execution means that every successive partial map is defined. Its endpoint, when it exists, is
$x+\sum_{i=1}^t v_{a_i}$. The endpoint equation alone is insufficient: intermediate guards can fail even if the final vector is nonnegative.

A weighted Petri transition with consumption $c_a\in\N^p$ and production $d_a\in\N^p$ is obtained by taking
\[
 \ell_a=c_a,\qquad v_a=d_a-c_a,\qquad u_a=\infty.
\]
Upper guards also permit exact-value tests, in particular a zero test $\ell_{aj}=u_{aj}=0$. This elementary model is sufficiently broad to expose the distinction between resource feasibility and mere net displacement.

\subsection{An exact and computable independence relation}
We call distinct actions independent when they commute \emph{as partial maps}, including equality of domains:
\begin{equation}\label{cdc:ct:eq:commute}
 F_b\circ F_a=F_a\circ F_b.
\end{equation}
This is stronger than equality of final translations on the common domain. All translations have commuting displacements; their guards need not commute.

\begin{proposition}[Exact box test; manuscript~01, second route to \cref{cdc:thm:interval}]\label{cdc:ct:prop:boxes}
The domain of ``first $a$, then $b$'' is the box with bounds
\begin{align}
 L_{ab,j}&=\max(\ell_{aj},\ell_{bj}-v_{aj}),\label{cdc:ct:eq:pair-lower}\\
 U_{ab,j}&=\min(u_{aj},u_{bj}-v_{aj}).\label{cdc:ct:eq:pair-upper}
\end{align}
It is empty if $L_{ab,j}>U_{ab,j}$ for some $j$. Two actions commute if and only if both composition boxes are empty, or both are nonempty and their lower and upper bounds agree.
\end{proposition}
\begin{proof}
The two enabling conditions are $\ell_a\le x\le u_a$ and $\ell_b\le x+v_a\le u_b$. Intersecting them gives the displayed box. Whenever both compositions are defined, their outputs are the same vector $x+v_a+v_b$. Thus equality of partial maps is exactly equality of their domains. Nonempty boxes in $\N^p$ of this form have uniquely determined finite lower endpoints and finite or infinite upper endpoints, so equality is detected coordinatewise. Empty boxes must be compared as empty sets, not by their possibly different displayed bounds.
\end{proof}
Testing one pair requires $O(p)$ integer comparisons and additions. Testing all unordered pairs gives the largest static irreflexive independence relation in $O(pm^2)$ such operations. Any symmetric irreflexive subrelation of this relation can be used below, at the cost of retaining more distinct traces.

\begin{remark}[Not physical simultaneous firing]\label{cdc:ct:rem:not-step}
Our independence concerns interchangeable \emph{sequential} schedules. It is not the Petri-net step convention that reserves all input tokens simultaneously. For example, two named transitions each consuming one token and producing none commute as partial maps: either order is defined exactly at markings at least two. They are both individually enabled at marking one, but their two-event word is not executable there. The resource theorem below retains this distinction.
\end{remark}

\subsection{A closed classification for one Petri coordinate}
The box test has a useful explicit specialization. Let $(c,d)$ and $(e,f)$ be the consumption--production pairs of two actions on one coordinate, and put $v=d-c$, $w=f-e$.
\begin{proposition}[Scalar Petri commutation; the one-coordinate case of \cref{cdc:thm:commute}]\label{cdc:ct:prop:petri-commute}
The two lower-guarded translations commute precisely in the cases in Table~\ref{cdc:ct:tab:scalar}. For vector-valued Petri transitions the criterion must hold in every coordinate.
\end{proposition}
\begin{table}[ht]
\centering
\begin{tabular}{ll}
\toprule
Displacement signs & Necessary and sufficient condition \\\midrule
$v=w=0$ & Always\\
$v>0$, $w>0$ & $c=e$\\
$v<0$, $w<0$ & $d=f$\\
$v=0$, $w\ne0$ & $c\le\min(e,f)$\\
$w=0$, $v\ne0$ & $e\le\min(c,d)$\\
$vw<0$ & Never\\
\bottomrule
\end{tabular}
\caption{Commutation as equality of partial sequential maps, not concurrent token reservation.}\label{cdc:ct:tab:scalar}
\end{table}
\begin{proof}
There is no upper bound, so the criterion is
\begin{equation}\label{cdc:ct:eq:scalar-max}
 \max(c,e-v)=\max(e,c-w).
\end{equation}
If $v,w>0$ and $c>e$, the left side equals $c$, while both arguments on the right are less than $c$. The case $e>c$ is symmetric; $c=e$ works.

If $v=-r<0$ and $w=-s<0$, substitute $c=d+r$ and $e=f+s$. Both sides of~\eqref{cdc:ct:eq:scalar-max}, after subtracting $r+s$, become respectively $\max(d-s,f)$ and $\max(f-r,d)$. If $d>f$, the second is $d$ and the first is strictly smaller; the reverse inequality is symmetric. Equality of $d,f$ is sufficient.

For $v=0$,~\eqref{cdc:ct:eq:scalar-max} reads $\max(c,e)=\max(e,c-w)$. With $w>0$ it holds exactly when $c\le e$; with $w<0$ it holds exactly when $c\le e+w=f$. The other zero-displacement case is symmetric. If $v>0$ and $w=-s<0$, the right side is $\max(e,c+s)$, strictly greater than both $c$ and $e-v$, so equality fails. Finally, two zero displacements give the identical threshold $\max(c,e)$.

For vectors, all composition domains are nonempty upward boxes. Equality of their coordinate bounds is necessary and sufficient by Proposition~\ref{cdc:ct:prop:boxes}.
\end{proof}
This classification recognizes shared-resource commutation that a disjoint-coordinate test would miss. A read-only transition can commute with a consuming transition when its read threshold is no greater than the resource retained after the latter fires.

\section{The exact resource envelope of a multiset\srctag{01}}\label{cdc:ct:sec:envelope}
Write
\[
 v_{aj}^{+}=\max(v_{aj},0),\qquad v_{aj}^{-}=\max(-v_{aj},0).
\]
For multiplicities $n=(n_a)_{a\in A}\in\N^m$, let
\begin{equation}\label{cdc:ct:eq:loss-gain}
 S_j^- =\sum_a n_a v_{aj}^-,\qquad
 S_j^+ =\sum_a n_a v_{aj}^+,
 \qquad A_n=\{a:n_a>0\}.
\end{equation}
For each action-coordinate pair define the fixed constants
\begin{equation}\label{cdc:ct:eq:retained-cap}
 r_{aj}=\ell_{aj}-v_{aj}^-\ge0,
 \qquad s_{aj}=u_{aj}+v_{aj}^+\quad(u_{aj}<\infty).
\end{equation}
The extrema below are coordinatewise. The maximum of an empty set of nonnegative lower margins is zero; the minimum of an empty set of upper margins is infinity.

\begin{theorem}[Universal-serialization envelope]\label{cdc:ct:thm:envelope}
For arbitrary actions satisfying~\eqref{cdc:ct:eq:guard-valid}, the following are equivalent:
\begin{enumerate}[label=(\roman*)]
\item Every word with multiplicity vector $n$ is executable from $x\in\N^p$.
\item For every coordinate $j$,
\begin{align}
 x_j&\ge S_j^-+\max_{a\in A_n}r_{aj},\label{cdc:ct:eq:envelope-lower}\\
 x_j&\le \min_{\substack{a\in A_n\\u_{aj}<\infty}}s_{aj}-S_j^+.
 \label{cdc:ct:eq:envelope-upper}
\end{align}
\end{enumerate}
When these conditions hold, every such word ends at
\begin{equation}\label{cdc:ct:eq:envelope-output}
 y=x-S^-+S^+.
\end{equation}
No pairwise commutation assumption is needed for this theorem.
\end{theorem}
\begin{proof}
Fix a coordinate $j$ and distinguish one occurrence of an active action $a$. Just before that occurrence, let $\delta$ be the sum of displacements contributed by the preceding occurrences in this coordinate. Its guard is
\[
 x_j\ge\ell_{aj}-\delta,\qquad
 x_j\le u_{aj}-\delta\quad(u_{aj}<\infty).
\]
Among all permutations and all choices of this distinguished occurrence, the smallest possible preceding displacement is
\[
 \delta_{\min}=-S_j^-+v_{aj}^-.
\]
Indeed, put all other negative-displacement occurrences before it and all positive-displacement occurrences after it; zero-displacement occurrences are immaterial. No subset of the other occurrences has smaller total displacement. The strongest lower requirement due to $a$ is consequently
\[
 \ell_{aj}-\delta_{\min}
   =S_j^-+\ell_{aj}-v_{aj}^-
   =S_j^-+r_{aj}.
\]
Similarly the largest possible preceding displacement is
$\delta_{\max}=S_j^+-v_{aj}^+$, obtained by putting all other positive occurrences first. The strongest finite upper requirement due to $a$ is
\[
 u_{aj}-\delta_{\max}=s_{aj}-S_j^+.
\]
Taking the maximum or minimum over active $a$ proves that the displayed box is the intersection of every guard constraint at every position of every ordering.

For sufficiency, fix any ordering. Each of its prefix displacements lies between the two extremal values just used for its current action. The assumed inequalities therefore imply that every formal intermediate state satisfies its next guard. Induction on the word length then makes these formal states an actual execution in $\N^p$; condition~\eqref{cdc:ct:eq:guard-valid} ensures the next output is natural.

For necessity, if all orderings execute, in particular the extremal ordering for each action and coordinate executes. Thus each of its extremal guard inequalities holds. Different coordinates may use different extremal permutations: this is legitimate because the hypothesis quantifies over \emph{all} permutations. If $n=0$, there are no occurrences, the empty word is always executable, and both formulas give the identity domain and endpoint. Summing displacements gives~\eqref{cdc:ct:eq:envelope-output}.
\end{proof}

\begin{corollary}[Exact acceleration under commutation]\label{cdc:ct:cor:some-every}
If the actions in $A_n$ commute pairwise as partial maps, then \emph{some} ordering of the multiset is executable from $x$ if and only if the envelope conditions hold. Thus some ordering, every ordering, and the explicit resource box are equivalent.
\end{corollary}
\begin{proof}
Any two multiset words are connected by adjacent transpositions. Swapping identical letters changes nothing; swapping distinct active letters replaces two equal composite partial maps. Pre- and post-composition preserve that equality. All complete words therefore denote the same partial map, with the same domain. Apply Theorem~\ref{cdc:ct:thm:envelope}.
\end{proof}

\begin{corollary}[Petri resource formula]\label{cdc:ct:cor:petri-envelope}
For Petri transitions, every ordering of multiplicities $n$ executes from $x$ exactly when, for every $j$,
\begin{equation}\label{cdc:ct:eq:petri-envelope}
 x_j\ \ge\ \sum_a n_a(c_{aj}-d_{aj})_+
           +\max_{a\in A_n}\min(c_{aj},d_{aj}).
\end{equation}
Under pairwise commutation of the active transitions, the same condition characterizes existence of an ordering. The endpoint is $x+\sum_a n_a(d_a-c_a)$.
\end{corollary}
\begin{proof}
Here $v^-=(c-d)_+$ and
$c-(c-d)_+=\min(c,d)$; there are no upper guards. Substitute into Theorem~\ref{cdc:ct:thm:envelope} and Corollary~\ref{cdc:ct:cor:some-every}.
\end{proof}
The two terms in~\eqref{cdc:ct:eq:petri-envelope} have different meanings. The sum measures total irreversible loss over the block. The maximum measures the largest resource that must remain available while a transition fires. Shared read-only requirements are maximized, not summed.

\begin{example}[Why existential and universal feasibility differ]
Take $a:x\mapsto x+1$ with lower guard zero and $b:x\mapsto x-1$ with lower guard one. From zero, $ab$ executes and $ba$ does not. The envelope for one occurrence of each requires $x\ge1$. This is the correct universal condition and an incorrect existential condition without commutation. The two partial maps fail the box test.
\end{example}
\begin{example}[A repeated read versus repeated loss]
For a read $x\mapsto x$ guarded by $x\ge r$, any positive number of repetitions requires only $x\ge r$. For a unit loss guarded by $x\ge1$, $n$ repetitions require $x\ge n$. A test that reserves $r$ fresh tokens for each repeated read overestimates requirements; a test that only checks individual enabling underestimates repeated loss.
\end{example}


% ---- manuscript 06, lines 309-534 ----
\section{The exact resource-summary algebra\srctag{06}}
\label{cdc:cs:sec:resources}

\subsection{Every word is a principal partial translation}

For $U,V\in\N^p$, let $[U,V]$ denote the partial map
\begin{equation}
\label{cdc:cs:eq:partialmap}
U+r\longmapsto V+r\qquad(r\in\N^p).
\end{equation}
Its domain is $U+\N^p$. We compose such maps in execution order: the left
factor executes first.

\begin{lemma}[Faithful summaries]
\label[lemma]{cdc:cs:lem:faithful}
Each word $w$ has a unique pair $(U_w,V_w)\in\N^p\times\N^p$ such that
\begin{equation}
\label{cdc:cs:eq:resourcecontract}
m\xrightarrow{w}n
\quad\Longleftrightarrow\quad
\exists r\in\N^p\;(m=U_w+r\ \wedge\ n=V_w+r).
\end{equation}
For fixed $m,n$, the slack $r$, if it exists, is unique.
\end{lemma}

\begin{proof}
The empty word has summary $(0,0)$ and a letter $a$ has summary $(u_a,v_a)$.
Closure under concatenation is proved in \cref{cdc:cs:lem:resourceproduct}, so these
base cases give existence by induction on length. For uniqueness, the domain
of $[U,V]$ has the unique least element $U$, and the value at that element is
$V$. Thus equality of the partial maps forces equality of both coordinates.
Finally, $m=U_w+r$ determines $r$ uniquely.
\end{proof}

An independent closed description makes the enabling information explicit.
Set $\Delta_i=v_{a_i}-u_{a_i}\in\Z^p$. For each coordinate $j$,
\begin{equation}
\label{cdc:cs:eq:prefixformula}
(U_w)_j=
\max\left(0,\max_{1\leq i\leq L}
\left((u_{a_i})_j-\sum_{h<i}(\Delta_h)_j\right)\right),
\qquad V_w=U_w+\sum_{i=1}^L\Delta_i.
\end{equation}
The inner maximum is absent when $L=0$. Indeed, the inequalities saying that
all steps are enabled are exactly the inequalities whose largest right-hand
side occurs in \eqref{cdc:cs:eq:prefixformula}. Nonnegativity of $V_w$ follows from
executability at $U_w$. In particular, every finite word is executable from
\emph{some} sufficiently large marking, although not necessarily from the
marking of interest.

\subsection{Composition by cancellation}

\begin{lemma}[Resource product]
\label[lemma]{cdc:cs:lem:resourceproduct}
For $x=(u,v)$ and $y=(u',v')$, set $c=\min(v,u')$. Then
\begin{equation}
\label{cdc:cs:eq:product}
x\star y=(u+u'-c,\;v+v'-c)
\end{equation}
is the summary of executing $x$ followed by $y$.
\end{lemma}

\begin{proof}
It suffices to work in one coordinate. Write $v=c+s$ and $u'=c+t$, where
$s,t\geq0$ and at least one of $s,t$ is zero. Starting at $u+r$, the first
partial map produces $v+r$. The second is enabled precisely when
$c+s+r\geq c+t$. Since $st=0$, this is equivalent to $r\geq t$: if $t=0$ it
is automatic, and if $t>0$ then $s=0$. Writing $r=t+r'$ therefore gives the
composite action
\[
u+t+r'\longmapsto v'+s+r'.
\]
Here $u+t=u+u'-c$ and $v'+s=v+v'-c$, which proves the formula.
\end{proof}

\begin{corollary}
The operation $\star$ is associative with identity $(0,0)$. For $p$ counters
its algebra is the direct product of $p$ copies of the one-counter
partial-translation algebra.
\end{corollary}

\begin{proof}
Composition of partial maps is associative. The summary map is faithful by
\cref{cdc:cs:lem:faithful}, so the corresponding pairs are equal. The identity and
coordinatewise product assertions follow directly from the formula.
\end{proof}

In one coordinate this is a concrete realization of the bicyclic monoid:
first decrement $u$ times, then increment $v$ times. An increment followed by
a decrement cancels on all inputs, whereas a decrement followed by an
increment remains undefined at zero. We use the explicit pairs rather than
an abstract presentation because they display exactly which natural-number
witnesses a compiler needs. The broader connection between bicyclic storage,
partially blind counters, and graph monoids is classical \cite{zetzsche}.

\begin{lemma}[Unique minimum gadget]
\label[lemma]{cdc:cs:lem:min}
For fixed $x,y\in\N$, the equations
\begin{equation}
\label{cdc:cs:eq:mingadget}
c+s=x,\qquad c+t=y,\qquad st=0
\end{equation}
have the unique solution
$c=\min(x,y)$, $s=x-c$, $t=y-c$ in $\N^3$.
\end{lemma}

\begin{proof}
The displayed values satisfy the equations. Conversely, $c\leq x,y$. Since
$st=0$, either $s=0$ and $c=x\leq y$, or $t=0$ and $c=y\leq x$. These cases
agree when $x=y$ and determine all three variables.
\end{proof}

For composition we apply \eqref{cdc:cs:eq:mingadget} to $x=v$, $y=u'$ and set the
new coordinates to $u+t$ and $v'+s$. Every defining residual has degree at most
two, and every auxiliary coordinate is forced. This is stronger than merely
asserting that a minimum is Diophantine.

\subsection{Exact powers, including zero repetitions}

\begin{theorem}[Power formula]
\label{cdc:cs:thm:power}
Let $x=(u,v)$, and put $c=\min(u,v)$, $\alpha=u-c$, $\beta=v-c$
coordinatewise. Then
\begin{equation}
\label{cdc:cs:eq:power}
x^{\star k}=
\begin{cases}
(0,0),&k=0,\\
(c+k\alpha,\;c+k\beta),&k\geq1.
\end{cases}
\end{equation}
In every coordinate $\alpha\beta=0$.
\end{theorem}

\begin{proof}
In one coordinate, either $u\geq v$, so $\beta=0$, or $v\geq u$, so
$\alpha=0$. For $k=1$ the formula recovers $(u,v)$. Suppose $k\geq1$. If
$\beta=0$, the output after $k$ copies is $c$ and its cancellation against the
next input $c+\alpha$ is $c$. The composite is
$(c+(k+1)\alpha,c)$. If $\alpha=0$, the input stays $c$ and the output becomes
$c+(k+1)\beta$. This proves the positive-power formula by induction. The
zero-power formula is the identity of the monoid.
\end{proof}

\begin{example}[A catalyst cannot be retained at the zeroth power]
The command decrement-then-increment on one counter has summary $(1,1)$.
Every positive power still requires one token, but its zeroth power is the
empty word and is enabled at zero. Substituting $k=0$ into the positive-power
line of \eqref{cdc:cs:eq:power} would incorrectly retain the requirement $(1,1)$.
\end{example}

A two-way zero test can itself have unique polynomial witnesses. Let
$e,\ell\in\N$ satisfy
\begin{equation}
\label{cdc:cs:eq:twoway}
e(e-1)=0,\qquad k=e+\ell,\qquad (1-e)\ell=0.
\end{equation}
If $e=0$ these equations force $k=\ell=0$; if $e=1$ they force $k\geq1$ and
$\ell=k-1$. Thus $e=\mathbf1_{k>0}$ uniquely. The power coordinates are then
\begin{equation}
\label{cdc:cs:eq:powergadget}
U'=ec+k\alpha,\qquad V'=ec+k\beta.
\end{equation}
Together with \eqref{cdc:cs:eq:mingadget}, all residuals remain quadratic. Notice
that the child summary and its minimum witnesses are still constrained when
$k=0$. Disabling their equations would create free witnesses and destroy
single-foldness.

\subsection{A sharp global commutation test}

Disjoint resource footprints provide one sufficient commutation relation,
but are not necessary. Exact summaries give the largest relation appropriate
for exchanging adjacent labels on \emph{all} markings.

\begin{theorem}[Global adjacent-swap criterion; the counter form of \cref{cdc:thm:commute}]
\label{cdc:cs:thm:commutation}
For two transition labels $a,b$, the following conditions are equivalent:
\begin{enumerate}[label=\textnormal{(\roman*)},leftmargin=2em]
\item The partial actions of $ab$ and $ba$ are equal, including their domains.
\item Their minimum input vectors are equal.
\item Coordinatewise,
\begin{equation}
\label{cdc:cs:eq:commutation}
\min(u_b,v_a)=\min(u_a,v_b).
\end{equation}
\end{enumerate}
Consequently, the relation $I_*$ consisting of the distinct pairs satisfying
\eqref{cdc:cs:eq:commutation} is the largest symmetric, irreflexive label relation
for which every permitted adjacent swap preserves executability and result
at every starting marking.
\end{theorem}

\begin{proof}
By \eqref{cdc:cs:eq:product}, the minimum inputs of $ab$ and $ba$ are respectively
\[
u_a+u_b-\min(v_a,u_b),\qquad
u_a+u_b-\min(v_b,u_a).
\]
Equality is exactly \eqref{cdc:cs:eq:commutation}. Both words have the same total
increment $v_a-u_a+v_b-u_b$, so equality of minimum inputs also implies
equality of minimum outputs. Faithfulness proves the equivalence. The last
assertion follows by applying the equivalence to each pair separately.
\end{proof}

Two pure consumers of the same resource satisfy the criterion, as do two
pure producers. Pure increment and pure decrement of one counter do not.
The criterion is about equality of the \emph{two-step partial maps}; it is
not the usual stronger assertion that two individually enabled transitions
can always both execute. Two consumers may compete for a marking that enables
each separately but not their composition. This distinction is important
when relating the construction to other notions of concurrency.

Fix henceforth any symmetric, irreflexive $I\subseteq I_*$. Define
$\equiv_I$ as the congruence on $A^*$ generated by $ab\equiv_I ba$ for
$(a,b)\in I$. Every class is finite, since swaps preserve length and letter
multiplicities. Equality of partial actions is preserved under composition,
so
\begin{equation}
\label{cdc:cs:eq:traceinvariance}
w\equiv_I w'\quad\Longrightarrow\quad
\Summ(w)=\Summ(w').
\end{equation}
The converse need not hold: distinct traces can have the same partial action.
Our later counting theorem counts these trace classes, not all semantic
partial maps.


% ---- manuscript 07, lines 166-217 ----
\section{A first-failure product for resource legality\srctag{07}}
\subsection{Formal histories may continue after failure}
Fix $p$ resources, an initial vector $r_0\in\N^p$, and a candidate sequence of $T$ events. Event $i$ consumes $c_i\in\N^p$ and produces $a_i\in\N^p$. Define its \emph{formal} history over $\Z^p$ by
\[
 r_{i+1}=r_i-c_i+a_i\qquad(0\leq i<T),
\]
regardless of whether event $i$ is enabled. The sequence is legal when
\[
 r_{ij}\geq c_{ij}\quad\text{for every }i<T\text{ and }j<p.
\]
Let $\ff{x}{d}=\prod_{h=0}^{d-1}(x-h)$, with $\ff{x}{0}=1$.

\begin{theorem}[First-failure product]\label{cdc:qc:thm:first-failure}
The integer
\[
 G=\prod_{i=0}^{T-1}\prod_{j=0}^{p-1}\ff{r_{ij}}{c_{ij}}
\]
is positive for a legal event sequence and is exactly zero for an illegal one. Hence legality is equivalent to
\[
 \exists u\in\N\quad G=u+1,
\]
and the witness $u$ is unique whenever it exists.
\end{theorem}
\begin{proof}
If the sequence is legal, every factor in a nonempty falling factorial satisfies $r_{ij}-h\geq c_{ij}-(c_{ij}-1)=1$. Empty products equal one, so $G\geq1$.

Suppose instead that the sequence is illegal and choose its first disabled event $i$. Every earlier event is legal. Induction on the earlier events gives $r_i\in\N^p$: subtracting an enabled consumption and adding a nonnegative production cannot create a negative resource. There is a coordinate $j$ with
\[
 0\leq r_{ij}<c_{ij}.
\]
The falling factorial for this coordinate includes the factor with $h=r_{ij}$, which is zero. Thus the entire product is zero. Formal resource values at later events can be negative, but they cannot undo this zero factor.

The final assertion follows because the positive integers are exactly the numbers $u+1$ with $u\in\N$. Such a $u$ must equal $G-1$.
\end{proof}

\subsection{Why this is not a generic positivity trick}
Replacing an arbitrary collection of strict inequalities by positivity of their product is unsound: two negative factors have positive product. The theorem avoids this problem for a structural reason. The first disabled state is still nonnegative, and the falling factorial vanishes at every forbidden consumption level $0,\ldots,c_{ij}-1$. A negative factor arising later is always preceded by a zero factor.

For a queue deleting one symbol per event, put $p=1$ and $c_i=1$. If event $i$ appends $q_i\geq0$ symbols, then
\[
 \ell_i=n-i+\sum_{h<i}q_h,
 \qquad G=\prod_{i<T}\ell_i.
\]
Legality is exactly $G\geq1$. For deletion of $d$ symbols, use $\prod_{i<T}\ff{\ell_i}{d}$ instead.

\begin{remark}[Polynomial scope]
The theorem is stated for candidate \emph{integer} consumptions and productions. It becomes a fixed polynomial formula directly when the falling-factorial orders are fixed, as in the queue constructions below. An unrestricted variable consumption would make the displayed product length variable; that notation alone would not be an ordinary polynomial. Finite choices of consumption can be treated by separate finite-table constructions, but are not silently included in the degree claims here.
\end{remark}

\subsection{Separating content from resources}
The product sees only resource availability. It does not check which symbol is read, whether an instruction is selected correctly, or whether a produced word is the prescribed one. Its role is complementary: a content equation will identify the candidate symbols, and the resource product will guarantee that these symbols are available in the correct temporal order. This separation is the mechanism that permits us to eliminate intermediate queues.


\part{Canonical trace classes}\label{cdc:part:traces}

\paragraph{Source and scope.} \wtag{} Four manuscripts prove that the natural zeros of an explicit polynomial are in bijection with commutation classes of bounded executions. Manuscript~05 is the base: its Sections 4--8 (normal forms with linear Boolean memory, raw runs, the quartic and the convex quadratic certificates, costs and a ballot formula, and counting hardness) are printed first, together with its corollary for interval-guarded programs, moved here from its section on other substrates (\cref{cdc:cor:guarded}). Manuscript~04's Sections 3--5 (normal-form monitor, Petri quartic, counting and Parikh refinements), manuscript~06's Sections 4 and~7 (transfer signatures for trace normality and the bounded trace-class bijection) and manuscript~01's Sections 6--7 (Cartier--Foata layers and the causal-history quartic) follow.

\paragraph{The shared theorem and its second routes.} \wtag{} For exact length~$T$ the bijection is \cref{cdc:thm:quartic} (degree four) and \cref{cdc:thm:main} (degree two, convex, with a rational polytope as real zero set), from manuscript~05. It needs only a sound independence relation $I\subseteq I_{\max}$ and extends to interval guards and zero tests (\cref{cdc:cor:guarded}); manuscript~04's \cref{cdc:wf:thm:petri} needs independence by \emph{disjoint supports} \eqref{cdc:wf:eq:disjoint}, and its variable count $(p+r)(2T+1)$ is the count $(2T+1)(d+m)$ of \cref{cdc:thm:quartic} in manuscript~04's letters. Manuscript~06's \cref{cdc:cs:thm:tracebijection} (quartic, $q(2T+1)+p(7T-4)$ witnesses, resource summaries on a binary concatenation tree) is a third construction. Both are kept, with their proofs, as marked second routes. Manuscript~01's \cref{cdc:ct:thm:history} is a different theorem: it bounds the Cartier--Foata height~$H$ of a trace rather than its length, and its box-guarded actions carry upper guards.

\paragraph{Other duplicates.} \wtag{} Manuscript~04's monitor (\cref{cdc:wf:thm:monitor}) is the flag recurrence of \cref{cdc:prop:flags}, and its forbidden-factor lemma, like manuscript~06's \cref{cdc:cs:lem:forbidden}, is \cref{cdc:lem:forbidden}; all are kept with their proofs. The exponential automaton lower bound is proved twice with different families: \cref{cdc:thm:memory} uses $2k$ labels, and \cref{cdc:wf:thm:memory} uses $2k+1$ letters and realizes the relation by disjoint supports in a $1$-safe net. $\sharpP$-completeness of class counting is proved for the maximal relation with unit weights (\cref{cdc:thm:counting}) and for empty independence (\cref{cdc:wf:thm:counting}); these are different quotients and both are kept. The Parikh and additive-cost refinement is printed in manuscript~05's generating-function form, \cref{cdc:cor:weighted}; manuscripts~01 (\cref{cdc:ct:cor:counts}), 04 (\cref{cdc:wf:cor:parikh}) and~06 (the corollary after \cref{cdc:cs:thm:tracebijection}) state it for their own certificates.

\paragraph{Setting and hypotheses.} \wtag{} Manuscript~05: a Petri net with $d$ places and $m$ distinctly labelled transitions, any symmetric irreflexive $I\subseteq I_{\max}$, optional interval guards. Manuscript~04: $p$ places, $r$ transitions, $I$ restricted to disjoint supports. Manuscript~06: consume/produce transitions on $p$ counters, $q$ labels, $I$ a relation of globally commuting labels, no upper guards. Manuscript~01: box-guarded translations with lower and upper guards, $I$ a valid relation of commuting pairs, a fixed height bound $H$.

% ---- manuscript 05, lines 267-527 ----
\section{Canonical schedules with linear Boolean memory\srctag{05}}
\label{cdc:sec:normal}
Fix any symmetric irreflexive $I\subseteq I_{\max}$. On words over $\A$, let $\equiv_I$ be the equivalence relation generated by adjacent swaps $ab\leftrightarrow ba$ for $(a,b)\in I$. Equivalent words have equal length and equal multiplicities of all labels. Every class of words of length $T$ is finite and has a unique lexicographically least word, denoted its normal form.

\begin{lemma}[Execution invariance]\label{cdc:lem:invariance}
If $w\equiv_I w'$, then $w$ is enabled from $M$ if and only if $w'$ is, and their final markings agree whenever they are enabled.
\end{lemma}
\begin{proof}
Apply \cref{cdc:thm:commute} to each generating swap and compose the resulting partial-function equalities.
\end{proof}

\subsection{The forbidden-factor test}
The next characterization is a classical trace-language fact \cite{barthelemy}; we record a proof to make its use in the compiler explicit.

\begin{lemma}[Lexicographic forbidden factors]\label{cdc:lem:forbidden}
A word is not lexicographically least in its $I$-class if and only if it has a factor
\begin{equation}\label{cdc:eq:forbidden}
bua,\qquad a<b,\quad(a,b)\in I,\quad
(a,c)\in I\text{ for every letter }c\text{ of }u.
\end{equation}
\end{lemma}
\begin{proof}
Such a factor allows the final $a$ to move left across $u$ and $b$, making the word smaller. Conversely, order the occurrences of a word by the transitive closure of $i<j$ whenever their labels do not commute. Equivalent words are precisely the labelled linear extensions of this occurrence order: adjacent incomparable swaps preserve it, and moving the first occurrence of a target linear extension to the front, then inducting, connects any two extensions.

If the given extension is not lexicographically least, consider its first nongreedy choice. A later occurrence with smaller label $a$ was already available. It is incomparable with every occurrence skipped before reaching it, since a predecessor would contradict availability and a successor cannot occur earlier in the given extension. Hence all these skipped labels commute with $a$. The first skipped label is some $b>a$, giving \eqref{cdc:eq:forbidden}.
\end{proof}

\begin{example}[Adjacent inversion checks miss normal-form violations]\label{cdc:ex:adjacent}
Let $a<b<c$, with $aIb$ and $bIc$ but not $aIc$. The word $cab$ has no adjacent decreasing independent pair: $ca$ is dependent and $ab$ is increasing. Nevertheless
\[
cab\equiv_I cba\equiv_I bca,
\]
and $bca$ is lexicographically smaller. Normalization must remember more than the immediately preceding label.
\end{example}

\subsection{A forbidden-set recurrence}
For each candidate next label $a$, keep one flag $F_a$. After a prefix $w$, set $F_a(w)=1$ exactly when $w$ has a suffix $bu$ with
\[
b>a,\qquad aIb,\qquad\hbox{every letter of }u\hbox{ independent of }a.
\]
The empty suffix $u$ is permitted. All flags initially vanish. Appending $a$ is forbidden precisely when $F_a=1$.

\begin{proposition}[Exact flag update]\label{cdc:prop:flags}
On appending a label $b$, the flags change by
\begin{equation}\label{cdc:eq:flag-update}
F'_a=\ind_{aIb}\bigl(\ind_{a<b}\mathbin{\lor}F_a\bigr).
\end{equation}
A word is a normal form if and only if no step appends a label whose old flag is one.
\end{proposition}
\begin{proof}
If $a$ and $b$ are dependent, any suffix whose letters must commute with $a$ is broken, and $b$ itself cannot start a new forbidden suffix. If $aIb$ and $a<b$, the one-letter suffix $b$ starts such a suffix regardless of the old flag. If $aIb$ and $b<a$, the new letter can only prolong a previously existing suffix, not start one. These are exactly the three cases in \eqref{cdc:eq:flag-update}. The last assertion is \cref{cdc:lem:forbidden} applied at the right endpoint of every forbidden factor.
\end{proof}

For arithmetic use, at step $i$ let $e_{i,b}$ be one-hot selectors. Write
\begin{equation}\label{cdc:eq:CH}
C_{i,a}=\sum_{b:aIb}e_{i,b},\qquad
H_{i,a}=\sum_{b>a:aIb}e_{i,b},\qquad
W_{i,a}=C_{i,a}-H_{i,a}.
\end{equation}
On a natural one-hot assignment, $C,H,W$ are Boolean, $H\le C$, and $H$ and $W$ cannot both equal one. The recurrence is
\begin{equation}\label{cdc:eq:flag-poly}
F_{i+1,a}=H_{i,a}+W_{i,a}F_{i,a}.
\end{equation}
It uses only $m$ bits of state. All indices range over $i=0,\ldots,T-1$ unless indicated otherwise.

\subsection{Why an exponentially large automaton is avoidable but real}
\begin{theorem}[Exponential state lower bound]\label{cdc:thm:memory}
For every $k\ge1$ there is an independence relation on $2k$ ordered labels whose language of lexicographic normal forms requires at least $2^k$ nonrejecting states in any deterministic finite automaton. Thus any deterministic finite-state implementation requires at least $k$ bits of memory in the worst case. The $m$ flags of \cref{cdc:prop:flags} are optimal in order of growth, up to a constant factor.
\end{theorem}
\begin{proof}
Order the alphabet as
\[
a_1<b_1<a_2<b_2<\cdots<a_k<b_k.
\]
Let $a_iIb_j$ exactly when $j\le i$, impose symmetry, and allow no independence between two $a$'s or between two $b$'s. For $S\subseteq\{1,\ldots,k\}$, let $w_S$ consist of the letters $b_j$ with $j\in S$, in decreasing index order. Each $w_S$ is normal, since all $b$-letters are dependent.

After this prefix, $F_{a_i}=1$ exactly when $i\in S$. Any selected $b_j$ with $j>i$ is dependent with $a_i$ and occurs before all smaller indices; it cannot leave a final flag. The letter $b_i$, if present, sets the flag. Every later $b_j$ has $j<i$, commutes with $a_i$, and is smaller than $a_i$, so it preserves but does not set the flag. Therefore $w_Sa_i$ is normal exactly when $i\notin S$.

For distinct $S,S'$, choose $i$ in their symmetric difference. The continuation $a_i$ is accepted after exactly one of $w_S,w_{S'}$. A deterministic automaton must assign distinct states to these $2^k$ prefixes. They are all accepted prefixes, so these are nonrejecting states. A device with $b$ bits has at most $2^b$ states, giving $b\ge k$.
\end{proof}

By \cref{cdc:thm:graph}, these difficult relations occur as exact semantic commutation graphs of unit-weight Petri nets. This lower bound is not a lower bound on the number of unbounded arithmetic variables: one integer can store many bits. It specifically explains why expanding the recognizer into one variable per automaton state would be the wrong compiler architecture.

\section{A parsimonious quadratic encoding of raw executions\srctag{05}}
\label{cdc:sec:raw}
Fix a net, natural endpoints $M,N\in\N^d$, and an exact horizon $T\in\N$. Introduce the variables
\[
x_{i,p}\quad(0\le i\le T),\qquad
r_{i,p}\quad(0\le i<T),\qquad
e_{i,a}\quad(0\le i<T),
\]
where $1\le p\le d$ and $1\le a\le m$. They all range over $\N$. The equations are
\begin{align}
x_{0,p}&=M_p,&x_{T,p}&=N_p,\label{cdc:eq:endpoints}\\
\sum_{a=1}^m e_{i,a}&=1,\label{cdc:eq:onehot}\\
x_{i,p}&=\sum_a c_{a,p}e_{i,a}+r_{i,p},\label{cdc:eq:consume}\\
x_{i+1,p}&=r_{i,p}+\sum_a v_{a,p}e_{i,a}.\label{cdc:eq:produce}
\end{align}
Let $P_{\mathrm{raw}}$ be the sum of squares of the left-minus-right residuals of these equations.

\begin{theorem}[Exact raw-run fibres]\label{cdc:thm:raw}
The natural zeros of $P_{\mathrm{raw}}$ are in explicit bijection with the enabled length-$T$ transition words taking $M$ to $N$. The polynomial has degree at most two. The displayed, unoptimized presentation has
\begin{align*}
V_{\mathrm{raw}}&=(2T+1)d+Tm,\\
R_{\mathrm{raw}}&=2(T+1)d+T
\end{align*}
variables and squared residuals, respectively.
\end{theorem}
\begin{proof}
A natural solution of \eqref{cdc:eq:onehot} has exactly one coordinate equal to one and all others zero. Let $a_i$ be its selected label. Equations \eqref{cdc:eq:consume} and the nonnegativity of $r$ give $x_i\ge c_{a_i}$, so the step is enabled. They also force
\[
r_i=x_i-c_{a_i},\qquad x_{i+1}=x_i-c_{a_i}+v_{a_i}.
\]
Together with \eqref{cdc:eq:endpoints}, this decodes a valid execution.

Conversely, a valid word fixes the selectors. Starting from $x_0=M$, it uniquely determines each marking $x_i$ and then each remainder $r_i$. Enabledness guarantees that these are natural numbers. They satisfy every equation, and there are no other variables. This proves both existence and uniqueness of extension. All residuals are affine, so their squared sum has degree at most two. Counting the three variable arrays and the equations gives the formulas.
\end{proof}

\begin{lemma}[Uniform coordinate bounds]\label{cdc:lem:height}
Every natural solution of \eqref{cdc:eq:endpoints}--\eqref{cdc:eq:produce}, and every nonnegative real solution of these same equations, satisfies
\begin{equation}\label{cdc:eq:height}
0\le x_{i,p},r_{i,p}\le B_p:=M_p+T\max_a v_{a,p},
\qquad 0\le e_{i,a}\le1.
\end{equation}
\end{lemma}
\begin{proof}
The selector sum bounds each nonnegative selector by one. From \eqref{cdc:eq:consume}, $0\le r_{i,p}\le x_{i,p}$, and from \eqref{cdc:eq:produce},
\[
x_{i+1,p}\le x_{i,p}+\max_a v_{a,p}.
\]
Induction from $M_p$ proves the sharper bound with $i$ in place of $T$ for $x_i$, and hence the stated common bound.
\end{proof}

The raw-run encoding already avoids a frequent error: gating a separate transition equation by a selector while leaving inactive-branch auxiliaries free. There is only one enabling remainder per place and time, and it belongs to the selected step.

\section{Canonical traces: compact quartic and convex quadratic forms\srctag{05}}
\label{cdc:sec:main}
\subsection{A small quartic presentation}
Add the flag variables $f_{i,a}\in\N$ for $0\le i\le T$. Besides the raw-run equations, impose
\begin{align}
f_{0,a}&=0,\label{cdc:eq:fzero}\\
f_{i+1,a}&=H_{i,a}+W_{i,a}f_{i,a},\label{cdc:eq:frecur}\\
\sum_a e_{i,a}f_{i,a}&=0.\label{cdc:eq:reject}
\end{align}
The linear forms $H,W$ are those of \eqref{cdc:eq:CH}. Let $P_4$ be the sum of the squares of all displayed residuals.

\begin{theorem}[Compact canonical certificate]\label{cdc:thm:quartic}
Natural zeros of $P_4$ are in bijection with the $I$-commutation classes of enabled length-$T$ words from $M$ to $N$. It has degree at most four, with
\begin{align}
V_4&=(2T+1)(d+m),\label{cdc:eq:V4}\\
R_4&=2(T+1)d+(T+1)m+2T.\label{cdc:eq:R4}
\end{align}
At every zero all selectors and flags are Boolean; no separate equations $f(f-1)=0$ are required.
\end{theorem}
\begin{proof}
Decode a raw run using \cref{cdc:thm:raw}. Since each selector row is one-hot, the recurrence \eqref{cdc:eq:frecur}, starting from zero, forces $f_i$ to be exactly the Boolean flags of \cref{cdc:prop:flags}. In \eqref{cdc:eq:reject}, every summand is nonnegative and only the selected label can contribute. Thus the selected old flag vanishes at every step, so the word is normal.

Conversely, each normal enabled word uniquely determines its raw witness and all flag rows. It satisfies the rejection equations. By \cref{cdc:lem:invariance}, every enabled commutation class has its normal representative among the executions with the same endpoints. Uniqueness of this representative and of all its auxiliaries proves the bijection. The extra flag array contributes $(T+1)m$ variables; its initial and recurrence equations contribute $(T+1)m$ residuals, and rejection contributes $T$. Recurrence and rejection are at most quadratic, so the squared sum is at most quartic.
\end{proof}

\subsection{Eliminating products without introducing multiplicity}
We next replace each recurrence and rejection test by affine constraints. At each $(i,a)$ introduce five variables
\[
u_{i,a},\quad s^1_{i,a},\quad s^2_{i,a},\quad
s^3_{i,a},\quad s^4_{i,a}\in\N.
\]
Replace \eqref{cdc:eq:frecur}--\eqref{cdc:eq:reject} by
\begin{align}
f_{i,a}&=u_{i,a}+s^1_{i,a},\label{cdc:eq:q1}\\
W_{i,a}&=u_{i,a}+s^2_{i,a},\label{cdc:eq:q2}\\
1+u_{i,a}&=f_{i,a}+W_{i,a}+s^3_{i,a},\label{cdc:eq:q3}\\
f_{i+1,a}&=H_{i,a}+u_{i,a},\label{cdc:eq:q4}\\
1&=e_{i,a}+f_{i,a}+s^4_{i,a}.\label{cdc:eq:q5}
\end{align}
Keep the raw-run equations and \eqref{cdc:eq:fzero}. Let $L_1,\ldots,L_{R_2}$ be all their affine residuals, including \eqref{cdc:eq:q1}--\eqref{cdc:eq:q5}, and define
\begin{equation}\label{cdc:eq:P2}
P_2(z)=\sum_{j=1}^{R_2}L_j(z)^2.
\end{equation}
This is an explicit polynomial specified without an expensive expansion into monomials. The implementation also supports expansion for small examples.

\begin{theorem}[Canonical trace polytope theorem]\label{cdc:thm:main}
For every $\Net,M,N,T$ and symmetric irreflexive $I\subseteq I_{\max}$ as above, the polynomial \eqref{cdc:eq:P2} has all the following properties.
\begin{enumerate}[label=\textup{(\roman*)}]
\item It belongs to $\Z[z_1,\ldots,z_{V_2}]$, has degree at most two, and is convex on $\R^{V_2}$.
\item Its natural zero set is in explicitly computable bijection with the $I$-commutation classes of enabled length-$T$ executions from $M$ to $N$.
\item Its nonnegative real zero set is a bounded rational polytope. Its lattice points are exactly the natural zeros in \textup{(ii)}.
\item The unoptimized presentation has
\begin{align}
V_2&=(2T+1)d+(7T+1)m,\label{cdc:eq:V2}\\
R_2&=2(T+1)d+(5T+1)m+T.\label{cdc:eq:R2}
\end{align}
\item At every natural zero, each selector, flag, $u$, and $s^j$ is at most one, and all markings and enabling remainders satisfy \eqref{cdc:eq:height}. Thus the polynomial has finitely many natural zeros, with an explicit coordinate bound.
\end{enumerate}
\end{theorem}
\begin{proof}
We separate arithmetic soundness, uniqueness, and real geometry.

\emph{Natural soundness.} The raw constraints decode an enabled run. Each $W_{i,a}$ and $H_{i,a}$ is Boolean and they are disjoint. Start with the forced Boolean row $f_0=0$. If $f_i$ is Boolean, \cref{cdc:lem:and} applied to \eqref{cdc:eq:q1}--\eqref{cdc:eq:q3} forces
\[
u_{i,a}=f_{i,a}W_{i,a}.
\]
Then \eqref{cdc:eq:q4} gives the correct next flag. It is Boolean because $u\le W$ and $H+W=C\le1$. Induction proves the Boolean flag interpretation at every layer. Equation \eqref{cdc:eq:q5} is exactly $e_{i,a}+f_{i,a}\le1$. It rules out selecting a forbidden label. The decoded run is consequently normal.

\emph{Existence and uniqueness.} Starting from a canonical enabled word, set the raw variables and flags as before. Each $u$ and each slack must be
\begin{align*}
u&=fW,&s^1&=f-u,&s^2&=W-u,\\
s^3&=1+u-f-W,&s^4&=1-e-f.
\end{align*}
All these values are zero or one. Nonnegativity of the last expression is precisely the no-forbidden-selection property. They satisfy every affine equation. Conversely, the equations and the preceding induction force exactly these values. Thus the extension is unique, and the class bijection follows as in \cref{cdc:thm:quartic}.

\emph{Real geometry.} Write the affine system as $Az=b$ with integer $A,b$. Scalarization gives
\[
\{z\in\R_{\ge0}^{V_2}:P_2(z)=0\}
=\mathcal C:=\{z\in\R_{\ge0}^{V_2}:Az=b\}.
\]
This is a rational polyhedron, possibly empty. The selectors lie in simplices, so $0\le H,W$ and $H+W=C\le1$. Inductively assume $0\le f_i\le1$. Equations \eqref{cdc:eq:q1}--\eqref{cdc:eq:q3} give $0\le u\le f_i,W$. Hence
\[
0\le f_{i+1}=H+u\le H+W=C\le1.
\]
The first two slacks are at most one. The third is $1+u-f_i-W\le1$ because $u\le f_i$ and $W\ge0$; the fourth is $1-e-f_i\le1$. All are nonnegative by the variable domains. The raw variables obey \cref{cdc:lem:height}. Every coordinate of $\mathcal C$ is therefore bounded. A bounded rational polyhedron is a rational polytope, with the empty case allowed.

Finally, $P_2=\|Az-b\|_2^2$ has Hessian $2A^{\mathsf T}A$, which is positive semidefinite. The counts are the raw counts, plus $(T+1)m$ flag variables and $5Tm$ local auxiliaries; the added residuals are $m+5Tm$. This gives \eqref{cdc:eq:V2} and \eqref{cdc:eq:R2} and completes all assertions.
\end{proof}

\begin{remark}[What is fixed and what is a parameter]
The net, the alphabet order, the horizon, and the relation $I$ are fixed when this polynomial is formed. The endpoints can alternatively remain free polynomial parameters; they enter only affinely, so the joint degree is still at most two. Treating arbitrary arc weights as additional variables would create products with selectors and would change this degree accounting. Treating $T$ as a variable in one fixed polynomial is a different problem altogether.
\end{remark}

The number of variables and equations is linear in $T(d+m)$. In a sparse residual representation, the coefficient occurrences require at most order
$T(dm+|I|+d+m)$ space, up to the endpoint terms and bit lengths. This is not a claim that a fully expanded square has the same sparsity: squaring a dense affine residual introduces cross terms. Nor are the displayed constants claimed optimal; obvious deterministic or redundant arrays can be eliminated in particular instances.

\subsection{The real relaxation is not automatically integral}
\begin{example}[A nonempty polytope with no lattice points]\label{cdc:ex:gap}
Take two places, two transitions, and
\[
c_a=(2,0),\quad c_b=(0,2),\quad v_a=v_b=(0,0),
\quad M=(1,1),\quad N=(0,0),\quad T=1.
\]
The two pure consumers commute by \cref{cdc:thm:commute}. Neither is enabled at $M$, so there are no integer solutions. Nevertheless set
\[
e_{0,a}=e_{0,b}=\tfrac12,\quad x_0=(1,1),
\quad x_1=r_0=(0,0).
\]
These satisfy the raw affine constraints. In the order $a<b$, take $f_0=0$, $u=0$, and $f_{1,a}=1/2$, $f_{1,b}=0$. For each label set
\[
s^1=0,\qquad s^2=W,\qquad s^3=1-W,\qquad s^4=1-e.
\]
All are nonnegative and satisfy \eqref{cdc:eq:q1}--\eqref{cdc:eq:q5}. Thus $\mathcal C$ is nonempty but $\mathcal C\cap\Z^{V_2}$ is empty.
\end{example}

Convexity of $P_2$ is not an algorithm for its natural zeros. Linear programming can find the fractional point in this example and still decide nothing affirmative about a genuine execution.

\subsection{A concrete scheduling example}
Let $a$ produce one token, let $b$ be a no-op, and let $c$ consume one token, on one place. In the order $a<b<c$, the maximal commutation relation has precisely the pairs $aIb$ and $bIc$ and their reverses. The word $cab$ is enabled from marking one but not canonical, as in \cref{cdc:ex:adjacent}. Its normal representative is $bca$; its marking trace is
\[
1\ \xrightarrow{b}\ 1\ \xrightarrow{c}\ 0\ \xrightarrow{a}\ 1.
\]
After the first $b$, the flag for $a$ is one. The subsequent $c$, which is dependent with $a$, clears it before the final $a$ is selected. The construction records that flag history uniquely; it does not merely impose a global sorting of all labels.

For this net, horizon six and endpoints $M=N=2$, there are 140 enabled words and 28 commutation classes. The quadratic certificate has 142 variables and 113 residuals; the quartic one has 52 variables and 47 residuals. These counts are reproduced by the included program, not inferred from a floating-point solver.


\subsection{Canonical quadratic certificates for guarded programs}\label{cdc:sec:guarded}
\srcnote{\wtag{} This corollary and the paragraph after it are moved from manuscript~05's section on other substrates (Part~VIII), where they follow \cref{cdc:thm:interval} (printed in Part~I, \cref{cdc:sec:interval}); the bound $B_p=M_p+T\max_a v_{a,p}$ is that of \eqref{cdc:eq:bigM} there.}

% ---- manuscript 05, lines 694-714 ----
\begin{corollary}[Canonical quadratic certificates for guarded programs]\label{cdc:cor:guarded}
Fix an interval-guarded translation system as in \cref{cdc:thm:interval}, numerical endpoints, and horizon $T$. Let $g$ be the number of finite upper-guard entries among all transition coordinates. Its executions modulo the maximal commutation relation in \cref{cdc:thm:interval} have a parsimonious convex quadratic certificate with
\[
V_2+Tg\quad\hbox{variables},\qquad R_2+Tg\quad\hbox{residuals},
\]
where $V_2,R_2$ are given by \eqref{cdc:eq:V2}--\eqref{cdc:eq:R2}. Its nonnegative real zero set is bounded.
\end{corollary}
\begin{proof}
Use the underlying Petri updates $c_a=l_a$ and
$v_a=l_a+\delta_a$, which are nonnegative by assumption. Add, for every finite $u_{a,p}$ and each time step, a unique natural slack $g_{i,a,p}$ with
\begin{equation}\label{cdc:eq:upper-guard}
x_{i,p}+g_{i,a,p}
=u_{a,p}e_{i,a}+B_p(1-e_{i,a}),
\end{equation}
where $B_p=M_p+T\max_a v_{a,p}$. At a natural one-hot selector, this enforces the selected upper guard; each inactive guard enforces only the already established bound $x_{i,p}\le B_p$. Hence each enabled guarded word has exactly one raw extension, including all upper-guard slacks, and no other words have extensions.

Use the sound relation from \cref{cdc:thm:interval} in the same flag constraints. The proof of the class bijection needs only partial-map commutation, now supplied by that theorem, so it applies without change. In the real relaxation the new right-hand side is between $B_p$ and $u_{a,p}$, and the added slack is bounded by their maximum. The old coordinate bounds remain valid. There is one new variable and one affine residual per finite guard and time step.
\end{proof}

The distinction between the two maximal relations matters. A pair can commute vacuously because both guarded orders have empty domains, even if the underlying unguarded Petri transitions do not commute. It is the \emph{guarded} execution set that is then invariant. The reference compiler implements these optional upper guards and recomputes their exact commutation relation; it does not silently reuse the unguarded relation.


% ---- manuscript 05, lines 528-629 ----
\section{Costs, exact counts, and a closed-form example\srctag{05}}
\label{cdc:sec:cost}
\subsection{Which observations survive the quotient?}
For label costs $\kappa_a\in\N^q$, define
\[
K(w)=\sum_{i=0}^{T-1}\kappa_{a_i}.
\]
This quantity is constant on commutation classes because adjacent swaps preserve the multiplicity of each label. It can be retained by adding the affine equations
\begin{equation}\label{cdc:eq:cost}
y_j=\sum_{i,a}\kappa_{a,j}e_{i,a}\qquad(1\le j\le q).
\end{equation}
The new output variables $y_j$ are unique. Prescribing a target value of $y$, or imposing $y_j+s_j=B_j$ for a budget, retains unique extension. These additions do not increase the polynomial degree.

\begin{corollary}[Generating functions are preserved]\label{cdc:cor:weighted}
Let $\mathcal R_T(M,N)/I$ be the finite set of commutation classes represented by $P_2$. Then, with commuting indeterminates $t_1,\ldots,t_q$,
\[
\sum_{[w]\in\mathcal R_T(M,N)/I}t^{K(w)}
=\sum_{z\in\Sol(P_2)}t^{\sum_{i,a}\kappa_a e_{i,a}(z)}.
\]
Here $t^v=\prod_jt_j^{v_j}$. In particular, the number of natural zeros is the number of classes, and taking one cost coordinate per label retains the entire Parikh-vector generating polynomial.
\end{corollary}
\begin{proof}
Apply the bijection of \cref{cdc:thm:main}; each matched pair has the same displayed cost. Both sums are finite by the explicit bounds.
\end{proof}

Neither state-dependent costs nor maximum intermediate resource usage are automatically preserved. For the one-place producer $a$ and no-op $b$, the orders $ab$ and $ba$ commute, yet a cost charged to $b$ equal to the current marking has different values. More generally, any observation depending on the prefix state must be included in the partial-map semantics before commutation is certified. A label-additive cost is a sufficient condition, not a universal description of observable equivalence.

A stochastic interpretation also requires care. Counting one representative per class is not the same as summing the probability of all schedules in that class. The latter requires a separate class-weight calculation, potentially including the number and probabilities of its linear extensions.

\subsection{A ballot formula for the producer--idle--consumer net}
The small example from \cref{cdc:sec:main} admits an exact family of counts. This gives an independent combinatorial check on both the semantics and the enumerator.

\begin{theorem}[Class and schedule counts for one resource]\label{cdc:thm:ballot}
For the one-place net with producer $a$, idle transition $b$, and consumer $c$, fix $M,N,T\in\N$. Let
\[
\mathcal L=\{\ell:0\le\ell\le T,\quad
r_\ell=(\ell+N-M)/2\in\{0,1,\ldots,\ell\}\}.
\]
Use $\binom nk=0$ when $k$ is an integer outside $0\le k\le n$. Set
\[
A_\ell=\binom{\ell}{r_\ell}
       -\binom{\ell}{r_\ell+M+1}.
\]
Then the number of maximal-commutation classes and the number of raw schedules are respectively
\begin{equation}\label{cdc:eq:ballot}
\#\bigl(\mathcal R_T(M,N)/I_{\max}\bigr)=\sum_{\ell\in\mathcal L}A_\ell,
\qquad
\#\mathcal R_T(M,N)=\sum_{\ell\in\mathcal L}\binom T\ell A_\ell.
\end{equation}
The class Parikh generating polynomial is
\[
\sum_{\ell\in\mathcal L}
 A_\ell U^{r_\ell}V^{\ell-r_\ell}B^{T-\ell},
\]
where $U,V,B$ mark producer, consumer, and idle occurrences.
\end{theorem}
\begin{proof}
Deleting all idle letters leaves a word over $a,c$, and adjacent allowed swaps change neither this reduced word nor the number of idle letters. Conversely, all placements of a fixed number of idle letters around the same reduced word are equivalent, since $b$ commutes with both other labels. Thus a class is exactly an enabled reduced word of some length $\ell$, together with $T-\ell$ idle letters.

A reduced word ending at $N$ has $r_\ell$ up-steps and $\ell-r_\ell$ down-steps. Of the $\binom\ell{r_\ell}$ unrestricted words, the ones that go below zero are counted by $\binom\ell{r_\ell+M+1}$. To see this, reflect the part of the path from its starting height $M$ through its first visit to $-1$. This maps it bijectively to an unrestricted path from $-M-2$ to $N$ of length $\ell$. Such a path has $r_\ell+M+1$ up-steps; reflecting through its first crossing of $-1$ reverses the construction. Hence $A_\ell$ is exactly the enabled reduced-word count.

There are $\binom T\ell$ choices of positions for a fixed reduced word in a raw schedule. Idle steps cannot alter enabledness. Summing proves both counts and recording the three multiplicities proves the generating polynomial.
\end{proof}

For $M=N=2,T=6$, the nonzero terms $A_\ell$ for $\ell=0,2,4,6$ are $1,2,6,19$. Thus the class count is $28$, whereas the raw count is
\[
1+15\cdot2+15\cdot6+19=140.
\]
The formula also explains why a canonical class count cannot generally be obtained by dividing the raw count by one uniform symmetry factor: the factors $\binom T\ell$ differ among classes.

\section{Counting hardness survives maximal commutation\srctag{05}}
\label{cdc:sec:counting}
We now let the net and endpoints vary as finite input data. Arc weights and markings are encoded in binary; the exact horizon $T$ is encoded in unary. The latter convention is essential for a polynomial-size unrolling. The complexity class $\sharpP$ counts accepting branches of polynomial-time nondeterministic computations; Valiant's work is a foundational source for this counting framework \cite{valiant}.

\begin{theorem}[Parsimonious counting completeness]\label{cdc:thm:counting}
Counting maximal-static-commutation classes of enabled length-$T$ Petri-net executions from $M$ to $N$ is $\sharpP$-complete under parsimonious polynomial-time reductions. Hardness holds even for nets with only unit arc weights that are $1$-safe from the specified initial marking. Consequently, counting natural zeros of the quadratic certificates produced by \cref{cdc:thm:main} is $\sharpP$-complete on this promised family of certificates.
\end{theorem}
\begin{proof}
\emph{Membership.} Compute $I_{\max}$ by \cref{cdc:thm:commute}. Guess a word of length $T$ using a fixed-width binary encoding of each label, rejecting unused binary codes. Every word has exactly one guessing branch. Simulate it, check the endpoints, and run the flag recognizer. Exactly the canonical enabled words accept. Since $T$ is unary, the word length is polynomial in the input size. By \cref{cdc:lem:height}, intermediate marking bit lengths are polynomial as well. The arithmetic and flag checks run in polynomial time. Hence each commutation class contributes exactly one accepting branch.

\emph{Hardness construction.} Start with a Boolean circuit on $n$ input wires, with $g$ gates of fan-in at most two and designated output wire $o$. Gates are topologically ordered. Write $w=n+g$. We construct a net that guesses the input, evaluates every gate deterministically, tests that the output is one, and finally erases the stored values. The erasure is important: all accepting inputs must lead to the same final marking.

For each wire $j$, create two value places $V_j^0,V_j^1$, initially empty. Create control places $Q_0,\ldots,Q_{2w+1}$ with a single token initially at $Q_0$. Every transition moves the control token from one stage to the next. The data operations at the stages are as follows.
\begin{enumerate}[label=\textup{\arabic*.},leftmargin=2em]
\item At input stage $j<n$, two distinct transitions write a token to $V_j^0$ or $V_j^1$.
\item At gate stage $j$, include one transition for each assignment to its distinct predecessor wires. There are at most four. The transition reads the corresponding predecessor value places catalytically: a read place has both a unit input and a unit output arc. It writes one token to the correct value place of the new output wire. For repeated predecessors, consolidate duplicate reads rather than consuming two tokens from one place.
\item At stage $w$, the only transition requires a token at $V_o^1$, retains it, and advances the control. No transition accepts an output zero.
\item At the next $w$ stages, two possible transitions for each wire erase its zero or one token. Exactly one is enabled. At the end all value places are empty.
\end{enumerate}
The common final marking has its only token at $Q_{2w+1}$, and the horizon is
\[
T=2w+1.
\]
Every reachable marking is $1$-safe: each wire is written once, read without change, and then erased once; only one control token exists. All arcs have weight zero or one. After the input choices, exactly one gate or erasure transition is enabled at each stage. Thus satisfying input assignments are in bijection with accepting runs of this common length and common endpoint. The construction is polynomial in the circuit size.

\emph{The quotient does not merge assignments.} Distinct input assignments select different labelled input transitions. Their Parikh vectors are different, and all commutations preserve Parikh vectors. Thus no two accepting runs belonging to different assignments can become equivalent, even when the relation is the maximal one computed from the constructed net. There is already only one run for each assignment. Soundness of the relation also ensures that any word equivalent to an accepting run is an accepting run. Consequently each such class has exactly one run.

For completeness, counting satisfying assignments to circuits is itself a parsimonious universal $\sharpP$ problem: write a polynomial-time accepting predicate as a Boolean circuit on a fixed-length certificate, using deterministic padding to make the certificate representation unique. Its satisfying input strings are exactly the accepting certificates. Applying the construction above proves hardness. Finally, \cref{cdc:thm:main} preserves the count when passing to the quadratic certificate.
\end{proof}

This theorem pinpoints the limitation of the geometric representation. A compact convex quadratic polynomial can encode a hard \emph{integer counting} problem because its real zero polytope need not be integral. It does not imply hardness for solving the corresponding unconstrained real least-squares problem. The theorem also does not say that every family with a large commutation relation is hard; graph and net structure can still support useful parameterized algorithms.


% ---- manuscript 04, lines 149-353 ----
\section{A symbolic normal-form monitor for concurrent traces\srctag{04}}\label{cdc:wf:sec:monitor}
\subsection{Words, independence, and the forbidden-factor criterion}
Let $\A=\{0,\ldots,r-1\}$ be a totally ordered alphabet and let $I\subseteq\A^2$ be symmetric and irreflexive. Letters in $I$ are called independent. Write $u\equiv_I v$ if $u$ and $v$ are related by finitely many swaps of adjacent independent letters. This is a length-preserving congruence; each class is finite and therefore has a unique lexicographically least word, denoted $\NF_I(u)$.

\begin{lemma}[Classical normal-form criterion]\label{cdc:wf:lem:forbidden}
A word is lexicographically least in its $I$-class if and only if it has no factor $bva$ such that $a<b$ and $a$ is independent of $b$ and of every letter of $v$.
\end{lemma}
\begin{proof}
Such a factor is an immediate obstruction: swap its last $a$ successively left through $v$ and $b$. The resulting equivalent word is lexicographically smaller.

For the converse, distinguish repeated occurrences by their order among occurrences of the same letter. Equivalent words have the same occurrences, and the relative order of every pair of dependent occurrences is unchanged by an allowed swap. Suppose $u'$ is an equivalent, smaller word. Delete the common prefix, matching the distinguished occurrences there. Let $b$ be the next letter of $u$, and $a<b$ the next letter of $u'$. The occurrence of $a$ used by $u'$ occurs later in $u$, so $u$ has a factor $bva$ at that position. Every occurrence in $bv$ is still unmatched. If one were dependent on $a$, its order before $a$ would persist in $u'$, contradicting that $a$ is the next unmatched occurrence of $u'$. Thus $a$ is independent of every letter of $bv$, giving the forbidden factor.
\end{proof}

This proof uses only the invariance of dependent-occurrence order; it does not require a construction of a full automaton. In particular, checking only adjacent inversions would be inadequate. The middle word $v$ can contain letters that prevent a local sorting rule from exposing the obstruction immediately.

\subsection{Flags with a direct arithmetic update}
For a prefix $u$, define $f_a(u)\in\{0,1\}$ to record whether there exists an occurrence of some $b>a$ in $u$ such that $b$ is independent of $a$ and every letter after that occurrence in $u$ is independent of $a$. Initially all flags are zero. On appending $c$, the update is
\begin{equation}\label{cdc:wf:eq:flagcases}
 f_a(uc)=
 \begin{cases}
 0,&(c,a)\notin I,\\
 1,&(c,a)\in I\text{ and }c>a,\\
 f_a(u),&(c,a)\in I\text{ and }c<a.
 \end{cases}
\end{equation}
The equality case belongs to the first line because $I$ is irreflexive. An appended letter $c$ creates a forbidden factor precisely when $f_c(u)=1$.

Encode the letter at time $t$ by one-hot variables $x_{t,c}$. The case distinction becomes the quadratic equation
\begin{equation}\label{cdc:wf:eq:flagpoly}
 f_{t+1,a}=
 f_{t,a}\sum_{\substack{c<a\\(c,a)\in I}}x_{t,c}
 +\sum_{\substack{c>a\\(c,a)\in I}}x_{t,c}.
\end{equation}
Normality is imposed by the single equation per time step
\begin{equation}\label{cdc:wf:eq:normaltest}
 \sum_{a\in\A}x_{t,a}f_{t,a}=0.
\end{equation}

\begin{theorem}[Quadratic normal-form monitor; the flag recurrence of \cref{cdc:prop:flags}]\label{cdc:wf:thm:monitor}
For a fixed word of length $T$, equations $f_{0,a}=0$ and \eqref{cdc:wf:eq:flagpoly} force a unique flag tableau. Its entries lie in $\{0,1\}$ without separate Boolean equations. Adding \eqref{cdc:wf:eq:normaltest} at every time accepts exactly the lexicographically least words of their trace classes.
\end{theorem}
\begin{proof}
Because each row of selectors is one-hot, exactly one case of \eqref{cdc:wf:eq:flagcases} applies. Induction on $t$ proves both uniqueness and Booleanity, and proves that the computed flags have the defining suffix-obstruction meaning. Equation~\eqref{cdc:wf:eq:normaltest} is exactly $f_{t,c}=0$ for the selected $c$. Therefore all those equations hold exactly when no forbidden factor ever appears. Apply Lemma~\ref{cdc:wf:lem:forbidden}.
\end{proof}

The monitor uses $r$ bits of working memory, plus a rejection state in an ordinary automaton implementation. The displayed arithmetic description uses $O(T(r+|I|))$ terms, where $|I|$ counts ordered independent pairs. The number of \emph{equations} is $O(Tr)$, even for dense $I$; this must not be confused with an $O(Tr)$ bound on a fully expanded description for every graph.

\subsection{Why a symbolic representation matters}
\begin{theorem}[Exponential-state lower-bound family; a variant of \cref{cdc:thm:memory}]\label{cdc:wf:thm:memory}
For every $k\geq1$ there is an independence relation on $2k+1$ ordered letters such that every deterministic finite automaton recognizing its lexicographic trace normal forms has at least $2^k$ states. Thus, in the worst case, a deterministic streaming recognizer needs $\Omega(r)$ bits of memory. This remains true for independence relations realized by disjoint supports in an always-enabled, $1$-safe Petri net.
\end{theorem}
\begin{proof}
Order the alphabet as
\[
 r_1<\cdots<r_k<a_1<\cdots<a_k<b.
\]
Let the independent pairs, together with their reverses, be $(b,a_i)$ for all $i$ and $(r_i,a_j)$ for $i\ne j$. There are no other independent pairs.

For $S\subseteq\{1,\ldots,k\}$ let $p_S$ be $b$ followed by the $r_i$ with $i\in S$, in increasing order. Each $p_S$ is a normal form. After reading $b$, every flag $f_{a_i}$ is one. Reading $r_i$ clears $f_{a_i}$, while each other $r_j$ preserves it. Consequently
\begin{equation}\label{cdc:wf:eq:distinguish}
 p_Sa_i\text{ is a normal form}\quad\Longleftrightarrow\quad i\in S.
\end{equation}
If a deterministic automaton reached the same state after distinct prefixes $p_S$ and $p_{S'}$, choose $i$ in their symmetric difference. Appending $a_i$ would have to produce the same acceptance outcome, contradicting \eqref{cdc:wf:eq:distinguish}. Hence all $2^k$ prefixes require distinct states.

For the Petri-net realization, create one place for every unordered pair of distinct dependent letters. Put one token in every such place. Each transition has a consume-one/produce-one self-loop at precisely the places corresponding to pairs containing that transition. Every transition is always enabled and every marking remains unchanged and $1$-safe. Two distinct transition supports overlap exactly when their letters are dependent. Thus disjoint support realizes the prescribed $I$.
\end{proof}

The lower bound concerns deterministic recognition of normal forms. It is \emph{not} a lower bound on all Diophantine encodings, nondeterministic certificates, or the number of variables in a polynomial. Moreover, the net in the proof uses structural resource independence; transitions with the same overall effect are not automatically identified.

\section{Quartic certificates for Petri-net trace classes\srctag{04}}
\subsection{Operational semantics and trace closure}
A Petri net here has $p\geq1$ places and $r\geq1$ distinctly labeled transitions. Let $U_{ia},V_{ia}\in\N$ be the numbers of tokens consumed and produced by transition $a$ at place $i$. At marking $m\in\N^p$, transition $a$ is enabled if $m_i\geq U_{ia}$ for every $i$; its successor is
\[
 m'_i=m_i-U_{ia}+V_{ia}.
\]
Its support is $\supp(a)=\{i:U_{ia}+V_{ia}>0\}$. Fix a symmetric irreflexive $I$ satisfying
\begin{equation}\label{cdc:wf:eq:disjoint}
 (a,b)\in I\quad\Longrightarrow\quad
 \supp(a)\cap\supp(b)=\varnothing.
\end{equation}
One may use all disjoint-support pairs, or any symmetric subset.

\begin{lemma}[Enabled-swap property]\label{cdc:wf:lem:swap}
If a firing word contains adjacent independent transitions $ab$, swapping them gives another enabled firing word with the same final marking. Hence the set of firing words of length $T$ from $s$ to $z$ is a union of complete $I$-classes.
\end{lemma}
\begin{proof}
Suppose $a$ fires at $m$ and $b$ fires after $a$. Since $a$ changes no place in the support of $b$, all enabling inequalities for $b$ already hold at $m$. After $b$, the enabling inequalities for $a$ still hold. The effects add coordinatewise and have disjoint supports, so both orders reach $m-U_a+V_a-U_b+V_b$. The rest of the firing word is unchanged. Iterating this argument along any finite sequence of allowed swaps proves closure.
\end{proof}

This quotient is by \emph{the specified transition commutations}. It does not identify arbitrary executions with equal endpoints, equal Parikh vectors, indistinguishable token identities, or isomorphic occurrence processes. These are different semantic quotients.

\subsection{The complete residual system}
Fix a horizon $T\in\N$ and endpoint markings $s,z\in\N^p$. Introduce the following natural variables:
\begin{align*}
 m_{t,i}&\quad(0\leq t\leq T,\ 1\leq i\leq p),\\
 x_{t,a}&\quad(0\leq t<T,\ 0\leq a<r),\\
 e_{t,i}&\quad(0\leq t<T,\ 1\leq i\leq p),\\
 f_{t,a}&\quad(0\leq t\leq T,\ 0\leq a<r).
\end{align*}
The $m$ variables are markings, $x$ selects a transition, $e$ is the enabling slack, and $f$ is the monitor. The system consists of
\begin{align}
 m_{0,i}&=s_i,&m_{T,i}&=z_i,\label{cdc:wf:eq:petri-end}\\
 x_{t,a}(x_{t,a}-1)&=0,&\sum_a x_{t,a}&=1,\label{cdc:wf:eq:petri-onehot}\\
 m_{t,i}&=\sum_a U_{ia}x_{t,a}+e_{t,i},\label{cdc:wf:eq:petri-enable}\\
 m_{t+1,i}+\sum_a U_{ia}x_{t,a}
 &=m_{t,i}+\sum_a V_{ia}x_{t,a},\label{cdc:wf:eq:petri-update}\\
 f_{0,a}&=0,\label{cdc:wf:eq:petri-f0}\\
 f_{t+1,a}&=f_{t,a}\sum_{\substack{c<a\\(c,a)\in I}}x_{t,c}
 +\sum_{\substack{c>a\\(c,a)\in I}}x_{t,c},\label{cdc:wf:eq:petri-fupdate}\\
 \sum_a x_{t,a}f_{t,a}&=0.\label{cdc:wf:eq:petri-normal}
\end{align}
Unspecified indices range over all values for which their variables are defined. Every residual has degree at most two. Let $P_{\mathcal N,I,T,s,z}$ be the sum of their squares.

\begin{theorem}[Witness-faithful concurrent compilation; manuscript~04, second route to \cref{cdc:thm:quartic} under disjoint supports]\label{cdc:wf:thm:petri}
The natural zeros of $P_{\mathcal N,I,T,s,z}$ are in canonical bijection with the $I$-classes of enabled length-$T$ firing words from $s$ to $z$. The polynomial has degree at most four. The displayed construction uses exactly
\begin{align}
 V_{\mathrm{Petri}}&=(p+r)(2T+1),\label{cdc:wf:eq:petri-V}\\
 E_{\mathrm{Petri}}&=2p+r+T(2p+2r+2)\label{cdc:wf:eq:petri-E}
\end{align}
variables and quadratic-or-linear residuals, respectively. Every solution coordinate is at most
\begin{equation}\label{cdc:wf:eq:petri-height}
 B=\max\left\{1,\ \max_i s_i+T\max_{i,a}V_{ia}\right\}.
\end{equation}
If all markings reachable from $s$ are $1$-safe, every solution coordinate is Boolean.
\end{theorem}
\begin{proof}
\emph{From a class to a solution.} By Lemma~\ref{cdc:wf:lem:swap}, the unique least word $w=a_0\cdots a_{T-1}$ in the class is still an enabled run from $s$ to $z$. Let $m_t$ be its successive markings, set $x_{t,a}=1$ exactly when $a=a_t$, and put $e_{t,i}=m_{t,i}-U_{i,a_t}$. Enabling makes every $e_{t,i}$ natural. Use the unique monitor tableau of Theorem~\ref{cdc:wf:thm:monitor}. All equations hold, including the normality tests.

\emph{From a solution to a class.} Because the domain is $\N$, the equations in \eqref{cdc:wf:eq:petri-onehot} select exactly one transition $a_t$ at each step. Equation~\eqref{cdc:wf:eq:petri-enable} proves its enabling inequalities. Equation~\eqref{cdc:wf:eq:petri-update} is precisely its successor marking. The endpoint equations therefore produce a length-$T$ run from $s$ to $z$. The flag equations and Theorem~\ref{cdc:wf:thm:monitor} force this word to be least in its class.

\emph{Uniqueness.} For a fixed word, the initial marking and update equations force all markings; the enabling equations force all slacks; and the monitor equations force all flags. The selectors are the word itself. Thus there is exactly one solution per normal word, and hence per class.

\emph{Counts.} The four variable families contribute $(T+1)p$, $Tr$, $Tp$, and $(T+1)r$, giving \eqref{cdc:wf:eq:petri-V}. Endpoints contribute $2p$ residuals and initial flags contribute $r$. Each step contributes $r$ Boolean equations, one selector sum, $p$ enabling equations, $p$ updates, $r$ monitor updates, and one normality equation, giving \eqref{cdc:wf:eq:petri-E}.

\emph{Height.} At each step, any place increases by at most $\max_{i,a}V_{ia}$. Hence every marking coordinate satisfies \eqref{cdc:wf:eq:petri-height}. Each enabling slack is between zero and the corresponding marking; all other coordinates are bits. If every reachable marking is $1$-safe, this also makes every slack a bit. The sum-of-squares conversion proves the degree statement.
\end{proof}

For $T=0$, the same formulas apply: there is exactly one solution if $s=z$ and none otherwise. Some equations then refer to the same marking variables. They are still counted as displayed residuals. Counts throughout this paper are construction counts, not minimality claims.

\subsection{Why token-balance equations are insufficient}
An incidence equation
\[
 z=s+(V-U)\nu
\]
records only a transition-count vector $\nu\in\N^r$. It does not certify enabling or an order in which transitions can fire. For example, take two places and transitions moving a token in opposite directions. From the zero marking, neither transition is enabled, but $\nu=(1,1)$ satisfies the zero-to-zero incidence equation. Equations~\eqref{cdc:wf:eq:petri-enable} rule out this spurious object. The normal-form monitor then removes a different kind of redundancy: multiple valid orders of independent transitions.

\begin{example}[A fork and a join]\label{cdc:wf:ex:fork}
Use five places and three transitions $a<b<c$:
\[
\begin{array}{c|cc}
 &\text{Consumed vector}&\text{Produced vector}\\\hline
 a&(1,0,0,0,0)&(0,1,0,0,0)\\
 b&(0,0,1,0,0)&(0,0,0,1,0)\\
 c&(0,1,0,1,0)&(0,0,0,0,1).
\end{array}
\]
From $s=(1,0,1,0,0)$ to $z=(0,0,0,0,1)$ in three steps, the enabled words are $abc$ and $bac$. Only $a$ and $b$ are independent, so there is one trace class, represented by $abc$. The certificate has $56$ variables and $67$ residuals, with a unique Boolean solution. The accompanying file \path{data/04-witness-faithful-petri_fork_join.json} contains the complete residuals and that solution.
\end{example}

\section{Counting, weights, and the limits of low degree\srctag{04}}
\subsection{Exact counting and Parikh refinements}
Write $\mathcal C_{T,s,z}$ for the class set in Theorem~\ref{cdc:wf:thm:petri}. Then
\begin{equation}\label{cdc:wf:eq:count}
 \#\mathcal C_{T,s,z}
 =\#\{\mathbf y\in\N^{V_{\mathrm{Petri}}}:
       P_{\mathcal N,I,T,s,z}(\mathbf y)=0\}.
\end{equation}
Both sides are finite, and the height bound supplies a finite search box. The count is at most $r^T$. This is a parsimonious representation: no correction factor for permutations, unused branches, or arithmetic coding choices is needed.

\begin{corollary}[Parikh and additive-cost refinement; compare \cref{cdc:cor:weighted}]\label{cdc:wf:cor:parikh}
Adjoin variables $n_a$ and equations $n_a=\sum_{t<T}x_{t,a}$. The extension remains witness-faithful, and $n_a$ is the number of occurrences of transition $a$ in the represented class. For fixed natural transition costs $c_a$, a further equation $C=\sum_a c_a n_a$ gives a unique cost coordinate. Thus coefficient counts in
\[
 G_{T,s,z}(\mathbf u)=\sum_{[w]\in\mathcal C_{T,s,z}}
                 \prod_a u_a^{\Par_a(w)}
\]
are exact solution counts of quartic certificates with the corresponding linear constraints.
\end{corollary}
\begin{proof}
The new variables are uniquely determined by the selectors. Independent swaps preserve each transition count and hence every additive cost. Imposing prescribed values filters precisely the desired classes. The new residuals are linear, so degree at most four is retained after summing squares.
\end{proof}

This does not imply that equal Parikh vectors determine equal trace classes. The polynomial $G$ can have coefficients larger than one for exactly that reason.

\subsection{A complexity consequence for succinct nets}
The following consequence explains why the small degree should not be interpreted as an efficient generic solution method. It is a standard counting-reduction phenomenon instantiated by our explicit compiler, rather than a new general complexity classification of traces.

\begin{theorem}[Counting hardness for the bounded-net family; empty independence, compare \cref{cdc:thm:counting}]\label{cdc:wf:thm:counting}
Counting the classes represented in Theorem~\ref{cdc:wf:thm:petri}, with the net and a unary horizon as input, is in $\#\mathrm P$. It is $\#\mathrm P$-hard even for $1$-safe nets in which all transitions share a self-loop place, so that disjoint-support independence is empty. Consequently the corresponding structured family of quartic zero-counting problems is $\#\mathrm P$-complete, even on this Boolean-witness subclass.
\end{theorem}
\begin{proof}
For membership, nondeterministically choose a length-$T$ word using a fixed binary encoding of each letter, rejecting unused encodings. There is exactly one choice string per word. Simulate the net, check the endpoint, and run the monitor. Marking bit lengths are bounded by the input bit length plus $O(\log(T+1))$ using \eqref{cdc:wf:eq:petri-height}. With $T$ unary, these operations take polynomial time. Exactly one word from each class accepts.

For hardness, start with a Boolean circuit over fan-in-two AND and unary NOT gates. Counting its satisfying assignments is the defining standard circuit form of $\#\mathrm P$. Let it have $n$ input bits and $g$ gates in topological order. Build a net with one control token moving through successive stages and a dual-rail pair of places for each input and gate value. At an input stage, two transitions choose the bit and put a token on its zero or one rail. At each gate stage, transitions for its possible input valuations read the relevant rails by self-loops, put one token on the fresh correct output rail, and advance control. If an input wire occurs twice, its rail is read only once in the enabling pattern; this avoids an unintended demand for two tokens. Exactly one gate transition is enabled for each assigned input valuation.

An additional stage advances only when the designated circuit output is true. Finally, erase the $n+g$ stored bits in a fixed order, with a zero and a one transition at each stage. Precisely one erasure transition is enabled. The target is the final control token with all value rails empty. Add a common consume-one/produce-one lock place to every transition. It makes disjoint-support independence empty without changing the executions.

All stages are visited once, all value rails receive at most one token, and the net is $1$-safe. Every satisfying assignment determines exactly one accepting word, including its deterministic gate evaluations and erasures. No other accepting words exist. Every such word has
\[
 T=(n+g)+1+(n+g)=2(n+g)+1
\]
steps and the same target marking. The construction is polynomial in circuit size, and because $I$ is empty, words are classes. Theorem~\ref{cdc:wf:thm:petri} then gives the same count as natural quartic zeros, all of whose coordinates are Boolean.
\end{proof}

The theorem concerns a succinct operational representation, not an explicitly enumerated reachability automaton. This matters when comparing it with algorithms whose input size is the number of automaton states. It also shows that exact uniqueness of the auxiliary variables removes encoding redundancy, not the intrinsic combinatorial difficulty of choosing a computation.


% ---- manuscript 06, lines 535-721 ----
\section{A compact signature for lexicographic trace normality\srctag{06}}
\label{cdc:cs:sec:normality}

Let $<$ be the fixed total order on $A$. A word is \emph{normal} when it is
the lexicographically least word in its $I$-trace class. Equal-length words
are being compared, so no convention about comparing a proper prefix with a
longer word affects the definition. We first recall the classical local
obstruction to normality and then derive a small block signature.

\subsection{The forbidden-factor criterion}

\begin{lemma}[Lexicographic obstruction]
\label[lemma]{cdc:cs:lem:forbidden}
A word $w$ is nonnormal if and only if it has a factor $bua$ such that
$a<b$ and every letter in $bu$ is independent of $a$.
\end{lemma}

\begin{proof}
If such a factor occurs, the last $a$ can be commuted left across all of $bu$.
This replaces $bua$ by $abu$, producing a lexicographically smaller equivalent
word.

For the converse, distinguish occurrences even when their labels agree. On
the occurrences of $w$, place an order relation from position $i$ to position
$j$ whenever $i<j$ and their labels are not independent, and take transitive
closure. This is a partial order. The words in the trace class are exactly
the labelings of its linear extensions: an adjacent independent swap
preserves all these precedence constraints, and any two linear extensions
can be connected by swapping adjacent incomparable occurrences. The latter
claim follows by moving the first occurrence of a target linear extension
leftward, then applying induction to the remaining occurrences. Every
occurrence passed in this move must be incomparable.

The smallest labeling of a linear extension is obtained greedily by choosing
the least label among the available minimal occurrences. Equal-labeled
occurrences are dependent and retain their order, so ties do not introduce
a different labeled word. If $w$ is not the smallest labeling, let its first
greedy failure be at a position labeled $b$. An available occurrence with a
smaller label $a$ appears later in the displayed word. Availability means
that none of the intervening remaining occurrences is a predecessor of that
$a$; in particular, each intervening label, including $b$, is independent of
$a$. They form a factor $bua$ as asserted.
\end{proof}

The normal form and the underlying sorting interpretation are classical
\cite{anisimov-knuth}. We include the proof because the exact obstruction,
not just regularity of the normal-form language, determines the compiler.

\subsection{One bit of memory per possible smaller letter}

For each $a\in A$, maintain a bit $F_a$. After reading a prefix, the meaning of
$F_a=1$ is that the prefix contains an occurrence of some $b>a$ such that
this occurrence and every subsequent letter are independent of $a$. Thus
reading $a$ next would complete a forbidden factor. Initially $F_a=0$.

When the next letter is $b$, reject if $F_b=1$. The memory update is
\begin{equation}
\label{cdc:cs:eq:dfaupdate}
F'_a=d_{ab}F_a+h_{ab}(1-F_a),\qquad
 d_{ab}=\mathbf1_{(a,b)\in I},\quad
 h_{ab}=\mathbf1_{(a,b)\in I\,\wedge\,b>a}.
\end{equation}
If $a$ is dependent on $b$, every old candidate is blocked and the new bit is
zero. If $b>a$ and $a$ is independent of $b$, the letter $b$ itself is a new
candidate and the bit becomes one. In the remaining independent case the bit
is unchanged. These alternatives prove the memory invariant by induction.
Together with \cref{cdc:cs:lem:forbidden}, they prove that this deterministic
recognizer accepts exactly the normal words.

It is convenient to continue applying the memory update even after a
rejection. Acceptance and memory evolution are separate: a block may have a
well-defined final memory even though some earlier position caused failure.

\begin{definition}[Block signature]
\label[definition]{cdc:cs:def:signature}
A signature is a tuple
\[
S=(C,R,\alpha,\beta),\qquad C\in\B,\quad
R,\alpha,\beta\in\B^q,\quad \alpha_a\beta_a=0.
\]
It describes a block $w$ when, for every incoming bit vector $F$,
\begin{align}
\label{cdc:cs:eq:signaturetransform}
f_w(F)_a&=\alpha_a F_a+\beta_a,\\
\label{cdc:cs:eq:signatureaccept}
w\text{ causes no rejection from }F
&\quad\Longleftrightarrow\quad
C=1\ \text{ and }\ F_aR_a=0\ \text{ for all }a.
\end{align}
Here $f_w$ is the memory transformation obtained by processing the whole block.
\end{definition}

The disjointness condition means that each coordinate transformation is
exactly one of the identity map, the constant-zero map, and the constant-one
map. The bit $C$ records acceptance from the all-zero incoming memory.
The vector $R$ records which incoming bits can trigger a rejection. When
$C=0$, the acceptance predicate is already false and different $R$ vectors
can describe the same predicate. We do not require a minimal semantic
signature. We use the particular signature determined by the following
recursive formulas; that choice is fully deterministic.

\begin{lemma}[Empty and one-letter signatures]
\label[lemma]{cdc:cs:lem:atomsignature}
The empty block has signature
\[
\mathbf{1}_{\mathcal S}=(1,\zeros,\ones,\zeros).
\]
For the one-letter block $b$, a signature is
\begin{equation}
\label{cdc:cs:eq:atomsignature}
C=1,\quad R_a=\mathbf1_{a=b},\quad
\alpha_a=\mathbf1_{(a,b)\in I\,\wedge\,b<a},\quad
\beta_a=\mathbf1_{(a,b)\in I\,\wedge\,b>a}.
\end{equation}
\end{lemma}

\begin{proof}
The empty word does nothing and never rejects. For one letter $b$, rejection
is precisely the incoming condition $F_b=1$. Formula
\eqref{cdc:cs:eq:dfaupdate} is constant one in the independent $b>a$ case, identity
in the independent $b<a$ case, and constant zero otherwise. Irreflexivity of
$I$ excludes a fourth case with $a=b$ independent of itself.
\end{proof}

\subsection{Composition without enumerating subsets}

For bits $x,y$, write $x\lor y=x+y-xy$. All vector operations below are
coordinatewise except for the displayed scalar product.

\begin{theorem}[Signature product]
\label{cdc:cs:thm:sigproduct}
Suppose $S=(C,R,\alpha,\beta)$ describes $x$ and
$T=(D,S_0,\gamma,\delta)$ describes $y$. Then the following signature
$S\odot T=(C',R',\alpha',\beta')$ describes $xy$:
\begin{align}
\label{cdc:cs:eq:sigproduct}
\alpha'_a&=\gamma_a\alpha_a,&
\beta'_a&=\gamma_a\beta_a+\delta_a,\\
R'_a&=R_a\lor\bigl((S_0)_a\alpha_a\bigr),&
C'&=CD\prod_{a\in A}\bigl(1-(S_0)_a\beta_a\bigr).
\notag
\end{align}
All resulting coordinates are bits and $\alpha'_a\beta'_a=0$.
\end{theorem}

\begin{proof}
Substitution in the final-memory formula gives
\[
f_y(f_x(F))_a
=\gamma_a(\alpha_aF_a+\beta_a)+\delta_a
=(\gamma_a\alpha_a)F_a+(\gamma_a\beta_a+\delta_a).
\]
The summands in the last constant term cannot both be one because
$\gamma_a\delta_a=0$. If $\gamma_a\alpha_a=1$, then
$\gamma_a=\alpha_a=1$, forcing $\beta_a=\delta_a=0$ and hence disjointness.

For acceptance, both blocks must accept. Their conditions are
\[
C=D=1,\qquad F_aR_a=0,\qquad
(\alpha_aF_a+\beta_a)(S_0)_a=0\quad(a\in A).
\]
All terms in the last product are nonnegative. It vanishes precisely when
$\beta_a(S_0)_a=0$ and $F_a\alpha_a(S_0)_a=0$. The conditions independent of
$F$ give $C'$, and the remaining conditions forbid an incoming one whenever
$R_a$ or $\alpha_a(S_0)_a$ is one. This is the stated $R'_a$.
\end{proof}

These formulas can also be viewed as an associative product on all such
signatures, not merely as a rule on signatures arising from words. To verify
associativity, the memory coefficients associate by composition of affine
maps. For three factors, the forbidden incoming bits are the coordinatewise
union
\[
R\lor(\alpha S_0)\lor(\alpha\gamma T_0),
\]
independent of bracketing. The conditions on constant terms say, in either
bracketing, that none of the constant-one contributions at a later rejection
test is present. Distributivity of Boolean conjunction and disjunction makes
the scalar expressions equal. The empty signature is a two-sided identity.
Thus recursive multiplication agrees with reading letters one at a time.

There are at most $2\cdot 6^q$ possible signatures: $2$ choices for $C$, $2^q$
for $R$, and $3^q$ for the disjoint pairs $(\alpha_a,\beta_a)$. A compiler need
not list them. Each product requires only $O(q)$ Boolean gates. This is the
useful distinction between a small \emph{description of a transformation}
and an explicit transition table on all $2^q$ incoming memory states.


% ---- manuscript 06, lines 1071-1200 ----
\section{A bijection between bounded trace classes and polynomial roots\srctag{06}}
\label{cdc:cs:sec:counting}

The compressed compiler fixes a schedule template and its parameters. For a
stronger statement about concurrent \emph{behaviors}, we now allow every word
of a fixed length $T$. The witness tuple will include the chosen word itself,
and normality will remove exactly the redundant commutation representatives.

\subsection{One-hot word selection and deterministic memory}

Fix $T\geq1$. For each position $0\leq i<T$ and label $a\in A$, introduce
$z_{i,a}\in\N$ with
\begin{equation}
\label{cdc:cs:eq:onehot}
z_{i,a}(z_{i,a}-1)=0,\qquad
\sum_{a\in A}z_{i,a}=1.
\end{equation}
These variables select exactly one word $w$ of length $T$. Its $i$th letter
summary has coordinates
\begin{equation}
\label{cdc:cs:eq:selectedsummary}
U_i=\sum_a u_a z_{i,a},\qquad V_i=\sum_a v_a z_{i,a}.
\end{equation}
Use any full binary concatenation tree on these $T$ leaves, with the resource
minimum gadget at each internal node. Impose the root endpoint equations
\eqref{cdc:cs:eq:endpoint}.

For canonicality, introduce memory variables $F_{a,i}$ for
$a\in A$ and $0\leq i\leq T$, beginning with $F_{a,0}=0$. Define the following
\emph{linear expressions}, without additional variables:
\[
D_{a,i}=\sum_{b:(a,b)\in I}z_{i,b},\qquad
H_{a,i}=\sum_{b:(a,b)\in I,\ b>a}z_{i,b}.
\]
Impose the quadratic residuals
\begin{align}
\label{cdc:cs:eq:wordmemory}
F_{a,i+1}&=D_{a,i}F_{a,i}+H_{a,i}(1-F_{a,i}),\\
\label{cdc:cs:eq:wordreject}
\sum_{a\in A} z_{i,a}F_{a,i}&=0\qquad(0\leq i<T).
\end{align}
For a selected word, the recurrence forces the unique memory evolution of
\eqref{cdc:cs:eq:dfaupdate}. Its values are bits by induction, so separate Boolean
constraints on these variables are unnecessary. Equation
\eqref{cdc:cs:eq:wordreject} is exactly the rejection test for the selected letter.
Even without using its one-hot interpretation, all its summands are
nonnegative and therefore cannot cancel.

Let $P_T(m,n,z,\eta)$ be the sum of the squares of all residuals in this
construction, where $z$ denotes the word selectors and $\eta$ all other
witnesses.

\begin{theorem}[Bounded trace-class root bijection; manuscript~06, second route to \cref{cdc:thm:main}]
\label{cdc:cs:thm:tracebijection}
For fixed $m,n\in\N^p$ and $T\geq1$, the natural solutions of $P_T=0$ are in
bijection with the $I$-trace classes of words $w$ satisfying
\[
|w|=T,\qquad m\xrightarrow{w}n.
\]
The polynomial has degree at most four and exactly
\begin{equation}
\label{cdc:cs:eq:tracecount}
q(2T+1)+p(7T-4)
\end{equation}
witness variables in the allocation described above. All witnesses are at
most $\max(1,WT,\norm m)$.
\end{theorem}

\begin{proof}
A root determines a unique word by \eqref{cdc:cs:eq:onehot}. The resource equations
force its exact summaries and endpoint slack, and the memory equations force
its exact recognition run. Therefore the selected word is executable with
the required endpoints and is normal. Conversely, any such normal word gives
one selector assignment and, by the uniqueness of the resource and memory
computations, exactly one full root.

Every executable trace class has a unique normal word. Adjacent swaps in $I$
preserve the entire partial action by \eqref{cdc:cs:eq:traceinvariance}; hence that
normal word is executable from $m$ to $n$, whether or not it was the original
representative used to specify the class. Thus normal executable words are in
bijection with precisely the stated trace classes.

For the variable count, selectors contribute $qT$. A full binary tree on
$T$ leaves has $2T-1$ nodes, each with $2p$ summary coordinates, and $T-1$
internal nodes, each with $3p$ cancellation witnesses. The endpoint contributes
$p$ slacks and memory contributes $q(T+1)$. The total is
\[
qT+2p(2T-1)+3p(T-1)+p+q(T+1)
=q(2T+1)+p(7T-4).
\]
Every residual is quadratic or linear, proving the quartic bound. The resource
height proof with subwords of length at most $T$, together with Boolean
selectors and memory, gives the height assertion.
\end{proof}

For $T=0$, no word-choice or internal witnesses are needed. The polynomial
$\sum_{j=1}^p(m_j-n_j)^2$ has the single empty witness tuple when $m=n$ and
none otherwise, exactly matching the empty trace.

\begin{corollary}[Prescribing the Parikh vector; compare \cref{cdc:cor:weighted}]
Given $\nu\in\N^q$, add the linear equations
\[
\nu_a=\sum_{i=0}^{T-1}z_{i,a}\qquad(a\in A).
\]
The same construction, without additional witnesses, counts exactly the
executable trace classes with that Parikh vector.
\end{corollary}

A Parikh vector is generally coarser than a trace: dependent labels retain
order information. Conversely, two different traces may induce the same
resource summary. The three quotients---letter counts, trace classes, and
partial actions---should not be substituted for one another in a counting
claim.

\subsection{Why the length parameter is not free here}

The polynomials $P_T$ form a family. Their number of variables grows with $T$.
There is no assertion that one fixed polynomial receives $T$ as a free input
and has exactly one root per trace for all $T$. Passing from the family to a
single unbounded trace representation would require an additional coding
construction. Ordinary Diophantineness of an encoded trace does not guarantee
that the additional coding witnesses are unique or finite in number.

For each fixed $T$ the root set is finite and has at most $q^T$ members. The
height bound gives a finite search box, although enumerating that entire box
is not proposed as an efficient algorithm. The value of the theorem is its
exact combinatorial meaning, not a speedup over direct word enumeration.
The compressed compiler addresses a different efficiency question: avoiding
expansion of one highly structured schedule.


% ---- manuscript 01, lines 413-556 ----
\section{Causal traces and canonical layers\srctag{01}}\label{cdc:ct:sec:traces}
\subsection{The quotient to be represented}
Fix a symmetric irreflexive relation $I\subseteq A\times A$ consisting of commuting action pairs. Put $D=(A\times A)\setminus I$; in particular $(a,a)\in D$. Two words are trace-equivalent when one can be obtained from the other by swapping adjacent $I$-independent letters. Write $\Tr(A,I)$ for this quotient monoid.

By equality of partial maps, equivalent words have the same execution domain and endpoint. Accordingly, a trace is executable from $x$ when any, equivalently every, representative is executable. This equivalence is why we required equality of domains in the independence test.

A trace is not the same thing as an endpoint. Distinct labels with identical effects remain distinct labels, and dependent orders remain distinct traces even when their particular endpoints happen to coincide. The quotient removes precisely the swaps specified by $I$, no more.

\subsection{A self-contained normal-form lemma}
The underlying normal form is classical Cartier--Foata theory \cite{cartier-foata,krattenthaler}. We prove the exact version needed by the polynomial compiler to make the padding and predecessor conditions explicit.

For a word with individually distinguished occurrences, generate a partial order by putting occurrence $i$ below occurrence $j$ whenever $i<j$ in the word and their labels are $D$-dependent. Take the transitive closure. The resulting labelled dependence poset has an intrinsic order on repeated occurrences of each letter because self-dependence is included.

\begin{lemma}[Posets and swaps]\label{cdc:ct:lem:poset-swaps}
Two words represent the same trace exactly when their labelled dependence posets agree, with the $k$th occurrence of each label matched to the $k$th occurrence. The representatives of a trace are exactly the label words of the linear extensions of its dependence poset.
\end{lemma}
\begin{proof}
Swapping adjacent independent occurrences changes the order of no dependent pair and so preserves all generating comparisons and their transitive closure. Conversely, fix a dependence poset and two linear extensions. Locate the first occurrence in the second extension within the first extension. All occurrences before it in the first extension are incomparable with it: a predecessor would also have to precede it in the second extension, while a successor could not precede it in the first. Incomparable occurrences must have independent labels, since every dependent pair was made comparable. Move the selected occurrence left by adjacent allowed swaps. Remove this common first occurrence and repeat inductively. Thus any two linear extensions are connected by allowed swaps.

The original word is a linear extension. A linear extension orders every dependent pair in the prescribed direction, so reconstructs the same dependence poset. This proves both assertions.
\end{proof}

Give each occurrence the level
\begin{equation}\label{cdc:ct:eq:level}
 h(e)=1+\max\{h(f): f\text{ is an earlier occurrence with }(\lambda(f),\lambda(e))\in D\},
\end{equation}
with the maximum of an empty set equal to zero. Let $C_i$ be the set of labels of level $i$.

\begin{theorem}[Cartier--Foata layer characterization]\label{cdc:ct:thm:foata}
Every trace has a unique sequence of nonempty layers $C_1,\ldots,C_h$ satisfying:
\begin{enumerate}[label=(\alph*)]
\item each $C_i$ is a clique for $I$, so distinct members commute and no label is repeated;
\item every $a\in C_{i+1}$ has a predecessor $b\in C_i$ with $(b,a)\in D$.
\end{enumerate}
Concatenating the layers, in any order within each layer, represents the trace. The number $h$ is the length of a longest chain in the dependence poset.
\end{theorem}
\begin{proof}
The level of an occurrence is its longest-chain rank in the dependence poset. Indeed, every generating predecessor has a lower rank, and any longest chain can be refined to use generating comparisons between successive levels; the recurrence~\eqref{cdc:ct:eq:level} computes exactly these ranks in a topological order. Thus the levels are intrinsic to the trace by Lemma~\ref{cdc:ct:lem:poset-swaps}.

Two dependent occurrences cannot have the same level, because the earlier is a generating predecessor of the later. Hence each layer has distinct pairwise independent labels. If an occurrence has level $i+1$, the finite maximum in~\eqref{cdc:ct:eq:level} is attained by a dependent earlier occurrence of level $i$. This is condition (b). Sorting occurrences by level respects the poset, so concatenating the layers yields a linear extension and hence the same trace.

Conversely, start from a sequence satisfying (a) and (b), and concatenate its layers in arbitrary internal orders. Induct on $i$. All letters of the first layer have level one, since no two of them are dependent. A letter in layer $i+1$ is independent of its same-layer predecessors, has no predecessor from a later layer, and has a dependent predecessor in layer $i$. By induction the maximum predecessor level is exactly $i$, so its computed level is $i+1$. Thus the proposed layers are precisely the intrinsic levels of the resulting trace. This also proves uniqueness. Finally, the maximum rank equals the longest-chain length by the rank recurrence.
\end{proof}

\subsection{Unique trailing padding}
For a fixed bound $H$, append empty layers to a normal form until it has $H$ layers. The local conditions above characterize these padded forms if empty layers are allowed: if $C_i=\varnothing$, condition (b) forces $C_{i+1}=\varnothing$, and hence all later layers are empty. No choice of a padding length needs to be encoded. The all-empty sequence represents the empty trace.

This small observation is important for fibre preservation. A separate ``finished'' flag with unconstrained transitions, or arbitrary idle letters, can create several certificates for the same finite history. Trailing empty layers with the predecessor condition do not.

\section{An explicit quartic for causal histories\srctag{01}}\label{cdc:ct:sec:history}
\subsection{Parameters and witnesses}
Fix the action alphabet $A$, a valid independence relation $I$, and a natural causal-height bound $H$. Only $x,y\in\N^p$ are endpoint parameters. For $H\ge1$ introduce Boolean-intended variables $z_{ha}$ for $1\le h\le H$ and $a\in A$; they indicate membership in the $h$th layer.

Let $e$ be the number of unordered \emph{distinct} dependent pairs and set
\[
 D(a)=\{b:(b,a)\in D\},\qquad
 d=\sum_a |D(a)|=m+2e.
\]
Fix an ordering $b_{a1},\ldots,b_{a,d_a}$ of each $D(a)$, where $d_a=|D(a)|\ge1$.

The witnesses are the layer bits, internal states $M_h\in\N^p$ for $1\le h<H$, totals $L_{hj},G_{hj}$, lower slacks $\alpha_{haj}$, finite-upper slacks $\beta_{haj}$, and blocker products $U_{hak}$ for $1\le h<H$, $1\le k\le d_a$. Put $M_0=x$, $M_H=y$, and use $U_{ha0}=1$ as a constant, not a witness.

\subsection{The complete residual catalogue}
Use the following residuals, for all indices in their specified ranges:
\begin{align}
 &z_{ha}(z_{ha}-1),\label{cdc:ct:eq:hist-bool}\\
 &z_{ha}z_{hb}\quad\text{for each unordered distinct dependent pair }\{a,b\},
 \label{cdc:ct:eq:hist-clique}\\
 &L_{hj}-\sum_a z_{ha}v_{aj}^-,\qquad
 G_{hj}-\sum_a z_{ha}v_{aj}^+,\label{cdc:ct:eq:hist-total}\\
 &M_{hj}+L_{hj}-M_{h-1,j}-G_{hj},\label{cdc:ct:eq:hist-end}\\
 &z_{ha}(M_{h-1,j}-L_{hj}-r_{aj}-\alpha_{haj}),\qquad
 (1-z_{ha})\alpha_{haj},\label{cdc:ct:eq:hist-low}\\
 &z_{ha}(s_{aj}-G_{hj}-M_{h-1,j}-\beta_{haj}),\qquad
 (1-z_{ha})\beta_{haj}\quad(u_{aj}<\infty),\label{cdc:ct:eq:hist-up}\\
 &U_{hak}-U_{ha,k-1}(1-z_{h,b_{ak}})\quad(h<H),\label{cdc:ct:eq:blockers}\\
 &z_{h+1,a}U_{ha,d_a}\quad(h<H).\label{cdc:ct:eq:hist-foata}
\end{align}
Let $Q_{A,I,H}$ be the sum of their squares. All residuals have degree at most two, including the blocker recurrences; no high-degree conjunction is left unlifted.

\begin{theorem}[Parsimonious causal-history compiler; bounded causal height, not length]\label{cdc:ct:thm:history}
For fixed $A,I,H\ge1$, the nonnegative roots in all witness coordinates of
\[
 Q_{A,I,H}(x,y;w)=0
\]
are in an explicit bijection with the executable $I$-traces from $x$ to $y$ whose Cartier--Foata height is at most $H$. The polynomial has degree at most four, with
\begin{align}
 N_H&=H\bigl(m(1+p)+f+2p\bigr)+(H-1)(d+p),\label{cdc:ct:eq:history-N}\\
 E_H&=Hm+He+3Hp+2H(pm+f)+(H-1)(d+m)\label{cdc:ct:eq:history-E}
\end{align}
witnesses and squared residuals, respectively. In particular, every such trace gives exactly one complete natural solution, including all auxiliary coordinates.
\end{theorem}
\begin{proof}
\emph{Extraction.} A zero of $Q$ makes all residuals vanish. Equation~\eqref{cdc:ct:eq:hist-bool} makes each $z_{ha}$ Boolean. The corresponding subsets $C_h=\{a:z_{ha}=1\}$ are cliques by~\eqref{cdc:ct:eq:hist-clique}.

Induction on $k$ in~\eqref{cdc:ct:eq:blockers} gives
\begin{equation}\label{cdc:ct:eq:block-product}
 U_{hak}=\prod_{t=1}^k(1-z_{h,b_{at}})\in\{0,1\}.
\end{equation}
Therefore~\eqref{cdc:ct:eq:hist-foata} says that an action present in layer $h+1$ has at least one dependent action in layer $h$. Section~\ref{cdc:ct:sec:traces} implies that these are the uniquely padded layers of a trace of height at most $H$.

The totals~\eqref{cdc:ct:eq:hist-total} are the losses and gains for the layer multiplicities $z_{h\cdot}$. The guard slacks~\eqref{cdc:ct:eq:hist-low}--\eqref{cdc:ct:eq:hist-up} assert its envelope at $M_{h-1}$. Within a layer all actions are pairwise independent and hence commute, so Corollary~\ref{cdc:ct:cor:some-every} makes every internal ordering executable. Equation~\eqref{cdc:ct:eq:hist-end} sets the endpoint to $M_h$. Concatenating all layers yields an execution from $x$ to $y$. Its trace is independent of the chosen internal order.

\emph{Construction.} Take an executable trace of height at most $H$. Its unique padded normal form fixes $z$. Any representative has the same partial map, so the layer-concatenated representative is executable. Start at $M_0=x$ and set the boundary states by the cumulative displacements. Set $L,G$ by their aggregate formulas. Set each active slack to its exact nonnegative guard margin and every inactive slack to zero. Define $U$ by the product~\eqref{cdc:ct:eq:block-product}. The resource theorem and normal-form characterization verify every residual.

\emph{Uniqueness.} The trace fixes its padded normal form, hence all $z$. From $x$ and $z$, the displacement recurrence fixes all states; the aggregate equations fix all totals. The complementary active/inactive equations fix every slack, and the blocker recurrences fix every $U$. Thus there is exactly one full root over that trace. Extraction and construction are inverse because they use these same formulas.

\emph{Counts.} There are $Hm$ bits, $(H-1)p$ internal-state entries, $2Hp$ totals, $Hpm$ lower slacks, $Hf$ upper slacks, and $(H-1)d$ blockers. Summing gives~\eqref{cdc:ct:eq:history-N}. The residual families contribute, in order, $Hm$, $He$, $2Hp$, $Hp$, $2Hpm$, $2Hf$, $(H-1)d$, and $(H-1)m$, which gives~\eqref{cdc:ct:eq:history-E}. Squaring residuals of degree at most two proves the degree bound.
\end{proof}
For $H=0$ there are no witnesses and the separate polynomial
\[
 Q_{A,I,0}(x,y)=\sum_{j=1}^p(y_j-x_j)^2
\]
represents the single empty trace when $x=y$, and no trace otherwise.

\subsection{Counting, costs, and finite fibres}
\begin{corollary}[Exact trace counts; compare \cref{cdc:cor:weighted}]\label{cdc:ct:cor:counts}
For each fixed $x,y,H$, the number of natural roots of $Q_{A,I,H}$ is the number of executable trace classes between these endpoints with height at most $H$, and is at most $2^{Hm}$.
If prescribed letter counts $n_a$ are desired, adjoining the squared linear residuals
\[
 \sum_{h=1}^H z_{ha}-n_a
\]
preserves degree at most four and selects precisely those traces with these counts. No new witness is required.
\end{corollary}
\begin{proof}
Use the bijection in Theorem~\ref{cdc:ct:thm:history}. There are at most $2^{Hm}$ Boolean layer matrices, each with at most one completion. The additional residuals record exactly the number of occurrences of each action in that matrix.
\end{proof}
The same construction records exact total length or any fixed additive integer action cost by one further linear equation. It does not encode a scheduling-dependent cost that varies between equivalent words; such a cost is not a function on the chosen quotient and would require a different semantic contract.

\begin{proposition}[An explicit bound on all history witnesses]\label{cdc:ct:prop:witness-bound}
Suppose $W\ge1$ bounds every finite guard constant and every $|v_{aj}|$, and $X=\max_j x_j$. Every root of $Q_{A,I,H}$ has all witness coordinates at most
\begin{equation}\label{cdc:ct:eq:witness-bound}
 B=X+(H+1)mW+2W+1.
\end{equation}
\end{proposition}
\begin{proof}
A layer contains at most $m$ actions, so every executable boundary state is at most $X+HmW$ in each coordinate. Each loss or gain total is at most $mW$. An active lower slack is $M_{h-1,j}-L_{hj}-r_{aj}$, which is nonnegative and no greater than its state coordinate because $r_{aj}\ge0$. An active upper slack is at most $s_{aj}\le2W$. Inactive slacks are zero; all bits and blockers are zero or one. These bounds are each no greater than~\eqref{cdc:ct:eq:witness-bound}.
\end{proof}
Consequently the bounded-height construction permits complete finite enumeration in principle, with witness bit length logarithmic in $B$. This is a representation and size statement, not a claim that enumerating $B^{N_H}$ points is a practical solver.

\subsection{A parameter that must not be hidden}
The bound $H$ determines the number of variables and equations. Thus
$Q_{A,I,H}$ is an effectively generated \emph{family} of ordinary polynomials, not one fixed-arity polynomial with $H$ as an input variable. Taking the union over all $H$ yields an effective reachability relation, and MRDP can represent that union existentially. Nothing in that observation preserves the preceding root bijections or finite-foldness.


\part{Compressed repetition and accelerators}\label{cdc:part:accel}

\paragraph{Source and scope.} \wtag{} This Part collects the fixed-arity constructions whose witness count does not grow with repetition counts. Manuscript~01's Sections 4--5 give the single-fold quartic accelerator for a multiset of box-guarded actions (\cref{cdc:ct:thm:accelerator}), fixed block programs, exact powers of a fixed word (\cref{cdc:ct:prop:word-power}) and rank and progress criteria for endpoint-level single- and finite-foldness; its Section~8 works the examples (eleven witnesses for multiplicities $10^{100},10^{80},10^{90}$). Manuscript~06's Sections 5--6 prove that nested normality only sees repetition counts $0$, $1$ and at least~$2$ (\cref{cdc:cs:thm:squaresaturation,cdc:cs:thm:cap}) and build the single-fold quartic compiler for parametrized nested repetition expressions (\cref{cdc:cs:thm:compiler}); its Section~8 locates the expressiveness boundary (flat schemes are Presburger-definable, separated depth-two schemes represent arbitrary quadratic systems).

\paragraph{Relation between the two.} \wtag{} Neither construction contains the other. Manuscript~06's nested expressions generalize manuscript~01's single-level word power to arbitrarily nested repetitions and add the normality test; manuscript~01's actions carry upper guards, and its accelerator characterizes \emph{every} serialization of a multiset, and \emph{some} serialization under commutation. The single-word case appears in both forms, \cref{cdc:ct:prop:word-power} and \cref{cdc:cs:thm:power} (Part~I), and both are credited.

\paragraph{Setting and hypotheses.} \wtag{} Manuscript~01: a fixed finite alphabet of box-guarded translations on $\N^p$; multiplicities, endpoints and phase counts are parameters of the polynomial, not witnesses, unless a corollary says otherwise. Manuscript~06: consume/produce transitions on $p$ counters with lower guards only, $q$ labels, and a fixed repetition expression with $N$ syntax nodes whose repetition parameters are inputs.

% ---- manuscript 01, lines 281-412 ----
\section{A fixed-arity single-fold quartic accelerator\srctag{01}}\label{cdc:ct:sec:accelerator}
\subsection{The complete polynomial system}
Fix the action descriptions once and for all. Treat $x,y\in\N^p$ and $n\in\N^m$ as \emph{parameters}. Let
\[
 f=\#\{(a,j):u_{aj}<\infty\}.
\]
Introduce natural witnesses $z_a,k_a$, two totals $L_j,G_j$, lower slacks $\alpha_{aj}$ for all $a,j$, and upper slacks $\beta_{aj}$ exactly for the finite upper guards. Define the following residuals:
\begin{align}
 &z_a(z_a-1),\qquad n_a-z_a(1+k_a),\qquad (1-z_a)k_a;
 \label{cdc:ct:eq:acc-support}\\
 &L_j-\sum_a n_a v_{aj}^-,\qquad
 G_j-\sum_a n_a v_{aj}^+,\qquad
 y_j+L_j-x_j-G_j;\label{cdc:ct:eq:acc-total}\\
 &z_a(x_j-L_j-r_{aj}-\alpha_{aj}),\qquad
 (1-z_a)\alpha_{aj};\label{cdc:ct:eq:acc-lower}\\
 &z_a(s_{aj}-G_j-x_j-\beta_{aj}),\qquad
 (1-z_a)\beta_{aj}\quad(u_{aj}<\infty).
 \label{cdc:ct:eq:acc-upper}
\end{align}
Here and throughout the residuals are \emph{integer} expressions evaluated on natural variables. Subtraction is not truncated. Let $P_A$ be the sum of the squares of all listed residuals.

\begin{theorem}[Single-fold quartic accelerator]\label{cdc:ct:thm:accelerator}
For every fixed action alphabet $A$, $P_A$ has total degree at most four and
\begin{align}
 N_{\rm acc}&=2m+2p+pm+f \quad\text{witnesses},\label{cdc:ct:eq:acc-N}\\
 E_{\rm acc}&=3m+3p+2pm+2f \quad\text{squared residuals}.
 \label{cdc:ct:eq:acc-E}
\end{align}
For every fixed $(x,y,n)$, its natural zero set has exactly one point if every serialization of $n$ executes from $x$ to $y$, and is empty otherwise. If the active actions commute, ``every'' can be replaced by ``some.'' Neither arity nor degree depends on the magnitudes of the multiplicities.
\end{theorem}
\begin{proof}
A sum of integer squares vanishes exactly when every residual vanishes. The first residual in~\eqref{cdc:ct:eq:acc-support} forces $z_a\in\{0,1\}$. If $n_a=0$, the second forces $z_a=0$, and the third then forces $k_a=0$. If $n_a>0$, the second forces $z_a=1$ and $k_a=n_a-1$, after which the third is automatic. Thus the support encoding has one and only one solution for every natural $n_a$.

Equations~\eqref{cdc:ct:eq:acc-total} uniquely set $L=S^-$ and $G=S^+$. If $z_a=1$, the lower residual forces
$\alpha_{aj}=x_j-S_j^--r_{aj}$; the naturalness of this uniquely determined slack is exactly the relevant lower bound. If $z_a=0$, the inactive equation uniquely sets $\alpha_{aj}=0$, and no guard is imposed. The same argument for the finite upper guard gives
$\beta_{aj}=s_{aj}-S_j^+-x_j$ when active and zero otherwise.

The existence of all these natural slacks is therefore exactly the envelope box. The remaining endpoint residual is~\eqref{cdc:ct:eq:envelope-output}. Theorem~\ref{cdc:ct:thm:envelope} gives soundness and completeness; Corollary~\ref{cdc:ct:cor:some-every} gives the existential interpretation under commutation. Every witness was uniquely determined in the preceding argument, proving single-foldness.

There are $2m$ support witnesses, $2p$ totals, $pm$ lower slacks, and $f$ upper slacks. The three support residuals per action, three total/endpoint residuals per coordinate, and two residuals per slack give the stated counts. Every residual has degree at most two, including degrees in the parameter variables, so $P_A$ has degree at most four.
\end{proof}

\subsection{Why the inactive equations are indispensable}
The equation $z\alpha=0$ by itself is not a canonical conditional constraint: when $z=0$ it permits every natural $\alpha$. In our construction the active residual imposes the inequality and the complementary residual $(1-z)\alpha=0$ fixes unused storage to zero. A single omitted inactive equation can turn a singleton fibre into an infinite one.

Similarly, $n=z(1+k)$ does not determine $k$ when $n=z=0$. The third residual in~\eqref{cdc:ct:eq:acc-support} is not an optimization detail; it is part of the uniqueness proof. These are small examples of a general rule: a semantic proof that ignores inactive branches need not be a fibre-preserving proof.

\subsection{What is compressed and what is not}
The construction compresses arbitrarily many repetitions into fixed-arity arithmetic. It does not promise small \emph{numbers}: the witness $k_a=n_a-1$ has essentially the same bit length as $n_a$, as it should. The totals are computed by linear forms in the multiplicities. No powers, factorials, prime coding, or CRT witnesses occur.

The sparse residual presentation has $O(pm+f)$ terms for a fixed action description, counting the terms in all aggregate sums. Expanding the sum of squares can produce $O(pm^2)$ terms from those sums, before cancellation. Thus a short constraint list, a short arithmetic circuit, and a small expanded monomial list are different size measures. The attachment exposes both representations for the worked examples, rather than silently equating them.

\begin{remark}[Not a new universal variable bound]
The arity in~\eqref{cdc:ct:eq:acc-N} is a bound for the specific envelope relation of a fixed finite alphabet. It is not a universal Diophantine-equation bound, and it does not improve ProveIt's independently reported universal certificate operation count \cite{repo-h10}.
\end{remark}

\section{Fixed block programs with unbounded repetition\srctag{01}}\label{cdc:ct:sec:blocks}
A single commuting block is not the only unbounded relation obtained without sequence coding. Fix an ordered list of $b\ge1$ phases. In phase $s$, a fixed alphabet $A_s$ of $m_s$ actions is pairwise commuting. Multiplicities $n^{(s)}$ may be arbitrarily large. Only the number and order of the phases are fixed.

\begin{theorem}[Single-fold composition with explicit counts]\label{cdc:ct:thm:block-composition}
The relation asserting that the successive phases with multiplicities
$n^{(1)},\ldots,n^{(b)}$ carry $x$ to $y$ has a single-fold quartic representation. It has
\begin{equation}\label{cdc:ct:eq:block-arity}
 (b-1)p+\sum_{s=1}^b(2m_s+2p+pm_s+f_s)
\end{equation}
witnesses, in addition to the endpoint and multiplicity parameters. Its arity is independent of every repetition count.
\end{theorem}
\begin{proof}
Introduce natural boundary states $X_1,\ldots,X_{b-1}$, set $X_0=x$ and $X_b=y$, and use the residual system of Section~\ref{cdc:ct:sec:accelerator} for each phase between $X_{s-1}$ and $X_s$, with disjoint local witnesses. The sum of all residual squares is still a quartic.

For fixed input and all phase counts, the endpoint of each phase is forced by its net displacement. Consequently every intermediate state is unique. Once those states are fixed, Theorem~\ref{cdc:ct:thm:accelerator} gives unique local witnesses whenever the phase executes. Conversely, a zero provides a legal execution of each phase; their concatenation is a legal block program. Counting the boundary states and local witnesses gives~\eqref{cdc:ct:eq:block-arity}.
\end{proof}
The phase counts are semantic data in this theorem. If the same label can occur in adjacent phases, one untagged trace could have more than one decomposition into phase counts. The theorem does not claim to eliminate that ambiguity. One may retain phase tags as part of the semantic object, require counts as input, or establish a separate unique-decomposition result.

\subsection{Fixed-word macros: noncommuting loop bodies}
The same model also compresses repetition of a fixed word whose individual actions need not commute. This is a closure consequence, not a new universality claim.

\begin{proposition}[Exact powers of a fixed word]\label{cdc:ct:prop:word-power}
Let $w=a_1\cdots a_t$ be fixed, put $\delta_i=\sum_{k=1}^i v_{a_k}$ with $\delta_0=0$, and set $\Delta=\delta_t$. Its partial map has domain given by
\begin{align}
 \ell_{w,j}&=\max\bigl(0,\max_{1\le i\le t}(\ell_{a_i j}-\delta_{i-1,j})\bigr),\label{cdc:ct:eq:macro-lo}\\
 u_{w,j}&=\min_{1\le i\le t}(u_{a_i j}-\delta_{i-1,j}).\label{cdc:ct:eq:macro-up}
\end{align}
If this box is nonempty, it defines a valid guarded-translation macro with displacement $\Delta$. For $n>0$, the word $w^n$ executes precisely when
\begin{equation}\label{cdc:ct:eq:word-power}
 x_j\ge\ell_{w,j}+(n-1)\Delta_j^-,\qquad
 x_j\le u_{w,j}-(n-1)\Delta_j^+
\end{equation}
for every coordinate, with infinite upper bounds omitted. Its endpoint is $y=x+n\Delta$. This relation, including $n=0$, has a single-fold quartic representation with $2+3p+f_w$ witnesses, where $f_w$ counts the finite macro upper bounds. The arity is independent of both $n$ and the length of the already compiled fixed word.
\end{proposition}
\begin{proof}
Before letter $a_i$, the formal state is $x+\delta_{i-1}$. Pulling every guard back to $x$ gives~\eqref{cdc:ct:eq:macro-lo}--\eqref{cdc:ct:eq:macro-up}. The same induction on formal prefix states as in Theorem~\ref{cdc:ct:thm:envelope} proves that their intersection is exactly the execution domain. If nonempty, its lower corner itself lies in the domain. Its final state is natural, so $\ell_{w,j}+\Delta_j\ge0$, proving the required validity condition for the macro.

Apply Theorem~\ref{cdc:ct:thm:envelope} to this one macro with multiplicity $n$. Its lower envelope is $n\Delta_j^-+\ell_{w,j}-\Delta_j^-$, and its upper envelope is $u_{w,j}+\Delta_j^+-n\Delta_j^+$, giving~\eqref{cdc:ct:eq:word-power}. Theorem~\ref{cdc:ct:thm:accelerator} with $m=1$ supplies the stated witness count, includes the zero-repetition case, and proves unique completion.

If the macro domain is empty, only $n=0$ can execute; that special case is instead represented without witnesses by $n^2+\sum_j(y_j-x_j)^2$. For the empty word, use the identity macro with lower bound zero and no upper bound.
\end{proof}
This extends the fixed-phase construction to expressions such as $w_1^{n_1}\cdots w_b^{n_b}$ for fixed words $w_s$, even if the letters within a loop body are highly dependent. The macro description is computed in time linear in the total word length times the state dimension; its guard and displacement constants can grow with that length. Thus constant witness arity does not mean that the source word has disappeared from the coefficient description. An empty-domain macro simply forces its repetition count to zero. Multiplicities remain specified parameters unless a separate argument such as the next two corollaries controls their fibres.

\subsection{Endpoint-level single-foldness by displacement rigidity}
There is a useful sufficient criterion for making the counts existential without creating extra roots. List all phase-labelled action occurrences as $q=\sum_s m_s$ action types. Let $V\in\Z^{p\times q}$ have their displacements as columns.

\begin{corollary}[Rigid endpoint certificates]\label{cdc:ct:cor:rank}
If $V$ has full column rank over $\Q$, then the endpoint relation
\[
 R(x,y)\iff\text{some phase multiplicities carry $x$ to $y$}
\]
has a single-fold quartic representation obtained from Theorem~\ref{cdc:ct:thm:block-composition} by turning the $q$ count parameters into witnesses.
\end{corollary}
\begin{proof}
Every root satisfies $y-x=Vn$. Full column rank makes $n$ unique over $\Q$, hence at most one natural count vector can occur. For that vector, Theorem~\ref{cdc:ct:thm:block-composition} makes every remaining witness unique. Existence is unchanged by treating the counts existentially.
\end{proof}
Full column rank is sufficient, not necessary. It is deliberately easy to verify and identifies the source of uniqueness: the endpoint recovers the semantic counts, while the compiler recovers their auxiliary certificate.

\subsection{Endpoint-level finite-foldness from progress}
\begin{corollary}[Progress-bounded fibres]\label{cdc:ct:cor:progress}
Suppose there is a vector $w\in\Q^p$ and $\varepsilon>0$ such that
$w\mathbin{\cdot}v_i\ge\varepsilon$ for every phase-labelled action type. The same existential-count quartic is finite-fold at each fixed pair $(x,y)$.
\end{corollary}
\begin{proof}
Every root obeys
\[
 w\mathbin{\cdot}(y-x)=\sum_i n_i(w\mathbin{\cdot}v_i)
                 \ge\varepsilon\sum_i n_i.
\]
If the left side is negative there are no roots. Otherwise all count vectors lie in the finite set $\{n\in\N^q:\sum_i n_i\le w\cdot(y-x)/\varepsilon\}$. For each such vector the remaining fibre is empty or a singleton by Theorem~\ref{cdc:ct:thm:block-composition}. Hence the complete fibre is finite.
\end{proof}
These corollaries are restricted positive results, not consequences of a general finite-fold theorem. Their hypotheses are explicit linear-algebraic or progress conditions. In particular, the number of phases is fixed in the polynomial, and arbitrary branching between an unbounded number of phases is outside their scope.

\subsection{A directly linear description of each support case}
For a fixed support set of each phase, all active-action tests in the envelope become finitely many linear inequalities in the endpoints, boundary states, and multiplicities. Inactive counts are zero and active counts are positive. Thus these block relations have a finite union of such integer linear descriptions, with boundary states eliminable by the displacement equations. The polynomial construction is valuable because it supplies an explicit unique completion and a uniform degree bound, not because it proves previously unavailable decidability for these restricted relations.


% ---- manuscript 01, lines 557-620 ----
\section{Worked examples\srctag{01}}\label{cdc:ct:sec:examples}
\subsection{Eleven witnesses for an enormous shared-resource block}
Consider three one-place Petri actions:
\begin{equation}\label{cdc:ct:eq:example-actions}
 a:(c,d)=(2,1),\qquad b:(c,d)=(3,1),\qquad r:(c,d)=(1,1).
\end{equation}
The first two decrease the marking by one and two respectively; the third reads the shared retained token. By Proposition~\ref{cdc:ct:prop:petri-commute}, all three commute. For $a,b$, both final production thresholds equal one; the read threshold is at most both retained production thresholds.

Put $L=n_a+2n_b$. If at least one multiplicity is positive, the exact condition is
\begin{equation}\label{cdc:ct:eq:example-condition}
 x\ge L+1,\qquad y=x-L.
\end{equation}
If all multiplicities vanish, the condition is instead $y=x$, including $x=0$. The empty-support case cannot be silently replaced by the positive-support inequality.

In this example the eleven witnesses are
\[
 z_a,z_b,z_r,\quad k_a,k_b,k_r,\quad L,G,\quad\alpha_a,\alpha_b,\alpha_r.
\]
For $i\in\{a,b,r\}$ the support residuals are
\[
 z_i(z_i-1),\quad n_i-z_i(1+k_i),\quad(1-z_i)k_i.
\]
The remaining nine residuals are
\[
 L-n_a-2n_b,\quad G,\quad y+L-x-G,
\]
\[
 z_i(x-L-1-\alpha_i),\quad(1-z_i)\alpha_i
       \qquad(i=a,b,r).
\]
Their eighteen squares define one quartic. The generic construction retains the forced-zero gain witness $G$ for a uniform interface; it could be eliminated in this particular example, so eleven is not asserted to be minimal.

Set
\[
 n_a=10^{100},\quad n_b=10^{80},\quad n_r=10^{90},\qquad
 x=10^{100}+2\cdot10^{80}+1,\quad y=1.
\]
The unique witness is $z_i=1$, $k_i=n_i-1$, $L=n_a+2n_b$, $G=0$, and all three lower slacks zero. These formulas evaluate directly, without traversing the astronomical number of events. The attached expanded polynomial has total degree four and 76 nonzero monomials before specializing the parameters. Its JSON file includes all eighteen residuals, the full expanded coefficients, and the exact natural certificate.

\subsection{Many schedules, one trace, one polynomial root}
For the same three actions, take the normal form
\[
 C_1=\{a,b,r\},\qquad C_2=\{a,r\},\qquad C_3=\{a\}.
\]
It has counts $(3,1,2)$. Since all distinct labels are independent, all
\[
 \frac{6!}{3!\,1!\,2!}=60
\]
words with these counts are one trace. Its height is three, not six, because repeated copies of a label remain dependent.

From $x=9$, the boundary markings are $9,6,5,4$, so $y=4$. The compiler with $m=3$, $p=1$, $H=3$, $f=0$, $e=0$, and $d=3$ produces
\[
 N_H=32,\qquad E_H=48.
\]
There is exactly one root for this trace. There may be other roots for other traces with the same endpoints and height bound; adding the three count equations isolates this particular trace. This qualification is essential: endpoint equality alone is not a specification of the history being counted.

\subsection{Upper guards and an identity block}
Let $a$ be the increment on $0\le x\le4$ and let $r$ be a read on $0\le x\le5$. They commute: the domain of either two-action order is $0\le x\le4$. For multiplicities $n_a,n_r$, the upper envelope includes $x\le5-n_a$ whenever $n_a>0$, and the same inequality from the read whenever it is active; the lower envelope is $x\ge0$. Thus with three increments and any positive number of reads, the allowed starts are exactly $0,1,2$, and the endpoint is $x+3$.

If no increment occurs but the read does occur, the bound is $x\le5$. If neither occurs, there is no upper bound at all. The support variables in the quartic distinguish all three cases without an arbitrary branch-choice witness.

\subsection{A diagnostic for spurious concurrency}
Let $a,b$ be distinct labels, each consuming one token and producing none. At $x=1$, the naive Boolean layer conditions ``both selected actions are individually enabled'' would accept $\{a,b\}$. The exact layer envelope has loss two and retained margin zero, so rejects it. At $x=2$, both sequential orders execute and the layer is valid. This is precisely the operational interpretation of Remark~\ref{cdc:ct:rem:not-step}; the compiler does not confuse a clique with a feasibility certificate.


% ---- manuscript 06, lines 722-1070 ----
\section{Saturation: nested normality only sees zero, one, and two\srctag{06}}
\label{cdc:cs:sec:saturation}

\begin{theorem}[Square saturation]
\label{cdc:cs:thm:squaresaturation}
Every signature $S=(C,R,\alpha,\beta)$ satisfies
\begin{equation}
\label{cdc:cs:eq:square}
S^{\odot2}=
\left(C\prod_{a\in A}(1-R_a\beta_a),\ R,\ \alpha,\ \beta\right),
\qquad S^{\odot3}=S^{\odot2}.
\end{equation}
Consequently $S^{\odot k}=S^{\odot2}$ for every $k\geq2$.
\end{theorem}

\begin{proof}
Apply \eqref{cdc:cs:eq:sigproduct} with both factors equal to $S$.
Since the entries are bits,
$\alpha_a^2=\alpha_a$ and
$R_a\lor(R_a\alpha_a)=R_a$. Disjointness gives
$\alpha_a\beta_a+\beta_a=\beta_a$. The scalar coordinate is the one shown in
\eqref{cdc:cs:eq:square}. Multiplication by $S$ once more changes only the scalar
coordinate and multiplies it by the same Boolean factor
$C\prod_a(1-R_a\beta_a)$. That factor is idempotent, so nothing changes.
Induction proves the assertion for every larger power.
\end{proof}

There are two separate stabilization phenomena. A block's final-memory map
is already idempotent after one application. A second copy can nevertheless
expose a forbidden pattern straddling the boundary between the copies. Once
that possibility has been checked, a third copy cannot expose a new kind of
failure. The signature calculation records both facts, including rejecting
blocks.

Define $\capTwo(k)=\min(k,2)$ and apply it componentwise to parameter tuples.

\begin{theorem}[Cap-at-two theorem for nested schemes]
\label{cdc:cs:thm:cap}
For every expression $E$ from \cref{cdc:cs:def:expression} and every parameter tuple
$\mathbf{k}\in\N^r$,
\begin{equation}
\label{cdc:cs:eq:cap}
\sig\bigl(\llbracket E\rrbracket_{\mathbf{k}}\bigr)
=
\sig\bigl(\llbracket E\rrbracket_{\capTwo(\mathbf{k})}\bigr).
\end{equation}
In particular, the two expanded words are either both normal or both
nonnormal. Parameters may be reused, and repetitions may be nested to any
fixed finite depth.
\end{theorem}

\begin{proof}
Proceed by structural induction. The empty and letter cases do not involve
parameters. For concatenation, the induction hypotheses give equal child
signatures, and the product formula gives equal parent signatures. At a
repetition node $F^{k_j}$, the induction hypothesis first identifies the
signature $S$ of the body under the two parameter assignments. The two parent
signatures are then $S^{\odot k_j}$ and
$S^{\odot\capTwo(k_j)}$, which agree by \cref{cdc:cs:thm:squaresaturation}, including
the identity case $k_j=0$. Reuse of a parameter causes no difficulty because
the same componentwise capping is applied everywhere. Finally, acceptance
from initial memory zero is exactly the scalar condition $C=1$.
\end{proof}

\begin{corollary}[Finite parameter cells for normality]
\label[corollary]{cdc:cs:cor:normalcells}
For a fixed expression with $r$ distinct parameters, its normality predicate
is a union of some of the $3^r$ cells
\[
X_1\times\cdots\times X_r,
\qquad X_j\in\bigl\{\{0\},\{1\},\{2,3,4,\ldots\}\bigr\}.
\]
It can be evaluated from an $O(qN)$-size circuit over the three-way tests on
parameters, without constructing an expanded word or enumerating all cells.
\end{corollary}

\begin{proof}
The theorem shows constancy on each cell. Its inductive proof supplies the
circuit: one signature product per concatenation, and at a repetition choose
the identity, the child signature, or its square according to the three-way
parameter test. Each signature product has $O(q)$ gates.
\end{proof}

\begin{example}[Why the threshold two is sharp]
Take $a<b$ and $(a,b)\in I$. The word $ab$ is normal, but $(ab)^2=abab$ is not:
its middle $ba$ can be swapped. Thus capping all positive powers at one would
be false, even for one unnested repetition node. The theorem does not claim
that normality is unchanged when zero is identified with one, either.
\end{example}

\begin{remark}[What saturation does not preserve]
The words $a^2$ and $a^{1000}$ need not have the same resources, output,
length, Parikh vector, or trace class. The signature is a homomorphic image
of the free \emph{word} monoid, not of the trace monoid: the independent words
$ab$ and $ba$ can have different normality signatures. By contrast, the
resource-summary map factors through the trace relation. Cap-at-two is a
theorem about the canonicality test only, not a compression of the actual
resource action to bounded powers.
\end{remark}

\section{A single-fold quartic compiler\srctag{06}}
\label{cdc:cs:sec:compiler}

We now compile the two summaries together. The construction is elementary:
it uses ordinary polynomial equations, nonnegative unknowns, and no
exponentiation predicate, bounded universal quantifier, or invocation of
MRDP.

\subsection{Quadratic gates with forced values}

A Boolean wire $b$ is constrained by $b(b-1)=0$. With Boolean inputs, the
following equations uniquely determine their Boolean outputs:
\begin{equation}
\label{cdc:cs:eq:booleangates}
t=xy\quad(\mathrm{AND}),\qquad
 t=x+y-xy\quad(\mathrm{OR}),\qquad
 t=1-x\quad(\mathrm{NOT}).
\end{equation}
One may add the output bit equation explicitly, although its value already
follows from the inputs. Long conjunctions are implemented by a chain or a
balanced tree of binary AND gates. Expressions such as
$R_a\lor((S_0)_a\alpha_a)$ must use an intermediate wire for the inner
product; substituting it directly into the OR equation would create a cubic
residual. Likewise the product defining $C'$ in \eqref{cdc:cs:eq:sigproduct} is not
left as a single high-degree residual.

To share a repetition switch with both branches of the compiler, introduce
$z_0,z_1,z_2,h\in\N$ and impose
\begin{align}
\label{cdc:cs:eq:threeway}
z_i(z_i-1)&=0\quad(0\leq i\leq2),&
 z_0+z_1+z_2&=1,\\
k&=z_1+2z_2+h,&(1-z_2)h&=0.
\notag
\end{align}
Exactly one selector is one. If it is $z_0$, then $h=0$ and $k=0$. If it is
$z_1$, then $h=0$ and $k=1$. If it is $z_2$, then $k=2+h\geq2$. Conversely,
each $k$ chooses exactly one case and a unique $h$. Thus the selector gadget
has no surplus solutions.

For resources the active bit is $e=z_1+z_2$, used in
\eqref{cdc:cs:eq:powergadget}. For a signature, compute the child's square and then
select, coordinate by coordinate, among the identity, child, and square:
\begin{equation}
\label{cdc:cs:eq:signatureselection}
s'=z_0 s_{\mathrm{id}}+z_1s+z_2s_{\mathrm{square}}.
\end{equation}
The three terms are at most quadratic, since each $s$ is a wire rather than
an expanded polynomial. Exactly one term is active. Both the child and its
square remain fully constrained in inactive cases.

\subsection{The polynomial and its correctness contract}

At each syntax node allocate its $2p$ resource coordinates. For an empty
node they are zero; for a letter node they are the transition weights. At a
concatenation node impose the minimum gadget coordinatewise and the two
output equations following \cref{cdc:cs:lem:resourceproduct}. At a repetition node
impose the child minimum, the selector gadget, and the power output equations.
For the canonical version, also allocate the signature at each node and the
intermediate Boolean wires described above.

At the root impose, with a new slack vector $r\in\N^p$,
\begin{equation}
\label{cdc:cs:eq:endpoint}
m=U_E+r,\qquad n=V_E+r.
\end{equation}
In the canonical version also impose $C_E=1$. List all integer residuals as
$r_1,\ldots,r_s$. Every residual has total degree at most two in all free
inputs and witnesses. Define
\begin{equation}
\label{cdc:cs:eq:compiledpoly}
Q_E(m,n,\mathbf{k},z)=\sum_{i=1}^s r_i(m,n,\mathbf{k},z)^2.
\end{equation}
This is one ordinary polynomial in $\Z[m,n,\mathbf{k},z]$ of degree at most
four. Its sum-of-squares presentation is simply a convenient exact way to
store it; it is not an extension of polynomial syntax.

\begin{theorem}[Compressed single-fold compilation]
\label{cdc:cs:thm:compiler}
For every fixed expression $E$, the construction above effectively produces
polynomials $Q_E^{\mathrm{res}}$ and $Q_E^{\mathrm{can}}$ of degree at most four.
For every fixed $m,n\in\N^p$ and $\mathbf{k}\in\N^r$,
\begin{align}
\#\{z:Q_E^{\mathrm{res}}(m,n,\mathbf{k},z)=0\}&=
\mathbf1_{m\xrightarrow{\llbracket E\rrbracket_{\mathbf{k}}}n},
\label{cdc:cs:eq:compilerres}\\
\#\{z:Q_E^{\mathrm{can}}(m,n,\mathbf{k},z)=0\}&=
\mathbf1_{m\xrightarrow{\llbracket E\rrbracket_{\mathbf{k}}}n\;\wedge\;
\llbracket E\rrbracket_{\mathbf{k}}\text{ normal}}.
\label{cdc:cs:eq:compilercan}
\end{align}
The resource-only compiler uses $O(pN)$ witnesses and equations. The canonical
compiler uses $O((p+q)N)$ witnesses and equations. These bounds are independent
of the expanded execution length and of the numerical parameter values.
\end{theorem}

\begin{proof}
All residuals are integer-valued. Therefore $Q_E=0$ is equivalent to every
residual being zero: a sum of nonnegative integer squares vanishes only when
each square vanishes.

We first show that all internal witnesses are uniquely determined by the
parameters. At an empty or letter node the summary coordinates are fixed
constants. If the child summaries at a concatenation are fixed,
\cref{cdc:cs:lem:min} uniquely determines the cancellation variables, and the
output equations uniquely determine the parent coordinates. At a repetition,
the child minimum is likewise unique, \eqref{cdc:cs:eq:threeway} uniquely determines
the selectors and remainder, and \eqref{cdc:cs:eq:powergadget} uniquely determines
the parent summary. Structural induction proves uniqueness everywhere in the
resource branch, as well as agreement with the actual word summary by
\cref{cdc:cs:lem:resourceproduct,cdc:cs:thm:power}.

For the canonical branch, the base signatures are fixed. Each elementary
Boolean gate has a unique value for fixed inputs; concatenation therefore
computes exactly \eqref{cdc:cs:eq:sigproduct}. The same argument computes the child
square and the selected power signature, which agrees with the full word
power by \cref{cdc:cs:thm:squaresaturation}. All inactive branches are still
computed, so zero or one repetitions do not release any unconstrained
variables. This proves uniqueness of every signature and intermediate wire.

The only remaining witnesses are the root slacks. By
\cref{cdc:cs:lem:faithful}, they exist exactly for an execution with the stated
endpoints and are unique when they exist. In the canonical version, the
additional equation $C_E=1$ holds exactly when the word accepts from zero
memory, hence exactly when it is normal. This proves both counting equalities,
including nonexistence on rejected inputs.

For the size estimate, each resource node needs a constant number of vectors
of dimension $p$ and each signature node needs $O(q)$ wires. A repetition
switch has constant size. Thus summing over $N$ syntax nodes gives the stated
bounds. The use of intermediate wires ensures the quadratic residual bound,
and \eqref{cdc:cs:eq:compiledpoly} doubles the degree at most.
\end{proof}

\subsection{An explicit resource-only variable bound}

The asymptotic estimate does not hide a dependence on the numerical
repetition counts. A more precise bound is available for a streamlined
resource-only implementation using the two-way switch \eqref{cdc:cs:eq:twoway}.
Let $C$ be the number of concatenation nodes and $R$ the number of repetition
nodes. Allocate all $2p$ summary coordinates at every node, three minimum
witnesses per coordinate at every internal node, two switch witnesses per
repetition, and $p$ endpoint slacks. This gives
\begin{equation}
\label{cdc:cs:eq:resourcecount}
2pN+3p(C+R)+2R+p\ \leq\ (5p+2)N+p
\end{equation}
witnesses. Some are constant or could be shared, but retaining them makes the
inductive proof transparent. The supplied prototype intentionally uses the
three-way switch in both modes and some redundant Boolean bit equations; its
exact counts therefore need not equal this optimized resource-only count.

\subsection{Certificate height and bit complexity}

Let $W$ be the largest transition weight, $K$ the largest supplied repetition
parameter, and $h$ the repetition depth. Put $K=0$ when no parameter occurs.
The next bound concerns a syntax tree, not a recursively defined program.

\begin{proposition}[Witness height]
\label[proposition]{cdc:cs:prop:height}
For an accepted input to the compiler in \cref{cdc:cs:thm:compiler}, its unique
witness can be chosen---and in the displayed construction is forced---with
all coordinates at most
\begin{equation}
\label{cdc:cs:eq:height}
H=\max\left(1,K,\norm{m},\,WN(1+K)^h\right).
\end{equation}
In particular, exponentially long expansions need not require exponentially
many witness bits per coordinate.
\end{proposition}

\begin{proof}
For any subexpression, each letter occurrence in its syntax contributes at
most a product of $h$ repetition counts to its expanded length. Summing over
at most $N$ letter leaves gives length at most $N(1+K)^h$ for every
subexpression, including ones discarded by an outer zero repetition.

For a word of length $L$, its summary satisfies $U_w,V_w\leq WL$
coordinatewise. One way to see this is to start with the bound $W$ at leaves
and use \eqref{cdc:cs:eq:product}, which never exceeds the sum of the child bounds;
an equivalent direct bound follows from \eqref{cdc:cs:eq:prefixformula}. Minimum
witnesses and their residual differences are bounded by the corresponding
child coordinates. The intermediate nonnegative terms in a power output
are bounded by that output. Selector remainders are at most $K$, Boolean
wires are at most one, and the endpoint slack is at most $m$. These bounds
give \eqref{cdc:cs:eq:height}.
\end{proof}

Thus a safe per-coordinate binary-length bound is
\[
O\!\left(1+\log(1+K)+\log(1+\norm m)+\log(1+W)+\log(1+N)
+h\log(1+K)\right).
\]
The root expansion length alone would not provide a correct bound: a zero
outer power can erase an enormous child, whose uniquely constrained internal
certificate is still present. This is another reason to track all syntax
nodes, rather than only the visible final word.

A directed acyclic expression graph can share child summaries and retain the
$O((p+q)N)$ variable bound when $N$ counts distinct graph nodes. The tree-height
estimate must then be replaced by a bound on expansion of shared nodes. A
coarse universal estimate is $2^N(1+K)^N$ for their lengths, obtained by
induction over an ordering of the acyclic graph. Its logarithm is still
polynomial in $N$ and the input bit lengths. The prototype in this package
uses trees, not graph sharing.

\subsection{Worked examples and a scope warning}

\begin{example}[A nested multiplication schedule]
\label[example]{cdc:cs:ex:multiplication}
Let $+$ increment a counter and $-$ decrement it when positive. The scheme
\begin{equation}
\label{cdc:cs:eq:multexample}
E_{x,y,z}=\bigl((++)^x\bigr)^y\; +++\; (-)^z
\end{equation}
executes from zero to zero exactly when
\[
z=2xy+3.
\]
Its expansion first makes $2xy+3$ increments and then $z$ decrements. Thus
there is no intermediate enabling obstruction when the equality holds.
For $x=3,y=4,z=27$, the concrete word has length $54$. For
$x=10^{100}$, $y=10^{120}$ and $z=2xy+3$, the same fixed compiler verifies a
word whose length $2z$ has $221$ decimal digits, without expanding it.
There is no independence between $+$ and $-$ for a single counter, so every
word is normal for the resulting empty independence relation.
\end{example}

\begin{example}[Shared-resource commutation]
Suppose $a$ consumes one token and $b$ consumes two tokens from the same
counter, producing none, and $a<b$. The global commutation criterion includes
$(a,b)$: either order has summary $(3,0)$. From marking $3$, the words $ab$ and
$ba$ are two executable representatives of one trace. The canonical compiler
accepts the former and rejects the latter even though their resource
certificates describe the same partial action.
\end{example}

\begin{remark}[A canonicality filter is not a normalizer]
\label[remark]{cdc:cs:rem:filter}
For a restricted expression family, adding the canonicality condition can
remove every feasible instance. If a fixed expression is $ba$ with
$a<b$ and $(a,b)\in I$, its normal representative $ab$ need not occur anywhere
in that expression family. Therefore \eqref{cdc:cs:eq:compilercan} tests whether the
\emph{given expansion} is canonical; it does not preserve existential
reachability for every restricted schedule language. The trace-class
bijection in the next section avoids this problem by ranging over all words
of the fixed length, a trace-closed universe.
\end{remark}


% ---- manuscript 06, lines 1201-1447 ----
\section{Expressiveness and an exact nesting boundary\srctag{06}}
\label{cdc:cs:sec:hierarchy}

The compiler gives short certificates \emph{once the repetition counts are
fixed}. The existence of suitable counts is a separate problem. We now show
that sharing parameters between nested repetitions is enough to recover
Diophantine universality even with exceptionally simple counter commands.

\subsection{Flat schemes are Presburger-definable}

\begin{theorem}[Flat-scheme decidability]
\label{cdc:cs:thm:flat}
For every fixed token system and every expression of repetition depth at most
one, the relation
\[
\{(m,n,\mathbf{k}):m\xrightarrow{\llbracket E\rrbracket_{\mathbf{k}}}n\}
\]
is effectively definable in Presburger arithmetic. The same holds with the
additional requirement that the expanded word is normal. Consequently,
existence of parameters for fixed endpoints is decidable in this fragment.
\end{theorem}

\begin{proof}
A repetition body has no repetition nodes and hence has a constant resource
summary $(u,v)$. In the zero case its power is $(0,0)$; in the positive case
\eqref{cdc:cs:eq:power} is affine in the repetition parameter, because
$c,\alpha,\beta$ are constants. This case split is expressible with $k=0$ or
$k\geq1$.

Concatenation combines summaries by addition and coordinatewise minimum.
The graph of $c=\min(x,y)$ has the Presburger definition
\[
(x\leq y\ \wedge\ c=x)\ \lor\ (y\leq x\ \wedge\ c=y).
\]
Thus structural induction yields a Presburger formula for all summaries and
endpoints. Reusing a parameter only reuses the same variable in affine
expressions; it does not introduce a product of two variable quantities at
a depth-one repetition node. Normality is also Presburger-definable, either
by the finite cells of \cref{cdc:cs:cor:normalcells} or by the finite Boolean circuit
and the tests $k=0,k=1,k\geq2$. Existential quantification over parameters
preserves Presburger definability, and Presburger arithmetic is decidable.
\end{proof}

This proof deliberately does not use the quadratic equation $st=0$ to
define the minimum inside Presburger arithmetic. It replaces that gadget by
a linear case distinction. The nonlinear certificate language and the
simpler decision procedure for a restricted fragment are compatible but
serve different purposes.

\subsection{Separated one-counter schemes and polynomial degree}

Consider one counter with commands $+$ and $-$. Call a scheme
\emph{separated} if it is a concatenation $E_+E_-$ where $E_+$ uses only $+$
and $E_-$ uses only $-$. Empty expressions are allowed. For a fixed parameter
tuple, let $A(\mathbf{k})$ and $B(\mathbf{k})$ be their respective lengths.

\begin{lemma}[Length polynomial]
\label[lemma]{cdc:cs:lem:lengthpoly}
For any expression using a single letter, its expansion length is a polynomial
in its repetition parameters with nonnegative integer coefficients. Its total
degree is at most the repetition depth.
\end{lemma}

\begin{proof}
The length polynomials of $\eps$ and a letter are $0$ and $1$. Concatenation
adds length polynomials, while repetition by $k_j$ multiplies the body length
by $k_j$. Induction proves both nonnegativity of coefficients and the degree
bound, since repetition increases nesting depth by one and concatenation
takes the maximum of child depths.
\end{proof}

\begin{theorem}[Solution-preserving degree--depth correspondence]
\label{cdc:cs:thm:degree_depth}
For each fixed $d\geq0$, the following two problem presentations are
effectively interconvertible, preserving the parameter solution set exactly:
\begin{enumerate}[label=\textnormal{(\roman*)},leftmargin=2em]
\item one polynomial equation $P(\mathbf{k})=0$ over $\N$, with
$P\in\Z[\mathbf{k}]$ and $\deg P\leq d$;
\item a separated one-counter scheme of repetition depth at most $d$,
required to execute from zero to zero.
\end{enumerate}
No extra existential parameters are needed in either direction.
\end{theorem}

\begin{proof}
A separated scheme executes first $A(\mathbf{k})$ increments and then
$B(\mathbf{k})$ decrements. It executes from zero to zero if and only if
$A(\mathbf{k})=B(\mathbf{k})$: equality supplies enough tokens for every
decrement and leaves none. By \cref{cdc:cs:lem:lengthpoly}, $P=A-B$ has degree at
most $d$ and gives exactly the same parameter solutions.

Conversely split an integer polynomial coefficientwise as $P=P_+-P_-$,
where $P_+,P_-$ have nonnegative coefficients. For a monomial
$c\prod_j k_j^{e_j}$, begin with $c$ copies of the appropriate sign command
and place them under $\sum_j e_j$ nested repetitions, using each $k_j$
exactly $e_j$ times. Its length is the monomial and its repetition depth is
the monomial degree. Concatenate the positive monomial blocks to form $E_+$
and the negative ones to form $E_-$. The resulting scheme has depth at most
$d$ and accepts exactly when $P=0$.
\end{proof}

The direct construction writes a coefficient $c$ as $c$ letter occurrences.
It is an effective solution-set correspondence, not a claim of polynomial
size in a binary encoding of coefficients. Constant blocks could instead be
represented by shared straight-line expressions. We do not need that
optimization for any undecidability or multiplicity conclusion.

At $d=1$, this recovers a linear Diophantine fragment. At $d=2$, the separated
\emph{one-counter} problem is exactly natural solvability of one quadratic
equation; decision procedures for positive integral solutions of arbitrary
quadratic equations, and more general linear-region restrictions, are due to
Grunewald and Segal \cite{grunewald-segal}. Translating variables by one
interconverts positive and nonnegative domains. At $d=3$, the theorem states
an exact equivalence with cubic Diophantine solvability; we make no additional
classification of that problem here. At $d=4$, undecidability follows from
the next degree-reduction lemma and MRDP.

\subsection{Unique quadratic systems and quartic equations}

\begin{lemma}[Multiplicity-preserving degree reduction]
\label[lemma]{cdc:cs:lem:quartic}
For every $P(x,y)\in\Z[x,y]$, one can effectively construct quadratic or
linear polynomials $r_1(x,y,g),\ldots,r_s(x,y,g)$ such that, for each fixed
$(x,y)\in\N^{a+b}$, the system $r_i=0$ has exactly one natural solution $g$
if $P(x,y)=0$, and no solution otherwise. Hence
\[
Q(x,y,g)=\sum_{i=1}^s r_i(x,y,g)^2
\]
has degree at most four and its roots project bijectively to the roots of $P$.
\end{lemma}

\begin{proof}
Split $P=P_+-P_-$ into nonnegative-coefficient polynomials. Evaluate each by a
finite arithmetic circuit using only nonnegative integer constants, original
input variables, addition, and multiplication. Introduce one new variable for
each gate output. An addition gate gives $g-u-v=0$ and a multiplication gate
gives $g-uv=0$, with $u,v$ either previously available variables or constants.
Finally equate the two output wires, or the corresponding constants when a
circuit has no gates.

Topological induction on the circuit shows that every gate variable is
uniquely forced to a nonnegative integer for any fixed original inputs.
The final equation then holds exactly when $P_+=P_-$. This proves the
unique-lifting assertion for the system. A sum of its squares is zero
precisely when every gate equation and the final equality are zero, and its
degree is at most four.
\end{proof}

Avoiding signed intermediate witnesses is important. Representing each signed
integer as an arbitrary difference of two naturals would introduce infinitely
many encodings. Splitting the original polynomial into two nonnegative
circuits avoids that problem entirely. This is a standard arithmetic-circuit
normal form, used here with its multiplicity contract made explicit.

\begin{corollary}[Depth-four one-counter universality]
\label[corollary]{cdc:cs:cor:depthfour}
Every computably enumerable relation $R\subseteq\N^a$ has a representation
\[
x\in R\quad\Longleftrightarrow\quad
\exists k\in\N^b\;
0\xrightarrow{\llbracket E\rrbracket_{(x,k)}}0
\]
by a fixed separated one-counter scheme $E$ of repetition depth at most four.
The scheme uses only increment and enabled decrement. Its parameters may be
globally reused.
\end{corollary}

\begin{proof}
Apply MRDP to obtain a polynomial representation of $R$. Apply
\cref{cdc:cs:lem:quartic}, and then \cref{cdc:cs:thm:degree_depth} at $d=4$. Treat the input
coordinates $x$ as prescribed repetition parameters and all the other
coordinates as existential repetition parameters.
\end{proof}

This is a normal form for \emph{existential schedule parameterization}. It
does not say that an ordinary one-counter machine with freely chosen steps
can compute an arbitrary partial recursive function. The numerical
relationships between globally shared loop counts carry the Diophantine
expressiveness.

\subsection{Depth two is already universal with several resources}

For each counter $j$, let $+_j$ and $-_j$ increment or decrement that counter
and leave all others unchanged. A vector-separated scheme consists of a
positive block using only the $+_j$, followed by a negative block using only
the $-_j$.

\begin{theorem}[Quadratic-system substrate and sharp flat/nested boundary]
\label{cdc:cs:thm:depthtwo}
A finite system of polynomial equations $P_j(\mathbf{k})=0$ of degrees at
most $d$ can be represented, without new parameters, by a vector-separated
scheme of depth at most $d$ on one counter per equation, executing from
$\zeros$ to $\zeros$. Conversely, every such scheme gives a system of
polynomial equations of degrees at most its depth.

In particular, every computably enumerable relation has a fixed
vector-separated representation of depth at most two. Therefore parameter
existence for depth-two schemes in this class is undecidable, whereas it is
decidable for all depth-one token schemes by \cref{cdc:cs:thm:flat}.
\end{theorem}

\begin{proof}
Write $P_j=A_j-B_j$ coefficientwise with nonnegative polynomials. In the
positive phase, use the monomial construction in
\cref{cdc:cs:thm:degree_depth} to perform $A_j$ increments of counter $j$. Do this
for every $j$. In the negative phase, perform $B_j$ decrements of counter $j$
for every $j$. Counter $j$ starts and ends at zero exactly when $A_j=B_j$.
If all equalities hold, every negative phase is enabled. If the run ends at
zero, its net increment in every coordinate is zero, so all equalities hold.
The depth bound and the absence of extra parameters follow from the monomial
construction. The converse takes the increment and decrement counts in each
coordinate and applies the length-polynomial induction.

For a computably enumerable relation, MRDP supplies one polynomial. Replace
it by the quadratic system of \cref{cdc:cs:lem:quartic}, without first summing its
squares. The preceding construction gives a depth-two scheme, with the gate
variables among its existential parameters. Taking a computably enumerable
undecidable relation proves undecidability. Since a single universal relation
has one fixed polynomial and one fixed quadratic system, even the number of
counters and the schedule template can be fixed for that relation; only its
input parameters vary. The depth-one decision theorem proves the sharp
boundary in the stated multi-resource class.
\end{proof}

\begin{corollary}[Universality can coexist with an already canonical word]
The universal schemes in \cref{cdc:cs:thm:depthtwo} can be chosen so that every
expansion is lexicographically normal for the global commutation relation.
\end{corollary}

\begin{proof}
In the positive phase, group all occurrences of $+_1$, then $+_2$, and so on.
Group the negative phase by $-_1,-_2,\ldots$ in the same way. Choose the total
alphabet order
\[
+_1<\cdots<+_p<-_1<\cdots<-_p.
\]
Every expanded word is globally nondecreasing in that order. It is therefore
lexicographically least even among all permutations with the same letter
multiplicities, hence also among the smaller set of words in its trace class.
Grouping does not change the coordinatewise positive and negative counts or
the depth of any monomial block.
\end{proof}

The nontrivial search thus does not come from interleaving ambiguity. It
persists for sorted, phase-separated schedules whose commands merely change
one counter by one.


\part{Polynomial trajectories}\label{cdc:part:poly}

\paragraph{Source and scope.} \wtag{} This Part is manuscript~02, ``Canonical Diophantine Certificates for Polynomial Trajectories'': its statement of the problem and of the restricted result (the first two subsections of its introduction), then its Sections 3--9 (the finite-difference sign certificate, the single-fold quartic theorem and its cost, Boolean guards, acceleration of polynomial trajectories, the computational substrates covered, nontermination certificates, and worked examples), and its conclusion. Its canonical signed coordinates and total circuits (its Section~2) are printed with the conventions, in \cref{cdc:pt:sec:arithmetic}; its remaining introduction is in \cref{cdc:sec:manuscripts}; its other substrates are in Parts~VII and~VIII and its boundary results in Part~IX.

\paragraph{Setting and hypotheses.} \wtag{} Coefficient parameters are \emph{signed} integers, $\boldsymbol c\in\Z^{d+1}$, with $d$ a nominal degree bound; existential coordinates range over $\N$. $T$ is the horizon of the index range $\{0,\ldots,T\}$, not a number of execution steps. The main theorem is a family $P_d$ indexed by the degree bound, uniform in coefficients and horizon.

\section{The problem addressed and the precise restricted result\srctag{02}}\label{cdc:pt:sec:problem}
\srcnote{\wtag{} The first two subsections of manuscript~02's introduction; the rest of it is in \cref{cdc:sec:ms02}.}

% ---- manuscript 02, lines 97-160 ----
\subsection{Existence is not multiplicity preservation}
A computational model can be represented arithmetically in several different
senses. An existence-preserving representation says that a computation
exists exactly when an equation has a solution. A counting-faithful
representation must additionally explain which solutions correspond to
which computations, and whether auxiliary arithmetic choices create extra
solutions. Even a deterministic computation can acquire infinitely many
Diophantine witnesses through a noncanonical trace encoding.

The distinction is central to Matiyasevich's finite-fold and single-fold
questions. Classical results give both polynomial Diophantine
representations of enumerable sets and single-fold \emph{exponential}
Diophantine representations. Combining the two properties is a separate
problem, not a consequence of the ordinary MRDP theorem
\cite{mat2010}. The present paper does not assume that combination.

Our starting question is more specific:
\begin{quote}
Can a long, deterministic, polynomial trajectory with nonmonotone guards be
replaced by a fixed-dimensional polynomial certificate that has exactly one
witness when the claimed run is valid?
\end{quote}
Here ``fixed-dimensional'' means independent of the numerical duration of
the run; the dimension may depend on the fixed program and on fixed degree
bounds. The answer proved below is affirmative.

\subsection{The precise restricted result}
For $d\in\N$, put
\[
 p_{\boldsymbol c}(t)=\sum_{j=0}^d c_jt^j,
 \qquad \boldsymbol c\in\Z^{d+1},\quad T\in\N.
\]
Coefficients may vanish, so $d$ is a \emph{nominal} degree bound, not
necessarily the actual degree. Define
\[
 \PolyPos_d(\boldsymbol c,T)
 \quad\Longleftrightarrow\quad
 p_{\boldsymbol c}(i)\geq0\quad\text{for every integer }0\leq i\leq T.
\]
The main theorem constructs an effective family of polynomials $P_d$ with
\begin{equation}
 \PolyPos_d(\boldsymbol c,T)
 \quad\Longleftrightarrow\quad
 \exists\boldsymbol w\in\N^{M_d}\;
 P_d(\boldsymbol c,T,\boldsymbol w)=0,
 \label{cdc:pt:eq:intro-main}
\end{equation}
where $\deg P_d\leq4$, $M_d=O((d+1)^3)$, and every true instance has
\emph{exactly one} complete witness $\boldsymbol w$.

The family is uniform in coefficients and horizon. It is not a single
polynomial accepting an arbitrarily long coefficient list as one encoded
input. Varying $d$ changes the finite set of variables and the polynomial
itself. This distinction prevents an unjustified leap from our theorem to
arbitrary bounded arithmetic.

The main mathematical development proceeds in three stages. First, a finite
tower of differences yields a unique short sign table. Second, a completely
evaluated arithmetic and Boolean circuit checks this table, including its
padding and all comparisons. Third, canonical signed-number and sign
relations turn this circuit into one quartic equation. Only after uniqueness
of the entire witness has been proved do we project away an execution time,
and then only for a first-exit relation where that time is unique.


% ---- manuscript 02, lines 341-1251 ----
\section{The finite-difference sign certificate\srctag{02}}
\label{cdc:pt:sec:sign-tower}

\subsection{Differences and sign runs}
For $p\in\Z[t]$, define
\[
 \Delta p(t)=p(t+1)-p(t),\qquad q_k=\Delta^kp\quad(0\leq k\leq d).
\]
Then $\deg q_k\leq d-k$, and $q_d$ is constant. In ascending coefficient
order, the difference transformation is explicit:
\begin{equation}
 [t^j]\Delta p(t)=\sum_{r=j+1}^d\binom rj c_r.
 \label{cdc:pt:eq:coefficient-difference}
\end{equation}
Thus every coefficient of every $q_k$ is a fixed integer linear
combination of the input coefficients.

A \emph{sign run} on $\{0,\ldots,T\}$ is a maximal nonempty interval of
consecutive integers on which the sign is constant. We encode the interval
as $[a,b)\cap\Z$, with $0\leq a<b\leq T+1$. The signs $-1,0,1$ are
encoded by natural tags $0,1,2$, respectively.

\begin{lemma}[A bound on the number of runs]
\label{cdc:pt:lem:runs}
A real polynomial of degree at most $m$ has at most $2m+1$ sign runs on any
nonempty finite interval of consecutive integers. The zero polynomial has
one run.
\end{lemma}
\begin{proof}
Suppose the polynomial is nonzero. Let $Z$ be the number of distinct
integer zeros in the interval. Let $C$ be the number of edges between
consecutive sampled integers at which the signs are strictly opposite,
with neither value zero. The intermediate value theorem puts a noninteger
real root in the open interval of each such edge. These roots are distinct
from one another and from the $Z$ integer roots. Therefore $C+Z\leq m$.

Every change of sign tag occurs either at one of these $C$ edges or at an
edge incident with an integer zero. There are at most $2Z$ edges of the
second kind. Thus the number of tag changes is at most
$C+2Z\leq2m$. There is one more run than changes. Consecutive zero values
cause no difficulty: edges joining two zeros are not changes, and our bound
may only overcount them.
\end{proof}

\subsection{Canonical padding}
Set
\[
 L_k=2(d-k)+1\qquad(0\leq k\leq d).
\]
Level $k$ has $L_k$ slots, with cuts and tags
\[
 b_{k,0},\ldots,b_{k,L_k},\qquad
 \tau_{k,0},\ldots,\tau_{k,L_k-1}.
\]
The outside cuts are fixed:
\[
 b_{k,0}=0,\qquad b_{k,L_k}=T+1.
\]
All other cuts and all tags are natural primary witness coordinates. The
following finite conditions specify the permitted shape:
\begin{align}
 b_{k,j}&\leq b_{k,j+1},& 0&\leq\tau_{k,j}\leq2,
 \label{cdc:pt:eq:shape-order}\\
 b_{k,j}=b_{k,j+1}&\Longrightarrow
 \bigl(b_{k,j}=T+1\ \wedge\ \tau_{k,j}=1\bigr),
 \label{cdc:pt:eq:shape-padding}\\
 b_{k,j}<b_{k,j+1}<b_{k,j+2}&\Longrightarrow
 \tau_{k,j}\neq\tau_{k,j+1}.
 \label{cdc:pt:eq:shape-maximal}
\end{align}
The last condition is imposed only when the displayed indices exist.
A slot is \emph{active} when its two cuts are different. Condition
\eqref{cdc:pt:eq:shape-padding} forces all inactive slots to be terminal padding;
their tag is the fixed zero-sign tag. Condition
\eqref{cdc:pt:eq:shape-maximal} prohibits splitting a constant-sign run into two
active slots. Every genuine maximal-run table therefore has exactly one
padded encoding.

The count of primary coordinates is exact:
\begin{equation}
 \sum_{k=0}^d\bigl((L_k-1)+L_k\bigr)
 =\sum_{m=0}^d(4m+1)
 =(d+1)(2d+1).
 \label{cdc:pt:eq:primary-count}
\end{equation}
The $2(d+1)$ outside cuts are not counted, because they are fixed
expressions rather than witnesses.

\subsection{The local endpoint condition}
The top level has one slot. Require
\begin{equation}
 \tau_{d,0}=1+\sgn(q_d(0)).
 \label{cdc:pt:eq:top-level}
\end{equation}
Now let $k<d$. Consider any active parent slot
$[a,b)=[b_{k,j},b_{k,j+1})$ and active child slot
$[c,e)=[b_{k+1,\ell},b_{k+1,\ell+1})$. Define
\begin{equation}
 u=\max(a,c),\qquad v=\min(b-1,e).
 \label{cdc:pt:eq:intersection}
\end{equation}
Whenever $u\leq v$, require
\begin{equation}
 1+\sgn(q_k(u))=\tau_{k,j}
 =1+\sgn(q_k(v)).
 \label{cdc:pt:eq:endpoint-condition}
\end{equation}
The child right endpoint in \eqref{cdc:pt:eq:intersection} is $e$, not $e-1$.
This is intentional: information about $\Delta q_k$ on $[c,e)$ makes
$q_k$ monotone on the integer interval $[c,e]$. Using $e-1$ would also
cover the domain if handled consistently, but \eqref{cdc:pt:eq:intersection}
gives a convenient closed monotonicity interval and is the convention used
throughout the implementation.

All pairs of slots are considered. An implication disables the endpoint
assertion for an inactive slot or empty intersection. Its arithmetic
computations are nevertheless total, as in Lemma~\ref{cdc:pt:lem:circuit}.
Maxima and minima are implemented by comparison bits and selection gates,
so these conditions are a fixed quantifier-free polynomial formula for
fixed $d$.

\begin{lemma}[Endpoint propagation]
\label{cdc:pt:lem:propagation}
Assume every child tag at level $k+1$ is correct. If a parent slot satisfies
\eqref{cdc:pt:eq:endpoint-condition} for every active child slot, then its tag is
the actual sign tag of $q_k$ at every integer point in that parent slot.
\end{lemma}
\begin{proof}
On an active child interval $[c,e)$, the difference
$q_k(i+1)-q_k(i)=q_{k+1}(i)$ has a constant sign. Hence $q_k$ is
nondecreasing, nonincreasing, or constant on $[c,e]\cap\Z$.
Its restriction to the intersection
\[
 [a,b-1]\cap[c,e]\cap\Z=[u,v]\cap\Z
\]
is also monotone. If both endpoint values are positive, every intervening
value is positive; the analogous statement holds for negative values. If
both endpoint values are zero, monotonicity forces every intervening value
to be zero. Thus the parent tag is correct throughout each nonempty
intersection.

Every integer $t$ in the parent slot belongs to an active child slot,
because the child slots partition $\{0,\ldots,T\}$. For that child,
$t\in[c,e-1]\subseteq[c,e]$, so the relevant intersection is nonempty and
contains $t$. This proves correctness at every point of the parent.
\end{proof}

\begin{theorem}[Existence and uniqueness of the sign tower]
\label{cdc:pt:thm:tower}
For every $\boldsymbol c\in\Z^{d+1}$ and $T\in\N$, there is exactly one
primary tuple satisfying
\eqref{cdc:pt:eq:shape-order}--\eqref{cdc:pt:eq:endpoint-condition}. It is the
canonically padded sequence of maximal sign-run tables of
$q_0,\ldots,q_d$ on $\{0,\ldots,T\}$.
\end{theorem}
\begin{proof}
For existence, take the actual maximal sign runs of each $q_k$.
Lemma~\ref{cdc:pt:lem:runs} provides at most $L_k$ runs, so terminal padding is
possible. The shape conditions hold by construction. The top difference
is constant, so \eqref{cdc:pt:eq:top-level} holds. Each endpoint appearing in
\eqref{cdc:pt:eq:endpoint-condition} is an integer point in the parent slot;
its tag is therefore correct.

For soundness, start at the top level, whose unique active slot is correct
by \eqref{cdc:pt:eq:top-level}. Descend through $k=d-1,\ldots,0$ and apply
Lemma~\ref{cdc:pt:lem:propagation}. Every active slot at every level is then a
constant-sign interval with its true tag.

Finally, the shape conditions imply that the active slots form an ordered
partition of the entire domain, and adjacent active slots have different
tags. Such a partition must be the maximal sign-run partition: a boundary
inside a constant-sign run is prohibited, and a missing boundary at a sign
change would make a slot nonconstant. Terminal empty slots and their tags
are prescribed by \eqref{cdc:pt:eq:shape-padding}. Consequently every coordinate
is uniquely determined.
\end{proof}

\begin{remark}[Degenerate cases are part of the theorem]
The proof includes $T=0$, the zero polynomial, a nominal degree larger than
the actual degree, and runs consisting of one integer zero. No genericity,
simple-root, nonzero-leading-coefficient, or sufficiently-large-horizon
assumption is used. A zero polynomial at any level has one active zero-sign
slot followed by uniquely fixed padding.
\end{remark}

\section{The single-fold quartic theorem and its cost\srctag{02}}
\label{cdc:pt:sec:main}

\begin{theorem}[Canonical bounded polynomial positivity]
\label{cdc:pt:thm:main}
There is an effective family of integer polynomials $P_d$, indexed by
$d\in\N$, and an absolute constant $C$ such that
\[
 \deg P_d\leq4,\qquad M_d\leq C(d+1)^3,
\]
and \eqref{cdc:pt:eq:intro-main} is a single-fold representation for every signed
coefficient tuple and natural horizon. Exactly $(d+1)(2d+1)$ primary
coordinates suffice before auxiliary arithmetic and Boolean wires are
introduced. Both the number of quadratic residuals and a bounded-arity
straight-line description length are $O((d+1)^3)$.
\end{theorem}
\begin{proof}
Take the finite sign-tower conditions and add
\begin{equation}
 b_{0,j}<b_{0,j+1}\Longrightarrow \tau_{0,j}\neq0
 \quad(0\leq j<L_0).
 \label{cdc:pt:eq:nonnegative-tags}
\end{equation}
By Theorem~\ref{cdc:pt:thm:tower}, the primary tuple is unique before these
conditions are added. They hold precisely when $q_0=p$ has no negative
value in the domain. Apply Lemma~\ref{cdc:pt:lem:circuit} to the complete
quantifier-free checking formula. It gives one quartic equation, and its
internal wires have a unique extension for every fixed primary tuple.
Thus the complete existential tuple is unique whenever the positivity
relation holds, and no tuple exists otherwise.

It remains to justify the cubic size bound, rather than merely asserting
that the formula is finite. At a level where the parent degree bound is
$m\geq1$, there are $(2m+1)(2m-1)=4m^2-1$ parent--child pairs.
Their intersection and tag comparisons take boundedly many gates per
pair. The potentially expensive operations are polynomial evaluations.

Do not evaluate $q_k$ afresh at the selected endpoint of every pair.
The possible left endpoints are parent starts and child starts; the
possible right endpoints are parent stops minus one and child stops.
There are $O(m)$ such candidate expressions. Evaluate $q_k$ once at each
candidate by Horner's rule, using $O(m)$ arithmetic gates per evaluation,
and cache the sign tags. A comparison then selects the appropriate
precomputed endpoint tag in $O(1)$ gates. The total cost at this level is
$O(m^2)$.

The difference coefficient tower can be constructed in $O((d+1)^3)$
scalar arithmetic operations using \eqref{cdc:pt:eq:coefficient-difference}.
The shape tests use $O((d+1)^2)$ gates. Summing the $O(m^2)$ level costs
therefore gives $O((d+1)^3)$ gates overall. The canonical circuit lemma
introduces only a constant number of variables and residuals per gate.
The exact primary count is \eqref{cdc:pt:eq:primary-count}.
\end{proof}

\subsection{An explicit specification, not only an existence proof}
The preceding construction completely specifies $P_d$: unroll the stated
slot conditions, evaluate the stated finite circuit, introduce the signed
wire and sign equations, and sum the squared residuals. A factored
sum-of-squares description is an exact description of an ordinary
polynomial; expanding all monomials is unnecessary.

The accompanying compiler makes this description executable. Its exported
JSON records every residual as a sparse polynomial with integer
coefficients and named variables. The equation represented by the file is
\[
 P_d=Q_1^2+\cdots+Q_N^2=0.
\]
There is no exponentiation primitive, bounded quantifier, floor function,
black-box Diophantine predicate, or implicit oracle left in this output.
The only powers are the fixed squares and products appearing in the
polynomial itself.

\subsection{Generating the unique witness without scanning the horizon}
\begin{proposition}[Generation and verification complexity]
\label{cdc:pt:prop:complexity}
The sign tower can be generated using
$O((d+1)^3\log(T+2))$ exact arithmetic operations, and its canonical
endpoint checker can be implemented using $O((d+1)^3)$ such operations.
With dense binary integer coefficients and binary $T$, generation and
verification have polynomial bit complexity in
\[
 d+1,\quad \log(H+2),\quad\log(T+2),
 \qquad H=\max_j|c_j|.
\]
The arithmetic auxiliary witnesses also have polynomially bounded bit
length in these quantities.
\end{proposition}
\begin{proof}
The constant top row is immediate. Suppose the child row has been
constructed. On each child interval $[c,e)$, $q_k$ is monotone. For a
nondecreasing sequence, find the first index with value at least zero and
the first index with value strictly greater than zero by binary search.
These two indices split the interval into its negative, zero, and positive
subintervals. Empty subintervals are discarded. For a nonincreasing
sequence, apply the same procedure to $-q_k$ and reverse the sign labels.
A constant sequence is included in the nondecreasing case. Merge adjacent
subintervals with equal tags when passing from one child interval to the
next.

There are $O(m)$ child intervals at a level of degree bound $m$. Each
requires at most two binary searches, each with $O(\log(T+2))$
evaluations. Horner evaluation costs $O(m)$ arithmetic operations. This
gives the generation bound after summing over the levels. The verification
bound follows from the cached endpoint implementation in the proof of
Theorem~\ref{cdc:pt:thm:main}.

For bit complexity, the explicit identity
\[
 \Delta^kp(t)=\sum_{\ell=0}^k(-1)^{k-\ell}\binom{k}{\ell}p(t+\ell)
\]
shows that the coefficient heights of all differences have bit length
$O(\log(H+2)+d\log(d+2))$. One can also obtain this bound by expanding
$(t+\ell)^r$ and bounding its binomial coefficients. Evaluation at an
integer of magnitude at most $T+1$ adds
$O(d\log(T+2))$ bits. All candidate cuts lie in $[0,T+1]$, signs and
tags are bounded, and the sign-test slack is at most the magnitude of its
argument. Polynomially many additions and multiplications on numbers with
these polynomial bit bounds establish the claim using ordinary integer
arithmetic. The same estimates apply to stored arithmetic outputs and
comparison slacks. No assertion that unit-cost arithmetic is unit-cost
\emph{bit} computation is needed.
\end{proof}

\begin{remark}[Cubic size is an upper bound, not an optimum]
The supplied compiler favors transparency and total evaluation over small
constants. The theorem does not claim minimal witness count, optimal
quartic sparsity, or an improved universal Diophantine equation. The
specialized quartics should not be compared numerically with universal
witness-count records as if they represented the same class of relations.
\end{remark}

\section{Boolean guards and empty execution prefixes\srctag{02}}
\label{cdc:pt:sec:boolean}

\begin{theorem}[Bounded Boolean polynomial guards]
\label{cdc:pt:thm:boolean}
Fix a Boolean formula $\Phi$ in signs of $r\geq1$ polynomials
$p_1(t),\ldots,p_r(t)$, with respective fixed degree bounds $d_j\leq D$.
The relation
\[
 \forall i\in\{0,\ldots,T\}\;
 \Phi\bigl(\sgn p_1(i),\ldots,\sgn p_r(i)\bigr)
\]
has an effective single-fold quartic representation, uniform in all
coefficients and in $T$. A witness and circuit size bound is
\[
 O\!\left(r(D+1)^3+r^2(D+1)^2+|\Phi|r(D+1)\right).
\]
A guard with no polynomial atoms is simply a constant Boolean guard and is
handled directly.
\end{theorem}
\begin{proof}
Construct the unique difference sign tower for every $p_j$. Take the union
of the starts of all active level-zero runs. It includes zero. Between two
consecutive distinct starts in this union, every polynomial has constant
sign, hence $\Phi$ has constant truth value.

It is unnecessary to existentially sort or deduplicate this union. For
every active start of every polynomial, evaluate all $r$ polynomials at
that start and require $\Phi$ to be true there. Repeated starts merely
repeat a true assertion and do not add choices. These tests are sufficient
because every common sign cell begins at one of the tested starts; they
are necessary because each tested start belongs to the domain.

There are $O(r(D+1))$ candidate starts. Evaluating $r$ polynomials of
degree at most $D$ at each uses $O(r^2(D+1)^2)$ gates. Evaluating the
Boolean formula at these starts uses $O(|\Phi|r(D+1))$ gates. The sign
towers contribute $O(r(D+1)^3)$ gates. Every primary tower is unique and
every extra evaluation is total and deterministic, so
Lemma~\ref{cdc:pt:lem:circuit} gives the claimed single-fold quartic.
\end{proof}

Strict comparisons, equations, disequalities, and arbitrary negations are
all included. For example, a guard may be
\[
 (p_1>0\wedge p_2\neq0)\ \vee\ (p_3=0\wedge\neg(p_4\geq0)).
\]
There is no monotonicity assumption on this entire guard or on $p_j$.
Monotonicity is used only on the certified difference intervals.

\subsection{Why a zero-length run requires explicit treatment}
In execution semantics the usual condition is $\forall i<n\;G(i)$.
For $n=0$ this is true, even if $G(0)$ is false. Replacing it by a test on
$[0,\max(n-1,0)]$ without an additional condition would therefore be wrong.
Use instead the uniquely determined number
\[
 h=\max(n-1,0).
\]
Always construct the full sign towers on $[0,h]$, but enable the final
universal-guard assertions only when $n>0$. The tables and their checking
wires remain fully determined when $n=0$; only the demand that the guard be
true is disabled. The graph of this maximum is a quantifier-free
piecewise-polynomial relation, and the circuit lemma supplies its unique
arithmetic implementation. This method preserves both vacuity and
single-foldness.

\section{Accelerating polynomial trajectories\srctag{02}}
\label{cdc:pt:sec:trajectories}

\begin{definition}[Polynomially iterable map]
A fixed integer-state map $f:\Z^s\to\Z^s$ is \emph{polynomially iterable}
when there is a fixed vector
$F(t,\boldsymbol x)\in\Q[t,\boldsymbol x]^s$ such that
\[
 F(0,\boldsymbol x)=\boldsymbol x,
 \qquad F(t+1,\boldsymbol x)=f(F(t,\boldsymbol x))
\]
for every natural $t$ and integer $\boldsymbol x$. In particular,
$F(t,\boldsymbol x)=f^t(\boldsymbol x)$ is integer-valued at these inputs.
The trajectory formula, not just its existence, is supplied as part of the
fixed program specification.
\end{definition}

Let $G(\boldsymbol x)$ be a fixed Boolean combination of polynomial
comparisons with integer coefficients. The loop tests $G$ \emph{before}
performing each simultaneous update. Define
\begin{align}
 \Run_f^G(\boldsymbol x,n,\boldsymbol y)
 \quad\Longleftrightarrow\quad&
 \boldsymbol y=F(n,\boldsymbol x)
 \ \wedge\ \forall i<n\;G(F(i,\boldsymbol x)),
 \label{cdc:pt:eq:run}\\
 \Exit_f^G(\boldsymbol x,n,\boldsymbol y)
 \quad\Longleftrightarrow\quad&
 \Run_f^G(\boldsymbol x,n,\boldsymbol y)
 \ \wedge\ \neg G(\boldsymbol y).
 \label{cdc:pt:eq:exit}
\end{align}
The run relation allows stopping while the guard is still true; the exit
relation does not. This semantic distinction is essential when an
existential execution time is introduced.

\begin{theorem}[Single-fold polynomial acceleration]
\label{cdc:pt:thm:acceleration}
For every fixed polynomially iterable $f$ and fixed polynomial Boolean
guard $G$, both \eqref{cdc:pt:eq:run} and \eqref{cdc:pt:eq:exit} have effective
single-fold representations by ordinary integer quartics with finitely
many natural witnesses. Their witness dimensions are independent of $n$
and of the values of the initial and final states.
\end{theorem}
\begin{proof}
Substitute $F(t,\boldsymbol x)$ into each polynomial atom of $G$.
The result is a polynomial in $t$ whose coefficients are rational
polynomials in $\boldsymbol x$. Multiply each atom's polynomial difference
by a fixed positive common denominator. This preserves its sign and gives
an integer polynomial in $(t,\boldsymbol x)$. Its degree in $t$ is fixed
by $f$ and $G$, not by $n$ or by the state values.

Compute the coefficient values from $\boldsymbol x$ using a fixed
arithmetic circuit and canonical signed wires. Apply
Theorem~\ref{cdc:pt:thm:boolean} to these polynomials, with the empty-prefix
convention from Section~\ref{cdc:pt:sec:boolean}. Add the equations
$\boldsymbol y=F(n,\boldsymbol x)$ after clearing fixed denominators.
These equations can be checked by the same canonical circuit lemma, so
no polynomial degree increase survives the introduction of gate variables.
For exit, add the Boolean test $\neg G(\boldsymbol y)$.

Coefficient computations and endpoint computations are unique. The primary
sign tables are unique. All their auxiliary wires are unique.
Consequently the complete witness is unique for each valid
$(\boldsymbol x,n,\boldsymbol y)$. Every invalid triple fails at least one
of the equivalent checking conditions. Finally sum the squared quadratic
residuals.
\end{proof}

\begin{corollary}[Projecting the first-exit time]
\label{cdc:pt:cor:first-exit}
The halting domain
\[
 H_f^G(\boldsymbol x)\iff
 \exists n\in\N\;\exists\boldsymbol y\in\Z^s\;
 \Exit_f^G(\boldsymbol x,n,\boldsymbol y)
\]
and the partial output graph
\[
 O_f^G(\boldsymbol x,\boldsymbol y)\iff
 \exists n\in\N\;\Exit_f^G(\boldsymbol x,n,\boldsymbol y)
\]
have effective single-fold quartic representations with natural witnesses.
\end{corollary}
\begin{proof}
For a fixed initial state, there is at most one first-exit time. Indeed, if
$n<m$ were two such times, the exit at $n$ would give
$\neg G(F(n,\boldsymbol x))$, while the valid prefix for $m$ would give
$G(F(n,\boldsymbol x))$. The final state is then uniquely
$F(n,\boldsymbol x)$. Thus existentially projecting these quantities
cannot introduce additional choices. Represent every projected signed
state coordinate by the normalized pair of Lemma~\ref{cdc:pt:lem:signed}; merely
using an unrestricted difference of two naturals would be insufficient.
The existing certificate for the unique exit triple has a unique extension.
The added normalization equations are quadratic.
\end{proof}

\begin{remark}[Reachability is not automatically single-fold]
Replacing $\Exit$ by $\Run$ in this corollary would be unjustified. For the
identity update and the guard \emph{true}, the same state is reachable at
every duration. Even a perfect single-fold certificate for each fixed
$n$ then gives infinitely many witnesses after $n$ is projected away.
The theorem does not rule out a different single-fold representation of
that trivial relation; it identifies why this particular projection
argument would fail.
\end{remark}

\section{Computational substrates covered by the theorem\srctag{02}}
\label{cdc:pt:sec:substrates}

\subsection{Unipotent affine state transformations}
Consider
\[
 f(\boldsymbol x)=A\boldsymbol x+\boldsymbol b,
 \qquad A\in M_s(\Z),\quad\boldsymbol b\in\Z^s.
\]
Introduce the augmented matrix
\[
 M=\begin{pmatrix} A&\boldsymbol b\\0&1\end{pmatrix},
 \qquad \widetilde{\boldsymbol x}=(\boldsymbol x,1)^{\mathsf T}.
\]
Then $M^t\widetilde{\boldsymbol x}$ is the augmented state at time $t$.

\begin{proposition}[Unipotent affine acceleration]
\label{cdc:pt:prop:unipotent}
Suppose $(M-I)^r=0$ for a fixed $r\geq1$. Then
\begin{equation}
 M^t=\sum_{j=0}^{r-1}\binom tj(M-I)^j\qquad(t\in\N).
 \label{cdc:pt:eq:unipotent}
\end{equation}
Consequently this affine update is polynomially iterable and every fixed
polynomial Boolean guard admits the single-fold execution, first-exit,
and output representations above.
\end{proposition}
\begin{proof}
Expand $(I+(M-I))^t$ by the finite binomial theorem. All powers of
$M-I$ of exponent at least $r$ vanish. For a fixed $j$, the expression
\[
 \binom tj=\frac{t(t-1)\cdots(t-j+1)}{j!}
\]
is a rational polynomial in $t$, with the usual value $1$ for $j=0$.
For natural $t<j$ it is zero, so \eqref{cdc:pt:eq:unipotent} also covers small
$t$. Matrix multiplication by the initial augmented state gives the
required polynomial trajectory. The identities are valid over integers
after clearing fixed denominators.
\end{proof}

This includes translations, counter transfers with nilpotent coupling,
and accumulation chains. State coordinates may increase or decrease;
nonnegativity requirements for a natural-counter interpretation can be
inserted as guard atoms. There is no assumption that the full state or the
guard value is monotone.

\subsection{Unitriangular nonlinear updates}
Consider simultaneous assignments
\begin{align}
 x_1'&=x_1+c,\nonumber\\
 x_i'&=x_i+R_i(x_1,\ldots,x_{i-1})\qquad(2\leq i\leq s),
 \label{cdc:pt:eq:triangular}
\end{align}
where $c\in\Z$ and every $R_i$ is an integer polynomial. The word
\emph{simultaneous} means that all right-hand sides use the old state.

\begin{lemma}[Discrete antidifferentiation]
\label{cdc:pt:lem:antidifference}
For every $q(t)\in K[t]$ over a characteristic-zero field $K$, there is a
unique polynomial $S(t)\in K[t]$ satisfying
\[
 S(0)=0,\qquad S(t+1)-S(t)=q(t).
\]
It has degree at most $\deg q+1$ when $q\neq0$.
\end{lemma}
\begin{proof}
The polynomials $\binom tj$, $j\geq0$, form a basis of $K[t]$, since
they have distinct degrees and nonzero leading coefficients. Write
$q(t)=\sum_j\alpha_j\binom tj$ and set
$S(t)=\sum_j\alpha_j\binom t{j+1}$. Pascal's identity proves the
difference equation and $S(0)=0$. If two such polynomials existed, their
difference would have zero forward difference; a nonconstant polynomial
in characteristic zero cannot have that property, because taking a
difference lowers its degree by exactly one. Thus their difference is a
constant, and its value at zero is zero.
\end{proof}

\begin{theorem}[Nonlinear unitriangular substrates]
\label{cdc:pt:thm:triangular}
Every map \eqref{cdc:pt:eq:triangular} is polynomially iterable. Therefore its
exact guarded execution, first exit, halting domain, and partial output
graph admit the single-fold quartic representations of
Theorem~\ref{cdc:pt:thm:acceleration} and Corollary~\ref{cdc:pt:cor:first-exit}.
\end{theorem}
\begin{proof}
The first coordinate is $F_1(t,\boldsymbol x)=x_1+ct$. Suppose the
trajectories of the first $i-1$ coordinates have already been expressed as
rational polynomials in $(t,\boldsymbol x)$. Then
\[
 R_i(F_1(t,\boldsymbol x),\ldots,F_{i-1}(t,\boldsymbol x))
\]
is another such polynomial. Apply discrete antidifferentiation in the
variable $t$, treating the initial coordinates as coefficients, and add
$x_i$ to obtain $F_i$. The resulting polynomial satisfies exactly the
update recurrence and the initial condition. Induction on the integer time
then shows that its values coincide with the actual integer computation,
so no separate integer-valuedness assumption is needed. Continue through
all coordinates.
\end{proof}

Degree growth in this construction can be substantial. For example, if
$\deg_t F_j\leq D_j$ for $j<i$, the degree in $t$ of $F_i$ is bounded by
one plus the largest weighted degree
$\sum_{j<i}\alpha_jD_j$ among monomials
$x_1^{\alpha_1}\cdots x_{i-1}^{\alpha_{i-1}}$ of $R_i$, with the
initial constant added. The resulting certificate is fixed-dimensional
for a \emph{fixed} program, but the theorem does not promise a small
polynomial bound in the source-code length of an arbitrary nested
triangular expression.

\subsection{A certified eventually quasi-polynomial affine extension}
Polynomial trajectories exclude sign-alternating linear modes. A fixed
residue decomposition handles a wider class without introducing variable
exponentiation.

\begin{theorem}[Eventually quasi-polynomial affine acceleration]
\label{cdc:pt:thm:quasi}
Let $M$ be the augmented integer matrix of an affine update. Suppose fixed
integers $h\geq0$, $p\geq1$, and $r\geq1$ are supplied and satisfy
\begin{equation}
 M^h(M^p-I)^r=0.
 \label{cdc:pt:eq:quasi-certificate}
\end{equation}
For every fixed polynomial Boolean guard, exact guarded execution and
first exit have single-fold quartic representations of fixed dimension.
Their halting domains and partial output graphs are also single-fold.
\end{theorem}
\begin{proof}
For $0\leq s<p$ and $k\in\N$, all the matrices involved commute,
and the binomial expansion gives
\begin{equation}
 M^{h+s+pk}
 =\sum_{j=0}^{r-1}\binom kj M^{h+s}(M^p-I)^j.
 \label{cdc:pt:eq:quasi-form}
\end{equation}
Terms of exponent at least $r$ vanish by
\eqref{cdc:pt:eq:quasi-certificate}. Thus the trajectory is a polynomial in $k$
on each of the finitely many residue classes after the fixed prefix of
length $h$.

In a claimed run of length $n$, the prefix tests at times $i<h$ are
required only when $i<n$. In residue class $s$, the relevant times are
$h+s+pk<n$. Their number is the uniquely determined natural number
\[
 N_s=\left\lceil\frac{\max(n-h-s,0)}p\right\rceil.
\]
The graph of this expression has a canonical polynomial certificate:
compute the positive part by a sign test, divide its sum with $p-1$ by
the fixed positive integer $p$, and impose the unique Euclidean remainder
condition $0\leq u<p$. Comparisons are converted using the circuit lemma.
This determines $N_s$ and all its auxiliary values uniquely, including
when $N_s=0$.

Substituting \eqref{cdc:pt:eq:quasi-form} into the guard gives a polynomial
Boolean guard in $k$. Apply the bounded guard theorem on the prefix
$0\leq k<N_s$ for each $s$, with the zero-length convention already
proved. The terminal state is selected from the fixed prefix or the
unique residue decomposition of $n-h$. These are finite, deterministic
piecewise-polynomial computations; every inactive branch is evaluated on
uniquely prescribed arguments or explicitly padded by zero. Thus all
auxiliaries are unique. The circuit lemma again reduces the finite
combined system to a quartic. First exit and its projections follow by the
same uniqueness argument as before.
\end{proof}

Condition \eqref{cdc:pt:eq:quasi-certificate} is an exact integer matrix identity
that can be checked directly. It is sufficient; this paper does not require
an algorithm for finding optimal $h,p,r$ or a spectral classification.
Taking $p=2$ accommodates appropriate alternating modes, while nilpotent
transient modes can be absorbed into $h$. The period is fixed as part of
the substrate certificate, not an unbounded existential choice.

\subsection{Composition of deterministic phases}
\begin{proposition}[Single-fold composition of partial functions]
\label{cdc:pt:prop:composition}
For fixed integer or natural tuple domains $A,B,C$, if the graphs of
partial functions $f:A\rightharpoonup B$ and $g:B\rightharpoonup C$
have single-fold polynomial representations, then
so does the graph of $g\circ f$. It can be presented by a quartic after
introducing canonical circuit variables.
\end{proposition}
\begin{proof}
Use the conjunction of the two graph representations and existentially
quantify the intermediate value $b$. For fixed input $a$, functionality
of $f$ makes $b$ unique if it exists. The two auxiliary tuples are then
unique by hypothesis. Encode signed components of $b$ canonically.
Flatten the finite polynomial expressions into arithmetic gates and sum
squared quadratic residuals.
\end{proof}
A fixed sequence of terminating phases from our loop classes therefore
inherits single-foldness. This is not closure under arbitrary recursive
calls or a variable number of phases. A deterministic branch controlled
by a polynomial test can also be handled, provided every inactive
certificate tuple is explicitly normalized. Unrestricted nondeterministic
choice requires a separate multiplicity analysis.

\section{Canonical nontermination certificates and input bounds\srctag{02}}
\label{cdc:pt:sec:nontermination}
The exact-run theorem treats a supplied duration. For polynomial
trajectories, a canonical stabilization bound also handles all future
times without leaving ordinary polynomial arithmetic.

\begin{lemma}[An elementary canonical root bound]
\label{cdc:pt:lem:root-bound}
Let $q(t)=\sum_{j=0}^m a_jt^j$ be an integer polynomial, with zero leading
coefficients permitted in the displayed list. Set
\[
 B_q=2+\sum_{j=0}^m|a_j|.
\]
If $q$ is nonzero, it has no real root $t\geq B_q$. Its sign is constant
on $[B_q,\infty)$. The zero polynomial has constant zero sign there.
\end{lemma}
\begin{proof}
Constant polynomials require no argument. Otherwise discard zero leading
coefficients and let $a_m\neq0$, $m\geq1$. Since $a_m$ is a nonzero
integer, $|a_m|\geq1$. Put
$A=\max_{j<m}|a_j/a_m|\leq\sum_j|a_j|$.
For a real $t>1+A$, the finite geometric sum gives
\[
 \left|\sum_{j<m}a_jt^j\right|
 \leq |a_m| A\frac{t^m-1}{t-1}
 < |a_m|t^m.
\]
Thus the lower terms cannot cancel the leading term, and $q(t)$ has the
sign of $a_m$. Every $t\geq B_q$ satisfies the strict inequality needed
above. The extra constant $2$ removes any endpoint ambiguity.
\end{proof}

\begin{theorem}[Specified-input nontermination is single-fold]
\label{cdc:pt:thm:nontermination}
Let $f,F,G$ be as in Theorem~\ref{cdc:pt:thm:acceleration}. The relation
\[
 \Nonterm_f^G(\boldsymbol x)
 \quad\Longleftrightarrow\quad
 G(F(t,\boldsymbol x))\text{ holds for every }t\in\N
\]
has an effective single-fold quartic representation. Moreover, an explicit
canonical bound $B(\boldsymbol x)$ can be computed such that every
terminating input first exits at some time at most $B(\boldsymbol x)$.
\end{theorem}
\begin{proof}
After clearing fixed positive denominators as before, write the
polynomials of the guard along the trajectory as
\[
 q_j(t,\boldsymbol x)=\sum_{\ell=0}^{D_j}
 a_{j\ell}(\boldsymbol x)t^\ell,
 \qquad a_{j\ell}\in\Z[\boldsymbol x].
\]
Define the exact, not merely upper-bounded, number
\begin{equation}
 B(\boldsymbol x)=2+\sum_{j,\ell}|a_{j\ell}(\boldsymbol x)|.
 \label{cdc:pt:eq:canonical-B}
\end{equation}
For each specified input, Lemma~\ref{cdc:pt:lem:root-bound} makes every atom's
sign constant from time $B(\boldsymbol x)$ onward. The Boolean combination
$G$ is therefore constant from that time onward as well. It follows that
\begin{equation}
 \Nonterm_f^G(\boldsymbol x)
 \iff \forall\,0\leq t\leq B(\boldsymbol x)\;
 G(F(t,\boldsymbol x)).
 \label{cdc:pt:eq:nonterm-finite}
\end{equation}
Indeed, if the right-hand side holds, its truth at $B$ extends to the
entire tail; the converse is immediate. If the loop exits at all, some
false guard value must consequently occur no later than $B$.

The coefficient values and their absolute values are computed uniquely
using canonical signed pairs. Equation \eqref{cdc:pt:eq:canonical-B} determines
$B$ uniquely. Apply the bounded Boolean guard theorem to
\eqref{cdc:pt:eq:nonterm-finite}. Its primary tables and all auxiliary arithmetic
values are unique, so the complete nontermination certificate is
single-fold. Flattening the finite combined circuit gives a quartic.
\end{proof}

\begin{corollary}[Bit-efficient specified-input classification]
\label{cdc:pt:cor:input-complexity}
For a fixed polynomial trajectory and fixed polynomial Boolean guard,
termination or nontermination from a binary-encoded integer initial state
can be decided in polynomial time in that state's bit length. The
corresponding canonical certificates have polynomial bit length. When the
loop terminates, its first-exit time and output can be produced within the
same polynomial-time bound.
\end{corollary}
\begin{proof}
Let $H=\max(1,\|\boldsymbol x\|_\infty)$. Since the program and the
coefficient polynomials are fixed, there are fixed constants $C,K$ with
$B(\boldsymbol x)\leq C(H+1)^K$. Coefficient computation therefore uses
polynomial bit time and yields $O(\log(H+2))$-bit values, with constants
allowed to depend on the program.

Generate the fixed-degree sign towers on $[0,B]$ by
Proposition~\ref{cdc:pt:prop:complexity}. Their running time is polynomial in
$\log(H+2)$. Inspect the finitely many common sign-cell starts used in
Theorem~\ref{cdc:pt:thm:boolean}. If all guard values are true, the loop does not
terminate. Otherwise the smallest false cell start is exactly the first
exit time. Evaluate the fixed polynomial trajectory there to obtain the
output. The same magnitude estimates bound the complete certificate's
bit length.
\end{proof}

This is an acceleration statement about binary arithmetic, not a bound on
the literal number of loop iterations. A loop can run for a very large
number of steps and still have an efficiently produced compressed
certificate. The constants in the complexity bound depend on the fixed
trajectory and its degrees; this corollary is not a polynomial bound for
arbitrary programs presented as input.

\subsection{Why this does not decide global termination}
The quantifier over the initial state is crucial. For an arbitrary integer
polynomial $Q(\boldsymbol x)$, consider the identity update with guard
$Q(\boldsymbol x)=0$. There is a nonterminating initial state exactly when
$Q$ has an integer zero. With natural initial states, the same example
expresses natural solvability. Hence a uniform procedure deciding whether
\emph{every} initial state terminates for all such supplied loops would
decide a Hilbert-tenth solvability problem. Ordinary MRDP supplies the
undecidability background \cite{lf2022}.

Our theorem assigns a finite single-fold certificate to each specified
initial state. It does not decide the existential assertion that some
state has a nontermination certificate. That assertion is itself a
Diophantine solvability problem. Thus the polynomial-time specified-input
corollary is compatible with the broader termination distinctions in the
triangular-loop literature \cite{hfg2024,lmg2024}.

\section{Worked examples\srctag{02}}
\label{cdc:pt:sec:examples}

\subsection{A guard that fails only in the interior}
Let
\[
 p(t)=(t-2)(t-4)=t^2-6t+8,\qquad T=6.
\]
Its values are
\[
 (p(0),\ldots,p(6))=(8,3,0,-1,0,3,8).
\]
Both outer endpoint values are positive, but bounded nonnegativity is
false. The difference tower is
\[
 q_0=t^2-6t+8,\qquad q_1=2t-5,\qquad q_2=2.
\]
The actual active runs are shown in Table~\ref{cdc:pt:tab:runs}. Their tags are
one plus the signs shown in the table.

\begin{table}[H]
\centering
\begin{tabular}{cl}
\toprule
Level & Active half-open sign intervals on $\{0,\ldots,6\}$\\
\midrule
$q_0$ & $[0,2):+$, $[2,3):0$, $[3,4):-$, $[4,5):0$, $[5,7):+$\\
$q_1$ & $[0,3):-$, $[3,7):+$\\
$q_2$ & $[0,7):+$\\
\bottomrule
\end{tabular}
\caption{A unique short certificate exposing a negative interior value.}
\label{cdc:pt:tab:runs}
\end{table}

The nominal level capacities are $5,3,1$. Thus $q_1$ gets one inactive
slot $[7,7)$ with tag $1$, while the other levels need no padding. The
sign-table certificate is valid, but the additional positivity assertion
rejects the negative level-zero run.

For the loop $x'=x+1$ with guard $(x-2)(x-4)\geq0$, started at $x=0$,
the first exit is $n=3$, $x=3$. A procedure checking only the guard at the
start and at time $6$ would incorrectly infer that the entire prefix was
valid. The sign-table method does not make that inference.

\subsection{A nonlinear accumulation chain}
Consider the simultaneous update
\begin{equation}
 x'=x+1,\qquad y'=y+x,\qquad z'=z+y^2,
 \label{cdc:pt:eq:example-loop}
\end{equation}
with initial state $(a,b,c)\in\Z^3$. Its exact trajectory is
\begin{align}
 x_t&=a+t,\nonumber\\
 y_t&=b+at+\binom t2,\label{cdc:pt:eq:xyz-trajectory}\\
 z_t&=c+b^2t+(2ab+a^2)\binom t2\nonumber\\
 &\quad +(2a^2+2b+4a+1)\binom t3
 +(6a+6)\binom t4+6\binom t5.\nonumber
\end{align}
To verify the last formula, write $y_j=b+aj+\binom j2$ and use
\begin{align*}
 j^2&=\binom j1+2\binom j2,\\
 j\binom j2&=2\binom j2+3\binom j3,\\
 \binom j2^2&=\binom j2+6\binom j3+6\binom j4.
\end{align*}
Expanding $y_j^2$ in this basis gives the coefficients in
\eqref{cdc:pt:eq:xyz-trajectory} after summation with
$\sum_{j=0}^{t-1}\binom jk=\binom t{k+1}$. This is a symbolic proof;
the accompanying tests also compare the formula with $2{,}500$ explicitly
iterated states.

An example guard is
\[
 (z^2-3xy+\lambda\geq0)\ \wedge\ (y\neq\mu),
\]
where $\lambda,\mu$ are fixed parameters or additional unchanged input
coordinates. Along the trajectory, the first atom has degree at most ten
in $t$, and the second has degree at most two. The bounded Boolean theorem
therefore supplies a fixed-dimensional certificate for any duration. The
result does not depend on the first atom being monotone; it may pass
through several sign intervals.

\subsection{An enormous horizon with a small certificate}
The test suite includes
\[
 A=10^{50},\qquad T=10^{100},\qquad p(t)=(t-A)^2.
\]
The three difference levels have respectively three, two, and one active
sign runs. Their total active run count is six. The supplied generator
used $832$ distinct exact polynomial evaluations to obtain this
certificate; the horizon was not enumerated. The degree-two quartic
compiler then verified its generated natural witness against every
quadratic residual.

The numbers $10^{50}$ and $10^{100}$ are simply concrete test inputs,
not new primitive operations in the output equation. The compiled
polynomial is the same degree-two template used for small coefficients
and horizons; only its free parameter assignment changes.


% ---- manuscript 02, lines 1627-1646 ----
\section{Manuscript 02: Conclusion\srctag{02}}
A long computation need not be represented by a long explicit trace.
For polynomial trajectories, the signs of a guard can be certified by a
unique hierarchy of maximal sign intervals of successive finite
differences. The hierarchy is finite because degree decreases; its
correctness is local because differences control monotonicity; and its
uniqueness is global because all runs are maximal and all padding is fixed.

Combined with total canonical arithmetic circuits, this yields an
ordinary quartic equation with a unique natural witness and a dimension
independent of the run duration. It applies to genuinely nonlinear
unitriangular dynamics, not just monotone linear counters. First-exit
semantics permits safe projection of time, and a canonical polynomial
root bound yields unique specified-input nontermination certificates.

The construction gives a concrete addition to a research program in which
ordinary Diophantine reachability is already established: a
multiplicity-preserving accelerator with an executable polynomial output.
Its exact scope, proof obligations, and exponential boundary are explicit.


\part{Memory logs, RAM and graph evaluation}\label{cdc:part:memory}

\paragraph{Source and scope.} \wtag{} This Part is manuscript~03, ``Linear-Size Canonical Diophantine Certificates for Random-Access and Graph Computation'': its Sections 2--9 (canonical comparison and zero-test gadgets, a sorting network as a canonical polynomial circuit, the canonical random-access memory theorem, size and witness height, the linear-size compiler by exact large-base products, complete RAM executions, applications to counters, cellular automata, stacks, tapes and an immutable SK/SKI/Iota graph evaluator, and exact counting and first-halting probabilities), and its conclusion. Its conventions are in \cref{cdc:conv:sec}, its introduction in \cref{cdc:sec:manuscripts}, its boundary theorem in Part~IX.

\paragraph{Setting and hypotheses.} \wtag{} A log of length $L$ records reads and writes of a random-access memory with natural addresses and values; $N$ is the number of padded records and $C$ the number of comparators of the sorting network. ``Canonical certificate'' means a single-fold representation (\cref{cdc:conv:sec}); ``witness height'' bounds the magnitude of witness coordinates. The radix packing of \cref{cdc:mem:sec:linear} is the device that Part~VI uses for queue words (\cref{cdc:qc:lem:word-code}).

% ---- manuscript 03, lines 196-1254 ----
\section{Canonical arithmetic with quadratic residuals\srctag{03}}\label{cdc:mem:sec:gadgets}
\subsection{Comparison in two equations}
The elementary observation doing most of the work is a signed-gap identity.

\begin{lemma}[Canonical comparison]\label{cdc:mem:lem:compare}
For every $x,y\in\N$, the system in $b,d\in\N$
\begin{align}
 b(b-1)&=0,\label{cdc:mem:eq:cmp-bit}\\
 y-x-(2b-1)d-b+1&=0\label{cdc:mem:eq:cmp-gap}
\end{align}
has exactly one solution. It is
\begin{equation}\label{cdc:mem:eq:cmp-sol}
 (b,d)=
 \begin{cases}
 (1,y-x),&x\leq y,\\
 (0,x-y-1),&x>y.
 \end{cases}
\end{equation}
Thus $b$ is the indicator of $x\leq y$, with equality assigned to the first
branch, and $0\leq d\leq\max(x,y)$.
\end{lemma}
\begin{proof}
Equation~\eqref{cdc:mem:eq:cmp-bit} forces $b=0$ or $b=1$. If $b=1$,
\eqref{cdc:mem:eq:cmp-gap} becomes $y-x=d$, possible precisely when $x\leq y$, and
then $d=y-x$ is forced. If $b=0$, it becomes $x-y=d+1$, possible precisely
when $x>y$, and then $d=x-y-1$ is forced. The two cases are disjoint and
exhaustive. The stated bound follows in either case.
\end{proof}

The strict offset $1$ is important. Without it, a tie could choose either
branch. Equally importantly, there is only one slack coordinate: no unused
branch slack survives as an arbitrary witness. A naive disjunctive encoding
with two unnormalised slacks would generally have many auxiliary solutions.

\begin{example}
For $(x,y)=(3,7)$, the unique witness is $(b,d)=(1,4)$. For $(7,3)$ it is
$(0,3)$, not $(0,4)$, because the reverse branch uses a strict inequality.
For $(3,3)$ it is $(1,0)$. The formulas include zero without a separate
convention.
\end{example}

\subsection{Zero tests and conditional choice}
Putting $(x,y)=(\delta,0)$ in Lemma~\ref{cdc:mem:lem:compare} gives the following useful
form.
\begin{corollary}[Canonical zero test]\label{cdc:mem:cor:zero}
For $\delta\in\N$, the two equations
\begin{equation}\label{cdc:mem:eq:zero}
 z(z-1)=0,\qquad \delta+(2z-1)h+z-1=0
\end{equation}
have the unique natural solution
$(z,h)=(1,0)$ when $\delta=0$, and
$(z,h)=(0,\delta-1)$ when $\delta>0$.
\end{corollary}
The nonnegativity assumption on $\delta$ is part of this corollary. When we use
it on an address difference, that assumption will follow from the sorting
proof, not from an unproved premise in the polynomial system.

For a Boolean $b$ and natural inputs $x,y$, the residual
\begin{equation}\label{cdc:mem:eq:mux}
 u-bx-(1-b)y=0
\end{equation}
uniquely specifies a natural output. Addition, multiplication, copying, and
fixed constants are likewise represented by their defining equations. Boolean
AND is multiplication, NOT is $1-b$, and OR can be represented by
$u-b-c+bc=0$. The latter is nonnegative for Boolean inputs.

Natural truncated subtraction is also canonical without allowing a negative
intermediate. Compare $x$ with $y$ using $(b,d)$, and set
\begin{equation}
 u-(1-b)(d+1)=0.
\end{equation}
Then $u=\max(x-y,0)$. For fixed $B\geq2$, Euclidean division is specified by
\begin{equation}\label{cdc:mem:eq:division}
 x-Bq-r=0,\qquad r+h+1-B=0,
\end{equation}
with $q,r,h\in\N$. The second equation says $r<B$ and fixes $h$; division with
remainder proves uniqueness. No multiplicity is introduced by the bound.

\subsection{Composing canonical gadgets}
\begin{lemma}[Acyclic composition]\label{cdc:mem:lem:composition}
Suppose a finite directed acyclic circuit is built from gadgets whose output
and auxiliary coordinates are uniquely determined by their inputs. Use fresh
coordinates for different gadgets and identify only prescribed connecting
wires. Then the combined residual system has a unique assignment to every
internal coordinate for each external input. Adding predicates on the final
outputs either keeps that assignment or removes it; it never creates another.
\end{lemma}
\begin{proof}
Order the gadgets topologically. The first gadget has fixed inputs, so its
output and auxiliaries are forced. The same statement then holds for the next
gadget, and induction forces the entire assignment. Conversely, computing the
gadgets in that order constructs an assignment satisfying their residuals.
Additional final predicates merely test the already determined outputs.
\end{proof}

This elementary lemma is a stronger contract than closure of Diophantine
relations under existential quantification. It retains the size of each fiber.
It applies to acyclic certificate circuits, not to an arbitrary cyclic system
of equations. In the RAM theorem below, temporal memory correctness supplies a
separate induction for the apparent dependence of future states on read values.

\section{Deterministic sorting as a canonical polynomial circuit\srctag{03}}
\subsection{One four-field compare--exchange}
A record has four natural coordinates
\begin{equation}
 X=(\alpha,\kappa,w,v),
\end{equation}
where $\kappa$ is the sorting key. Given two records $X,Y$, apply
Lemma~\ref{cdc:mem:lem:compare} to their keys, obtaining $b,d$. For each field
$f\in\{\alpha,\kappa,w,v\}$ introduce outputs $\ell_f,u_f$ and require
\begin{align}
 \ell_f-bX_f-(1-b)Y_f&=0,\label{cdc:mem:eq:record-low}\\
 u_f-X_f-Y_f+\ell_f&=0.\label{cdc:mem:eq:record-high}
\end{align}
The first output record has the smaller key, with the left record first at a
tie; the second contains the other record. Descending comparators simply
interchange the two output ports. No additional variable is necessary for that
change of wiring.

\begin{proposition}[Comparator cost]\label{cdc:mem:prop:comparator-cost}
The four-field comparator introduces exactly ten natural variables and ten
residual equations, all of degree at most two. Its complete auxiliary and
output assignment is unique for every pair of input records. Its outputs are
exactly those records, not merely records with the same keys.
\end{proposition}
\begin{proof}
There are two variables and two equations for $(b,d)$, and eight output
variables with eight equations for the fields. The key comparison forces
$(b,d)$; \eqref{cdc:mem:eq:record-low} forces each lower field, and
\eqref{cdc:mem:eq:record-high} forces each upper field. In either Boolean case these
are the two original records in the required order.
\end{proof}

Consequently, a fixed sorting network with $C$ comparators becomes a canonical
arithmetic circuit with exactly $10C$ new coordinates and $10C$ residuals.
Unlike existential permutation routing, there is no choice of internal switch
settings once the input records are fixed.

\subsection{A completely explicit network}\label{cdc:mem:subsec:bitonic}
For $N=2^k$, use the iterative bitonic network whose complete generator is
\begin{lstlisting}
for block in 2, 4, ..., N:
    for stride in block/2, block/4, ..., 1:
        for i in 0, ..., N-1:
            j = i xor stride
            if j > i:
                compare wires i and j
                ascending exactly when (i bitwise-and block) == 0
\end{lstlisting}
All operations in this generator act on fixed wire indices at compilation time.
No bitwise arithmetic is needed inside the Diophantine equation.

For completeness, here is the mathematical sorting argument behind the
network. A finite sequence is bitonic when some cyclic rotation is
nondecreasing and then nonincreasing. Comparing positions separated by half the
length, and sending the smaller entries to the first half, is a bitonic split.
Both resulting halves are bitonic, and every entry of the first half is at most
every entry of the second half.

To verify this split property rigorously, threshold all entries at an arbitrary
value. A bitonic zero--one sequence has a single circular interval of ones. If
that interval has length at most half the total length, no opposite pair
contains two ones, so the lower half is all zero; the upper half has a circular
interval of ones. If the interval is longer, every opposite pair contains a
one, so the upper half is all one and the lower half again has a circular
interval of ones. This proves the split property for every threshold.
Comparison commutes with thresholding. Any failure of separation or bitonicity
in the original total order would appear at a threshold between offending
values. Thus the split property holds in the original order as well.

Recursively splitting each half sorts a bitonic sequence. To sort an arbitrary
sequence, sort one half ascending and the other descending, producing a bitonic
sequence, and merge. The displayed iterative generator performs precisely
these merges on successively doubled blocks, with the alternating block
orientations needed for the next merge. Its final block is ascending.

\begin{lemma}[Exact bitonic count]\label{cdc:mem:lem:bitonic-count}
For $N=2^k$, the displayed network sorts arbitrary totally ordered keys and has
\begin{equation}\label{cdc:mem:eq:bitonic-count}
 C_N=\frac{N}{4}k(k+1)
\end{equation}
comparators, with $C_1=0$.
\end{lemma}
\begin{proof}
Correctness was established above. At block size $2^r$ there are $r$ stride
layers, each containing $N/2$ disjoint comparisons. Therefore
$C_N=(N/2)\sum_{r=1}^k r=Nk(k+1)/4$.
\end{proof}

The proof also explains why duplicate keys cause no problem. The comparator's
left-biased tie convention makes its internal variables unique, while sorting
correctness only requires a nondecreasing output. In the memory application,
timestamps will in fact make all keys distinct.

\section{The canonical random-access memory theorem\srctag{03}}\label{cdc:mem:sec:memory}
\subsection{The semantic relation}
A finite chronological log is
\begin{equation}\label{cdc:mem:eq:log}
 E=((a_t,w_t,v_t))_{0\leq t<L},\qquad
 a_t,v_t\in\N,\quad w_t\in\{0,1\}.
\end{equation}
The memory is initially zero at every address. A write event ($w_t=1$) stores
$v_t$ at address $a_t$. A read event ($w_t=0$) asserts that the current value at
$a_t$ is $v_t$ and does not change memory. Write values are unrestricted natural
numbers; reads before any write must return zero. Let $\Valid_L(E)$ mean that
all events obey this semantics. Non-Boolean write flags make the log invalid.

We now construct one polynomial for each $L$, treating the $3L$ event
coordinates as free parameters. The polynomial's coefficients do not depend on
the numerical values of those parameters.

\subsection{Padding and chronological keys}
Set
\begin{equation}\label{cdc:mem:eq:padding}
 N=2^{\lceil\log_2\max(1,L)\rceil},\qquad B=N+1.
\end{equation}
For each genuine event introduce
\begin{equation}\label{cdc:mem:eq:keys}
 \alpha_t=a_t+1,\qquad \kappa_t=B\alpha_t+t,
\end{equation}
and the record $(\alpha_t,\kappa_t,w_t,v_t)$. The two definitions are enforced
by linear residuals; also include $w_t(w_t-1)=0$. For each $L\leq t<N$, append
the fixed dummy record
\begin{equation}\label{cdc:mem:eq:dummy}
 (\alpha_t,\kappa_t,w_t,v_t)=(0,t,1,0).
\end{equation}
These records are constants, not additional existential choices. They are zero
writes at a reserved shifted address. No genuine event has that shifted
address, so padding cannot affect the original memory.

\begin{lemma}[Key order]\label{cdc:mem:lem:key-order}
The keys of all $N$ records are distinct. Sorting by $\kappa$ first sorts by
shifted address $\alpha$ and, within an address, by original timestamp $t$.
All dummy records precede all genuine records.
\end{lemma}
\begin{proof}
All timestamps lie in $[0,N-1]$. If $\alpha<\alpha'$, then
\[
 B\alpha+t\leq B\alpha+N-1
 <B(\alpha+1)\leq B\alpha'+t'.
\]
At equal addresses, the key comparison is the timestamp comparison. Distinct
events have distinct timestamps, proving distinctness. Dummy addresses are
zero and genuine shifted addresses are positive.
\end{proof}

Apply any fixed sorting network on $N$ wires using the four-field comparator.
Write the sorted records as
\begin{equation}
 (\alpha_j^*,\kappa_j^*,w_j^*,v_j^*)\qquad(0\leq j<N).
\end{equation}
By the preceding lemmas, all these records and all internal sorting coordinates
are uniquely determined by $E$.

\subsection{Four equations per adjacent pair}
For the first sorted record require
\begin{equation}\label{cdc:mem:eq:first}
 (1-w_0^*)v_0^*=0.
\end{equation}
For each $1\leq j<N$, the difference
$\delta_j=\alpha_j^*-\alpha_{j-1}^*$ is nonnegative by
Lemma~\ref{cdc:mem:lem:key-order}. Introduce three natural variables $z_j,h_j,u_j$,
and the following four residuals:
\begin{align}
 z_j(z_j-1)&=0,\label{cdc:mem:eq:adj-bit}\\
 \alpha_j^*-\alpha_{j-1}^*+(2z_j-1)h_j+z_j-1&=0,\label{cdc:mem:eq:adj-gap}\\
 u_j-z_jv_{j-1}^*&=0,\label{cdc:mem:eq:adj-prev}\\
 (1-w_j^*)(v_j^*-u_j)&=0.\label{cdc:mem:eq:adj-read}
\end{align}
The first two are the canonical zero test of $\delta_j$. Consequently
$z_j=1$ exactly when the addresses agree; $h_j$ is the unique corresponding
gap witness. The third equation supplies the previous value when the address is
unchanged, and zero otherwise. The fourth checks a read but puts no condition
on the value of a write.

Introducing $u_j$ is necessary for our degree bound: writing
$(1-w_j^*)(v_j^*-z_jv_{j-1}^*)=0$ directly would create a cubic residual and a
sixth-degree square. The new coordinate is uniquely defined and therefore does
not spoil parsimony.

\begin{lemma}[Sorted memory criterion]\label{cdc:mem:lem:local-memory}
For sorted records arising from \eqref{cdc:mem:eq:keys}--\eqref{cdc:mem:eq:dummy},
\eqref{cdc:mem:eq:first}--\eqref{cdc:mem:eq:adj-read} have a natural solution exactly when the
original log is valid. When a solution exists, all their auxiliary coordinates
are unique.
\end{lemma}
\begin{proof}
Consider a maximal consecutive block of one shifted address. At its first
record, the value supplied as the predecessor is zero: this follows either
from \eqref{cdc:mem:eq:first}, or from $z_j=0$ and \eqref{cdc:mem:eq:adj-prev}. Thus its first
read must return zero, while its first write may set any value.

At every subsequent record of that block, $z_j=1$ and $u_j=v_{j-1}^*$. A read
is therefore required to agree with the preceding record's value; a write may
replace it. Inductively, the value in every record is the memory contents just
after that event. A read leaves that contents unchanged, and a write changes
it to its stated value. Because timestamp order is preserved within the block,
this is precisely the chronological memory semantics at that address.

Different addresses do not interact, and dummy zero writes are harmless. This
proves equivalence with the original log. For uniqueness, the sorted addresses
and values fix each $z_j,h_j$ by Corollary~\ref{cdc:mem:cor:zero}, then fix $u_j$ by
\eqref{cdc:mem:eq:adj-prev}. The read equations are only tests of those forced values.
\end{proof}

\subsection{The main theorem and its exact constants}
\begin{theorem}[Canonical quartic memory compiler]\label{cdc:mem:thm:memory}
Fix $L\geq0$ and $N$ as in \eqref{cdc:mem:eq:padding}. Let a sorting network on $N$
wires have $C$ comparators. One can explicitly construct
\[
 P_L\in\Z[a_0,w_0,v_0,\ldots,a_{L-1},w_{L-1},v_{L-1},\boldsymbol\eta]
\]
of degree at most four such that, for every natural event-parameter tuple $E$,
\begin{equation}\label{cdc:mem:eq:main-equivalence}
 \#\{\boldsymbol\eta\in\N^{W}:P_L(E,\boldsymbol\eta)=0\}
 =\begin{cases}1,&\Valid_L(E),\\0,&\text{otherwise}.
 \end{cases}
\end{equation}
The construction has exactly
\begin{align}
 W&=10C+2L+3(N-1),\label{cdc:mem:eq:witness-count}\\
 R&=10C+3L+4(N-1)+1\label{cdc:mem:eq:residual-count}
\end{align}
auxiliary variables and residuals, before omission of identically zero
residuals. It uses no exponentiation, divisibility predicate, bounded universal
quantifier, or appeal to MRDP.
\end{theorem}
\begin{proof}
Form the input residuals in \eqref{cdc:mem:eq:keys} and the write-bit residuals, the
$C$ comparator systems, and the final memory residuals. Let $P_L$ be their sum
of squares. Every residual is quadratic or linear, proving the degree bound.

If $P_L=0$, every residual vanishes. Write flags are Boolean. The input
coordinates force the shifted addresses and keys; the comparator lemma and
network correctness force the sorted permutation of complete records. The
sorted memory criterion then proves $\Valid_L(E)$. This proves soundness.

If the log is valid, compute the shifted records, run the deterministic
network, and choose every comparator and adjacent-pair variable by its explicit
formula. All residuals vanish by Lemma~\ref{cdc:mem:lem:local-memory}, giving a root.
This proves completeness. The same sequence of forcing arguments shows that
any root must be the one just constructed, proving uniqueness.

There are $2L$ input auxiliaries and $3L$ input residuals. Each comparator
contributes ten of each. Each of the $N-1$ adjacent pairs contributes three
auxiliaries and four residuals, while the first-record check contributes one
residual and no auxiliary. Adding these quantities gives
\eqref{cdc:mem:eq:witness-count}--\eqref{cdc:mem:eq:residual-count}.
\end{proof}

\begin{remark}[The empty log]
For $L=0$, take $N=1$. The sole record is the constant dummy record $(0,0,1,0)$.
There are no witnesses and just the identically zero first-record residual.
The empty tuple is the unique witness in $\N^0$. Thus the theorem includes the
empty computation without an exceptional existential parameter.
\end{remark}

\begin{remark}[Natural versus unrestricted integer witnesses]
The theorem does not assert uniqueness, or even the same relation, when the
witnesses are allowed to be arbitrary integers. A negative signed-gap witness
would destroy the meaning of comparison. Replacing natural variables by sums
of four integer squares restores existential representability, but generally
creates many representations and destroys the stated bijection. A formal
implementation must retain the natural-number domain explicitly.
\end{remark}

\section{Size, height, and the cost of canonical routing\srctag{03}}\label{cdc:mem:sec:resources}
\subsection{Explicit and asymptotic size}
Substituting \eqref{cdc:mem:eq:bitonic-count} into Theorem~\ref{cdc:mem:thm:memory} gives
\begin{align}
 W_L&=\frac{5N}{2}k(k+1)+2L+3(N-1),\label{cdc:mem:eq:explicit-w}\\
 R_L&=\frac{5N}{2}k(k+1)+3L+4(N-1)+1.\label{cdc:mem:eq:explicit-r}
\end{align}
These are scalar-coordinate and scalar-equation counts. For $L\geq1$ we have
$L\leq N<2L$ unless $L=1$, where $N=L$. Thus the explicit construction is
$O(L\log^2(L+2))$. For any sorting network with $C$ comparators the simpler
bounds $W_L\leq10C+5N$ and $R_L\leq10C+7N$ hold.

The theorem of Ajtai, Koml\'os, and Szemer\'edi supplies sorting networks of size
$C=O(N\log N)$ \cite{aks}. Substituting those networks in the \emph{same}
canonical comparator construction proves the following statement.
\begin{corollary}\label{cdc:mem:cor:optimal-size}
There is an effectively constructible family of quartic natural-number
polynomials representing valid length-$L$ memory logs with exactly one
auxiliary witness and $O(L\log(L+2))$ auxiliary coordinates and residuals.
\end{corollary}
This is an asymptotic existence/construction result using a known sorting
network theorem, not a practical improvement implemented in the companion
code. The simple bitonic network has transparent constants and a short checker.

Every displayed memory residual has at most eight monomials after collecting
terms. Therefore the expansion of $P_L$ has at most $64R_L$ monomial
\emph{contributions}, and hence at most that many distinct nonzero monomials.
The generous constant includes repeated or constant input wires. Expanding the
sum of squares does not entail an exponential expression-size explosion.
The residual coefficients have absolute value at most $2(N+1)$; a safe bound for
the absolute coefficients of the expanded polynomial is
$256R_L(N+1)^2$. Variable names or indices themselves require logarithmically
many bits in a serialized description. None of these statements identifies a
scalar count with the bit length of a complete file.

\subsection{Witness height}
\begin{theorem}[Height bound]\label{cdc:mem:thm:height}
For a valid log let
\[
 A=\max\bigl(\{a_t:0\leq t<L\}\cup\{0\}\bigr),\qquad
 V=\max\bigl(\{v_t:0\leq t<L\}\cup\{0\}\bigr).
\]
Every coordinate of its unique auxiliary witness is at most
\begin{equation}\label{cdc:mem:eq:height}
 H=\max\{1,V,(N+1)(A+1)+N-1\}.
\end{equation}
Consequently, if the addresses and values have at most $b$ binary digits, every
auxiliary coordinate has $b+O(\log(N+2))$ digits, with an absolute implied
constant.
\end{theorem}
\begin{proof}
Shifted addresses are at most $A+1$ and keys at most
$(N+1)(A+1)+N-1$. Every sorting output is a field of an input record, so it
obeys the same bound, as do Boolean flags. The comparator gap is bounded by
the larger input key by Lemma~\ref{cdc:mem:lem:compare}. In an adjacent-pair test,
$\alpha_j^*-\alpha_{j-1}^*\leq A+1$; its gap-minus-one coordinate is at most
$A$, and its previous-value coordinate is either zero or a log value.
These cover all auxiliary coordinates. Taking binary lengths proves the last
claim, also when the maximum input value is zero.
\end{proof}

The witness can be generated by the very comparisons and substitutions used
in the proof. In particular, its generation and checking are polynomial-time
in the input log's bit length and the displayed certificate size. Address
\emph{magnitude} does not enter the number of coordinates; address \emph{bit
length} necessarily enters the cost of reading and manipulating a coordinate.

This distinction becomes essential when certificates describe arithmetic
programs. The $T$-step program
\[
 x_0=2,\qquad x_{t+1}=x_t^2
\]
has $x_T=2^{2^T}$. Its local quadratic description has $O(T)$ coordinates, but
the last coordinate needs $2^T+1$ binary digits. Thus a small arithmetic
certificate, even a unique one, does not by itself supply an NP witness of
polynomial bit length. Memory-log height bounds and instruction-level growth
bounds must be stated separately.

\subsection{What is optimal, and what is not}
\begin{proposition}[Comparison-network lower bound]\label{cdc:mem:prop:lower}
An oblivious comparison network that sorts arbitrary $N$ distinct keys has at
least $\lceil\log_2(N!)\rceil=\Omega(N\log N)$ comparators. In particular, the
comparison order in Corollary~\ref{cdc:mem:cor:optimal-size} is optimal \emph{within
this sorting-network architecture}.
\end{proposition}
\begin{proof}
Give the input positions every possible permutation of $N$ distinct ranks.
The sorting process must route those positions into $N!$ different output
permutations. A sequence of $C$ binary comparison outcomes determines at most
$2^C$ routing patterns. Therefore $2^C\geq N!$. For example,
$N!\geq(N/2)^{N/2}$ for $N\geq2$, giving the asymptotic bound.
The memory-key format does not evade this argument: choosing distinct
addresses in an arbitrary order realizes every rank permutation of the keys.
\end{proof}
This is not a lower bound for arbitrary Diophantine encodings. A different
certificate architecture could avoid explicit comparison sorting. Nor does it
say that Batcher's extra logarithmic factor is necessary.

\section{A linear-size alternative: exact large-base products\srctag{03}}\label{cdc:mem:sec:linear}
The comparison lower bound is not a lower bound on Diophantine certificates.
We now remove the sorting network entirely. The key is to certify a proposed
sorted list by one exact product identity, at a uniquely determined integer
base large enough to prevent every possible carry collision. Unlike a random
fingerprint over a finite field, this check has no soundness error. It exchanges
small witness height for a linear number of scalar coordinates.

\subsection{A deterministic multiset lemma}
\begin{lemma}[Carry-free product fingerprint]\label{cdc:mem:lem:fingerprint}
Let $C\geq2$, $L\geq1$, $\beta=C^L+1$, and let
$c_1,\ldots,c_L,d_1,\ldots,d_L$ belong to $\{0,\ldots,C-1\}$. Then
\begin{equation}\label{cdc:mem:eq:fingerprint}
 \prod_{i=1}^L(\beta+c_i)=\prod_{i=1}^L(\beta+d_i)
\end{equation}
holds if and only if the two lists have the same multiset of entries.
\end{lemma}
\begin{proof}
Form $f(X)=\prod_i(X+c_i)$ and $g(X)=\prod_i(X+d_i)$ in $\Z[X]$.
Their coefficients are nonnegative, and the sum of the coefficients of $f$ is
\[
 f(1)=\prod_i(1+c_i)\leq C^L<\beta.
\]
Thus every coefficient of $f$ is a base-$\beta$ digit; the same holds for $g$.
Equality $f(\beta)=g(\beta)$ forces equality of all digits, by repeated Euclidean
division by $\beta$. Hence $f=g$ as polynomials.

To recover the multiset, evaluate at $X=-c_1$. Some factor $-c_1+d_j$ must be
zero, since a product of integers is zero only if a factor is zero. Cancel
$X+c_1$ from the two polynomial factorizations and proceed by induction.
Cancellation is valid in the integral domain $\Z[X]$. This argument includes
repeated entries and zero. Conversely, equal multisets plainly give equal
products.
\end{proof}

The carry bound, rather than any unproved collision-resistance assumption, is
the entire reason this evaluation at one point is exact. Evaluating at a small
fixed integer without the bound would not justify the conclusion. The elementary
base-expansion argument is not itself claimed to be a new algebraic principle;
the point is its fully canonical use in the memory compiler.

\subsection{Packing complete records}
For $L\geq1$, use the genuine records
\[
 X_t=(\alpha_t,\kappa_t,w_t,v_t),\quad
 \alpha_t=a_t+1,\quad\kappa_t=(L+1)\alpha_t+t\quad(0\leq t<L).
\]
No padding is needed. Define uniquely
\begin{equation}\label{cdc:mem:eq:linear-radix}
 p_0=0,\quad p_{t+1}=p_t+\alpha_t+\kappa_t+w_t+v_t,\qquad D=p_L+2.
\end{equation}
Each input field is strictly less than $D$. Introduce $L$ proposed output
records $X'_j=(\alpha'_j,\kappa'_j,w'_j,v'_j)$ as witness coordinates. For each
of their four fields $f$ add a natural slack $s$ and the equation
\begin{equation}\label{cdc:mem:eq:digit-bound}
 f+s+1-D=0.
\end{equation}
This both asserts $f<D$ and fixes the bound witness uniquely.

The radix-$D$ code of a record is
\begin{equation}\label{cdc:mem:eq:record-code}
 \rho(X)=\alpha+D\kappa+D^2w+D^3v.
\end{equation}
For four digits in $[0,D-1]$, the code belongs to $[0,D^4-1]$ and is injective
in the complete tuple, by uniqueness of base-$D$ expansion. Represent it with
three quadratic gates, not the high-degree expression in
\eqref{cdc:mem:eq:record-code}:
\begin{equation}\label{cdc:mem:eq:horner}
 h_1-vD-w=0,\quad h_2-h_1D-\kappa=0,\quad \rho-h_2D-\alpha=0.
\end{equation}
Do this for both input and proposed output records.

Introduce $D_2,C,q_1,\ldots,q_L,\beta$ and impose
\begin{equation}\label{cdc:mem:eq:large-base}
 D_2=D^2,\quad C=D_2^2,\quad q_0=1,\quad
 q_i=q_{i-1}C\ (1\leq i\leq L),\quad\beta=q_L+1.
\end{equation}
These quadratic equations force $C=D^4$ and $\beta=C^L+1$ without putting
exponentiation into the polynomial language. The chain length depends on the
fixed parameter $L$, not an exponent variable. Binary powering could shorten
the chain, but the displayed version gives particularly simple exact counts.

Let $A_0=B_0=1$ be constants. The two product chains are
\begin{equation}\label{cdc:mem:eq:product-chains}
 A_{i+1}=A_i(\beta+\rho(X_i)),\qquad
 B_{i+1}=B_i(\beta+\rho(X'_i))\quad(0\leq i<L).
\end{equation}
Impose $A_L-B_L=0$. Every equation remains quadratic. Finally require the
proposed keys to be nondecreasing by introducing one natural gap per adjacent
pair:
\begin{equation}\label{cdc:mem:eq:ordered-keys}
 \kappa'_j-\kappa'_{j-1}-g_j=0\qquad(1\leq j<L).
\end{equation}
The product lemma and radix injectivity force the proposed records to be a
permutation of the complete input records. The keys are distinct, so the
nondecreasing order fixes that permutation uniquely. All gap variables are
then uniquely determined.

Attach the first-record and adjacent-record memory equations from
Section~\ref{cdc:mem:sec:memory}, now on these $L$ sorted records. Their address
differences are nonnegative by the same key-order argument, and the previous
proof applies unchanged.

\subsection{The linear compiler theorem}
\begin{theorem}[Linear-size canonical quartic memory]\label{cdc:mem:thm:linear-memory}
For every $L\geq0$ there is an explicit polynomial
\[
 P^{\mathrm{lin}}_L\in\Z[X_1,\ldots,X_{3L},Y_1,\ldots,Y_{24L}]
\]
of degree at most four such that each valid length-$L$ memory log has exactly
one natural witness, and each invalid log has none. It is the sum of squares
of exactly $22L+1$ quadratic residuals, each with at most eight monomials.
Its expanded sparse polynomial therefore has $O(L)$ monomial contributions.
The polynomial and its coefficients depend on $L$, not on the numerical log.
\end{theorem}
\begin{proof}
For $L=0$, take no witnesses and the single zero residual. Suppose $L\geq1$.
The input definitions and write-bit constraints force precisely the declared
records and Boolean write flags. The prefix chain uniquely fixes $D$, so the
input fields satisfy the digit bounds. The explicit output bound equations
force every proposed output field into the same digit range. The Horner and
power chains fix all codes and $\beta$. The product chains and their final
equality, by Lemma~\ref{cdc:mem:lem:fingerprint}, force equality of code multisets.
By radix injectivity this is equality of complete record multisets. Distinct
keys and \eqref{cdc:mem:eq:ordered-keys} therefore force the single correctly sorted
output list. The sorted memory criterion proves soundness.

For a valid log, take that sorted list, use its uniquely determined digit
slacks, and compute every chain as displayed. Multiset equality verifies the
product identity, and memory validity verifies the final checks. This gives a
root. Conversely, in any root the preceding argument forces the output list;
all remaining coordinates are then forced by their defining equations, bound
slacks, key gaps, and canonical memory tests. Thus there is exactly one root.

The complete accounting is as follows. ``Output fields'' are existential
coordinates whose semantics are forced globally by the product and order
checks; they do not need separate defining residuals.
\begin{center}
\begin{tabular}{@{}lrr@{}}
\toprule
Component & Witnesses & Residuals\\
\midrule
Input shifted addresses, keys, write bits & $2L$ & $3L$\\
Prefix sums and $D$ & $L+1$ & $L+1$\\
Proposed output fields & $4L$ & $0$\\
Output digit-bound slacks & $4L$ & $4L$\\
Input and output Horner chains & $6L$ & $6L$\\
$D_2,C,q_1,\ldots,q_L,\beta$ & $L+3$ & $L+3$\\
Two product chains and final equality & $2L$ & $2L+1$\\
Adjacent key gaps & $L-1$ & $L-1$\\
Memory checks & $3(L-1)$ & $4(L-1)+1$\\
\midrule
Total & $24L$ & $22L+1$\\
\bottomrule
\end{tabular}
\end{center}
Every residual has degree at most two and a bounded number of terms. The
sum-of-squares argument gives the degree and sparse-size claims. All generated
constants are fixed small integers or functions of $L$ such as timestamps and
$L+1$; the large base itself is a witness fixed by quadratic equations.
\end{proof}

\begin{theorem}[Height cost of linear size]\label{cdc:mem:thm:linear-height}
For a valid nonempty log, let $D$ have the value in
\eqref{cdc:mem:eq:linear-radix}. Every witness of the linear compiler is at most
\begin{equation}\label{cdc:mem:eq:linear-height}
 H_{\mathrm{lin}}=2^L D^{4L^2}.
\end{equation}
If the input addresses and values have at most $b$ bits, each witness has
$O(L^2(b+\log(L+2)))$ bits. Generation and verification of the full certificate
are polynomial-time in the bit length of the finite log and in $L$.
\end{theorem}
\begin{proof}
All fields, digit slacks, prefix sums, key gaps, and local memory coordinates
are at most $D$. Horner coordinates are below $C=D^4$; the powers are at most
$C^L$, and $\beta=C^L+1$. Each factor in a product chain is at most
$C^L+C\leq2C^L$. Every prefix product is therefore at most
$(2C^L)^L=2^LD^{4L^2}$, which also dominates the other coordinates. Further,
$D\leq2+4LM$, where $M$ is the largest input-record field, and
$\log_2 M=O(b+\log(L+2))$. Taking logarithms gives the bit-length bound.
The chains have linear length, and all their integer operands have the
polynomial bit lengths just bounded. Exact addition, multiplication, comparison,
and division by fixed bases can thus generate and verify the witness in
polynomial time.
\end{proof}

There is a genuine tradeoff, not a free improvement in every resource:
\begin{center}
\begin{tabularx}{\textwidth}{@{}p{0.25\textwidth}p{0.27\textwidth}X@{}}
\toprule
Backend & Scalar witnesses & Maximum witness bit length\\
\midrule
Explicit bitonic network & $O(L\log^2(L+2))$ & $O(b+\log(L+2))$\\
Optimal sorting network & $O(L\log(L+2))$ & $O(b+\log(L+2))$\\
Exact large-base products & $24L$ & $O(L^2(b+\log(L+2)))$\\
\bottomrule
\end{tabularx}
\end{center}
The linear construction is not a bounded-word online memory checker or a
finite-field proof system, and makes no claim to evade lower bounds for those
different models. It is an offline exact natural-number certificate whose
large coordinates encode the entire multiset check without carries.
For $L=4$ the simple linear version has 96 witnesses versus the network's 77;
for $L=16$ it has 384 versus 877. Thus even the scalar improvement is asymptotic,
not a claim of the smallest constant for every short trace.

\begin{corollary}[Two-backend compilation principle]\label{cdc:mem:cor:two-backends}
Any composition argument using only the memory checker's existence,
uniqueness, and decoding contracts may substitute the linear backend for the
network backend. The semantic bijection is unchanged. Only the coordinate
count and the height estimate change.
\end{corollary}
\begin{proof}
Both theorems characterize the identical relation on the same chronological
log, with a singleton fiber on every valid log. Replace the unique witness of
one checker by the unique witness of the other, leaving the semantic log and
all other coordinates unchanged.
\end{proof}

\section{From memory logs to complete RAM executions\srctag{03}}\label{cdc:mem:sec:ram}
\subsection{A precise local interface}
A \emph{normalized RAM} has a fixed finite tuple of natural registers and a
finite control state. One time step has an explicit choice label $c$ from a
fixed finite alphabet (a singleton for a deterministic machine). It computes,
from its current local state $s$ and $c$, a request
\[
 (a,w,d)=\mathsf{request}(s,c),\qquad w\in\{0,1\}.
\]
The request is computed \emph{before} the response is used. Here $d=0$ on a
read and $d$ is the value to store on a write. Let $r$ be the response,
canonically zero on a write. The event submitted to the memory checker is
\begin{equation}\label{cdc:mem:eq:request-event}
 (a,w,v),\qquad v=d+r,
\end{equation}
and the next local state is $\mathsf{update}(s,c,r)$. Include the residuals
\begin{equation}\label{cdc:mem:eq:request-normalization}
 (1-w)d=0,\qquad wr=0,\qquad v-d-r=0.
\end{equation}
On a read the memory checker forces $r$ to be the old value; on a write the
local normalization forces $r=0$. This convention removes an otherwise unused
response variable on write steps.

The functions $\mathsf{request}$ and $\mathsf{update}$ must have canonical
acyclic quadratic implementations, with a bounded number of coordinates per
step for the fixed program. Addition, multiplication, comparison, truncated
subtraction, Boolean logic, and division by a fixed positive base satisfy this
requirement by Section~\ref{cdc:mem:sec:gadgets}. A fixed finite control table is
compiled by equality tests and Boolean selection among these operations.
Compute every branch using total operations and select the chosen outputs;
then even an inactive branch has uniquely determined internal data. An
alternative guarded encoding must explicitly zero or otherwise normalize
inactive coordinates. Merely multiplying all its constraints by a branch bit
would leave arbitrary witnesses.

A fixed finite tuple of possible branch labels can be validated by comparing
$c$ with its allowed labels and testing the resulting Boolean disjunction.
Exactly one label is retained after halting. Conditional branches of a
deterministic program are computed by comparisons, not represented by free
choice labels.

We reserve physical address zero for a no-operation event $(0,1,0)$. All user
addresses are shifted by one before they enter this physical interface. A
local step not needing memory uses that event. This reserve is distinct from,
and compatible with, the memory compiler's own shifted dummy records.
A finitely supported initial memory is supplied by $I$ prescribed initialization
writes. Their addresses and values are fixed by the input, not arbitrary
witness choices. The initial local state is likewise fixed by the input.

\subsection{Halting without padding multiplicity}
After the first halt, leave the local state unchanged, force the choice label
to zero, set the response to zero, and submit the reserved no-operation write.
Thus a run that halts before the horizon has one, not many, padded extensions.
The accepting and rejecting halted states are absorbing. Acceptance by time
$T$ is tested on the final control state after $T$ steps. A run already halted
at time zero is included.

\begin{theorem}[Parsimonious RAM compilation]\label{cdc:mem:thm:ram}
Fix a normalized RAM program as above and a horizon $T$. There is an effectively
constructed integer polynomial $Q_{P,T}(x,Y)$ of degree at most four whose
natural-number solutions $Y$, for each fixed well-formed input $x$, are in
bijection with the labelled executions of $P$ that accept by time $T$.
The bijection is projection to the chronological local states and choice
labels; its inverse is canonical certificate generation.

If the initialization has $I$ writes, the linear backend gives
$O_P(T+I)$ coordinates and residuals. The bitonic construction uses
$O_P((T+I)\log^2(T+I+2))$ coordinates and quadratic residuals, counting the
trace coordinates as witnesses. An optimal sorting network improves this to
$O_P((T+I)\log(T+I+2))$. In particular, a deterministic program has either no
root or exactly one root at each fixed input and horizon.
\end{theorem}
\begin{proof}
Introduce local state, label, request, and response coordinates for each step.
Add the canonical request and update circuits, normalization
\eqref{cdc:mem:eq:request-normalization}, the initial-state constraints, and the final
acceptance constraint. Join the $I$ initialization events and the $T$ step
events into one chronological log, and attach the memory checker of
Theorem~\ref{cdc:mem:thm:memory}. Use fresh internal coordinates for every circuit and
for the checker. The sum of the squares of all residuals has degree at most
four.

Consider any root. The local circuits describe valid instruction steps for the
values supplied to them. The memory theorem proves that the entire event log
obeys sequential zero-initialized memory semantics, including initialization.
Start at the fixed initial state and proceed in chronological order. The
choice label and current state determine the request before any read. On a
read, sequential memory fixes the response uniquely; on a write, the response
is zero by \eqref{cdc:mem:eq:request-normalization}. The update circuit then fixes the
next state. This proves that the projected data form a genuine labelled run,
not a self-consistent but causally impossible circular solution. Its final
control state is accepting.

Conversely, take any labelled accepting run. Extend it by the prescribed
absorbing steps to length $T$. Its chronological events form a valid memory
log. Assign every local auxiliary using its canonical circuit, then the unique
memory witness. All residuals vanish.

Finally, fix the semantic run. Its local coordinates and event values are
fixed. Lemma~\ref{cdc:mem:lem:composition} fixes each local circuit's auxiliaries, and
Theorem~\ref{cdc:mem:thm:memory} fixes the memory auxiliaries. Therefore the constructed
root is the only one projecting to that run. Distinct labelled runs project
differently, proving the bijection. There are $O_P(T+I)$ local coordinates and
residuals, in addition to the memory checker on $T+I$ events. The size estimates
follow; substituting Theorem~\ref{cdc:mem:thm:linear-memory} gives the linear bound.
A deterministic program has a unique length-$T$ padded run, so its
accepting root set has size zero or one.
\end{proof}

The instruction interface is a hypothesis with a concrete implementation
method, not an appeal to MRDP for arbitrary local predicates. For a family of
programs the constants depend on their finite instruction circuits. A fixed
interpreter can instead make program data part of the initialized memory, but
one must count that data and its interpreted runtime.

\begin{corollary}[Guarded resource restrictions]
Imposing additional Boolean tests on the uniquely decoded trace preserves
parsimony. Imposing a coordinate bound $y\leq H$ using
$y+z-H=0$, $z\in\N$, also preserves it, because the new slack is forced.
\end{corollary}
\begin{proof}
The tests can only remove decoded runs, and the bound slack is necessarily
$z=H-y$. Neither can split a surviving root into several roots.
\end{proof}

\section{Applications to computational substrates\srctag{03}}\label{cdc:mem:sec:substrates}
\subsection{Counter machines and bounded cellular computation}
For a fixed counter machine, increment is linear and a zero/decrement branch
uses the canonical zero test and truncated subtraction. Control is finite.
There is no random-access memory consistency problem if all counters are
explicit local registers. Unrolling $T$ steps therefore yields $O_P(T)$
quadratic residuals and a parsimonious quartic equation. The memory theorem is
valuable precisely when replacing a large register array by actual indirect
addressing; it is not a necessary overhead for every substrate.

A finite-alphabet cellular automaton has a similarly direct description. Fix a
finite space--time region, its boundary convention (for example a periodic
finite lattice), and an initial configuration. Use one-hot Boolean coordinates
for each cell state. Enforce their sum to be one. For each output symbol,
express its local transition rule as the sum of the products of the appropriate
input indicators. The neighborhood size and alphabet are fixed, so each product
has a fixed-length chain of quadratic multiplication gates. The products and
sums are unique on valid one-hot inputs. Temporal induction fixes the next
layer. Thus a region with $V$ cell updates has an $O(V)$ canonical quartic
certificate. For a nondeterministic rule, retain an explicit permitted local
choice label; solutions then count labelled space--time histories, not only
distinct final configurations.

This construction does not represent an infinite configuration by one finite
list without specifying a coding. Infinite-lattice claims require a finite
causal cone or an explicit finite-support convention.

\subsection{Stacks and Turing tapes}
Let the stack alphabet have $b$ symbols and set $B=b+1$. Encode symbols by
$1,\ldots,b$, with zero representing the empty stack. Pushing digit $d$ onto
stack code $n$ produces $Bn+d$. Popping a nonempty valid stack uses the unique
quotient and remainder in \eqref{cdc:mem:eq:division}; its remainder is the top symbol
and its quotient the residual stack. A zero test handles the empty stack with
a fixed blank-symbol convention. Nonzero digits avoid leading-zero ambiguity.
When checking externally supplied stack codes, validate their finite symbol
sequence or assume a well-formed initial encoding; a random natural number need
not encode a valid stack in this scheme.

Two stacks and a scanned symbol describe a Turing tape: moving right pushes the
old scanned symbol onto the left stack and pops the right stack, with blank on
an empty stack; moving left is symmetric. Therefore a fixed Turing machine has
a canonical bounded-horizon quartic description by the local-gadget method.
After $T$ steps only $O(T)$ symbols beyond the input can have been visited, so
these radix codes have $O(T+n)$ digits for input length $n$. Alternatively,
representing stacks and tape cells in RAM uses the memory theorem without
large radix codes. These are two resource tradeoffs, not distinct proofs of
computability equivalence.

\subsection{An immutable graph evaluator for SK, SKI, and Iota}
\label{cdc:mem:subsec:graph}
The following concrete substrate is close to ProveIt's combinatory-logic work,
but its evaluation strategy must be specified. An immutable directed acyclic
graph has leaf nodes $S,K,I,\iota$ and binary application nodes
$\app(p,q)$. A fresh application node points only to already allocated nodes.
Number the initial nodes densely, including the four reserved leaf nodes,
and allocate subsequent nodes at successive fresh identifiers. This is the
input representation used in the bit bounds below. A runtime state
consists of a focus pointer and a stack of argument pointers. The stack's first
item is the next argument to apply.

If the focus is $\app(p,q)$, push $q$ and set the focus to $p$. For a leaf with
sufficient arguments, perform one of the following actions, leaving any
remaining argument stack unchanged:
\begin{align*}
 I;[x,\ldots]&\longmapsto x;[\ldots],\\
 K;[x,y,\ldots]&\longmapsto x;[\ldots],\\
 S;[x,y,z,\ldots]&\longmapsto
    \app(\app(x,z),\app(y,z));[\ldots],\\
 \iota;[x,\ldots]&\longmapsto \app(\app(x,S),K);[\ldots].
\end{align*}
The $S$ action allocates three application nodes, and the Iota action two. If
there are too few arguments, the machine is in weak-head normal form and halts.
There are no reductions inside arguments before they become a focus.

\begin{proposition}[Unfolding and immutable sharing]\label{cdc:mem:prop:graph}
Unfold the finite acyclic heap into ordinary combinator terms and apply the
focus term to the argument stack in order. An application-navigation step
leaves this term unchanged. Each leaf action realizes exactly the corresponding
head contraction of SKI/Iota. The evaluator is deterministic. If the initial
heap has $n$ nodes, after $m$ navigation or leaf actions it has at most $n+3m$
nodes and can be implemented by $O(n+m)$ fixed-size RAM microsteps, including
initialization.
\end{proposition}
\begin{proof}
Acyclicity gives an unambiguous recursively defined unfolding of each pointer.
For navigation, unfolding $\app(p,q)$ with residual arguments is precisely
unfolding $p$ with $q$ pushed in front. The $I,K,S,\iota$ identities are exactly
the displayed rules after replacing every pointer by its unfolding. In the
$S$ case both occurrences of $z$ refer to the same immutable node, but their
unfoldings are identical copies of the same term. No update through one
occurrence changes the other. Fresh nodes point to older nodes, preserving
acyclicity and the meaning of every old pointer.

Node tags and two pointer fields occupy a fixed number of memory words.
Push/pop and each leaf action require a bounded number of field accesses,
register operations, and allocations. Distinct affine address ranges encode
node fields and stack slots without collision. The stack length grows by at
most one per action. Initialization is linear in $n$. The node count follows
because no action allocates more than three nodes. Determinism follows from
inspection of the focus tag and the available stack length.
\end{proof}

\begin{corollary}[Canonical graph-evaluation certificates]\label{cdc:mem:cor:graph}
For a well-formed size-$n$ initial combinator DAG and a budget of $m$ spine
actions, reaching weak-head normal form within that budget has a quartic
natural-number certificate with a unique witness and
$O((n+m)\log^2(n+m+2))$ coordinates. Using optimal sorting networks replaces
the squared logarithm by one logarithm. For this pointer-only evaluator, the
certificate coordinates have $O(\log(n+m+2))$ bits. Alternatively, the linear
backend uses $O(n+m)$ coordinates with
$O((n+m)^2\log(n+m+2))$ bits per coordinate.
\end{corollary}
\begin{proof}
Choose the fixed microstep implementation just described, with a bounded
number of phases per spine action. Insert an action counter so that the budget
is tested on completed spine actions, rather than accepting a partial
contraction. Pad halted states canonically. Apply Theorem~\ref{cdc:mem:thm:ram} with
a constant-factor microstep horizon large enough for the initialization and
$m$ actions. Pointer identifiers, stack positions, control values, allocation
counters, and addresses are bounded by $O(n+m+1)$ when nodes are given dense
identifiers. Its local canonical circuits have the same order of magnitude
for intermediate values. Theorem~\ref{cdc:mem:thm:height} bounds the memory-checker
coordinates by $O((n+m+1)^2)$, which gives the claimed bit length. Substituting the linear backend and
Theorem~\ref{cdc:mem:thm:linear-height} gives the alternative bounds.
\end{proof}

Here $m$ counts navigation \emph{and} contraction actions. It is not merely the
number of source reductions. Well-formedness can be checked separately in
linear-size canonical circuits by verifying tags and that each edge points to
an earlier identifier. The certificate does not expand shared subterms: the
acyclic family $D_0=K$, $D_{j+1}=\app(D_j,D_j)$ has only $j+1$ graph nodes but
$2^j$ leaves in its unfolded term. Pointer sharing can therefore be meaningful
even before any reduction occurs.

ProveIt's existing development uses context-closed weak beta and combinator
relations, while its cited Coq equivalence uses a different call-by-value
lambda model \cite{repo-cl}. The strategy above is a particular deterministic
weak-head graph evaluator. The proof establishes its unfolding semantics; it
does not silently identify those different evaluation relations, prove strong
normalization, or complete the repository's outstanding exact-semantics bridge.
Likewise, injectivity of a source compiler and forward simulation alone do not
imply a bijection on complete target execution histories.

\section{Exact counting and first-halting probabilities\srctag{03}}\label{cdc:mem:sec:counting}
\subsection{Counting executions, not certificate accidents}
Let $\mathsf{AccPaths}_{P,T}(x)$ be the finite set of labelled accepting runs of
the normalized machine, with the unique halted padding included. The RAM
theorem states the exact equality
\begin{equation}\label{cdc:mem:eq:count}
 \#\{Y\in\N^{s(P,T)}:Q_{P,T}(x,Y)=0\}
 =\#\mathsf{AccPaths}_{P,T}(x).
\end{equation}
Finiteness follows from the finite branch alphabet and the fixed horizon,
although the coordinates themselves may be large. This equality would generally
fail for an existential encoding with arbitrary inactive slacks, arbitrary
permutation routes, or unconstrained choices after halting.

The same idea yields a useful bounded counting statement. Define the
\emph{bounded quartic counting problem} to take a sparse integer polynomial of
degree at most four and a binary-encoded bound for each natural variable, and
to count its roots inside that box.
\begin{proposition}[Bounded counting completeness]\label{cdc:mem:prop:counting}
Bounded quartic counting is $\#\mathrm P$-complete under parsimonious
polynomial-time reductions. Hardness already holds for polynomials that are
sums of squares of quadratic residuals with finite natural root sets.
\end{proposition}
\begin{proof}
Membership is immediate from the definition of $\#\mathrm P$: guess the fixed
binary representation of each bounded coordinate, reject out-of-range strings,
and evaluate the sparse polynomial using exact integer arithmetic. The total
number of guessed bits is the sum of the input bound lengths. Degree four
ensures that evaluation and its intermediate bit lengths are polynomial in
the instance size. Fixed-length binary representations with leading zeros
represent each coordinate exactly once.

For hardness, let a $\#\mathrm P$ function be the number of accepting labelled
paths of a polynomial-time nondeterministic Turing machine. The stack encoding
in Section~\ref{cdc:mem:sec:substrates}, or its RAM implementation, gives a canonical
quartic certificate for every path within its polynomial horizon. Pad after
halting canonically. Stack codes have polynomially many bits, because only a
polynomial number of fixed-alphabet symbols can occur; the local arithmetic
and the memory height bound therefore give a polynomial-length binary bound
for every certificate coordinate. The certificate's sparse description and
these bounds are computable in polynomial time. Equation~\eqref{cdc:mem:eq:count}
proves that the number of bounded roots is precisely the required accepting-path
count. The compiler already has finitely many roots for the fixed horizon, so
one may choose a box containing all of them.
\end{proof}
This proposition is a consequence of the parsimonious construction and the
definition of $\#\mathrm P$, not a claim that bounded counting hardness is new.
In particular, it does not place unrestricted quartic root counting in
$\#\mathrm P$; unrestricted equations can have infinitely many roots.

\subsection{Canonical stopping versus random padding}
Suppose each live step consumes exactly one independent fair bit. At a halted
state no bit is consumed; the recorded padding label is forced to zero. If an
accepting path first halts after $\tau\leq T$ live steps, its probability is
$2^{-\tau}$, not $2^{-T}$. Merely counting canonically padded roots and dividing
by $2^T$ would therefore be incorrect.

Let $e_t\in\{0,1\}$ indicate that the machine is live at the start of step $t$.
It is computed canonically from the control state. Attach coordinates
$u_0,\ldots,u_T$ and quadratic residuals
\begin{equation}\label{cdc:mem:eq:weight}
 u_0-1=0,\qquad u_{t+1}-u_t(2-e_t)=0\quad(0\leq t<T).
\end{equation}
These coordinates do not change the number of solutions: their values are
forced by the run.
\begin{theorem}[Exact path weights]\label{cdc:mem:thm:probability}
For the extended quartic certificate, the first-halting accepting probability
by horizon $T$ is
\begin{equation}\label{cdc:mem:eq:probability}
 \Pr[\text{accept by }T]
 =2^{-T}\sum_{Y:\,\widetilde Q_{P,T}(x,Y)=0}u_T(Y).
\end{equation}
Every accepting root whose first halt occurs at $\tau$ has
$u_T=2^{T-\tau}$. An input already accepting at time zero has one root of weight
$2^T$, contributing probability one.
\end{theorem}
\begin{proof}
Before the first halt, $e_t=1$ and the multiplier in \eqref{cdc:mem:eq:weight} is one.
After the halt, $e_t=0$ and the multiplier is two. There are $T-\tau$ such
padding steps, so $u_T=2^{T-\tau}$. The root's contribution to
\eqref{cdc:mem:eq:probability} is $2^{-\tau}$, the probability of its live choice
sequence. Distinct first-halting labelled paths describe disjoint cylinders
in the space of infinite fair-bit streams: no already halted live path is a
proper prefix of another live path. Summing their cylinder probabilities gives
the probability of acceptance by $T$.
\end{proof}

More generally, replace $e_t$ by a canonical bit indicating whether a random
bit is \emph{consumed} at that step. The same proof uses the number of consumed
bits, provided every consumed bit is retained in the labelled path, even when
both choices lead to the same next state. Non-random steps have one label.
The fixed common denominator remains $2^T$ when at most one bit is consumed
per step. This formulation covers deterministic navigation between probabilistic
instructions.

\begin{example}[Stop on the first one]
At each live step, toss a fair bit; accept and halt when it is one, and otherwise
continue. There are exactly $T$ accepting canonical paths by time $T$, namely
$1,01,\ldots,0^{T-1}1$. Their weights are
$2^{T-1},2^{T-2},\ldots,1$. Hence
\[
 \Pr[\text{accept by }T]
 =\frac{2^{T-1}+\cdots+1}{2^T}=1-2^{-T}.
\]
Counting the $T$ roots without weights would instead produce $T/2^T$, illustrating
the error caused by conflating canonical halted padding with random padding.
\end{example}


% ---- manuscript 03, lines 1628-1646 ----
\section{Manuscript 03: Conclusion\srctag{03}}
The main constructive result is an exact linear-size Diophantine memory
compiler: $24L$ natural witnesses, $22L+1$ quadratic residuals, and one sparse
quartic equation, with a singleton fiber for every valid log. The central
nonlocal check is deterministic carry-free product equality. The complementary
comparison-network compiler retains much smaller coordinate heights, exposing
a useful quantitative distinction between arithmetic size and bit resources.

Both backends preserve the same semantic object. Their canonical fibers allow
complete RAM executions to be compiled parsimoniously, support immutable graph
evaluation without unfolding shared terms, and make exact path counting and
first-halting probability formulas possible. Those statements are stronger than
bare existential representability, but remain finite-horizon statements.
The article gives complete mathematical proofs, executable memory compilers,
and a precise boundary between what has been obtained and the further
fixed-variable compression problem. Establishing literature-wide novelty and
formalizing the new proofs are separate tasks, not facts inferred from the
successful computations.


\part{Queues: tag systems and FIFO networks}\label{cdc:part:queues}

\paragraph{Source and scope.} \wtag{} This Part is manuscript~07, ``Causality Without Tableaux: Compact Diophantine Certificates for Queue Universality'', delivered in batch~61. It prints the manuscript's little-endian word codes (its Section~2.3), then its Sections 4--10 in their own order: cyclic tag systems without intermediate queues (\cref{cdc:qc:thm:stream}), a real computation and a self-starting fiction, the variable/degree/witness-height tradeoffs, general deletion-tag systems (\cref{cdc:qc:thm:tag}), the exact dimension of polynomial word memory (\cref{cdc:qc:thm:one-dimensional}), a quartic certificate from two-coordinate memory (\cref{cdc:qc:thm:quartic}), and FIFO networks and switched arithmetic dynamics (\cref{cdc:qc:thm:network-certificate}); it closes with the manuscript's conclusion. Its first-failure product, on which every certificate here depends, is \cref{cdc:qc:thm:first-failure} in Part~I; its conventions are in \cref{cdc:conv:sec}; its fixed-arity boundary section is in Part~IX. The word-memory section is kept before the quartic certificate, as in the manuscript, because the quartic stores each queue in the two coordinates of \cref{cdc:qc:prop:two-coordinates}.

\paragraph{Relation to the other Parts.} \wtag{} Manuscript~07 names manuscripts~04 and~06 as its companions and treats what they leave out: FIFO word content and the temporal availability of produced symbols. Manuscript~06 says that its counts cannot encode FIFO order and asks which queue restrictions can be added (\cref{cdc:q:storage}); this Part answers that question for bounded runs of fixed-read FIFO networks, while data-dependent read counts stay open. The bounded executions here have one witness each, as in Parts~II and~III, but on a different substrate: an ordered queue instead of a multiset of tokens, so none of the constructions below is a second route to \cref{cdc:thm:main}. Where Part~II spends $(2T+1)(d+m)$ or more coordinates on per-step markings, the stream certificate spends $T+1$ (degree $\max\{4,2T\}$), the quartic $3T-2$, the deletion-tag certificate $dT+1$, and a fixed-read network $D+1$. The radix word code is the device of Part~V's linear-size compiler (\cref{cdc:mem:sec:linear}).

\paragraph{Setting and hypotheses.} \wtag{} Queues are first-in first-out. Symbols lie in $\{0,\ldots,B-1\}$; the first written symbol is the head and the \emph{least} significant digit. Formal lengths are signed integers and may become negative after an illegal step. The program, the input word, the terminal word and the horizon $T\ge1$ are data of the construction, not unknowns; for $T=0$ the certificate is a constant polynomial. In fixed-read networks the number of symbols read in each phase does not depend on the data. Inside this Part, $p$ is the period of a cyclic-tag program or the number of queues, $m_i$, $a_i$, $c_i$ are the length, value and $2^{m_i}-1$ of an appendant, $d$ is the deletion number, $K=(B-1)!$, $J=((B-1)!)^C$, $D$ is the total number of consumed symbols, $G$ the guard with slack $u$, and $H$ an upper bound on queue length (\cref{cdc:conv:sec}).

\section{Little-endian words and faithful concatenation\srctag{07}}\label{cdc:qc:sec:words}
\srcnote{\wtag{} Section~2.3 of manuscript~07; its Sections~2.1--2.2 are in \cref{cdc:conv:sec} and its Section~3, the first-failure product \cref{cdc:qc:thm:first-failure} used below, is in Part~I.}

% ---- manuscript 07, lines 146-165 ----
Let $B\geq2$ and let words use digits $0,\ldots,B-1$. The first written symbol is the queue head. Define
\[
 \val_B(a_0\cdots a_{n-1})=\sum_{i=0}^{n-1}a_iB^i,
 \qquad \val_B(\eps)=0.
\]
When $B=2$ the subscript is omitted. The first symbol is thus the \emph{least} significant digit. Written terminal zero symbols still contribute to length.

\begin{lemma}[Word-code identities]\label{cdc:qc:lem:word-code}
For finite words $u,v$,
\[
 \val_B(uv)=\val_B(u)+B^{|u|}\val_B(v).
\]
For each fixed $n$, the value map is injective on words of length $n$. In particular, equality of length and value is equivalent to equality of words.
\end{lemma}
\begin{proof}
Split the defining sum for $uv$ at position $|u|$. For injectivity, reduce an equality of values modulo $B$ to recover the first digit, subtract that digit, divide by $B$, and repeat $n$ times. The length condition is essential: $\eps$, $\word{0}$, and $\word{00}$ all have value zero.
\end{proof}

A candidate execution will be represented by its sequence of consumed symbols. A global concatenation identity checks the content of this sequence. The main difficulty is to verify that each of those symbols existed \emph{before} it was consumed.


% ---- manuscript 07, lines 218-362 ----
\section{Cyclic tag systems without intermediate queues\srctag{07}}\label{cdc:qc:sec:stream}
\subsection{Operational semantics}
A cyclic-tag program is a nonempty finite list of binary words
\[
 \CT=(A_0,A_1,\ldots,A_{p-1}).
\]
The phase at time $i$ is $i\bmod p$, with initial phase zero. From a nonempty queue $bq$, where $b\in\{0,1\}$, one step produces
\[
 qA_{i\bmod p}\quad\text{if }b=1,\qquad q\quad\text{if }b=0.
\]
The phase advances in both cases. There is no transition from an empty queue. Empty appendants are allowed. This is Cook's cyclic-tag semantics with the two symbols renamed $0$ and $1$ \cite{cook}.

Fix an input word $w$, a terminal word $v$, and $T\geq1$. Set
\[
 n=|w|,\quad x=\val(w),\quad h=|v|,\quad y=\val(v),
\]
and abbreviate the phase-dependent program constants by
\[
 m_i=|A_{i\bmod p}|,\quad a_i=\val(A_{i\bmod p}),\quad c_i=2^{m_i}-1.
\]
These constants depend on $i$, but do not depend on the candidate run.

\subsection{The complete polynomial}
Introduce precisely $T+1$ unknowns
\[
 b_0,\ldots,b_{T-1},u\in\N.
\]
The $b_i$ are intended to be the consumed bits. Define the following \emph{polynomial expressions}, not new variables:
\begin{align}
 \ell_i&=n-i+\sum_{j<i}m_jb_j,\label{cdc:qc:eq:lengths}\\
 D_i&=1+c_i b_i, & R_i&=\prod_{j<i}D_j,\label{cdc:qc:eq:prefixes}\\
 L&=\ell_T-h,\label{cdc:qc:eq:L}\\
 F&=x+2^n\sum_{i<T}a_i b_iR_i
       -\sum_{i<T}2^i b_i-2^T y,\label{cdc:qc:eq:F}\\
 G&=\prod_{i<T}\ell_i.\label{cdc:qc:eq:G}
\end{align}
The representing polynomial is
\begin{equation}\label{cdc:qc:eq:main-polynomial}
 \boxed{\displaystyle
 P_{\CT,w,v,T}(\mathbf b,u)
 =\sum_{i<T}\bigl(b_i(b_i-1)\bigr)^2+L^2+F^2+(G-u-1)^2.}
\end{equation}
Every symbol on the right has been defined by integer polynomial operations. In particular, $R_i$ uses products of affine factors, not an exponent whose exponent is an unknown.

\begin{theorem}[Causal stream certificate]\label{cdc:qc:thm:stream}
For the fixed data above, natural zeros of \eqref{cdc:qc:eq:main-polynomial} correspond bijectively to legal $T$-step executions of $\CT$ from $w$ to $v$. The solution records their consumed bits, and its last coordinate is
\[
 u=\prod_{i<T}|w_i|-1,
\]
where $w_i$ is the queue immediately before step $i$.

The polynomial has exactly $T+1$ named variables, $T+3$ residuals, degree at most $\max\{4,2T\}$, and an $O(T)$ arithmetic circuit for a fixed program. Since cyclic-tag execution is deterministic, its natural zero set is either empty or a singleton. When $v=\eps$, a zero exists exactly when first halting occurs at step $T$.
\end{theorem}

\subsection{Proof: reconstructing the entire queue from one word identity}
\begin{proof}
Suppose first that $P(\mathbf b,u)=0$. All residuals vanish, so $b_i\in\{0,1\}$. Write $A_i^{[b_i]}$ for $A_{i\bmod p}$ when $b_i=1$ and for $\eps$ when $b_i=0$. This notation denotes a choice of words, not a variable-exponent polynomial operation.

On these Boolean values,
\[
 D_i=2^{m_i b_i},\qquad R_i=2^{\sum_{j<i}m_jb_j}.
\]
By Lemma~\ref{cdc:qc:lem:word-code}, the right side of
\begin{equation}\label{cdc:qc:eq:global-word}
 b_0b_1\cdots b_{T-1}v
   =wA_0^{[b_0]}A_1^{[b_1]}\cdots A_{T-1}^{[b_{T-1}]}
\end{equation}
has value $x+2^n\sum_{i<T}a_i b_iR_i$. Its length is $n+\sum_{i<T}m_i b_i$. The left side has value $\sum_{i<T}2^i b_i+2^Ty$ and length $T+h$. Consequently $F=0$ and $L=0$ prove the actual word equality \eqref{cdc:qc:eq:global-word}, not merely a congruence.

It remains to show that this static equality describes a causal computation. The equation $G=u+1$ makes $G$ positive. Theorem~\ref{cdc:qc:thm:first-failure}, with initial length $n$, unit consumptions, and productions $m_i b_i\geq0$, implies
\begin{equation}\label{cdc:qc:eq:positive-lengths}
 \ell_i\geq1\qquad(0\leq i<T).
\end{equation}
For each $i$, let
\[
 U_i=wA_0^{[b_0]}\cdots A_{i-1}^{[b_{i-1}]}.
\]
It is a prefix of the right side of \eqref{cdc:qc:eq:global-word}, and
\[
 |U_i|=n+\sum_{j<i}m_jb_j=i+\ell_i\geq i+1.
\]
Thus the first $i+1$ symbols of the global word already occur in $U_i$, before appendant $i$ is produced. In particular the symbol at position $i$ is $b_i$ and is already available.

Inductively, the actual queue before event $i$ is the suffix of $U_i$ obtained by deleting its first $i$ symbols. The assertion holds at $i=0$, where $U_0=w$. It is nonempty by the preceding inequality, its head is $b_i$, and deleting that head and appending $A_i^{[b_i]}$ gives the stated suffix description for $U_{i+1}$. This proves every transition. After $T$ steps, \eqref{cdc:qc:eq:global-word} says that the remaining suffix is exactly $v$.

Conversely, consider a legal $T$-step execution with consumed bits $b_i$. An induction on the operational semantics gives the same suffix description of its queues, and therefore \eqref{cdc:qc:eq:global-word}. This implies $L=F=0$. Its queue lengths are the values \eqref{cdc:qc:eq:lengths}, all positive before the $T$ steps. Set $u=G-1\in\N$. All residuals vanish.

For uniqueness, a fixed input and program determine the first consumed bit and the next queue. Repeating this observation determines the entire consumed sequence. The guard equation then determines $u$. If $v=\eps$, all earlier queues are nonempty while the last is empty, so $T$ is exactly the first halting time.

Finally, $\ell_i$ is affine, $R_i$ has degree at most $i$, and $b_iR_i$ has degree at most $i+1$. Hence $\deg F\leq T$, $\deg L\leq1$, and $\deg G\leq T-1$ because $\ell_0=n$ is constant. Combining with the Boolean residuals gives degree at most $\max\{4,2T\}$. There are $T$ Boolean residuals and three further residuals. Prefix recurrence $R_{i+1}=R_iD_i$, length recurrence, and running sums and products evaluate all expressions in $O(T)$ arithmetic operations, with constants from the fixed program.
\end{proof}

\subsection{What the theorem does not need}
The proof does not need a guessed queue at each time, a guessed time-varying power of two, a division witness for each read, or a general-purpose bounded universal quantifier elimination. Once the consumed bits are fixed, all such quantities are polynomial expressions or are reconstructed at the word level.

The condition that $v$ is fixed matters for the variable count. A terminal \emph{word parameter} can be supplied as data to the compiler without adding witnesses. Treating its numerical code and length as arbitrary unknowns would require a separate validity interface and a different resource statement.

For $T=0$, no step is taken. Use the constant polynomial $0$ if $w=v$ and $1$ otherwise, with no witness variables. The unique empty tuple is the zero of the former. The executable compilers deliberately restrict their interface to $T\geq1$ and report this convention rather than manufacture negative variable counts.

\section{Two examples: a real computation and a self-starting fiction\srctag{07}}
\subsection{A small polynomial with exactly one natural zero}
Take the period-two program
\[
 \CT=(\word{10},\eps),\qquad w=\word{1},\qquad v=\eps,\qquad T=3.
\]
Remember that the first written symbol is the head. The execution is
\[
 \word{1}\longrightarrow\word{10}\longrightarrow\word{0}
 \longrightarrow\eps.
\]
It consumes $(b_0,b_1,b_2)=(1,1,0)$ and has pre-step lengths $(1,2,1)$. Therefore $G=2$ and $u=1$.

The general formulas simplify to
\begin{align*}
 L&=-2+2b_0+2b_2,\\
 F&=1+b_0-2b_1-2b_2+6b_0b_2,\\
 G&=2b_0(2b_0-1).
\end{align*}
Thus the completely explicit polynomial
\begin{align}\label{cdc:qc:eq:small-example}
 P={}&\sum_{i=0}^{2}b_i^2(b_i-1)^2+(-2+2b_0+2b_2)^2\notag\\
 &+(1+b_0-2b_1-2b_2+6b_0b_2)^2\notag\\
 &+\bigl(2b_0(2b_0-1)-u-1\bigr)^2
\end{align}
has the single natural zero $(1,1,0,1)$. Here cancellation lowers the actual degree to four; the general theorem only promises degree at most six when $T=3$.

For a direct check of uniqueness, $L=0$ forces $b_0+b_2=1$ after the Boolean conditions. If $b_0=0$, the guard is zero and its residual cannot vanish. Hence $b_0=1$, $b_2=0$, and $F=0$ gives $b_1=1$. The guard then gives $u=1$. This agrees with the semantic proof and with the exported exact-integer certificate.

\subsection{Why global balance alone is false}
Now take
\[
 \CT=(\eps,\word{1}),\qquad w=\word{0},\qquad T=2,
\]
and propose the consumed word $\word{01}$. The global word equation says
\[
 \underbrace{\word{01}}_{\text{consumed}}
 =\underbrace{\word{0}}_{\text{initial}}
   \underbrace{\eps}_{\text{step 0}}
   \underbrace{\word{1}}_{\text{step 1}},
\]
which is true. The length equation is also true. Algebraically these two residuals are $L=b_1-1$ and $F=-b_0$, so $(b_0,b_1)=(0,1)$ satisfies both.

But the actual queue evolves as $\word{0}\to\eps$ and halts. The proposed second read tries to consume the very symbol that would be appended \emph{as a consequence of that read}. This is not a legal cyclic-tag transition. Here $\ell_0=1$ and $\ell_1=0$, so $G=0$ and no natural slack can solve $G=u+1$.

This example is more than an endpoint bookkeeping issue. It exhibits a genuinely circular dependency in a globally consistent production/consumption equation. The product guard breaks precisely that circularity. In the finite test suite, 268 candidate bitstrings satisfy the raw content and length equations but are rejected by causality; this count is evidence for the implementation, while Theorem~\ref{cdc:qc:thm:stream} supplies the proof for all horizons.

\begin{remark}[\textup{[write]} Relation to the project's tag-system notes]\label{cdc:qc:rem:tagnotes}
The suffix induction in the proof of Theorem~\ref{cdc:qc:thm:stream}, and
the observation that a globally balanced proposed stream can read a symbol
supplied by its own later appendant, also appear in exploratory notes of
ProveIt's Hilbert's-tenth-problem project for deletion-$\beta$ tag systems
(\path{Computability/HilbertTenthProblem/Papers/1980/EXPLORATION_TAG_HALTING_WITHOUT_PREFIX_GUARDS.md},
Sections~2--3; \path{EXPLORATION_FINITE_UNIVERSAL_HISTORY_VERIFIERS.md}),
which manuscript~07 did not consult. Those notes record the causal prefix
inequalities that an exact execution requires, and prove that the
eventual-halting language needs none of them: the sampled-symbol equations
and the final length window already force a genuine halt at the first short
prefix. There is no conflict. Theorem~\ref{cdc:qc:thm:stream} certifies an
exact $T$-step execution with a unique witness, and the guard is what buys
the uniqueness: without it, the notes' own example ($\beta=2$, $w=00$,
$h(0)=0$, $h(1)=100$) is accepted with both selectors $0$ and $001$. The
notes' scalar FIFO identity $D=I+WA$ for a queue of constant length
(\path{EXPLORATION_NATIVE_STREAM_RAW_QUEUE.md}, Section~1) is the global word
identity of this Part in a case where the guard is automatic, since every
queue length there equals the initial length. The project's Lean route for
Jones's 1980 tag certificate stores each queue as a content and a length
marker $3^{\ell}$ (\path{Computability/HilbertTenthProblem/Lean/Diophantine/Paper1980/TagSystem.lean},
\path{TagHistory.lean}: \path{Jones1980.length_step'},
\path{Jones1980.content_step'}), the two-coordinate word memory of
Proposition~\ref{cdc:qc:prop:two-coordinates} in base three.
These notes count the arithmetic operations of fixed-arity universal
certificates; no result of this Part concerns operation counts.
\end{remark}

% ---- manuscript 07, lines 363-624 ----

\section{Variable count, degree, and witness height\srctag{07}}\label{cdc:qc:sec:tradeoffs}
\subsection{The single guard witness has controlled bit height}
Let $M=\max_j|A_j|$, put $g=\max\{0,M-1\}$, and define
\begin{equation}\label{cdc:qc:eq:H}
 H=\max\{1,n+Tg\}.
\end{equation}
During any legal $T$-step run, queue length can increase by at most $g$ per step. Hence $1\leq\ell_i\leq H$ for every $i<T$.

\begin{proposition}[Witness size]\label{cdc:qc:prop:height}
Every zero of the causal stream polynomial satisfies
\[
 0\leq u=G-1\leq H^T-1.
\]
The full witness has a binary description of length
\[
 O\!\left(T+T\log_2\max\{2,H\}\right).
\]
\end{proposition}
\begin{proof}
The $T$ consumed-symbol coordinates are bits. Multiplying $T$ positive queue lengths bounded by $H$ gives $1\leq G\leq H^T$. A nonnegative integer below $H^T$ needs at most $1+T\log_2\max\{2,H\}$ bits under any usual convention for encoding zero.
\end{proof}
The product can be enormous as an integer, but its bit length is only logarithmic in that magnitude. This is one reason not to confuse witness-coordinate count with witness-value bounds. The statement is an upper bound, not a proof that this witness height is best possible.

The coefficients $2^n$ and $2^T$ have $n+1$ and $T+1$ bits, respectively. If every power $2^i$ is stored as a separate literal, their total textual size is quadratic in $T$. An arithmetic circuit can instead generate these powers by repeated doubling. The $O(T)$ arithmetic statement therefore does not assert that every literal-based serialization is bit-linear, nor that multiplication of all intermediate values takes unit time.

\subsection{Removing the slack by a larger polynomial guard}
For the positive integer $H$ of \eqref{cdc:qc:eq:H}, define
\[
 Z_H(z)=\prod_{r=1}^{H}(z-r).
\]
Its integer roots are exactly $1,\ldots,H$. Keep the Boolean, length, and content residuals from Section~\ref{cdc:qc:sec:stream}, but replace $G-u-1$ by $T$ residuals $Z_H(\ell_i)$.

\begin{theorem}[No-slack certificate]\label{cdc:qc:thm:no-slack}
The polynomial
\[
 \widetilde P(\mathbf b)
 =\sum_{i<T}\bigl(b_i(b_i-1)\bigr)^2+L^2+F^2
   +\sum_{i<T}Z_H(\ell_i)^2
\]
has natural zeros in bijection with exact $T$-step executions from $w$ to $v$. It uses exactly $T$ variables, $2T+2$ residuals, has degree at most $2\max\{2,T,H\}$, and admits an $O(T(1+H))$ arithmetic circuit for a fixed program.
\end{theorem}
\begin{proof}
Every legal execution has all pre-step lengths in $\{1,\ldots,H\}$, so it satisfies the added residuals. Conversely, a zero forces every $\ell_i$ into that positive range. The same global word identity and suffix induction used in Theorem~\ref{cdc:qc:thm:stream} then reconstruct the legal execution. No auxiliary choice remains once its consumed bits are fixed.

Since $\ell_i$ is affine, $Z_H(\ell_i)$ has degree at most $H$. The other residual degree bounds were already established. Each bounded-root guard uses $O(H)$ operations, and there are $T$ of them.
\end{proof}

For an empty terminal word, $n>T$ is immediately impossible, because a step removes only one symbol and appends a nonnegative number. When $n\leq T$ and the program is fixed, $H\leq T\max\{1,M\}$. Thus the no-slack version still has degree $O(T)$ and a quadratic-size circuit. It improves the number of witnesses without preserving the linear guard circuit.

\subsection{Why these are not universal optimality claims}
For a fully specified bounded computation, an unconstrained compiler could simulate all $T$ steps and output the constant polynomial $0$ or $1$. That would use no variables at all, but would not preserve the computation as a certificate and would tell us nothing about a structural reduction. Meaningful comparisons must specify which computational data remain visible in the solution and what the compiler is permitted to do.

Our constructions are symbolic: they are obtained from the program, input, endpoint, and horizon without needing to run the machine. The implementation optionally simulates the run \emph{after constructing the formula}, solely to attach an example witness. The command-line option \texttt{--no-witness} disables this extra operation. The theorem's uniqueness is uniqueness of the complete listed witness, not merely uniqueness of the final machine state.

The no-slack construction also warns against claiming that one guard variable is logically unavoidable. It is not. What that variable buys is a particularly small arithmetic circuit, a simple proof, and a compact witness-height estimate. The exact tradeoff among low degree, nearly minimal trace coordinates, and linear circuit size remains a substantive design problem.

\section{General deletion-tag systems\srctag{07}}\label{cdc:qc:sec:general-tag}
\subsection{Semantics and interpolation of the production table}
Fix an alphabet $\{0,\ldots,B-1\}$ with $B\geq2$, a deletion number $d\geq1$, and a word $A(a)$ for each alphabet symbol $a$. A step requires at least $d$ symbols. It consumes the first $d$ of them and appends $A(a)$, where $a$ was the first consumed symbol. A queue of length less than $d$ is halted. The construction applies to every such system; no universality assumption is required for an individual instance. The family of 2-tag systems with suitable finite alphabets is a standard universal substrate \cite{cook,woods-neary}.

We need polynomial table lookups without one-hot variables. For each $a\in\{0,\ldots,B-1\}$ define
\[
 \delta_a(z)=\prod_{\substack{0\leq b<B\\b\ne a}}\frac{z-b}{a-b}.
\]
For an alphabet digit $z$, this equals one when $z=a$ and zero otherwise. Let
\begin{align}
 M(z)&=\sum_{a<B}|A(a)|\delta_a(z),\label{cdc:qc:eq:Mtag}\\
 A(z)&=\sum_{a<B}\val_B(A(a))\delta_a(z),\label{cdc:qc:eq:Atag}\\
 D(z)&=\sum_{a<B}B^{|A(a)|}\delta_a(z).\label{cdc:qc:eq:Dtag}
\end{align}
Here $A(z)$ is a numerical interpolant; $A(a)$ in the operational semantics is a word. The argument or surrounding value notation distinguishes them. All three interpolants lie in $\Q[z]$ and have degree at most $B-1$.

\subsection{The residuals and explicit integer clearing}
Fix input $w$, terminal word $v$, and $T\geq1$. Write $n=|w|$ and $x=\val_B(w)$. The unknowns are
\[
 z_0,\ldots,z_{dT-1},u\in\N,
\]
where $z_0\cdots z_{dT-1}$ is the full consumed word, including the $d-1$ ignored symbols of each step. Put $h_i=z_{di}$ and define
\begin{align}
 E_B(z)&=\prod_{a=0}^{B-1}(z-a),\label{cdc:qc:eq:Etag}\\
 \ell_i&=n-di+\sum_{j<i}M(h_j), & R_i&=\prod_{j<i}D(h_j),\label{cdc:qc:eq:length-tag}\\
 L&=\ell_T-|v|,\label{cdc:qc:eq:Ltag}\\
 F&=x+B^n\sum_{i<T}A(h_i)R_i
       -\sum_{j<dT}B^jz_j-B^{dT}\val_B(v),\label{cdc:qc:eq:Ftag}\\
 G&=\prod_{i<T}\ff{\ell_i}{d}.\label{cdc:qc:eq:Gtag}
\end{align}

These expressions are rational polynomials, but denominators can be cleared explicitly. Set $K=(B-1)!$. The denominator of $\delta_a$ is, up to sign, $a!(B-1-a)!$, which divides $K$. Hence $KM$, $KA$, and $KD$ have integer coefficients. It follows that the residuals
\begin{equation}\label{cdc:qc:eq:integer-tag-residuals}
 E_B(z_j)\ (j<dT),\qquad KL,\qquad K^T F,
 \qquad K^{d(T-1)}(G-u-1)
\end{equation}
all lie in $\Z[\mathbf z,u]$. The last exponent suffices because $\ell_0=n$ is constant, so only $T-1$ falling factorials have a possible denominator. All clearing multipliers are nonzero integers. No new unknowns are introduced.

\begin{theorem}[Deletion-tag stream certificate]\label{cdc:qc:thm:tag}
The sum of squares of \eqref{cdc:qc:eq:integer-tag-residuals} has natural zeros in bijection with exact $T$-step executions from $w$ to $v$. It uses $dT+1$ variables and $dT+3$ residuals. Its degree is at most
\begin{equation}\label{cdc:qc:eq:tag-degree}
 2\max\{B,\ T(B-1),\ d(T-1)(B-1),\ 1\}.
\end{equation}
For fixed $B,d$, and production table, it has an $O(T)$ arithmetic circuit. Each successful instance has a unique witness. If $|v|<d$, it certifies first halting at exactly $T$.
\end{theorem}
\begin{proof}
The digit residuals force every $z_j$ to be an alphabet symbol. On those symbols, the interpolants have their prescribed natural values. Clearing denominators does not change any zero condition.

The length and content equations, together with Lemma~\ref{cdc:qc:lem:word-code}, give the word equality
\begin{equation}\label{cdc:qc:eq:tag-word}
 z_0\cdots z_{dT-1}v=wA(h_0)A(h_1)\cdots A(h_{T-1}).
\end{equation}
The guard equation and Theorem~\ref{cdc:qc:thm:first-failure}, with consumption $d$ at every step, imply $\ell_i\geq d$ for every $i<T$. Let $U_i=wA(h_0)\cdots A(h_{i-1})$. Its length is $di+\ell_i\geq d(i+1)$. Thus the whole next consumed block $z_{di}\cdots z_{di+d-1}$ is already in $U_i$, not in an appendant produced by this or a later step. The suffix induction now deletes $d$ symbols at a time. Its first symbol is $h_i$, so the selected appendant is correct. The terminal suffix is $v$.

Conversely, any actual execution satisfies \eqref{cdc:qc:eq:tag-word}, the length equation, the alphabet constraints, and the positive resource product. It therefore gives a zero with $u=G-1$. Deterministic semantics fixes every consumed symbol, including the ignored ones, and the guard fixes the slack.

For the degree estimate, digit residuals have degree $B$, length has degree at most $B-1$, and $A(h_i)R_i$ has degree at most $(i+1)(B-1)$. The nonconstant part of the guard has at most $d(T-1)$ factors of degree at most $B-1$. The term $-u-1$ contributes degree at most one. Squaring gives \eqref{cdc:qc:eq:tag-degree}. With fixed alphabet and deletion number, each table evaluation and falling factorial uses a fixed number of arithmetic operations; prefix products, lengths, and sums are accumulated in $O(T)$ operations. The clearing multipliers are constant coefficients with bit length linear in $T$ for fixed $B,d$.
\end{proof}

\paragraph{Integer-only circuit realization.}
Denominator clearing need not introduce division gates. Put $\widehat M=KM$, $\widehat A=KA$, $\widehat D=KD$, and
\[
 \lambda_i=K(n-di)+\sum_{j<i}\widehat M(h_j),\qquad
 \Pi_i=\prod_{j<i}\widehat D(h_j).
\]
Then the cleared content and guard residuals can be evaluated as
\begin{align*}
 K^TF={}&K^T\!\left(x-\sum_{j<dT}B^jz_j-B^{dT}\val_B(v)\right)
       +B^n\sum_{i<T}K^{T-i-1}\widehat A(h_i)\Pi_i,\\
 K^{d(T-1)}(G-u-1)={}&\ff{n}{d}
     \prod_{i=1}^{T-1}\prod_{r=0}^{d-1}(\lambda_i-Kr)
       -K^{d(T-1)}(u+1).
\end{align*}
All exponents in these expressions are fixed nonnegative integers. Precompute the powers of $K$ by multiplication; every other gate is an integer addition, subtraction, or multiplication. This gives the asserted linear arithmetic size without using rational constants internally.

\subsection{Halting tails are not necessarily empty}
A deletion-$d$ tag machine can halt with $1,\ldots,d-1$ symbols still present. Equating its terminal queue with $\eps$ would miss those computations. There are
\[
 N_{B,d}=1+B+\cdots+B^{d-1}
\]
possible halted words. The endpoint theorem handles each of them.

\begin{corollary}[Unspecified halted endpoint]\label{cdc:qc:cor:halting-tail}
For fixed $B,d$, program, input, and $T\geq1$, first halting at exactly $T$ has a witness-faithful polynomial with $dT+1$ variables and degree at most $N_{B,d}$ times \eqref{cdc:qc:eq:tag-degree}.
\end{corollary}
\begin{proof}
For each word $v$ of length less than $d$, let $P_v$ be the sum-of-squares polynomial of Theorem~\ref{cdc:qc:thm:tag}, using the same names for all consumed-symbol variables and the slack. Use $\prod_{|v|<d}P_v$. An integer product vanishes exactly when one factor vanishes. Each factor certifies an actual execution to its specified endpoint. Determinism permits at most one such endpoint and one complete witness. The common guard does not create extra choices. The degree of the product is at most the sum of the degrees of its factors.
\end{proof}
The product here is over a fixed finite collection determined by $B,d$, not over an unbounded family of words. Its size can be large as these parameters grow; the linear-in-$T$ statement keeps them fixed.

\subsection{A ternary example with a nonempty halted tail}
Let $B=3$, $d=2$, and choose productions
\[
 A(0)=\eps,\qquad A(1)=\word{20},\qquad A(2)=\word{0}.
\]
Starting from $\word{10}$, the execution is
\[
 \word{10}\longrightarrow\word{20}\longrightarrow\word{0}.
\]
At $T=2$ it halts with one symbol left. The consumed digits are $(1,0,2,0)$; both pre-step lengths are two, so
\[
 G=\ff{2}{2}\ff{2}{2}=4,\qquad u=3.
\]
The global identity is $\word{1020}\,\word{0}=\word{10}\,\word{20}\,\word{0}$. Its value is $19$ in base three, with the stated little-endian convention. This illustrates why the ignored symbols and the terminal length must both be retained.

\section{The exact dimension of polynomial word memory\srctag{07}}\label{cdc:qc:sec:memory}
\subsection{Two affine coordinates suffice}
The stream certificate eliminates queue states globally. For a local encoding, a useful state representation is
\begin{equation}\label{cdc:qc:eq:word-pair}
 \mathcal E_B(w)=(\val_B(w),B^{|w|})=(v,s)\in\N^2.
\end{equation}
The second coordinate records scale, not length itself. Define
\begin{equation}\label{cdc:qc:eq:star}
 (v,s)\star(w,r)=(v+sw,sr).
\end{equation}
Lemma~\ref{cdc:qc:lem:word-code} gives
\[
 \mathcal E_B(uv)=\mathcal E_B(u)\star\mathcal E_B(v).
\]
The operation is associative, as is clear either by direct expansion or by identifying $(v,s)$ with the matrix
\[
 \begin{pmatrix}s&v\\0&1\end{pmatrix}.
\]
The empty word has code $(0,1)$, the identity element. Insertion of a fixed symbol at either end is affine:
\begin{equation}\label{cdc:qc:eq:insertion-pair}
 L_a(v,s)=(Bv+a,Bs),\qquad R_b(v,s)=(v+bs,Bs).
\end{equation}
For every pair of symbols, $L_aR_b=R_bL_a$, reflecting the word identity $a(wb)=(aw)b$.

\begin{proposition}\label{cdc:qc:prop:two-coordinates}
The map $\mathcal E_B$ is a faithful natural-number encoding of all finite words, with polynomial concatenation and affine insertion at both ends.
\end{proposition}
\begin{proof}
The scale determines the length because $B^m=B^n$ implies $m=n$ for natural $m,n$. The value then determines the word by Lemma~\ref{cdc:qc:lem:word-code}. The displayed formulas establish the operation claims.
\end{proof}
The next result shows that two is the minimum number of natural coordinates for this insertion architecture. It does not assume that arbitrary-word concatenation is polynomial; two prepend maps and a single fixed append map already suffice for the obstruction.

\subsection{An elementary centralizer lemma}
Composition of polynomials is denoted by $\circ$. A polynomial commutes with $Q$ when $F\circ Q=Q\circ F$.

\begin{lemma}[Uniqueness at a fixed degree and leading coefficient]\label{cdc:qc:lem:centralizer-unique}
Let $Q\in\R[X]$ have degree $d\geq2$. Two nonconstant polynomials $H_1,H_2$ that commute with $Q$, have the same degree $m$, and have the same leading coefficient must be equal.
\end{lemma}
\begin{proof}
Suppose $\Delta=H_1-H_2$ is nonzero and has degree $k<m$. Let $c\ne0$ be the leading coefficient of $Q$ and $a\ne0$ the common leading coefficient of $H_1,H_2$. Subtract the two commutation identities to obtain
\[
 Q(H_1)-Q(H_2)=\Delta(Q).
\]
The right side has degree $dk$. On the left, factor the difference of the leading powers of $Q$. Its highest term has nonzero coefficient
\[
 dc\,a^{d-1}\lc(\Delta)
\]
and degree $(d-1)m+k$. Lower powers of $Q$ give smaller degrees. Thus equality would imply $(d-1)m+k=dk$, or $(d-1)(m-k)=0$, a contradiction.
\end{proof}

\begin{lemma}[Positive-leading centralizers are commutative]\label{cdc:qc:lem:positive-centralizer}
Let $Q$ be a nonconstant real polynomial with positive leading coefficient, and assume $Q\ne X$. Any two nonconstant real polynomials $F,G$ with positive leading coefficients that commute with $Q$ also commute with one another.
\end{lemma}
\begin{proof}
First suppose $\deg Q=d\geq2$ and let $c>0$ be its leading coefficient. If $\deg F=m$ and $a>0$ is the leading coefficient of $F$, comparison in $F(Q)=Q(F)$ gives
\[
 ac^m=ca^d,\qquad a^{d-1}=c^{m-1}.
\]
The positive solution is $a=c^{(m-1)/(d-1)}$. Similarly, if $\deg G=n$, its leading coefficient is $c^{(n-1)/(d-1)}$.

Both $F\circ G$ and $G\circ F$ commute with $Q$. They have degree $mn$, and their leading coefficients both equal $c^{(mn-1)/(d-1)}$. Lemma~\ref{cdc:qc:lem:centralizer-unique} makes them equal.

Now suppose $Q(X)=cX+e$ with $c>0$. If $c=1$, then $e\ne0$ because $Q\ne X$. Commutation gives
\[
 F(X+e)-F(X)=e.
\]
For a nonconstant polynomial the left side has degree $\deg F-1$, so $F$ has degree one. Comparing its slope shows that $F(X)=X+f$ for a constant $f$. The same holds for $G$, and translations commute.

If $c\ne1$, translate the fixed point $e/(1-c)$ to zero. In the translated coordinate, the commuting condition is $F(cX)=cF(X)$. For the coefficient of $X^k$, this requires $c^k=c$ whenever that coefficient is nonzero. Since $c>0$ and $c\ne1$, only $k=1$ is possible, including the exclusion of the constant coefficient. Thus all such maps are scalar multiplications in the same translated coordinate, and they commute.
\end{proof}

This argument is a small positive-real specialization of the subject of commuting polynomials. It is included in full rather than importing a classification theorem or presenting the general centralizer theory as new \cite{atela-hu}.

\subsection{One natural coordinate cannot support both ends}
\begin{theorem}[One-coordinate polynomial obstruction]\label{cdc:qc:thm:one-dimensional}
There is no injective map
\[
 e:\{0,1\}^{*}\longrightarrow\N
\]
and no real polynomials $P_0,P_1,Q$ satisfying, for every finite binary word $w$,
\begin{equation}\label{cdc:qc:eq:scalar-operations}
 P_0(e(w))=e(0w),\qquad P_1(e(w))=e(1w),\qquad Q(e(w))=e(w0).
\end{equation}
In particular, a faithful natural-coordinate word representation with polynomial insertion of arbitrary binary symbols at both ends requires at least two coordinates. Proposition~\ref{cdc:qc:prop:two-coordinates} attains that dimension with affine maps.
\end{theorem}
\begin{proof}
The image of $e$ is an infinite subset of $\N$ and is therefore unbounded. Each of $P_0,P_1,Q$ is nonconstant, because its outputs on that image include infinitely many distinct word codes. Each has positive leading coefficient: a nonconstant real polynomial with negative leading coefficient is eventually negative, contrary to its natural values on an unbounded subset of $\N$.

For $a\in\{0,1\}$ and every $w$,
\[
 (P_a\circ Q)(e(w))=e(a(w0))=e((aw)0)=(Q\circ P_a)(e(w)).
\]
Polynomials agreeing on an infinite set are equal, so $P_a$ commutes with $Q$ as a polynomial identity. Moreover $Q\ne X$, since $Q(e(w))=e(w)$ would imply $w0=w$ by injectivity, which contradicts word length.

Lemma~\ref{cdc:qc:lem:positive-centralizer} now gives $P_0\circ P_1=P_1\circ P_0$. Evaluating at $e(w)$ yields
\[
 e(01w)=e(10w).
\]
The two words are distinct, contradicting injectivity.
\end{proof}

\srcnote{\wtag{} \Cref{cdc:qc:thm:one-dimensional} is consistent with manuscript~01's injective polynomial pairing $\pi(a,b)=(a+b)^2+a$ (\cref{cdc:ct:lem:pair}) and with manuscript~05's pairing \eqref{cdc:eq:pair}: those are maps of \emph{two} arguments used to code finite terms and lists, while the theorem concerns \emph{one}-argument insertion maps on the code of a single word.}

% ---- manuscript 07, lines 625-784 ----

\subsection{The boundary is architectural, not computational}
The theorem does not say that one integer cannot encode a word. For example,
\[
 e(w)=2^{|w|}+\val(w)
\]
is a faithful sentinel encoding. Both prepends are polynomial:
\[
 e(aw)=2e(w)+a.
\]
Appending a zero, however, requires adding $2^{|w|}$, the leading power of two of the code. The theorem proves that no alternative scalar coding can make both prepends and this append simultaneously polynomial on all word codes.

The result does not exclude one-register computation with conditional operations, floor, division, prime-exponent tests, or auxiliary Diophantine witnesses. It does not apply to a finite set of word lengths, where arbitrary finite tables can be interpolated, or to a code whose operations depend on additional hidden control state. Nor does it claim a dimension bound for all possible continuous or rational encodings. Its conclusion is exactly about faithful natural-number word codes with fixed scalar polynomial insertion maps.

\section{A quartic certificate from two-coordinate memory\srctag{07}}\label{cdc:qc:sec:quartic}
\subsection{Encoding a single queue step}
For binary words let $(v,s)=\mathcal E_2(w)$. A nonempty word has $s=2^{|w|}\geq2$. Write its head as $b\in\{0,1\}$ and the value of its tail as $q$. Then
\[
 v=b+2q.
\]
At phase $i$, the update after the optional appendant is
\begin{equation}\label{cdc:qc:eq:local-step}
 v'=q+\frac{a_i b s}{2},\qquad
 s'=\frac{s\bigl(1+(2^{m_i}-1)b\bigr)}{2}.
\end{equation}
The denominator two is exact on a valid nonempty input scale. The updated scale is $2^{|w|-1+m_i b}$, and the updated value is precisely the value of the tail followed by the selected appendant.

\subsection{The full quadratic residual system}
Fix $w$ and $T\geq1$, and require the endpoint to be empty. Substitute the constant endpoints
\[
 (v_0,s_0)=(\val(w),2^{|w|}),\qquad (v_T,s_T)=(0,1).
\]
The natural unknowns are the $2(T-1)$ internal coordinates
\[
 (v_i,s_i)\quad(1\leq i<T)
\]
and $T$ quotients $q_i$ for $0\leq i<T$. In the following formulas,
\[
 b_i=v_i-2q_i
\]
is an abbreviation for an affine expression, not a new variable. For each $i<T$ use the three residuals
\begin{align}
 B_i&=b_i(b_i-1),\label{cdc:qc:eq:quadB}\\
 V_i&=2v_{i+1}-2q_i-a_i b_i s_i,\label{cdc:qc:eq:quadV}\\
 S_i&=2s_{i+1}-s_i-(2^{m_i}-1)b_i s_i.\label{cdc:qc:eq:quadS}
\end{align}
Let $Q_{\CT,w,T}=\sum_{i<T}(B_i^2+V_i^2+S_i^2)$.

\begin{theorem}[Local quartic certificate]\label{cdc:qc:thm:quartic}
The natural zero set of $Q_{\CT,w,T}$ is a singleton if the cyclic-tag run from $w$ first halts at time $T$, and is empty otherwise. The polynomial has $3T-2$ variables, $3T$ quadratic residuals, degree at most four, and $O(T)$ arithmetic-circuit size for a fixed program.
\end{theorem}
\begin{proof}
Every actual run first halting at $T$ gives the displayed internal word codes and quotients $q_i=\lfloor v_i/2\rfloor$. Equations~\eqref{cdc:qc:eq:quadB}--\eqref{cdc:qc:eq:quadS} are just the head decomposition and the update \eqref{cdc:qc:eq:local-step} with denominators cleared. Thus it gives a zero.

Conversely, suppose there is a zero. We induct from the valid initial word code. Assume $(v_i,s_i)$ is the code of an actual queue. Equation $B_i=0$ forces the signed integer $b_i=v_i-2q_i$ to be either zero or one. Since $q_i\in\N$, the decomposition $v_i=2q_i+b_i$ uniquely determines both $q_i$ and $b_i$.

If the queue is nonempty, its scale is even and at least two. The other residual equations uniquely force the actual next value and scale from \eqref{cdc:qc:eq:local-step}; hence they again form a valid word code. Notice that this establishes scale validity by induction. There is no separate existential predicate for being a power of two.

If the queue were empty before a required step, then $v_i=0$ and $s_i=1$. The Boolean residual and nonnegativity of $q_i$ force $q_i=b_i=0$. The scale residual would become $2s_{i+1}=1$, which has no natural solution. Therefore no zero can pass through an empty queue before step $T$.

Induction proves all $T$ actual transitions, and the final constants $(0,1)$ identify the empty word. The run consequently first halts at $T$. Determinism and quotient uniqueness fix every variable. The count is $2(T-1)+T=3T-2$, and all displayed residuals are quadratic, proving the remaining bounds.
\end{proof}

\subsection{The small example in quartic coordinates}
For the example $\CT=(\word{10},\eps)$, $w=\word{1}$, $T=3$, the internal states are
\[
 (v_1,s_1)=(1,4),\qquad (v_2,s_2)=(0,2),
\]
and the quotients are $(q_0,q_1,q_2)=(0,0,0)$. Thus the seven-coordinate witness is
\[
 (v_1,s_1,v_2,s_2,q_0,q_1,q_2)=(1,4,0,2,0,0,0).
\]
This is the witness in \path{data/07-queue-causality-quartic_example.json}. The stream construction uses only four coordinates for the same computation; the quartic construction guarantees degree four for every horizon instead of allowing degree to grow with $T$.

\subsection{Comparing information content}
A valid state of length at most $n+T\max\{0,M-1\}$ has value and scale requiring $O(n+TM+1)$ bits. The local witness therefore has the elementary upper bound $O(T(n+TM+1))$ bits. The stream witness instead stores the consumed bits and a product whose bit length is controlled by Proposition~\ref{cdc:qc:prop:height}.

These are different certificates for the same semantic object. One maintains a low-degree local recurrence by carrying explicit memory coordinates. The other eliminates those coordinates, retains one global content equation, and pays for the elimination in polynomial degree. The no-slack variant makes the independence of degree, circuit size, and variable count even more visible.

\section{FIFO networks and switched arithmetic dynamics\srctag{07}}\label{cdc:qc:sec:substrates}
\subsection{Causality for several channels}
The causal reconstruction argument is not inherently a one-queue argument. Consider $p$ FIFO queues. At event $i$, a candidate trace consumes the word $C_{ij}$ from queue $j$ and appends the word $A_{ij}$ to it. All reads occur before that event's writes. Local rules, if present, must separately validate the proposed input/output words. Let $W_j,V_j$ be the initial and final queues, and write
\[
 \ell_{ij}=|W_j|+\sum_{k<i}(|A_{kj}|-|C_{kj}|),\qquad
 G=\prod_{i<T}\prod_{j<p}\ff{\ell_{ij}}{|C_{ij}|}.
\]

\begin{theorem}[FIFO network reconstruction]\label{cdc:qc:thm:network-causality}
The candidate words form a legal FIFO execution between the specified endpoints if and only if $G>0$ and, for each queue $j$,
\begin{equation}\label{cdc:qc:eq:network-word}
 C_{0j}C_{1j}\cdots C_{T-1,j}V_j
 =W_jA_{0j}A_{1j}\cdots A_{T-1,j}.
\end{equation}
Here ``legal'' refers to queue availability and the stated consumption/production words; satisfaction of any additional local control rule is a separate premise.
\end{theorem}
\begin{proof}
Any legal execution satisfies the word identities by following each queue's FIFO order, and its resource product is positive by Theorem~\ref{cdc:qc:thm:first-failure}.

Conversely, the same resource theorem makes every queue long enough for the prescribed consumption at each event. Before event $i$, the already produced prefix on queue $j$ is $W_jA_{0j}\cdots A_{i-1,j}$. Its length is at least the number of symbols previously consumed on that queue plus $|C_{ij}|$. The corresponding next block in the global identity \eqref{cdc:qc:eq:network-word} therefore lies in this already produced prefix and equals $C_{ij}$. Simultaneously induct on all queues: the current queue is the suffix of its produced prefix after deleting all previous consumptions. Removing the next block and appending the designated output preserves that description. The final suffixes are exactly $V_j$.
\end{proof}

\subsection{A trace-symbol compiler for fixed-read networks}
A useful concrete class of networks has a fixed finite cyclic control schedule. Each phase specifies a consumption vector $(c_{ij})_{j<p}$, independent of the data. A finite table maps the tuple of consumed symbols to one output word on each queue. Such a table may branch arbitrarily on the consumed tuple, but the number of symbols read in that phase is fixed. A phase is disabled if some queue is too short.

Fix $T$, all endpoints, and the program. Let
\[
 D=\sum_{i<T}\sum_{j<p}c_{ij},\qquad
 C=\max_{i<T}\sum_{j<p}c_{ij},\qquad h=C(B-1).
\]
For $T\geq1$, introduce one alphabet-digit variable for each of the $D$ consumed symbols, and one natural slack $u$. Tensor-product Lagrange interpolation of each finite output table gives its output length, output value, and output scale as polynomials of degree at most $h$ in that event's digits. Give each queue its length and value equation from \eqref{cdc:qc:eq:network-word}; append prefixes are accumulated by multiplying their scale interpolants. Use the single global resource residual $G-u-1$.

\begin{theorem}[Fixed-read network certificate]\label{cdc:qc:thm:network-certificate}
The resulting residuals, after fixed integer denominator clearing, give a witness-faithful polynomial with $D+1$ variables, $D+2p+1$ residuals, and degree at most
\[
 2\max\{B,Th,hD,1\}.
\]
For a fixed network program with a fixed number of queues, it has $O(T)$ arithmetic-circuit size. Its natural zeros correspond bijectively to exact $T$-event executions between the fixed endpoints.
\end{theorem}
\begin{proof}
The $D$ digit residuals force the local table arguments to be alphabet symbols. Table evaluation then returns actual nonnegative lengths and actual finite words. The two residuals per queue imply \eqref{cdc:qc:eq:network-word}, by equal length and equal radix value. The resource residual makes the global guard positive, so Theorem~\ref{cdc:qc:thm:network-causality} reconstructs the execution. Conversely a legal execution provides all its consumed digits and the unique slack $G-1$. The fixed schedule and deterministic tables determine every consumed word from the endpoints' initial queues, hence the complete witness is unique.

Each digit residual has degree $B$. An output value times the product of preceding output scales has degree at most $Th$. Formal queue lengths have degree at most $h$. The product guard has $D$ linear-in-length factors and therefore degree at most $hD$. Including $u$ and squaring gives the stated bound.

For explicit clearing, put $J=((B-1)!)^C$. Every tensor-product table interpolant has denominator dividing $J$. Multiply each length equation by $J$ and each content equation by $J^T$. Writing $\lambda_{ij}=J\ell_{ij}$, the cleared guard is
\[
 \prod_{i<T}\prod_{j<p}\prod_{r<c_{ij}}(\lambda_{ij}-Jr)
      -J^D(u+1).
\]
It is an integer polynomial. Weighted integer prefix products give integer-only content circuits exactly as in Section~\ref{cdc:qc:sec:general-tag}. With the program fixed, table sizes, $p$, $C$, and $B$ are fixed, while $D=O(T)$; the arithmetic size follows.
\end{proof}

This class contains cyclic tag systems as the case of one queue, one symbol read per phase, and the two-entry append/skip table. It therefore includes universal computation, but the theorem does not claim that every individual network is universal. It also does not cover arbitrary data-dependent read counts or hidden zero tests without additional residuals. The general network theorem is proved here; the shipped executable compilers instantiate the one-queue cyclic-tag and deletion-tag cases, not the full network-table interface.

\subsection{A two-track arithmetic substrate}
On valid binary word codes, \eqref{cdc:qc:eq:local-step} becomes the following periodically switched affine rule:
\begin{equation}\label{cdc:qc:eq:affine-dynamics}
 (v,s)\longmapsto
 \begin{cases}
 (v/2,s/2),&v\text{ even},\\[2pt]
 ((v-1+a_i s)/2,\,2^{m_i-1}s),&v\text{ odd}.
 \end{cases}
\end{equation}
The state domain is
\[
 \mathcal W=\{(v,2^n):n\in\N,\ 0\leq v<2^n\},
\]
and the process stops at $s=1$. When $m_i=0$, the second multiplier is $1/2$, a rational coefficient, not an omitted case.

\begin{corollary}[Conjugate arithmetic computation]
The map $\mathcal E_2$ is a bijection from binary words to $\mathcal W$ and conjugates cyclic-tag transitions to \eqref{cdc:qc:eq:affine-dynamics}, including halting. Consequently this family of two-track, parity-switched rational-affine processes inherits cyclic-tag universality and the exact certificates of Theorems~\ref{cdc:qc:thm:stream} and \ref{cdc:qc:thm:quartic}.
\end{corollary}
\begin{proof}
Each integer $0\leq v<2^n$ has a unique length-$n$ base-two digit expansion, including terminal zeros. For a nonempty word, its head is the parity of $v$, and \eqref{cdc:qc:eq:local-step} is precisely the stated branch formula. Scale one corresponds only to the empty word. This establishes a step-for-step bijection, not merely a simulation with unspecified overhead.
\end{proof}
The switching condition is essential. This is not a claim that iteration of one everywhere-polynomial scalar function is universal. Nor does the two-coordinate lower bound apply to arbitrary parity-switched scalar dynamics.

\subsection{What transports to spatial computation}
Cook's construction connects cyclic tags to glider interactions and then to Rule~110 \cite{cook}. A certified simulation can transport a statement about an encoded cyclic-tag run to the corresponding spatial computation. But the transport has to preserve the chosen event and endpoint semantics. A tag step is not automatically one cellular-automaton time step, and an encoded infinite periodic background is not a finite queue state.

Thus our displayed counts measure tag events or fixed-read network events. They do not establish a direct $T+1$-variable polynomial for every arbitrary $T$-step Rule~110 space-time diagram. A direct glider or cellular-automaton compiler with certified clocks, boundary conditions, and observation conventions is a separate research target. This distinction keeps the unconventional-substrate connection concrete without conflating different notions of time and input.

\srcnote{\wtag{} The bounded cellular-automaton constructions of the other manuscripts are finite space--time tableaux with a fixed number of cells: manuscript~05's finite cellular automata and manuscript~04's cellular automata in a finite box, manuscript~01's circuit cylinders and manuscript~02's Boolean spacetime substrates (all in Part~VIII), and manuscript~03's bounded cellular computation (\cref{cdc:mem:sec:substrates}). The ``without tableaux'' of manuscript~07's title contrasts its queue certificates with such constructions; as just said, it does not claim a direct polynomial for arbitrary Rule~110 diagrams.}

% ---- manuscript 07, lines 785-785 ----


% ---- manuscript 07, lines 950-956 ----
\section{Manuscript 07: Conclusion\srctag{07}}
A queue computation has two logically different requirements: its consumed word must agree with what the computation produces, and each consumed symbol must already be present when it is read. A global radix identity addresses the first requirement. A first-failure product addresses the second without quantifying intermediate queues.

Together they give an exact $T+1$-variable cyclic-tag certificate, a $dT+1$-variable deletion-tag certificate, and a consumed-symbol-plus-one certificate for fixed-read FIFO networks. The no-slack and quartic alternatives show how circuit size, degree, and variable count can be traded without losing witness uniqueness. The polynomial-memory obstruction independently explains why two-coordinate local word memory cannot simply be replaced by one scalar with polynomial insertion at both ends.

These results answer concrete structural questions about Diophantine compilation. Their next useful tests are formalization, sharper resource bounds, and verified transport to genuinely spatial substrates. Their bounded uniqueness is a precise starting point for fixed-arity research, not a substitute for the additional argument that research requires.


\part{FRACTRAN, SKI and term rewriting}\label{cdc:part:rewriting}

\paragraph{Source and scope.} \wtag{} Manuscript~04's Sections 6--7 give a linear-size, unique-witness certificate for bounded FRACTRAN runs with first-applicable-fraction priority (\cref{cdc:wf:thm:fractran}) and a context-local certificate for SKI reduction along a prescribed schedule (\cref{cdc:wf:thm:ski}), extended to finite first-order rewriting (\cref{cdc:wf:cor:rewrite}). The last section of the Part collects three shorter treatments: manuscript~02's local single-fold FRACTRAN step interface, manuscript~01's local polynomial adapter for combinatory logic and its warning about forward simulations, and manuscript~05's bounded word rewriting and explicit term graphs.

\paragraph{Relations.} \wtag{} For FRACTRAN, manuscript~04's certificate covers a complete bounded run; manuscript~02's interface certifies one step and the terminal condition, and is printed after it as the local counterpart. For combinatory logic the three constructions are different and all are kept: manuscript~04 encodes whole subterms by integer codes and follows a fixed schedule, manuscript~01 gives a local head-rule adapter with the injective pairing of \cref{cdc:ct:lem:pair}, and manuscript~03's immutable graph evaluator is in Part~V (\cref{cdc:mem:subsec:graph}). The repository's combinatory-logic development distinguishes source syntax, runtime syntax and evaluation strategy (\cite{repo-cl}); none of these constructions is a formal bridge to it.

\paragraph{Setting and hypotheses.} \wtag{} Manuscript~04: an ordered list of $r$ fractions in lowest terms and a horizon of $T$ steps; SKI terms coded by $\App(\ell,r)$ with leaves $S,K,I$ coded $0,1,2$, a fixed schedule $\sigma$ of total context depth~$H$. Manuscript~01: the same pairing $\pi(a,b)=(a+b)^2+a$ with $\operatorname{App}(a,b)=3+\pi(a,b)$, validity of the input term as a domain promise. Manuscript~05: a fixed alphabet, substring rules, horizon~$T$ and an explicit word-length bound~$B$.

% ---- manuscript 04, lines 354-573 ----
\section{A linear-size, unique-witness FRACTRAN priority certificate\srctag{04}}
\subsection{Semantics and the issue with naive disjunction}
A FRACTRAN program is a finite ordered list of positive fractions
\[
 \frac{a_1}{b_1},\ldots,\frac{a_r}{b_r},\qquad r\geq1,
\]
which we first put in lowest terms. At a positive integer $n$, choose the least $j$ for which $b_j\mid n$ and replace $n$ by $a_jn/b_j$. If there is no such $j$, execution halts. Reduction to lowest terms is essential: applicability of an unreduced fraction $a/b$ is $b\mid an$, not generally $b\mid n$.

Diophantine treatment of this language is established prior work \cite{conway1987,lf2019}. Our question is the exact cost of a bounded, priority-correct, unique-witness certificate. An existential disjunction over all applicable fractions would lose priority. Giving every inactive branch an unrestricted remainder witness would also lose uniqueness. The following circuit handles both issues with a linear, rather than all-pairs, priority scan.

\subsection{Canonical division and zero detection}
Fix a positive denominator $b$. Introduce $q,\rho,s,u,z\in\N$ and impose
\begin{align}
 n&=bq+\rho,\label{cdc:wf:eq:frac-div}\\
 \rho+s&=b-1,\label{cdc:wf:eq:frac-bound}\\
 z(z-1)&=0,\label{cdc:wf:eq:frac-bit}\\
 z\rho&=0,\label{cdc:wf:eq:frac-zero}\\
 \rho&=(1-z)(u+1),\label{cdc:wf:eq:frac-positive}\\
 zu&=0.\label{cdc:wf:eq:frac-inactive}
\end{align}

\begin{lemma}[Single-valued divisibility gadget]\label{cdc:wf:lem:division}
For every $n\in\N$ and $b\geq1$, equations \eqref{cdc:wf:eq:frac-div}--\eqref{cdc:wf:eq:frac-inactive} have exactly one natural solution. It has $q=\lfloor n/b\rfloor$, $\rho=n\bmod b$, $s=b-1-\rho$, and
\[
 z=\mathbf 1_{b\mid n},\qquad
 u=\begin{cases}0,&\rho=0,\\\rho-1,&\rho>0.\end{cases}
\]
\end{lemma}
\begin{proof}
The first two equations imply $0\leq\rho<b$, so the quotient and remainder theorem forces $q,\rho$, and then $s$. If $\rho=0$, the fifth equation makes $z=0$ impossible, since it would say $0=u+1$. Hence $z=1$, and the last equation forces $u=0$. If $\rho>0$, the fourth equation forces $z=0$, and the fifth forces $u=\rho-1$. These values satisfy all six equations. The argument also covers $b=1$, where the remainder is always zero.
\end{proof}

Equation~\eqref{cdc:wf:eq:frac-inactive} is not cosmetic. Without it, the divisible case would permit every natural $u$, producing infinitely many encodings of the same division.

\subsection{A cumulative least-applicable-index circuit}
At time $t$, run the preceding gadget for every denominator $b_j$, obtaining $q_{t,j},\rho_{t,j},s_{t,j},u_{t,j},z_{t,j}$. Add natural variables $h_{t,0},\ldots,h_{t,r}$ and $x_{t,1},\ldots,x_{t,r}$ satisfying
\begin{align}
 h_{t,0}&=1,\label{cdc:wf:eq:frac-h0}\\
 h_{t,j}&=h_{t,j-1}(1-z_{t,j}),\label{cdc:wf:eq:frac-prefix}\\
 x_{t,j}&=h_{t,j-1}z_{t,j},\label{cdc:wf:eq:frac-first}\\
 \sum_{j=1}^r x_{t,j}&=1,\label{cdc:wf:eq:frac-nonhalt}\\
 n_{t+1}&=\sum_{j=1}^r a_jq_{t,j}x_{t,j}.\label{cdc:wf:eq:frac-output}
\end{align}
The prefix variables say that no applicable index has yet occurred. No separate Boolean equations for $h$ or $x$ are required: induction forces their values to be bits.

\begin{theorem}[Exact FRACTRAN compiler]\label{cdc:wf:thm:fractran}
Fix a reduced positive $r$-fraction program, positive integers $n_0^{\mathrm{in}}$ and $n^{\mathrm{out}}$, and $T\in\N$. There is an explicitly given polynomial of degree at most four with at most one natural zero; it has a zero if and only if the program executes exactly $T$ steps from $n_0^{\mathrm{in}}$ to $n^{\mathrm{out}}$. The construction uses exactly
\begin{align}
 V_{\mathrm{frac}}&=T(7r+2)+1,\label{cdc:wf:eq:frac-V}\\
 E_{\mathrm{frac}}&=T(8r+3)+2\label{cdc:wf:eq:frac-E}
\end{align}
variables and residuals. With $A=\max_j a_j$ and $B=\max_j b_j$, every solution coordinate is at most
\begin{equation}\label{cdc:wf:eq:frac-height}
 \max\{1,A^Tn_0^{\mathrm{in}},B\}.
\end{equation}
\end{theorem}
\begin{proof}
Use configurations $n_0,\ldots,n_T$ and endpoint equations $n_0=n_0^{\mathrm{in}}$, $n_T=n^{\mathrm{out}}$. At every $t<T$ impose the division gadgets and \eqref{cdc:wf:eq:frac-h0}--\eqref{cdc:wf:eq:frac-output}. Lemma~\ref{cdc:wf:lem:division} forces all the division data and all $z$ bits from $n_t$. The prefix recurrence then forces $h_{t,j}=1$ precisely when none of $1,\ldots,j$ is applicable. Consequently $x_{t,j}=1$ precisely at the first applicable index, if one exists; all other selectors are zero. Equation~\eqref{cdc:wf:eq:frac-nonhalt} rules out a halted step. The output equation is therefore exactly the FRACTRAN update.

Starting from the fixed positive input, this proves by induction that any solution is the actual run and all its auxiliaries are uniquely determined. In particular, each selected quotient is positive and every configuration remains positive. Conversely, an actual run supplies the uniquely determined division data and prefix bits and satisfies all equations. The final endpoint test is exact.

There are $T+1$ configuration variables. Each step adds $5r$ division variables, $r+1$ prefix variables, and $r$ selectors, giving $(T+1)+T(7r+1)$ variables. Each index contributes six division equations and two prefix/selection equations; each step adds an initial prefix, a nonhalt test, and an output equation. The two endpoint equations give \eqref{cdc:wf:eq:frac-E}. All residuals have degree at most two, including the output products $q_{t,j}x_{t,j}$.

Finally, each update satisfies $n_{t+1}\leq A n_t$, so $n_t\leq A^tn_0^{\mathrm{in}}$. Quotients are at most $n_t$, remainders and their slacks are below $B$, and all flags are bits. This proves \eqref{cdc:wf:eq:frac-height}. Sum the residual squares to obtain the polynomial.
\end{proof}

The theorem uses ``exactly $T$ steps'' to specify the length, not to require halting at time $T$. The next construction adds that requirement. For $T=0$, the basic certificate has one variable and two equations and asserts equality of the positive endpoints.

\subsection{First-halting certificates}
For each denominator $b_j$, add three natural variables $Q_j,U_j,V_j$ and two equations
\begin{align}
 n_T&=b_jQ_j+U_j+1,\label{cdc:wf:eq:halt-div}\\
 U_j+V_j&=b_j-2.\label{cdc:wf:eq:halt-positive}
\end{align}
They assert that division by $b_j$ has a strictly positive remainder. If $b_j=1$, the second equation has no natural solution, as it should: an integer fraction is always applicable.

\begin{corollary}[Unique first-halting witness]\label{cdc:wf:cor:halt}
Adding \eqref{cdc:wf:eq:halt-div}--\eqref{cdc:wf:eq:halt-positive} to the exact-$T$ certificate gives a polynomial with a unique natural zero precisely when the first halt occurs at step $T$ with the prescribed output. This adds $3r$ variables and $2r$ residuals, without increasing the degree bound.
\end{corollary}
\begin{proof}
For $b_j\geq2$, the equations force the remainder $U_j+1$ to lie in $\{1,\ldots,b_j-1\}$, so quotient and remainder are unique, as are $U_j,V_j$. They are solvable precisely when $b_j\nmid n_T$. All denominators must fail. Earlier times were explicitly required to have an applicable fraction by \eqref{cdc:wf:eq:frac-nonhalt}, proving first halting rather than mere terminality of an unrelated state.
\end{proof}

\begin{example}[Priority and a complete halt]
For the program $(3/2,5/3)$ on input $8$, the exact run is
\[
 8\longrightarrow12\longrightarrow18\longrightarrow27
 \longrightarrow45\longrightarrow75\longrightarrow125.
\]
At $12$ and $18$, both denominators divide the current integer, but the first fraction must be chosen. At $125$ neither is applicable. With $r=2$ and $T=6$, the basic certificate uses $97$ variables and $116$ residuals; the first-halting extension uses $103$ variables and $120$ residuals. Its largest witness coordinate is $125$. The exact example is supplied as \texttt{fractran\_first\_halt.json}.
\end{example}

The height estimate gives witness bit lengths
\[
 O\bigl(1+\log(n_0^{\mathrm{in}}+1)+T\log(A+1)+\log(B+1)\bigr).
\]
Thus the bounded certificate has polynomial bit-size in the input, the program, and a unary horizon. This quantitative statement is stronger than a variable count alone. It does not provide a bound on the first halting time as a computable function of an arbitrary program and input.

\section{Context-local Diophantine compilation of SKI rewriting\srctag{04}}\label{cdc:wf:sec:ski}
\subsection{A total natural coding of finite terms}
Let SKI terms be the finite binary application trees with leaves $S,K,I$, under the usual context-closed rules
\begin{equation}\label{cdc:wf:eq:ski-rules}
 Ix\to x,\qquad Kxy\to x,\qquad Sxyz\to xz(yz).
\end{equation}
Application associates to the left. We fix a concrete syntax; semantic equivalence of terms is not used as equality of their codes.

Define
\begin{align}
 \ulcorner S\urcorner&=0,&\ulcorner K\urcorner&=1,&\ulcorner I\urcorner&=2,\notag\\
 \App(\ell,r)&=3+\frac{(\ell+r)(\ell+r+1)}2+r.\label{cdc:wf:eq:app}
\end{align}
For natural $n,\ell,r$, let the application residual be
\begin{equation}\label{cdc:wf:eq:app-residual}
 C(n,\ell,r)=2n-6-(\ell+r)(\ell+r+1)-2r.
\end{equation}
It has degree two and integer coefficients.

\begin{lemma}[Unique natural tree decoding]\label{cdc:wf:lem:treecode}
Every natural number is the code of exactly one finite SKI term. For $n\geq3$, there is exactly one pair $(\ell,r)\in\N^2$ with $C(n,\ell,r)=0$, and both children are smaller than $n$. For $n<3$, no such pair exists.
\end{lemma}
\begin{proof}
Cantor's pairing function $\pi(\ell,r)=\frac12(\ell+r)(\ell+r+1)+r$ enumerates the diagonals $\ell+r=k$. On diagonal $k$, it takes the consecutive values $k(k+1)/2,\ldots,k(k+1)/2+k$, in increasing $r$. The next diagonal starts at the following integer. Thus it is a bijection $\N^2\to\N$.

Equation $C=0$ is exactly $n=3+\pi(\ell,r)$, giving the claimed unique pair and excluding $n<3$. If $k=\ell+r$, then $\pi(\ell,r)\geq k$ for $k\geq1$, and the assertion is also immediate for $k=0$. Hence $\ell,r<n$. Strong induction on $n$ now decodes a finite tree, with the three smaller codes interpreted as the leaves. The same induction and pairing uniqueness prove uniqueness.
\end{proof}

There is no separate well-formedness witness. Moreover, $\App(\ell,r)\geq\ell+r+3$ shows inductively that the decoded number of tree nodes is at most $n+1$. This is an effective bound, not a claim that decoding is polynomial in the \emph{binary} length of $n$.

\subsection{Prescribed schedules and context equations}
A reduction schedule is a finite list
\[
 \sigma=((R_0,\gamma_0),\ldots,(R_{T-1},\gamma_{T-1})),
 \qquad R_t\in\{S,K,I\},\quad\gamma_t\in\{L,R\}^*.
\]
The path $\gamma_t$ selects the redex in the term at time $t$. The schedule is \emph{external data}, not an existentially chosen variable-length path. Define its total context depth
\[
 H=\sum_{t<T}|\gamma_t|,
\]
and let $n_I,n_K,n_S$ count its three rule types, so $n_I+n_K+n_S=T$.

Use global configuration variables $c_0,\ldots,c_T$. During one descent through a context, suppose $u,v$ are the old and new whole-context codes. For a left descent introduce the old selected child $d$, the sibling $s$, and the new selected child $d'$, and impose
\begin{equation}\label{cdc:wf:eq:left-context}
 C(u,d,s)=0,\qquad C(v,d',s)=0.
\end{equation}
Continue the descent with the pair $(d,d')$. For a right descent use
\begin{equation}\label{cdc:wf:eq:right-context}
 C(u,s,d)=0,\qquad C(v,s,d')=0.
\end{equation}
Each context level therefore contributes three variables and two residuals. Sharing the sibling variable forces the off-path subtree to be unchanged, without inspecting its interior.

At the selected root, use the following equations, with $u,v$ now denoting old and new redex codes:
\begin{align}
 I:\quad & C(u,2,x)=0,\qquad v-x=0;\label{cdc:wf:eq:I-gadget}\\
 K:\quad & C(u,h,y)=0,\quad C(h,1,x)=0,\quad v-x=0;\label{cdc:wf:eq:K-gadget}\\
 S:\quad & C(u,h,z)=0,\quad C(h,g,y)=0,\quad C(g,0,x)=0,\notag\\
 & C(p,x,z)=0,\quad C(q,y,z)=0,\quad C(v,p,q)=0.\label{cdc:wf:eq:S-gadget}
\end{align}
The respective fresh-variable counts are $1$, $3$, and $7$; their residual counts are $2$, $3$, and $6$. Reusing the same argument code twice in the $S$ rule represents actual syntactic duplication, not an obligation to create two independently encoded copies.

\begin{theorem}[Context-local scheduled-reduction compiler]\label{cdc:wf:thm:ski}
Fix source and target codes $a,b\in\N$ and an external schedule $\sigma$. The context equations, root gadgets, and endpoint equations $c_0=a$, $c_T=b$ have exactly one natural solution if the schedule is valid and transforms $a$ into $b$, and none otherwise. Their sum of squares has degree at most four. The construction uses exactly
\begin{align}
 V_{\mathrm{SKI}}&=T+1+3H+n_I+3n_K+7n_S,\label{cdc:wf:eq:ski-V}\\
 E_{\mathrm{SKI}}&=2+2H+2n_I+3n_K+6n_S.\label{cdc:wf:eq:ski-E}
\end{align}
In particular, $V_{\mathrm{SKI}}\leq8T+3H+1$ and $E_{\mathrm{SKI}}\leq6T+2H+2$, independently of the number of nodes inside untouched subtrees.
\end{theorem}
\begin{proof}
Fix a source code for one step. At each level of the scheduled descent, Lemma~\ref{cdc:wf:lem:treecode} forces the old selected child and its sibling. If that level is a leaf, the corresponding application equation is impossible. At the selected root, the old-side equations force the argument codes and test the required leaf $I$, $K$, or $S$. Thus an incorrect redex kind also makes the system impossible.

For a correct redex, the root output is uniquely determined: it is $x$ for $I$ or $K$, and $\App(\App(x,z),\App(y,z))$ for $S$. All intermediate root variables are forced. Going back up the path, each new-side context equation uniquely forms an application with the already determined new child and old sibling. Hence the whole successor $v$ and all context auxiliaries are unique and equal to the scheduled contraction. This is a two-way argument: a valid contraction supplies exactly those values and satisfies all equations.

Induct on the steps beginning at $c_0=a$. A solution is precisely the prescribed valid reduction; its final code must equal $b$. Conversely, that reduction supplies its unique full witness. The counts follow by adding $T+1$ configurations, three variables and two equations per context level, the listed root-gadget costs, and two endpoint equations. Every residual has degree at most two, proving the quartic statement.
\end{proof}

\begin{example}[A two-step root computation]
The term $(S K I)I$ has code $635$. The root schedule $(S,\varepsilon),(K,\varepsilon)$ gives
\[
 (S K I)I\longrightarrow (K I)(I I)\longrightarrow I.
\]
Here $T=2$, $H=0$, $n_S=n_K=1$, so $V=13$ and $E=11$. The certificate \texttt{ski\_SKII.json} has a unique solution and maximum coordinate $635$.
\end{example}

\begin{example}[A nonempty context]
The term $I(IK)$ has code $91$. Contracting an $I$ at path $R$ and then at the root gives $I(IK)\to IK\to K$. Thus $T=2$, $H=1$, $n_I=2$, giving $V=8$ and $E=8$. The exact file is \texttt{ski\_context.json}.
\end{example}

\subsection{The unavoidable arithmetic-height qualification}
A small number of integer coordinates does not guarantee that those integers have small binary representations. The next bound makes the qualification explicit.

\begin{proposition}[A coarse height bound]\label{cdc:wf:prop:ski-height}
For a valid scheduled reduction with source code $a$, let $g=H+3n_S$. Every witness coordinate is at most a number $M$ satisfying
\begin{equation}\label{cdc:wf:eq:ski-height}
 \log_2(M+1)\leq
 2^g\bigl(\log_2(a+1)+4\bigr)-4.
\end{equation}
Thus the construction controls the number of coordinates linearly, but its general bit-height guarantee can be exponential in the amount of rebuilt context.
\end{proposition}
\begin{proof}
Decomposition into subterms never increases the current maximum code. An $I$ or $K$ contraction creates no new application at its root. An $S$ contraction creates three, and each context level creates one on reconstruction, so at most $g$ new applications are formed in temporal order.

For children at most $M$, formula~\eqref{cdc:wf:eq:app} gives
\[
 \App(\ell,r)+1\leq16(M+1)^2.
\]
This deliberately loose inequality follows directly from $\ell+r\leq2M$. If $L=\log_2(M+1)$, one new application increases the running maximum to a value whose $L$ is at most $2L+4$. Starting from $a$ and iterating that recurrence at most $g$ times gives \eqref{cdc:wf:eq:ski-height}. All witnesses are old subterms, reconstructed subterms, or configurations, so they are covered.
\end{proof}

Accordingly, Theorem~\ref{cdc:wf:thm:ski} is not a polynomial-time verification theorem in a measure that charges only for schedule length while treating large integer values as free. The code-based locality is genuine, but it trades structural expansion for potentially large arithmetic values. Better tree or DAG encodings are a concrete research target, not something supplied by the present pairing construction.

\subsection{Extension to finite first-order rewriting}
\begin{corollary}[Fixed-schedule finite-pattern rewriting]\label{cdc:wf:cor:rewrite}
Consider a finite first-order term signature with a disjoint, injective natural constructor coding, each constructor having a fixed-size unique-extension quadratic certificate. Fix finitely many rewrite rules in which every variable on the right occurs on the left, and fix a schedule of rule names and context paths. Then the scheduled reduction admits a unique-extension quadratic residual certificate whose size is linear in the total selected-context size plus the total sizes of the selected rule patterns and constructed right sides.
\end{corollary}
\begin{proof}
Decompose the source along the fixed context and then along the selected left pattern. Constructor injectivity makes all subterm codes unique. Repeated occurrences of a pattern variable use one shared code, thereby testing their equality; incompatible occurrences make the system unsatisfiable. Every right-side variable has now been determined. Construct the right side using the unique-extension constructor gadgets and rebuild the context. This is the same one-step uniqueness argument as in Theorem~\ref{cdc:wf:thm:ski}; iteration gives the result. Summing constant costs per visited pattern, constructor, and context node gives the stated bound.
\end{proof}

For a finite signature, such constructor codings can be obtained by tagged pairing, with fixed-arity tuples built from a fixed number of pairings. The corollary concerns ordinary first-order substitution at a fixed occurrence. It does not silently extend to binding syntax with capture-avoiding substitution, or identify different evaluation strategies of the lambda calculus.


\section{Local FRACTRAN and combinatory-logic interfaces, and bounded rewriting\srctag{02, 01, 05}}\label{cdc:sec:rewriting-misc}
\srcnote{\wtag{} Three subsections from the substrate surveys of manuscripts~02 (its Section~10.2), 01 (its Sections~9.3--9.4) and~05 (its Section~9.4). Manuscript~02's local FRACTRAN interface certifies one step and the terminal condition; the complete bounded run with first-applicable priority is manuscript~04's \cref{cdc:wf:thm:fractran} above.}

% ---- manuscript 02, lines 1278-1303 ----
\subsection{A local single-fold FRACTRAN interface}
For completeness, a second substrate illustrates how to avoid local
multiplicity before confronting the difficult global problem. Fix a
FRACTRAN list
\[
 \frac{p_1}{q_1},\ldots,\frac{p_m}{q_m},\qquad p_i,q_i>0.
\]
At positive integer state $n$, the first fraction producing an integer is applied.
For every $i$, introduce the \emph{unique} Euclidean quotient and remainder
\[
 p_i n=q_i a_i+r_i,\qquad 0\leq r_i<q_i.
\]
All quotients and remainders are evaluated, not just those before or at the
selected fraction. The applicable index $j$, if any, is characterized by
$r_j=0$ and $r_i\neq0$ for $i<j$; the next state is $a_j$.
The terminal condition is $r_i\neq0$ for every $i$.

These finite arithmetic conditions have unique quotient--remainder
witnesses and a unique selected index. The circuit construction therefore
gives single-fold local step and terminal-state certificates. This does
not imply a single-fold certificate for an unbounded FRACTRAN run.
FRACTRAN's role as a useful intermediate model in a full mechanized MRDP
proof is established by Larchey-Wendling and Forster
\cite{lf2022}; our discussion here isolates the distinct local
uniqueness issue.


% ---- manuscript 01, lines 638-671 ----
\subsection{A local polynomial adapter for combinatory logic}
The repository's combinatory-logic development makes careful distinctions between source syntax, runtime syntax, and evaluation strategy. In particular, its SKI-to-Iota simulation is not declared reduction-reflecting or fully abstract, and its exact Coq equivalence uses a weak-call-by-value lambda model distinct from the scoped contextual one \cite{repo-cl}. A new Diophantine adapter must not erase those distinctions.

There is nevertheless a useful elementary local construction. Define
\begin{equation}\label{cdc:ct:eq:pairing}
 \pi(a,b)=(a+b)^2+a,\qquad \operatorname{App}(a,b)=3+\pi(a,b).
\end{equation}
Give the three SKI leaves codes $0,1,2$. Codes of valid terms are generated inductively by these leaves and $\operatorname{App}$.

\begin{lemma}[Injective polynomial application coding]\label{cdc:ct:lem:pair}
The map $\pi:\N^2\to\N$ is injective. Moreover,
$\operatorname{App}(a,b)>\max(a,b)$, so validity of an alleged finite-term code is decidable by finite recursion on smaller codes.
\end{lemma}
\begin{proof}
For $s=a+b$, the value lies in the interval $[s^2,s^2+s]$. Consecutive such intervals are disjoint because $(s+1)^2>s^2+s$. Thus the value determines $s$, then determines $a=\pi(a,b)-s^2$, and then $b=s-a$. This proves injectivity. Since $s^2\ge s$ for natural $s$, the application code is greater than both children. To test validity, search finitely for a pair of children with the given application code and recursively test the unique pair if it exists; both recursive arguments are smaller.
\end{proof}
For the head rule $Sxyz\to xz(yz)$, fixed valid input code $u$ and proposed output code $v$ can be related by the six quadratic equations
\begin{align*}
 t_1&=\operatorname{App}(0,x),&
 t_2&=\operatorname{App}(t_1,y),&
 u&=\operatorname{App}(t_2,z),\\
 l&=\operatorname{App}(x,z),&
 r&=\operatorname{App}(y,z),&
 v&=\operatorname{App}(l,r).
\end{align*}
All seven natural witnesses $x,y,z,t_1,t_2,l,r$ are unique when they exist, because successive decompositions of the valid input determine the three arguments and then all constructions are functional. Their squared residuals give a quartic local rule. The same observation applies to the $K$ and $I$ head contractions.

Two restrictions must remain visible. First, the construction as stated has \emph{validity of the input term as a domain promise}. Applied indiscriminately to arbitrary natural codes, it could decompose a numeral with invalid subterms. Adding an arbitrary MRDP representation of validity need not preserve unique fibres. Second, contextual reduction needs a redex position and a surrounding context, and arbitrary-length reduction needs an additional history mechanism. The local rule therefore gives a concrete starting component, not a completed canonical representation of unrestricted SKI or Iota evaluation.

\subsection{What cannot be transported from a forward simulation alone}
Suppose a compiler maps source steps into nonempty target traces. This proves preservation of source execution, but by itself it does not prove that every target accepting trace is the image of a source trace. Nor does it prove that one source trace has a unique target realization. A parsimonious transport requires a decoding or reflection theorem and control of administrative target steps.

For ProveIt, this means that the existing operational embeddings are valuable reusable components, but a fibre-preserving Diophantine bridge must state the evaluation relation it preserves and prove the missing reflection and canonicalization obligations. The present guarded-translation construction avoids that ambiguity by defining its substrate and its represented histories explicitly.


% ---- manuscript 05, lines 724-730 ----
\subsection{Bounded word rewriting and explicit term graphs}
Fix a finite alphabet, a finite family of substring replacement rules, a horizon $T$, and an explicit maximum word length $B$. Store the word in $B$ cells with a canonical blank suffix and a uniquely encoded length. At each step choose a rule and a position using a fixed one-hot code, check the matched substring and size bound, and compute the uniquely padded next word. This is a finite deterministic verifier of the candidate choice history. Its circuit can be constructed by finite comparisons, multiplexers, and shifts of the $B$ cells. By \cref{cdc:thm:circuit} it has a quadratic witness-preserving arithmetic encoding.

The objects counted here are \emph{labelled derivations}, including the chosen rule and position, not merely their endpoint words. Distinct rules with identical replacements count separately unless the object definition first quotients them. For length-changing rules, the numerical positions of later redexes move; declaring two static position labels independent without accounting for that movement is unsound.

An analogous finite verifier is available for a specified bounded term-graph representation with deterministic allocation conventions. For tree representations of duplicating combinators, a bound on the number of reduction steps need not be a comparable bound on tree size. A size budget or a fully specified sharing representation is therefore essential. This article does not claim a linear-size arithmetic encoding of arbitrary tree expansion merely from a linear step bound.


\part{Other substrates}\label{cdc:part:substrates}

\paragraph{Source and scope.} \wtag{} Five manuscripts close with a survey of further computational substrates reached by their methods. This Part prints manuscript~05's (Boolean verifiers, counter and register programs with zero tests, finite cellular automata, exact time and padding), manuscript~04's (counter machines, stacks, Boolean circuits and cellular automata in a finite box, a compositional statement), manuscript~01's (counter machines with finite control, finite arithmetic and Boolean circuits), manuscript~02's (fixed-horizon counter machines, finite Boolean spacetime substrates) and manuscript~06's (what transfers without a further theorem, an abstract compiler contract, and why doubling, stacks and rewriting need more work). The rewriting subsections of these surveys are in Part~VII, and their boundary subsections in Part~IX.

\paragraph{Overlaps.} \wtag{} Boolean circuits are treated by four manuscripts; manuscript~05's \cref{cdc:thm:circuit} (affine equalities, convex quadratic, bounded real relaxation) is the strongest statement, and the circuit encodings of manuscripts~01, 02 and~04 are unique-extension encodings of the same kind with quadratic gate equations, hence quartic rather than quadratic polynomials. Counter machines appear in manuscripts~01, 02, 03 (Part~V), 04 and~05; manuscript~06 is the main treatment of counter programs (Parts~I--III), and the others are cross-referenced. Bounded cellular automata are compiled as finite space--time tableaux by manuscripts~01, 02, 03 (\cref{cdc:mem:sec:substrates}), 04 and~05; manuscript~07 does not compile Rule~110 diagrams directly and explains why its queue certificates do not transport to a $(T+1)$-variable space--time polynomial (the end of \cref{cdc:qc:sec:substrates}).

% ---- manuscript 05, lines 630-665 ----
\section{Other computational substrates\srctag{05}}
\label{cdc:sec:substrates}
The preceding theorem is most informative when the arithmetic front end is reusable. We give precise extensions and state which parts retain the Petri-net commutation theorem. In every case, size and time bounds refer to an explicitly unrolled finite computation, not a fixed-dimensional universal equation.

\subsection{A general Boolean verifier compilation lemma}
\begin{theorem}[Unique extension of a finite verifier]\label{cdc:thm:circuit}
Let $C$ be a finite deterministic Boolean circuit whose input strings encode candidate computational objects, with exactly one input encoding for each object. There is an effectively constructed system of affine integer equalities in natural variables, of size linear in the number of inputs and gates, whose full natural solutions are in bijection with the inputs accepted by $C$. Its sum-of-squares polynomial is convex quadratic, and its nonnegative real solution set is bounded.
\end{theorem}
\begin{proof}
For each input bit $p$, introduce its complement $\bar p$ and impose $p+\bar p=1$. This gives exactly one complement for each possible bit. Process gates topologically. A NOT gate is $p+y=1$. For an AND gate, use \eqref{cdc:eq:and}, with output $u$. For an OR gate, first compute $u=p\land q$ by that gadget and impose $y+u=p+q$. Constants have equations $y=0$ or $y=1$. Finally set the acceptance output to one.

On natural assignments the gate values are forced Boolean by induction, and all slacks are unique by \cref{cdc:lem:and}. Thus every accepted input has one full extension and every full solution decodes an accepted input. Over the reals the input bits are in $[0,1]$. The AND bounds keep the output in $[0,1]$, the NOT equation does the same, and the lower bound $u\ge p+q-1$ makes the OR output at most one. Its nonnegativity is part of its domain. All auxiliary slacks are bounded as in \cref{cdc:lem:and}. Hence the real feasible set is bounded. Scalarization proves the polynomial assertion.
\end{proof}

The unique input-encoding assumption must not be skipped. Padded records with unconstrained unused bits, multiple syntactic descriptions of the same object, or ignored nondeterministic choices would all change the solution count. The lemma is a compiler for a specified verifier and specified coding convention, not a theorem that arbitrary encodings are automatically parsimonious.

\subsection{Counter and register programs, including zero tests}
A finite counter program with natural registers and finite control can be unrolled using the raw variables of \cref{cdc:sec:raw}. Increment and decrement are consumption--production updates. To enforce a control edge with source $s(a)$ and target $t(a)$, add a control variable $q_i$ and equations
\[
q_i=\sum_a s(a)e_{i,a},\qquad
q_{i+1}=\sum_a t(a)e_{i,a},
\]
with prescribed initial and final control values. Equivalent one-hot control places can be used instead.

For a transition $a$ carrying a zero guard on register $p$, add
\begin{equation}\label{cdc:eq:zero-guard}
e_{i,a}x_{i,p}=0.
\end{equation}
This imposes the guard exactly on the selected branch. There is no auxiliary witness in the inactive branch, so uniqueness is unchanged. Nonzero/decrement guards are already expressed by consumption of a token. Thus a bounded counter-program execution has a direct parsimonious quartic sum-of-squares encoding, with finite control and every guard explicit.

For fixed numerical initial registers and a fixed horizon, the same extension can be made quadratic. Let $B_p=M_p+T\max_a v_{a,p}$, known from the raw induction. Replace \eqref{cdc:eq:zero-guard} by the affine equation
\begin{equation}\label{cdc:eq:bigM}
x_{i,p}+s_{i,a,p}=B_p(1-e_{i,a}),\qquad s_{i,a,p}\in\N.
\end{equation}
When $e_{i,a}=1$, it forces $x_{i,p}=0$. Otherwise it only enforces the already valid bound $x_{i,p}\le B_p$, and its slack is unique. The coefficient $B_p$ is a fixed integer for this instance. When the initial registers are free parameters, $B_pe_{i,a}$ is no longer jointly affine; the uniform-parameter presentation is then quartic unless another representation is supplied.


\srcnote{\wtag{} Manuscript~05 continues here with its exact commutation theorem for interval guards, printed in Part~I as \cref{cdc:thm:interval} (\cref{cdc:sec:interval}), and its corollary for guarded programs, printed in Part~II as \cref{cdc:cor:guarded} (\cref{cdc:sec:guarded}).}

% ---- manuscript 05, lines 715-723 ----
\subsection{Finite cellular automata}
Consider a one-dimensional deterministic cellular automaton on a ring of $W$ cells, alphabet size $q$, radius $r$, and fixed local rule. Give each cell a one-hot alphabet encoding, initially constrained by a sum equal to one. The local rule can be compiled by testing each of its $q^{2r+1}$ neighbourhood patterns with Boolean conjunctions and collecting patterns that produce each output symbol. Thus one time layer has a deterministic Boolean circuit of size
\[
O\bigl(W(r+1)q^{2r+1}\bigr),
\]
and $T$ layers have $T$ times this size. Apply \cref{cdc:thm:circuit}, treating the initial one-hot rows as the candidate data and their sums as affine constraints. Every valid initial configuration has exactly one space--time history and exactly one arithmetic extension. Constraints on the final pattern select its preimages without introducing extra encodings.

The same construction works on a finite interval with explicitly specified boundary conditions. For an infinite quiescent-background automaton, one must first specify and justify a sufficient finite space--time region; the ring theorem alone does not supply that justification. Also, synchronous local updates are not a sequence of freely swappable asynchronous transitions. A schedule quotient for an asynchronous automaton needs its own partial-map semantics.


\srcnote{\wtag{} Manuscript~05's subsection on bounded word rewriting and explicit term graphs, which stands here in the manuscript, is printed in Part~VII (\cref{cdc:sec:rewriting-misc}).}

% ---- manuscript 05, lines 731-735 ----
\subsection{Exact time, first halting, and padding}
All previous constructions use an exact horizon. For a deterministic program that halts in $t\le T$ steps, a unique length-$T$ record can be obtained by a distinguished stopped mode with a single mandatory padding action and frozen data. The first visit to stopped mode is the halting event; every later position is forced. The native time is then a uniquely determined observable.

Allowing arbitrary stuttering at arbitrary positions is different. It introduces multiple records of the same computation and can change the interaction with the commutation relation. Unused bits in a padded record must also be forced. A compiler for ``at most $T$'' needs a canonical padding convention as part of its correctness specification, not just an existential choice of a shorter length.


% ---- manuscript 04, lines 574-606 ----
\section{Other computational substrates through the same interface\srctag{04}}
The three main compilers illustrate separate difficulties: concurrent equivalence, priority, and context locality. The unique-extension viewpoint also gives a disciplined route to further substrates. The following constructions are elementary applications of Section~\ref{cdc:wf:sec:interfaces}, not claims of newly discovered universality.

\subsection{Counter machines and stacks}
An increment of a natural register has equation $R'=R+1$. A nonzero decrement has $R=R'+1$, with $R'\in\N$ supplying the guard. A zero branch has $R=0$ and the appropriate unchanged-register equations. Finite control and labeled instruction choices can be encoded by one-hot selectors and Lemma~\ref{cdc:wf:lem:activation}. A fixed horizon gives one solution per valid labeled instruction trace. In a deterministic program the guards, rather than an arbitrary disjunction, must force the chosen branch.

For a stack over digits $1,\ldots,B-1$, use zero for the empty stack and encode the top digit as the least significant digit. A push of digit $d$ and its inverse pop are
\[
 S'=BS+d,\qquad S=BS'+d,
\]
respectively. The selected digit is part of the labeled operation; it can alternatively be selected from the finite alphabet by a canonical circuit. Emptiness is $S=0$. Starting with a valid stack, these equations preserve the no-zero-digit representation and hence the intended stack semantics. For a run of $T$ pushes and pops, $S+1\leq B^T(S_0+1)$, so its code has $O(T\log B+\log(S_0+1))$ bits. Several stacks and finite control can be handled simultaneously.

A tape machine encoded by stacks additionally needs a convention for the current cell and implicit blank tails. Without that convention, distinct finite blank extensions would be distinct arithmetic objects representing the same tape. This is precisely the type of representational multiplicity that the compiler interface requires one to resolve explicitly.

\subsection{Boolean circuits and cellular automata in a finite box}
For Boolean inputs, AND and NOT gates have unique-extension equations
\[
 z=xy,\qquad z+x=1.
\]
Output Booleanity follows from the input Booleanity. Thus an acyclic circuit has one gate assignment per input, and every gate residual is quadratic. A finite-alphabet cellular automaton in a fixed finite box can be represented by a Boolean encoding of each cell and a fixed local update circuit. With $L$ cells, $T$ rounds, and $C$ gates per local update, a direct synchronous tableau uses $O(LTC)$ gate variables, together with the cell variables; all are bits. Boundary conditions and invalid finite-alphabet bit patterns must be specified, not inferred from an informal diagram.

Asynchronous local systems require more care. A proposed independence relation must guarantee that swapping actions preserves enabling and outcome. Disjoint complete read/write footprints are a sufficient condition. Disjoint write sets alone are not: an action may read a location changed by the other. Once a sound swap property is proved and the finite-action guards and updates have unique-extension certificates, the normal-form monitor of Section~\ref{cdc:wf:sec:monitor} applies independently of the particular state representation.

\subsection{A compositional statement}
\begin{proposition}[Trace-faithful composition]\label{cdc:wf:prop:composition}
Suppose a finitely labeled transition system has unique-extension polynomial certificates for each labeled step, and a fixed independence relation with the enabled-swap property. Suppose further that a length-$T$ operational tableau has no duplicate state encodings. Then adjoining the monitor equations yields a certificate whose solutions are in bijection with the independent-swap classes of the length-$T$ endpoint traces. Quadratization and the sum-of-squares construction preserve that bijection.
\end{proposition}
\begin{proof}
Use canonical branch activation to form the labeled-step tableau, ensuring that every labeled trace has exactly one auxiliary extension. Append the deterministic flag tableau for its action word. By the normal-form theorem, precisely the least words survive. The enabled-swap property makes the endpoint language trace-closed, so every class still has its least word. Unique-extension quadratization adds no multiplicity. This is the proof of Theorem~\ref{cdc:wf:thm:petri} with the specific token equations replaced by the assumed step certificates.
\end{proof}

The hypotheses are substantive. In particular, state-coding injectivity and canonical inactive variables are not automatic consequences of a simulation theorem or of ordinary Diophantine definability.


% ---- manuscript 01, lines 621-637 ----
\section{Other computational substrates and safe adapters\srctag{01}}\label{cdc:ct:sec:substrates}
\subsection{Counter machines with finite control and zero tests}
The guarded-translation model can describe a fixed finite counter program directly. Add a program-counter coordinate $q$ to the nonnegative registers. An instruction with source label $i$ and target label $j$ imposes the exact guard $q=i$ and displacement $j-i$ on that coordinate. An increment of register $t$ adds $+1$ to that register with lower guard zero. A positive decrement branch has lower guard one and displacement $-1$; its zero branch uses the exact guard zero and displacement zero. Other registers have zero displacement and unrestricted nonnegative guards.

For the program-counter coordinate, $\ell=i$ and $v=j-i$ satisfy~\eqref{cdc:ct:eq:guard-valid}, because $j\ge0$. For a decrement, its lower guard similarly guarantees natural output. Consequently the general envelope and history theorems apply to the resulting finite instruction alphabet.

This adapter is operational: the action labels retain which branch was taken, and each step agrees exactly with the specified counter instruction. It does not assert that all instructions commute. In a sequential program most useful adjacent instructions will be dependent, and the causal height may equal the ordinary execution length. Commuting-block acceleration applies only where the exact pair test certifies it.

To use first-halt semantics, the halt state has no outgoing action. Adding a free halt self-loop would create many finite histories after the same successful computation; neither determinism nor the earlier result would remove those genuinely different words. The classical use of counter or register machines in Diophantine arithmetization supplies the context for this adapter, not a new universality proof \cite{jm1984,lf2022}.

\subsection{Finite arithmetic and Boolean circuits}
A second elementary adapter concerns deterministic acyclic circuits. A gate equation such as $t=u+v$, $t=uv$, or a Boolean NAND equation $t=1-uv$ has degree at most two. Boolean wires can be constrained by $t(t-1)=0$. In a topological order the input values determine every wire value uniquely; hence the sum of squared gate residuals has degree at most four and a unique witness for the entire circuit evaluation.

More generally, any acyclic list of quadratic functional constraints with uniquely determined natural outputs can be lifted this way. The statement is proved by induction in the gate order, not by assuming that low degree alone forces uniqueness. A multiplication equation can have many factorizations if its two inputs are not already fixed.

This familiar circuit observation covers finite Boolean state-transition circuits and finite cellular-automaton cylinders with specified boundary conditions. It is a bounded circuit construction, not an unbounded cellular-automaton single-fold theorem. Boundary values, update order, and the exact cells represented must be part of the input specification. Unspecified boundaries would be additional semantic choices, not harmless auxiliary witnesses.


\srcnote{\wtag{} The remaining two subsections of this section of manuscript~01, a local polynomial adapter for combinatory logic and what cannot be transported from a forward simulation, are printed in Part~VII (\cref{cdc:sec:rewriting-misc}).}

% ---- manuscript 02, lines 1252-1277 ----
\section{Other substrates and the limit of the construction\srctag{02}}
\label{cdc:pt:sec:boundary}

\subsection{Fixed-horizon counter-machine certificates}
A fixed finite counter program has configurations consisting of a finite
control label and a fixed tuple of natural counters. Increment, decrement
with zero test, and control transfer are described by finite Boolean
combinations of polynomial equations and comparisons. The control and
zero-test branches are mutually exclusive for a deterministic program.

By Lemma~\ref{cdc:pt:lem:circuit}, each one-step relation has a single-fold
arithmetic certificate when its endpoint configurations are fixed.
For an \emph{externally fixed} horizon $T$, introduce the $T-1$
intermediate configurations and the one-step certificates. The number of
variables is $O(T)$ for a fixed program. A deterministic initial
configuration determines every intermediate configuration uniquely, so
this unrolled system is single-fold whenever the claimed run exists.
For a halting certificate, earlier configurations must also be required
not to be halting configurations, unless the chosen semantics already
prevents transitions from them.

This is a family of polynomials whose arity grows with $T$. It is not a
fixed polynomial with $T$ as a free parameter. The substantive gain in
Theorem~\ref{cdc:pt:thm:acceleration} is precisely the elimination of that
duration-dependent arity for a restricted, nontrivial class of updates.


\srcnote{\wtag{} Manuscript~02's local single-fold FRACTRAN interface, the second subsection of this section, is printed in Part~VII (\cref{cdc:sec:rewriting-misc}); its last three subsections, on the exponentiation boundary, universal first-halting computations and witness heights, are in Part~IX (\cref{cdc:pt:sec:limits}).}

% ---- manuscript 02, lines 1304-1318 ----
\subsection{Finite Boolean spacetime substrates}
A finite Boolean circuit has unique internal wire values once its inputs
are fixed. The same is true of a finite spacetime diagram obtained by
unrolling a deterministic cellular automaton with a fixed local rule,
specified boundary data, and an externally fixed time window. Boolean
gate equations turn such a diagram into a single-fold quartic system.
The witness count scales with the number of explicitly represented cells
and time steps.

Again, this elementary fact does not provide a fixed-dimensional
polynomial certificate for an unbounded computation. It clarifies a
useful design principle across substrates: local single-foldness is often
straightforward, whereas unique compression of variable-length execution
is the real obstacle.


% ---- manuscript 06, lines 1548-1632 ----
\section{Other computational substrates: transfers and limits\srctag{06}}
\label{cdc:cs:sec:substrates}

\subsection{What transfers without any additional theorem}

The same resource compiler applies to a fixed finite multiset-rewriting
system: a rule consumes the multiset $u_a$ and produces $v_a$. A discrete
chemical reaction network with stoichiometric integer updates is an instance
of these firing semantics. The claim concerns exact discrete firings and
token availability, not continuous concentrations, reaction rates, or
stochastic kinetics.

Finite control can be included by adding one place for each control state.
A transition consumes the token of its source control state and produces the
token of its target state. On markings with exactly one control token, this
reproduces the finite-state transition condition, while the remaining places
hold the original resources. The summary theorem then holds for the expanded
Petri system; semantic equivalence on all expanded markings is stronger than
one restricted to the one-control-token invariant. Accordingly the globally
computed independence relation can be conservative for the intended
invariant, but it is sound.

Bounded channels containing indistinguishable messages can also be expressed
by resource counts, with a complementary capacity place when needed. A FIFO
channel over distinct symbols is different: its state includes order, not
just multiplicity. The present proof does not encode that order by pretending
that each symbol count is sufficient.

\subsection{A reusable abstract compiler contract}

The resource argument suggests the following transfer principle.

\begin{proposition}[Local unique summaries imply global unique compilation]
\label[proposition]{cdc:cs:prop:abstract}
Suppose a class of word actions has finite-dimensional natural summaries with:
constant summaries for letters and the empty word; exact composition and
parameter-power relations given by fixed-size quadratic systems having a
unique output and unique local auxiliaries for each valid input summary;
and an endpoint-action relation with unique local witnesses. Assume the
relations preserve validity of summaries. Then every fixed parametrized
repetition expression in these actions has a single-fold quartic
representation of its endpoint relation, with a number of witnesses linear
in its syntax size, up to the fixed local summary dimensions.
\end{proposition}

\begin{proof}
Allocate all child summaries and local auxiliaries as in
\cref{cdc:cs:thm:compiler}. Structural induction on the syntax forces each valid
summary and every local auxiliary. The endpoint contract gives the final
zero-or-one solution condition. Every residual is quadratic, so its sum of
squares is quartic. Summing the fixed local allocations proves the size bound.
\end{proof}

This is a useful design contract, not a theorem that arbitrary machine models
satisfy its hypotheses. The difficult clause is usually the unique polynomial
representation of arbitrary powers, or the existence of a finite-dimensional
exact summary in the first place. If local residuals have degree at most
$D$ instead, the same argument yields degree at most $2D$.

\subsection{Why doubling, stacks, and rewriting require more work}

If a primitive command doubles a counter, then its $k$th power applied to
one produces $2^k$. A generalization of our arbitrary-power interface to that
command would in particular provide the corresponding witness-controlled
polynomial specification of exponentiation. Our translation-only resource
algebra avoids this issue; multiplication in \eqref{cdc:cs:eq:powergadget} multiplies
an externally supplied repetition count by a summary component, not the
runtime counter value by itself at each iteration. These are different
operations.

For stacks, a word of pushes and pops has cancellation information akin to
a required prefix and a residual suffix, but these summaries are words of
unbounded length, not fixed tuples of ordinary counters. Encoding the words
as integers then introduces base powers, length constraints, and canonical
coding questions. A proof of plain Diophantineness via MRDP would not preserve
single-foldness automatically.

Lambda terms and SKI or Iota reductions introduce another issue: sharing and
copying of syntax. A selected finite reduction sequence can be checked by an
interpreter, but an exact compressed summary of arbitrary iterated rewriting
need not have the affine principal-domain structure used here. Such models
are natural continuation targets for ProveIt, but they are not covered by the
current compiler. A reliable extension should state its own summary and
witness contracts before claiming transfer of our bounds.


\part{The boundary of fixed-arity compression}\label{cdc:part:boundary}

\paragraph{Source and scope.} \wtag{} Every construction of Parts~I--VIII is either a family of polynomials indexed by a horizon, a height, a length or a schedule, or a fixed-arity polynomial for a restricted relation. Every manuscript ends by explaining why this does not yield one fixed-arity single-fold or finite-fold polynomial for unbounded universal computation. This Part prints manuscript~05's boundary section first (base), then those of manuscripts~01, 02, 03, 04, 06 and~07.

\paragraph{Merged statements.} \wtag{} Four manuscripts prove that a single-fold (finite-fold) polynomial for universal halting exists exactly when every computably enumerable set has a single-fold (finite-fold) Diophantine representation: 05, 01, 03 and~04, with manuscript~05 adding a decidable trace-verifier form, manuscript~04 a degree-four form and manuscript~01 relations of any finite arity; manuscript~02 argues the first-halting form informally. The statement is printed once, as \cref{cdc:bd:thm:universal}, in manuscript~05's section; each other manuscript's proof stays in its own section, after a note giving its letters. Five manuscripts prove that no total computable function bounds a witness of a universal (or any undecidable) representation; that is printed once, as \cref{cdc:bd:prop:nobound}, in the most general form (manuscript~02's), with the proofs in place. Manuscript~06's finite-fold substrate equivalence (\cref{cdc:cs:thm:finitefold}) and its parameter-bound obstruction (\cref{cdc:cs:prop:nobound}) are different statements about separated one-counter schemes and are kept whole. Manuscript~07 proves neither theorem; it separates horizon-indexed from fixed-arity representations and remarks that its witness bound gives no computable bound on the unknown halting time.

\paragraph{What is and is not claimed.} \wtag{} The equivalence is a reduction, not a resolution: nothing in this report proves or refutes the existence of a single-fold or finite-fold universal polynomial. The single-fold \emph{exponential} Diophantine representation of every enumerable set (Jones--Matiyasevich 1984) is formalized in ProveIt (\cref{cdc:bd:rem:sfu}); the ordinary polynomial version is the open problem.

% ---- manuscript 05, lines 736-764 ----
\section{Fixed-dimensional compression and the finite-fold boundary\srctag{05}}
\label{cdc:sec:boundary}
The constructions above are a family of equations $P_T$, with a number of unknowns growing with $T$. A single polynomial with a fixed number of unknowns can represent reachability by MRDP or the repository's exact-iteration machinery, but preservation of witness multiplicity requires a separate argument. We make the precise logical strength of one natural proposed compression explicit.

\subsection{A decidable verifier with one successful history}
Fix a deterministic universal machine $U$, a canonical encoding of its programs and inputs, and a canonical encoding of finite configuration lists. Define
\[
\operatorname{Trace}_U(e,x,h)
\]
to mean that $h$ codes the complete execution of program $e$ on input $x$, beginning in its specified initial configuration and ending at the \emph{first} halting configuration. There is no padding and no irrelevant annotation. A run halting immediately has the singleton list containing its initial configuration.

This relation can be chosen primitive recursive. For specificity, a useful polynomial pairing injection is
\begin{equation}\label{cdc:eq:pair}
\pi(a,b)=(a+b)^2+a.
\end{equation}
It is injective: for $s=a+b$, its values lie in $[s^2,s^2+s]$, and these intervals for distinct $s$ are disjoint. Within an interval, $a$ determines $b$. Encode the empty list as zero and a nonempty list as
\[
\operatorname{cons}(a,h)=1+\pi(a,h).
\]
The tail is strictly smaller than the code. Decoding is by bounded search, rejecting numbers that are not valid codes, and checking the whole list requires only bounded iteration. Choose a standard machine encoding with primitive-recursive initial-configuration, one-step, and halting tests. The verifier checks all adjacent pairs and verifies that no earlier configuration halted. It is therefore primitive recursive.

Determinism and the absence of padding imply
\begin{equation}\label{cdc:eq:onehistory}
\text{for each }(e,x),\quad
\#\{h:\operatorname{Trace}_U(e,x,h)\}\le1.
\end{equation}
Existence of this unique history is equivalent to $U(e,x)$ halting. The history can be enormous; decidability of a proposed finite history says nothing about deciding whether some such history exists.

\subsection{An exact equivalence theorem}

\begin{theorem}[Universal halting and the single-fold and finite-fold problems; manuscripts 05, 01, 03 and~04]\label{cdc:bd:thm:universal}\label{cdc:thm:boundary}\label{cdc:ct:thm:universal}\label{cdc:mem:thm:compression}\label{cdc:wf:thm:universal}
Fix a universal programming system and let $K(e,x)$ mean that program~$e$ halts on input~$x$. The following statements are equivalent.
\begin{enumerate}[label=\textup{(\alph*)}]
\item Every computably enumerable subset of $\N$ has a finite-fold polynomial Diophantine representation over $\N$.
\item[\textup{(a$'$)}] Every computably enumerable relation of finite arity over $\N$ has a finite-fold polynomial Diophantine representation over $\N$.
\item There is one integer polynomial $H(e,x,z)$, with a fixed finite tuple $z$, representing universal halting $K(e,x)$ and having finitely many $z$-solutions for every fixed $(e,x)$.
\item[\textup{(b$'$)}] There is such a polynomial $H$ of degree at most four.
\item There is one integer polynomial $V(e,x,h,y)$, with a fixed finite tuple $y$, representing the decidable relation $\operatorname{Trace}_U$ of \cref{cdc:sec:boundary}, with finitely many $y$-solutions for every fixed $(e,x,h)$.
\end{enumerate}
The statements remain equivalent when ``finite-fold'' and ``finitely many'' are replaced throughout by ``single-fold'' and ``at most one''; in the single-fold form of~\textup{(b)}, the natural zero set of $H(e,x,\cdot)$ is a singleton when $K(e,x)$ holds and empty otherwise.
\end{theorem}

\paragraph{Sources, letters and proofs.} \wtag{} Manuscript~05 states (a)$\Leftrightarrow$(b)$\Leftrightarrow$(c) for a deterministic universal machine~$U$ with the canonical first-halting histories defined above; its proof follows. Manuscript~01 proves (a$'$)$\Leftrightarrow$(b) for a universal interpreter of register programs with first-halt semantics; its (i) is (a$'$) and its (ii) is~(b) (proof in \cref{cdc:ct:sec:boundary}). Manuscript~03 proves (a)$\Leftrightarrow$(b) for a universal deterministic interpreter of an acceptable enumeration; its (i) is (a) and its (ii) is~(b) (proof in \cref{cdc:mem:sec:boundary}). Manuscript~04 proves (a)$\Leftrightarrow$(b)$\Leftrightarrow$(b$'$) for an effective universal programming system with halting relation $H(e,x)$; its (i), (ii), (iii) are (a), (b), (b$'$) (proof in \cref{cdc:wf:sec:frontier}). The implications (a$'$)$\Rightarrow$(a) and (b$'$)$\Rightarrow$(b) are immediate, and (b)$\Rightarrow$(a$'$) is manuscript~01's argument; with the cycles proved by manuscripts~05 and~04, all five statements are equivalent. Manuscript~02 argues the first-halting form informally (\cref{cdc:pt:sec:firsthalt}), and manuscript~06 proves an equivalence for its own substrate (\cref{cdc:cs:thm:finitefold}). The theorem is a reduction: it does not say whether the statements are true. The formalized exponential counterpart is recorded in \cref{cdc:bd:rem:sfu}.

% ---- manuscript 05, lines 774-800 ----
\begin{proof}[Proof of (a)$\Leftrightarrow$(b)$\Leftrightarrow$(c), both versions (manuscript~05)]
\emph{(a)$\Rightarrow$(c).} The set of triple codes
\[
S=\{\pi(e,\pi(x,h)):\operatorname{Trace}_U(e,x,h)\}
\]
is computably enumerable, indeed decidable using the chosen bounded decoding. By (a), there is a polynomial $Q(n,y)$ with finite fibres representing $S$. Set
\[
V(e,x,h,y)=Q\bigl(\pi(e,\pi(x,h)),y\bigr).
\]
This is an integer polynomial because the code map is polynomial. Its natural $y$-fibre at a fixed triple is exactly the corresponding $Q$-fibre, so it has the required finiteness and represents exactly the trace relation.

\emph{(c)$\Rightarrow$(b).} Existentially quantify $h$ as one more witness and use $V$ itself as the halting polynomial. For fixed $(e,x)$, a nonhalting computation has no possible $h$, and a halting computation has exactly one possible $h$ by \eqref{cdc:eq:onehistory}. Above this unique value there are only finitely many $y$ by (c). Hence the full $(h,y)$-fibre is finite and is nonempty exactly for halting inputs.

\emph{(b)$\Rightarrow$(a).} For a computably enumerable set $A$, choose a program $e_A$ that halts on $x$ exactly when $x\in A$. Specializing the program parameter gives the polynomial $H(e_A,x,z)$ with finite fibres representing $A$.

Every step also preserves ``at most one'': polynomial substitution does not add witnesses, specialization does not add witnesses, and in the middle implication there is at most one admissible $h$ as well as at most one admissible $y$. This proves the single-fold version.
\end{proof}

This is a useful obstruction to an overambitious compiler contract. A fixed polynomial with finite fibres for a perfectly decidable first-halting trace verifier is already as strong as the general finite-fold conjecture. Calling the input a ``verified finite history'' does not make the required \emph{uniform arithmetic representation} harmless. The verifier is fixed while its history parameter is unbounded, and projecting that parameter yields universal halting.

The theorem is an equivalence and a design boundary, not a proof that such a polynomial cannot exist. The finite-fold and single-fold questions are the classical stronger representation problems discussed in \cite{mat2010,cco2024}. Nothing in our bounded family contradicts a positive resolution of either question. It simply does not furnish that resolution.

\subsection{Why obvious compression arguments lose multiplicity}
A common informal argument says: encode a long execution in one integer, represent the digit tests Diophantinely, and combine all equations. This can correctly prove \emph{existence}. It does not establish that a fixed history code has only one, or only finitely many, auxiliary arithmetic encodings. A representation of an exponentiation or sequence-decoding relation may have uncontrolled fibres. Each such relation needs its own finite-fold or single-fold theorem before this property can be composed.

Likewise, choosing ``the least witness'' at the semantic level is not a polynomial construction. Its minimality condition must itself be represented, and that representation can introduce further witnesses. The normal-form selection in \cref{cdc:sec:normal} avoids this particular problem because it is an explicit finite-state recurrence unrolled to a fixed length; the number of variables is allowed to grow.


\begin{proposition}[No computable witness-height bound; manuscripts 02, 01, 03, 04 and~05]\label{cdc:bd:prop:nobound}\label{cdc:prop:noheight}\label{cdc:ct:prop:no-bound}\label{cdc:pt:prop:height}\label{cdc:mem:prop:no-bound}\label{cdc:wf:thm:height-obstruction}
Let $A\subseteq\N$ be undecidable and suppose
\[
 x\in A\iff\exists\boldsymbol w\in\N^m\;P(x,\boldsymbol w)=0 .
\]
There is no total computable function $B(x)$ such that every positive instance has \emph{some} solution with all coordinates at most $B(x)$. In particular, let $P(e,x,z)$ be any fixed polynomial, with a fixed finite witness arity, representing universal halting, with no assumption on the multiplicity of its witnesses. There is no total computable function $B(e,x)$ with the following property: whenever $U(e,x)$ halts, there is a zero of $P(e,x,\cdot)$ whose every coordinate is at most $B(e,x)$.
\end{proposition}

\paragraph{Sources.} \wtag{} The first statement is manuscript~02's, for an arbitrary undecidable set; the second is the form stated by manuscripts~05 (whose proof follows), 01, 03 and~04, whose proofs are printed in their own sections below. The same finite-box search proves both. Manuscript~06 proves the corresponding statement for the parameters of a separated one-counter scheme (\cref{cdc:cs:prop:nobound}), and manuscript~07 records the point as a remark: its witness bound depends on~$T$ and supplies no computable bound on the unknown halting time (\cref{cdc:qc:sec:finite-fold}). As manuscripts~01, 02, 03, 04 and~05 each observe, finite-foldness alone gives no computable bound either, and the explicit bounds of the bounded constructions pose no conflict, since the horizon, height or log is given with the instance.

% ---- manuscript 05, lines 804-809 ----
\begin{proof}[Proof of \cref{cdc:bd:prop:nobound} for universal halting (manuscript~05)]
Such a function would decide halting. Compute $B(e,x)$, enumerate the finite box $\{0,\ldots,B(e,x)\}^{\dim z}$, and evaluate $P$ exactly at each tuple. A zero exists in the box exactly when the machine halts, by the assumed representation and bound. This contradicts the undecidability of universal halting.
\end{proof}

Even a bound on the \emph{least available} witness would be impossible in that sense; it is not necessary to bound every zero. Finite-foldness does not provide an effective bound of this kind. In contrast, the explicit bounds of \cref{cdc:thm:main} depend on the given horizon and therefore pose no conflict.


% ---- manuscript 01, lines 672-705 ----
\section{The boundary of unbounded canonicalization\srctag{01}}\label{cdc:ct:sec:boundary}
\subsection{Infinitely many histories can be semantically unavoidable}
\begin{proposition}[An infinite history fibre]\label{cdc:ct:prop:infinite-histories}
There is a two-action one-counter system with infinitely many distinct finite trace classes from a fixed state to itself. Therefore an unbounded representation in bijection with \emph{all} its finite histories cannot be finite-fold at that pair of endpoints.
\end{proposition}
\begin{proof}
Use an increment $a$ enabled at every natural marking and a decrement $b$ enabled at positive markings. They are dependent by Proposition~\ref{cdc:ct:prop:petri-commute}. Starting at zero, each word $(ab)^k$ executes and returns to zero. The words for distinct $k$ have different lengths and hence belong to different trace classes. A bijection with all these classes requires infinitely many natural roots at endpoints $(0,0)$.
\end{proof}
This does not disprove finite-fold representability of the endpoint relation. In this example an endpoint relation can have a very simple existential description. It proves only that ``preserve every finite history'' and ``have finitely many roots per endpoint'' are sometimes incompatible specifications.

\subsection{Semantic uniqueness does not canonicalize a coding base}
The next observation isolates a different source of multiplicity. It is an elementary warning about one coding architecture, not a no-go theorem for all possible codes.
\begin{proposition}[CRT base multiplicity]\label{cdc:ct:prop:crt}
Fix a nonempty finite list $a_0,\ldots,a_{t-1}\in\N$. There are infinitely many distinct pairs $(B,A_B)$ which encode that same list by residues modulo
\[
 m_i=1+(i+1)B,
\]
even if $A_B$ is required to be the least nonnegative simultaneous CRT representative.
\end{proposition}
\begin{proof}
Choose any sufficiently large positive multiple $B$ of $t!$ so that every $a_i<m_i$. For $i<j$, a common divisor $g$ of $m_i,m_j$ divides $(j-i)B$. But $\gcd(g,B)=1$, since $g$ divides $1+(i+1)B$. Thus $g$ divides $j-i$. Every positive divisor of $j-i$ divides $t!$ and hence $B$, so $g=1$. The moduli are pairwise coprime.

The Chinese remainder theorem therefore gives a unique $A_B$ with
$0\le A_B<\prod_i m_i$ and $A_B\equiv a_i\pmod{m_i}$ for all $i$. Infinitely many choices of $B$ satisfy the stated conditions. They give distinct pairs regardless of possible coincidences among the $A_B$ values.
\end{proof}
A canonical least residue fixes $A$ \emph{relative to} $B$, not $B$ itself. The example explains why extracting a unique run from a trace-coding equation is weaker than proving uniqueness of the complete arithmetic certificate. Our finite-height compiler avoids coding bases altogether; the accelerator avoids storing a list of repeated states.

\subsection{A sharp universal test for a single-fold compiler}
Let $U$ be a fixed deterministic universal interpreter for register programs, with first-halt semantics. Let
\[
 K(e,x)\iff U\text{ with program code }e\text{ and input }x\text{ eventually halts}.
\]
We use only the defining universal property: every computably enumerable relation has a recognizer simulated by some fixed program code. This is standard computability infrastructure, also underpinning the machine-to-Diophantine routes cited earlier \cite{lf2022,jm1984}.


\paragraph{Theorem of manuscript~01: ``Universal halting characterizes single-fold MRDP''.} \wtag{} Manuscript~01 states here, for the relation $K$ just defined, that (i) every computably enumerable relation of finite arity over $\N$ has a single-fold ordinary Diophantine representation if and only if (ii) one ordinary integer polynomial $P(e,x;w)$ with a fixed finite witness tuple has, at each $(e,x)$, a singleton natural zero set if $K(e,x)$ and an empty one otherwise; and the same with ``finite-fold'' and a nonempty finite zero set. This is (a$'$)$\Leftrightarrow$(b) of \cref{cdc:bd:thm:universal}, where the statement is printed once; manuscript~01's proof follows.

% ---- manuscript 01, lines 714-727 ----
\begin{proof}[Proof of (a$'$)$\Leftrightarrow$(b) of \cref{cdc:bd:thm:universal}, both versions (manuscript~01)]
The halting relation $K$ is computably enumerable, so (i) gives (ii).

Conversely, let $R\subseteq\N^k$ be computably enumerable. For $k\ge1$, iterating the polynomial injection $\pi$ of Lemma~\ref{cdc:ct:lem:pair} gives a polynomial injection $C_k:\N^k\to\N$ with a computable inverse on its image. A program can decode a valid encoded tuple and simulate a recognizer for $R$; on an invalid code it can diverge. Let $e_R$ be its fixed code. Then
\[
 R(a_1,\ldots,a_k)
 \iff K(e_R,C_k(a_1,\ldots,a_k))
 \iff\exists w\ P(e_R,C_k(a_1,\ldots,a_k);w)=0.
\]
Substituting a constant and a polynomial into an integer polynomial gives an integer polynomial. No new existential variable is introduced, so the entire fibre is exactly the fibre of $P$ at that substituted pair. A singleton remains a singleton and a finite set remains finite. Nullary relations are handled by a fixed dummy input. This proves the reverse implication in both versions.
\end{proof}
The theorem identifies the precise strength of a hoped-for universal fixed-arity canonical halting compiler. It would resolve the single-fold principle discussed by Matiyasevich, or its finite-fold version, rather than being a routine corollary of an ordinary exact-iteration interface \cite{mat2010,ccfo2019}. The explicit results of this report establish neither premise (ii). They instead provide fixed-arity compression for commuting blocks and a parsimonious family at bounded causal height.

\subsection{Witness bounds are a separate obstruction}

\paragraph{Proposition of manuscript~01: ``No computable global witness-height bound''.} \wtag{} For any integer polynomial $P(e,x;w)$ representing the halting relation~$K$, this is the universal-halting case of \cref{cdc:bd:prop:nobound}; manuscript~01's proof follows.

% ---- manuscript 01, lines 731-736 ----
\begin{proof}[Proof of \cref{cdc:bd:prop:nobound} for universal halting (manuscript~01)]
If such $B$ existed, compute it and test $P$ on the finite box
$\{0,\ldots,B(e,x)\}^{\dim w}$. A root would certify halting, and exhaustion without a root would certify nonhalting, contradicting undecidability. The argument uses only existence representation; it does not assume unique or finite fibres.
\end{proof}
There is no conflict with Proposition~\ref{cdc:ct:prop:witness-bound}: that proposition concerns a fixed bounded-height instance family, not a representation of unrestricted halting. Equally, finite-foldness by itself would not provide a computable bound on where its finitely many witnesses lie.


\section{The exponentiation boundary and universal first halting\srctag{02}}\label{cdc:pt:sec:limits}
\srcnote{\wtag{} The last three subsections of manuscript~02's section ``Other substrates and the limit of the construction'' (\cref{cdc:pt:sec:boundary}, Part~VIII).}

% ---- manuscript 02, lines 1319-1353 ----
\subsection{The exponentiation boundary}
The polynomial trajectory hypothesis cannot simply be dropped. Even the
single assignment
\[
 y'=2y,\qquad y_0=1
\]
has the exact-time graph $y=2^n$. This trajectory is exponential, not
polynomial in $n$.

\begin{proposition}[A precise implication of a stronger accelerator]
\label{cdc:pt:prop:exponential-boundary}
If a fixed-dimensional single-fold ordinary polynomial accelerator covered
this doubling loop uniformly in $n$, then the graph $y=2^n$ would have a
single-fold ordinary polynomial representation. Combined with the
classical single-fold exponential normal form, this would yield
single-fold ordinary polynomial representations of all enumerable sets.
\end{proposition}
\begin{proof}
The first assertion follows by specializing the accelerator's initial
state to $1$ and its guard to \emph{true}. For the second, use the
classical single-fold normal form in which exponentiation is isolated in
one relation $v=2^u$ \cite[p.~588]{cco2024}. Substitute the
hypothetical single-fold polynomial graph certificate for that relation.
The original outer tuple is unique, and the new graph certificate has a
unique extension for its fixed $(u,v)$. Thus the combined tuple remains
unique. Combine the resulting finite polynomial equations by a sum of
squares, or first lower their degrees by canonical arithmetic wires.
\end{proof}

The same issue arises for variable-base multiplication $y'=ay$, which
contains doubling as a specialization. The proposition is an implication,
not a claimed impossibility theorem. It identifies the strength of an
extension that the present paper has \emph{not} proved. The literature on
finite-fold representations explicitly distinguishes ordinary polynomial
and single-fold exponential results \cite{mat2010}.

\begin{remark}[\textup{[write]} The formalized exponential counterpart]\label{cdc:bd:rem:sfu}
The classical input used here --- every recursively enumerable set has a
single-fold unary exponential Diophantine representation
(Jones--Matiyasevich 1984) --- is formalized in ProveIt as
\path{JM1984.RM.re_sfu} and \path{JM1984.RM.re_single_equation}
(\path{Computability/HilbertTenthProblem/Lean/Diophantine/Paper1984/DPR.lean}),
with the single-fold predicate \path{JM1984.Exp.SFU}. The reductions in
this report are not formalized.
\end{remark}

% ---- manuscript 02, lines 1354-1373 ----

\subsection{Universal first-halting computations}\label{cdc:pt:sec:firsthalt}
A related conditional statement is useful when evaluating proposed
``substrate-independent'' encodings. Suppose one has a single-fold
polynomial representation of the input--time--output relation for the
\emph{first halt} of a fixed universal deterministic machine. For each
input program and data, there is at most one first-halting time and
output. Projecting those quantities, canonically encoding any signed
components, and specializing the program therefore gives a single-fold
polynomial representation of its halting set. Universality then gives
such a representation for every enumerable set.

Replacing single-fold by finite-fold gives the corresponding finite-fold
implication: there is at most one first-halting computation, and finitely
many auxiliary witnesses for it. Thus a universal finite-fold first-halt
compiler would settle a genuinely general representation question, rather
than merely repackage ordinary MRDP. Our accelerator is deliberately
restricted to supplied polynomial or fixed-residue polynomial trajectories.

\subsection{A height bound cannot be silently added to universality}

\paragraph{Proposition of manuscript~02: ``Witness-height obstruction''.} \wtag{} Manuscript~02's statement, for an arbitrary undecidable set $A\subseteq\N$ represented by $P(x,\boldsymbol w)$, is the first statement of \cref{cdc:bd:prop:nobound}, printed there once; its proof follows.

% ---- manuscript 02, lines 1383-1391 ----
\begin{proof}[Proof of \cref{cdc:bd:prop:nobound} (manuscript~02)]
Compute $B(x)$ and check the finitely many tuples in
$\{0,\ldots,B(x)\}^m$. A found zero certifies membership; if no zero is
found, the assumed bound certifies nonmembership. This would decide $A$.
\end{proof}
Finite-foldness alone does not give such a computable height bound. Our
height and bit-complexity bounds are possible because we treat a decidable
specified-input class, not an arbitrary universal halting predicate.


% ---- manuscript 03, lines 1255-1277 ----
\section{The exact boundary of uniform single-fold compression\srctag{03}}
\label{cdc:mem:sec:boundary}
\subsection{The quantifiers cannot be interchanged}
The construction proved so far has the form
\begin{equation}\label{cdc:mem:eq:family-order}
 \forall L\ \exists P_L\in\Z[X_1,\ldots,X_{3L},Y_1,\ldots,Y_{W_L}]
 \quad\text{with a canonical fiber at every fixed log.}
\end{equation}
The number of variables grows with $L$. For a fixed program we similarly have
one polynomial per horizon. This is not the assertion that one polynomial with
one fixed finite list of witnesses recognizes all halting computations while
preserving unique fibers.

For a set $A\subseteq\N$, a \emph{single-fold Diophantine representation} is a
polynomial $F(x,y_1,\ldots,y_s)$ such that $x\in A$ precisely when there is a
natural root in $y$, and there is at most one such tuple for every $x$. A
\emph{finite-fold representation} requires finitely many such tuples instead
of at most one. The general single-fold/finite-fold compression questions are
not settled by MRDP or by determinism of the represented computation
\cite{mat2010}. In particular, introducing a code for a unique trace does
not prove that the auxiliary variables used to decode and check that code are
unique.


\paragraph{Theorem of manuscript~03: ``A universal substrate is the exact bottleneck''.} \wtag{} Manuscript~03 fixes a universal deterministic interpreter $U(e,x)$ for an acceptable enumeration of partial computations and states that (i) every computably enumerable subset of $\N$ has a single-fold natural-number Diophantine representation if and only if (ii) one integer polynomial $F(e,x,Y_1,\ldots,Y_s)$, with fixed $s$, has natural roots in $Y$ that are unique when $U(e,x)$ halts and nonexistent otherwise; and the same with ``finite-fold/finitely many''. This is (a)$\Leftrightarrow$(b) of \cref{cdc:bd:thm:universal}, where the statement is printed once; manuscript~03's proof follows.

% ---- manuscript 03, lines 1292-1321 ----
\begin{proof}[Proof of (a)$\Leftrightarrow$(b) of \cref{cdc:bd:thm:universal}, both versions (manuscript~03)]
For (i)$\Rightarrow$(ii), encode pairs by the injective polynomial
\[
 \pi(e,x)=(e+x)^2+e.
\]
Its injectivity follows because for $n=e+x$ its possible values form the interval
$[n^2,n^2+n]$, and these intervals for successive $n$ are disjoint; within one
interval the value fixes $e$, then $x$. The set
\[
 K_\pi=\{\pi(e,x): U(e,x)\text{ halts}\}
\]
is computably enumerable. Apply (i) to obtain a single-fold polynomial
$G(z,Y)$ for that set and substitute $z=\pi(e,x)$. Injectivity ensures that the
result recognizes exactly the desired halting relation and preserves each
witness fiber.

For (ii)$\Rightarrow$(i), let $A$ be computably enumerable. By universality
there is a fixed program index $e_A$ such that $U(e_A,x)$ halts exactly when
$x\in A$. Specializing $e=e_A$ in $F$ gives a single-fold polynomial for $A$.
Both operations preserve finiteness of fibers as well, proving the finite-fold
version.
\end{proof}

The theorem is an elementary reduction exposing where an ambitious extension
would have to make genuinely new progress. It is not a solution of either
conjecture. The source determinism used in Theorem~\ref{cdc:mem:thm:ram} eliminates
multiple bounded executions, but says nothing by itself about the auxiliary
fibers of a fixed-variable arithmetization of arbitrary run length.

\subsection{An unavoidable witness-height obstruction}

\paragraph{Proposition of manuscript~03: ``No computable global witness bound''.} \wtag{} For any fixed polynomial $F(e,x,Y)$ representing the universal halting relation, without any assumption on multiplicity, and a candidate bound $h(e,x)$, this is the universal-halting case of \cref{cdc:bd:prop:nobound}; manuscript~03's proof follows.

% ---- manuscript 03, lines 1328-1349 ----
\begin{proof}[Proof of \cref{cdc:bd:prop:nobound} for universal halting (manuscript~03)]
Given $(e,x)$, compute $h(e,x)$ and test the finitely many natural tuples in
$[0,h(e,x)]^s$. Polynomial evaluation is decidable. By the supposed property,
the search finds a root exactly when $U(e,x)$ halts, deciding the universal
halting problem. This is impossible by the usual diagonal argument: a decider
for halting would allow a program that halts precisely when the decider predicts
that the same program on its own code does not halt.
\end{proof}

There is no conflict with Theorem~\ref{cdc:mem:thm:height}. That theorem takes a fixed
finite log as input, including the values already appearing in it. The machine
theorem fixes a horizon. Neither supplies a total computable bound on the
first halting time of an arbitrary universal computation. Even finite-fold
representability, by itself, would not make the maximum root computable.

A concrete formalization target is therefore a hierarchy of interfaces:
canonical local gadgets; canonical finite logs; canonical horizon-indexed
machine runs; ordinary fixed-variable Diophantine reachability by MRDP; and,
separately, the presently unproved finite-fold or single-fold upgrade of that
last interface. An implementation should not erase these distinctions behind
a common ``Diophantine'' label.


% ---- manuscript 04, lines 607-638 ----
\section{Unbounded computation and the fixed-arity frontier\srctag{04}}\label{cdc:wf:sec:frontier}
\subsection{What MRDP gives for the present substrates}
For any fixed effective substrate, the relation
\[
 \operatorname{Reach}(a,b)\quad\Longleftrightarrow\quad
 \text{some finite valid execution transforms $a$ into $b$}
\]
is computably enumerable: enumerate finite candidate histories and check each one. For SKI, every natural code is a finite term, and both redex checking and a finite contextual step are effective. For FRACTRAN and Petri nets the finite checks are explicit arithmetic procedures. MRDP therefore supplies a fixed finite-witness polynomial for each such reachability relation; program or net descriptions can also be included as natural inputs to a uniform interpreter.

This implication is classical, and the ProveIt exact-iteration interface is another route when the one-step relation has already been connected to its Diophantine API \cite{repo-trace}. Neither route automatically preserves the one-to-one correspondences constructed here.

\subsection{The difference between a family and one polynomial}
\begin{definition}[Finite-fold and single-fold]
A fixed polynomial $P(\mathbf a,\mathbf y)\in\Z[\mathbf a,\mathbf y]$ represents a relation $R$ if
\[
 R(\mathbf a)\iff\exists\mathbf y\in\N^m\ P(\mathbf a,\mathbf y)=0.
\]
It is finite-fold if the set of such $\mathbf y$ is finite for every fixed $\mathbf a$, and single-fold if that set has at most one element. The number $m$ is fixed independently of $\mathbf a$.
\end{definition}

The present FRACTRAN and scheduled-SKI polynomials are single-fold \emph{for each fixed horizon or schedule}. But their number of variables grows with that external data. They do not produce one fixed polynomial $P(T,a,b,\mathbf y)$ with fixed $m$ and the same uniqueness guarantee. The Petri-net polynomial intentionally permits multiple solutions, one for each distinct class, while giving exactly one solution per class.

The fixed-arity requirement is the essential frontier. Single-fold specification and its connection with eliminating exponentiation are discussed as substantive research questions by Cantone, Cuzziol, and Omodeo \cite{cco2024}. Nothing in the bounded constructions here supplies that missing elimination or settles the general finite-fold or single-fold conjectures.

\subsection{Three tempting but invalid shortcuts}
First, a deterministic machine has only one run, but a number-theoretic representation of that run may have many witnesses. The omitted equation $zu=0$ in the divisibility gadget gives a concrete example: the semantic branch remains unique, while the arithmetic fiber becomes infinite.

Second, representing a finite list by a variable base often introduces many choices of sufficiently large base and padding. One may still prove existential equivalence, but a uniqueness argument must specify and enforce a canonical base, length, and padding. Merely saying ``choose the least encoding'' is not a Diophantine construction: minimality introduces a statement excluding all smaller candidates, whose compilation and fiber behavior must themselves be justified.

Third, a forward operational simulation is not a reduction-reflecting equivalence. Extra target reductions can create arithmetic witnesses not corresponding to source computations. This matters for compiler chains involving different lambda-calculus evaluation strategies, or the distinction between a pure Iota source and a runtime syntax containing auxiliary combinators. The inspected repository explicitly distinguishes several of these semantics \cite{repo-cl}; our scheduled local theorem does not bypass those obligations.

\subsection{A precise universal reduction of the remaining problem}

\paragraph{Theorem of manuscript~04: ``Universal-halting scope equivalence''.} \wtag{} Manuscript~04 fixes an effective universal programming system with halting relation $H(e,x)$ and states, for either choice of ``finite-fold'' or ``single-fold'', the equivalence of (i) every computably enumerable subset of $\N$ has a representation of that kind; (ii) $H(e,x)$ has one by one fixed polynomial; (iii) $H(e,x)$ has one by a fixed polynomial of degree at most four. These are (a), (b) and (b$'$) of \cref{cdc:bd:thm:universal}, where the statement is printed once; manuscript~04's proof follows.

% ---- manuscript 04, lines 647-661 ----
\begin{proof}[Proof of (a)$\Leftrightarrow$(b)$\Leftrightarrow$(b$'$) of \cref{cdc:bd:thm:universal} (manuscript~04)]
Assume (i). Encode pairs by the Cantor pairing function $\pi$ from Section~\ref{cdc:wf:sec:ski}. The set $\{\pi(e,x):H(e,x)\}$ is computably enumerable, so it has the assumed representation. Adjoin a variable $z$ and the equation
\[
 2z-(e+x)(e+x+1)-2x=0.
\]
For fixed $e,x$, this determines exactly one natural $z$. Combining this equation with the representation of the paired set preserves its fiber cardinality and gives (ii).

Assume (ii). For a computably enumerable set $S$, choose the index $e_S$ of a program halting exactly on $S$. Substitute the fixed integer $e_S$ for the program coordinate. The remaining polynomial represents $S$ and has exactly the same fibers, proving (i).

Finally, (ii) implies (iii) by unique-extension quadratization and summing squares, while (iii) trivially implies (ii).
\end{proof}

This is a scope theorem, not a resolution. It identifies why a successful fixed-arity witness-faithful compiler for a universal substrate would be much stronger than any of the horizon-indexed families above. Restricting degree to four does not remove the main obstruction once an arbitrary finite number of variables is allowed.

\subsection{No computable global witness-height bound for halting}

\paragraph{Theorem of manuscript~04: ``Witness-height obstruction''.} \wtag{} For $P(e,x,\mathbf y)$ of fixed finite witness arity representing the halting relation of a universal programming system, this is the universal-halting case of \cref{cdc:bd:prop:nobound} (printed there as a proposition); manuscript~04's proof follows.

% ---- manuscript 04, lines 665-670 ----
\begin{proof}[Proof of \cref{cdc:bd:prop:nobound} for universal halting (manuscript~04)]
If such $B$ existed, compute it on $(e,x)$ and enumerate the finite box $\{0,\ldots,B(e,x)\}^m$. Evaluate the integer polynomial on every tuple. A zero would certify halting; if none existed, the assumed bound would certify nonhalting. This would decide the universal halting problem, a contradiction.
\end{proof}

Notice that the hypothesis asks only for \emph{one} small witness per positive instance, not for a bound on every witness. Finite-foldness does not imply such a computable bound: a finite set can have an extremely large maximum with no computably uniform control. Conversely, the bounded-horizon height estimates in this article are entirely compatible with the obstruction because the horizon is supplied as part of the compiled instance, and no computable universal bound on halting horizons is provided.


% ---- manuscript 06, lines 1448-1547 ----
\section{The exact finite-fold obstruction\srctag{06}}
\label{cdc:cs:sec:finitefold}

Single-foldness for fixed schedule parameters is strong enough to prevent
compiler-generated redundancy, but it cannot settle the general finite-fold
problem. The preceding normal forms turn this limitation into an exact
logical equivalence rather than a vague warning.

\begin{theorem}[Finite-fold substrate equivalence]
\label{cdc:cs:thm:finitefold}
For a relation $R\subseteq\N^a$, the following are equivalent:
\begin{enumerate}[label=\textnormal{(\roman*)},leftmargin=2em]
\item $R$ has a finite-fold Diophantine representation over $\N$.
\item There is a fixed separated one-counter scheme $E$ of depth at most four,
with input parameters $x$ and witness parameters $k$, such that
\[
x\in R\quad\Longleftrightarrow\quad
\exists k\in\N^b\;0\xrightarrow{\llbracket E\rrbracket_{(x,k)}}0,
\]
and for each fixed $x$ there are only finitely many successful parameter
tuples $k$.
\end{enumerate}
The same equivalence holds with ``finite-fold'' replaced by ``single-fold''
and ``finitely many'' replaced by ``at most one.''
\end{theorem}

\begin{proof}
From (i), apply \cref{cdc:cs:lem:quartic} to a witnessing polynomial. For each
original witness tuple there is exactly one tuple of circuit-gate values, so
finite or singleton fibers are preserved. Apply
\cref{cdc:cs:thm:degree_depth} to the resulting quartic. It preserves its parameter
solution set with no additional existential coordinates. This gives (ii)
with exactly the same fiber cardinalities.

Conversely, a separated depth-four scheme has acceptance polynomial
$A(x,k)-B(x,k)$ of degree at most four by
\cref{cdc:cs:thm:degree_depth}. Its roots are exactly the successful parameter
tuples, with no coding variables to eliminate. The assumed fiber cardinality
therefore gives precisely the required finite-fold or single-fold
Diophantine representation.
\end{proof}

The result is a normal-form equivalence, not a proof that every computably
enumerable relation satisfies these conditions. The general finite-fold
program and the role of single-fold specifications are discussed in
\cite{mat2010,cco2024}. We make no claim to settle either
conjecture. Instead, the theorem places their full difficulty inside a very
small substrate: one nonnegative counter, a block of increments followed by a
block of decrements, and four levels of globally shared repetition.

There is an analogous formulation using vector-separated depth-two schemes.
Use the quadratic-system version of \cref{cdc:cs:lem:quartic} in the forward
direction; in the reverse direction, sum the squares of the coordinate
balance equations. The same parameter tuples are preserved. Thus moving from
one counter to several trades nesting depth for simultaneous equations, not
for an automatic improvement in witness multiplicity.

\subsection{An exact accounting identity}

For fixed endpoints and input parameters, split the schedule parameters into
a prescribed tuple $x$ and a quantified tuple $k$. The compressed compiler
satisfies, with cardinalities understood in $\N\cup\{\infty\}$,
\begin{equation}
\label{cdc:cs:eq:fibercount}
\#\{(k,z):Q_E(m,n,x,k,z)=0\}
=
\#\{k:\llbracket E\rrbracket_{(x,k)}
\text{ meets the compiler's acceptance condition}\}.
\end{equation}
This follows by summing the zero-or-one fiber equality of
\cref{cdc:cs:thm:compiler}. There is no extra factor from arithmetic witnesses.
But the right side can still be infinite, and it counts parameter tuples,
not distinct words or traces. This is the exact point at which finite-fold
information must enter from somewhere other than the local compiler.

\subsection{No computable bound for a universal parameter search}

\begin{proposition}[Parameter-bound obstruction]
\label[proposition]{cdc:cs:prop:nobound}
Let $E$ be a fixed separated scheme whose existential parameter acceptance
set is an undecidable computably enumerable set of inputs. There is no total
computable function $B(x)$ such that every accepted input $x$ has a successful
witness parameter tuple all of whose coordinates are at most $B(x)$.
\end{proposition}

\begin{proof}
If such a function existed, compute $B(x)$ and test the finitely many witness
tuples in its box. For a separated scheme each test is an equality between
explicit integer polynomial values, hence decidable. An accepted input would
have a successful tuple in the box, and a rejected input would not. This would
decide the chosen undecidable set, a contradiction.
\end{proof}

The explicit height bound \eqref{cdc:cs:eq:height} does not conflict with this
proposition. It bounds \emph{auxiliary coordinates after the schedule
parameters are supplied}; it does not bound a successful choice of those
parameters from the original input. Even hypothetical single-fold
representations would not automatically provide a computable bound on their
unique witness.


% ---- manuscript 07, lines 786-815 ----
\section{The fixed-arity and finite-fold boundary\srctag{07}}\label{cdc:qc:sec:finite-fold}
\subsection{Three quantifiers that must not be interchanged}
For a fixed program and input, Theorem~\ref{cdc:qc:thm:stream} has the form
\[
 \text{for each }T\geq1,\quad
 \text{there is a polynomial }P_T\text{ in }T+1\text{ variables}.
\]
Unbounded halting is the existence of some successful horizon. This is not the same as one polynomial of fixed arity with $T$ as an ordinary input coordinate. Writing $\prod_{i<T}\ell_i$ when $T$ is an unknown does not perform that conversion; it merely writes a variable-length expression.

A fixed polynomial representation of a set $S\subseteq\N$ has the form
\[
 n\in S\quad\Longleftrightarrow\quad
 \exists\mathbf y\in\N^k\ P(n,\mathbf y)=0,
\]
where both $P$ and $k$ are chosen before $n$. It is finite-fold if each input has finitely many witnesses, and single-fold if each input has at most one. Research on obtaining finite-fold representations goes substantially beyond constructing a unique finite trace; Cantone, Cuzziol, and Omodeo give a concrete sufficient-condition approach to that problem \cite{cco-six}.

The constructions here solve the horizon-indexed witness problem. They do not provide the missing fixed-arity compression with controlled fibers. This is a statement of their exact logical strength, not a reason to abandon the stronger problem.

\subsection{What MRDP supplies and what uniqueness still requires}
An effective code for a finite queue computation has an effectively checkable transition relation. The set of successful codes, and the projection to halting inputs, can therefore be handled by the existing MRDP interface \cite{repo-mrdp}. That yields a fixed finite polynomial representation. It does not, by its stated interface alone, assert unique or finitely many auxiliary solutions.

A compression may introduce arithmetic witnesses for coding a list, extracting digits, evaluating powers, or eliminating bounded quantifiers. Even when the machine has one exact halting trace, these new witnesses must be separately controlled. Replacing a variable-length trace by one natural code does not prove that its polynomial validity predicate has unique witnesses.

Our treatment of first halting does remove one possible source of multiplicity: a halted run cannot be padded with fictitious empty-queue steps. The product guard and, independently, the local parity-of-scale obstruction both enforce that fact. This is useful but does not settle the multiplicity of a later fixed-arity arithmetization.

\subsection{No contradiction with bounded simulation}
All formulas are effectively constructible, and for a specified finite $T$ the actual deterministic run can also be simulated. Their purpose is not to outperform bounded simulation on this elementary decision problem. Their purpose is to preserve the run in a small, explicit algebraic witness and to support composition, formal verification, and comparison of computational substrates.

The witness bound depends on $T$. It supplies no computable bound on the unknown halting time from the input alone. Likewise, a circuit of size $O(T)$ is not necessarily small in the binary encoding length $\log T$. These distinctions are necessary when moving from exact finite certificates to unbounded computational questions.


\section{Implementation and validation}\label{cdc:sec:validation}
\wtag{} Every manuscript ships Python programs that import only the standard library, together with recorded outputs. Each subsection below is one manuscript's own account of what it implemented, what it tested and what the tests do not establish. The programs and data are shipped in \path{code/} and \path{data/} with the manuscript's prefix (\path{01-causal-traces-}, \path{02-poly-trajectories-}, \path{03-linear-memory-}, \path{04-witness-faithful-}, \path{05-trace-polytopes-}, \path{06-counter-schedules-}, \path{07-queue-causality-}); in the prose below the shipped names are printed. The displayed command lines are the manuscripts' own and use the delivered layout (for example \path{code/verify.py}, \path{examples/}, \path{tests/}), because the programs import one another and write their outputs by those names; the README gives rerun recipes that copy the shipped files into that layout, and warns which programs overwrite recorded outputs. The manuscripts' instructions for rebuilding their own PDFs do not apply to this report, whose build is described in the README. None of the finite checks is a proof; each manuscript says so, and its theorems rest on the proofs printed in Parts~I--IX.

% ---- manuscript 05, lines 810-877 ----
\subsection{Manuscript 05: Implementation and reproducible verification}
\label{cdc:sec:implementation}
\subsubsection{Reference implementation}
The companion package contains a Python implementation using only the standard library. It emits sparse residual polynomials and regards the final polynomial as their squared sum. Every natural assignment is checked using exact integer arithmetic. The one fractional-relaxation example uses exact rational numbers, not numerical tolerances.

The primary entry point is
\begin{lstlisting}[language=Python]
compile_certificate(net, initial, final, horizon,
                    form="quadratic", independence=None,
                    upper=None)
\end{lstlisting}
Here \code{form} can be \code{raw}, \code{quartic}, or \code{quadratic}. With no supplied independence relation the compiler computes the maximal sound relation. A supplied relation must be symmetric and contained in that maximal relation. The optional upper-guard matrix uses \code{None} for infinity and finite natural entries for upper bounds. One-step empty domains are rejected rather than silently retained; removing such transitions is a straightforward preprocessing step.

The returned certificate has methods for constructing the unique witness of a valid word, checking a whole natural assignment, decoding a zero, exporting the residual system as JSON, and expanding the polynomial for small instances. The code does not invoke an integer-programming solver or pretend to solve general Diophantine equations. It is a compiler and exact checker.

The source also implements the graph-realization construction of \cref{cdc:thm:graph}, the flag recognizer, a literal guarded-run simulator, and the circuit-to-safe-net reduction used in \cref{cdc:thm:counting}. The cellular-automaton and rewriting front ends in \cref{cdc:sec:substrates} are mathematical constructions; complete standalone translators for those front ends are not included.

\subsubsection{What was tested}
The full run of \texttt{python code/verify.py} produced \code{PASS}. The machine-readable record is \code{results/verification.json} \wtag{} (shipped as \path{data/05-trace-polytopes-verification.json}; the program is \path{code/05-trace-polytopes-verify.py}). The following are the exact finite scopes of the checks.
\begin{center}
\small
\begin{tabularx}{\textwidth}{@{}p{0.29\textwidth}X@{}}
\toprule
Check & Scope and result\\
\midrule
Petri demand criterion & 256 scalar transition pairs, arc weights $0$--$3$; 2,048 starting-marking comparisons.\\
Trace normal forms & All 75 independence graphs on ordered alphabets of sizes 1--4; all words of length at most 6: 358,509 words and 124,494 classes.\\
Boolean memory family & All 254 subset prefixes for $k=1,\ldots,7$, including distinguishing one-letter continuations.\\
Core certificates & 286 constructed full witnesses in the three presentations; 576 single-coordinate perturbations rejected.\\
Exhaustive zero search & All $2^{16}=65,536$ assignments in the full forced-height box of a tiny quadratic instance; exactly its two predicted zeros.\\
Real relaxation & Exact rational witness for \cref{cdc:ex:gap}, together with absence of an enabled integer step.\\
Circuit reduction & AND, OR, XOR, and a repeated-input example; accepting-run counts $1,3,2,1$, with safety and distinct Parikh vectors checked.\\
Ballot formula & 160 endpoint--horizon instances, with horizons through 7.\\
Interval guards & 729 scalar transition pairs and 6,561 marking comparisons; 45 bounded execution instances and 216 constructed witnesses.\\
\bottomrule
\end{tabularx}
\end{center}

The normal-form oracle is independent of the flag recurrence: it constructs the graph of legal adjacent swaps, computes its connected components with a disjoint-set data structure, and compares accepted words with the actual minimum of each component. The guarded commutation check compares literal two-step execution domains against the closed-form box criterion. The ballot formula provides a separate combinatorial count for the producer--idle--consumer examples.

For the core three-transition net, the computed endpoint examples are:
\begin{center}
\begin{tabular}{@{}rrrrrr@{}}
\toprule
$M$ & $N$ & $T$ & Raw runs & Classes & Quadratic $(V_2,R_2)$\\
\midrule
0&0&0&1&1&$(4,5)$\\
0&0&4&9&4&$(96,77)$\\
2&2&6&140&28&$(142,113)$\\
1&2&5&44&13&$(119,95)$\\
\bottomrule
\end{tabular}
\end{center}

A finite test does not prove the universal theorems. In particular, changing one coordinate of a constructed witness is not an exhaustive uniqueness proof, and the full box search concerns only its explicitly stated tiny instance. Universal correctness and uniqueness are supplied by the preceding mathematical arguments. The tests serve as independent checks of indexing, signs, degree accounting, boundary cases, and the executable realization of those arguments.

\subsubsection{Reproduction and polynomial format}
From the package directory \wtag{} (in the delivered layout; see the README, which also explains why the shipped \path{code/05-trace-polytopes-build.py} must not be run in the report directory), run
\begin{lstlisting}
python code/verify.py
pdflatex -interaction=nonstopmode -halt-on-error article.tex
pdflatex -interaction=nonstopmode -halt-on-error article.tex
pdflatex -interaction=nonstopmode -halt-on-error article.tex
\end{lstlisting}
The verification program should be run without Python's optimization flag, since its test assertions are part of the checks. The ordinary compiler's witness validation is an explicit runtime check. Python 3.10 or later is required; no third-party Python package or network access is needed for verification. A standard TeX installation with the packages listed in the source is needed to rebuild the PDF.

Each exported JSON residual is a list of coefficient--monomial pairs. A monomial is a list of variable names, repeated when a power is present. The represented polynomial is the sum of the squares of the residuals, not the sum of the residuals. The domain field explicitly says ``nonnegative integers.'' A separate small text file contains a fully expanded quadratic polynomial as a format sanity check.


% ---- manuscript 01, lines 737-801 ----
\subsection{Manuscript 01: Implementation, validation, and repository integration}\label{cdc:ct:sec:implementation}
\subsubsection{The reference compiler}
The accompanying \code{code/01-causal-traces-causal_diophantine.py} uses only the Python standard library. It contains an exact integer operational semantics, the composition-box test, the envelope, Cartier--Foata normalization, sparse integer polynomials, and four compiler entry points:
\begin{center}
\begin{tabular}{p{0.35\textwidth}p{0.52\textwidth}}
\toprule
Entry point & Mathematical contract\\\midrule
\code{compile_accelerator} & Theorem~\ref{cdc:ct:thm:accelerator}; all serializations, also some under active commutation.\\
\code{compile_history} & Theorem~\ref{cdc:ct:thm:history}; one root for each bounded-height executable trace.\\
\code{compile_blocks} & Theorem~\ref{cdc:ct:thm:block-composition}; fixed ordered commuting phases, optionally existential counts.\\
\code{compile_word_power} & Proposition~\ref{cdc:ct:prop:word-power}; arbitrary repetition of one fixed, possibly noncommuting word.\\
\bottomrule
\end{tabular}
\end{center}
A separate \code{code/01-causal-traces-verify_exports.py} checker reads the exported JSON without importing the compiler. It validates both natural certificates and compares the expanded polynomial with the residual-square expression on fifty seeded assignments per file; these comparisons are finite checks, not a symbolic identity proof.

The matching \code{complete_*} functions compute the canonical witness when a given semantic object is valid. They do not search for arbitrary polynomial roots. A \code{System} can check an assignment, evaluate the sum of squares, expand it, or export all coefficients and an optional verified certificate to JSON. The exporter refuses to attach an invalid certificate.

Natural-number domains are part of the specification. Polynomial evaluation uses ordinary signed integer arithmetic, and all numerical work is exact. The implementation deliberately uses no floating-point feasibility test, relaxed solver, modular substitute for integer equality, or fixed-width machine counter.

\subsubsection{What was actually tested}
Run
\begin{quote}\code{python code/verify.py}\end{quote}
from a machine with Python 3.10 or later. The recorded run used Python 3.13.5 and pseudorandom seed 20260930. The tests are reproducible and fail by raising an exception; they write their exact counts and examples to \code{data/} and \code{examples/}. \wtag{} (Shipped as \path{data/01-causal-traces-*}; the command runs \path{code/01-causal-traces-verify.py} in a copy with the delivered layout, as the README describes.)

\begin{table}[ht]
\centering\small
\begin{tabular}{p{0.64\textwidth}r}
\toprule
Finite check & Count\\\midrule
Pairs of scalar guarded actions, box test versus direct execution & 5,184\\
Petri pairs, closed sign classification versus box criterion & 1,296\\
State--multiset instances for the exact envelope & 9,600\\
Individual serialized executions within those instances & 102,768\\
Words checked against adjacent-swap connected components & 8,744\\
Boolean three-layer matrices across all eight independence graphs & 4,096\\
Accepted canonical matrices in that enumeration & 443\\
Random finite-history instances & 2,400\\
Compiler arity/equation-count instances & 200\\
Changed auxiliary coordinates correctly rejected & 43,804\\
Tiny history assignments exhaustively tested & 1,041\\
Tiny accelerator assignments exhaustively tested & 6,561\\
Fixed commuting-phase composition instances & 512\\
Serializations of the worked six-event trace & 60\\
Fixed-word power instances, including empty macro domains & 16,380\\
\bottomrule
\end{tabular}
\caption{Exact finite implementation checks. Counts are not a claim of formal verification.}\label{cdc:ct:tab:tests}
\end{table}

The domain-box tests use lower guards from zero through three, every displacement from the negative lower guard through two, and four upper-bound choices; direct execution is tested from zero through 25. The Petri classification uses every consumption and production value from zero through five. Random envelope instances have one through three coordinates, one through four action labels, multiplicities at most two, and total event count at most six. Every multiset serialization in each selected instance is executed independently of the envelope formula.

For the normal form, all eight independence graphs on three labels and every word of length at most six are checked. An independent union--find calculation of adjacent-swap components is compared with the computed layers. The tests also enumerate all $2^9$ Boolean three-layer matrices for each graph, using read-only actions so that every chosen independence subrelation is operationally valid.

The mutation checks change auxiliary witnesses while holding the semantic layer bits fixed. A different layer matrix could represent a different valid trace, so it would be incorrect to demand rejection of every change to a layer bit. In the accelerator, all witness coordinates are auxiliary relative to fixed $(x,y,n)$ and are tested accordingly.

The tiny exhaustive natural-root examples are a single read action at endpoint one with height bounds zero, one, and two. Their root counts are respectively one, two, and three. Every possible witness lies in the tested Boolean box for these examples, as follows directly from the residual formulas. A single decrement accelerator is similarly exhaustively checked for $x,y,n\le2$, with witness values zero through two. These finite tests exercise root multiplicities, not only production of successful certificates.

For fixed-word powers, every word of length at most five over a guarded increment, a decrement, and a zero test is checked at nine initial markings and five repetition counts. This includes impossible loop bodies and the zero-repetition exception.

The large worked accelerator is also evaluated as an expanded polynomial: its 76-monomial quartic evaluates to exactly zero at the supplied eleven-witness assignment. The 32-witness history example passes all 48 residuals. Full results, including the test ranges encoded in the script, accompany the paper.

\subsubsection{What the tests do not establish}
The tests are not a proof for all dimensions, counts, guards, or height bounds; those claims rest on the mathematical arguments above. They do not constitute a Lean kernel check, an independent source-level verification of the Python implementation, or a proof that the arity bounds are optimal. This report did not rebuild the ProveIt repository or rerun its published axiom audits. Repository facts are taken from the pinned source and documentation inspected for this study.


% ---- manuscript 02, lines 1392-1499 ----
\subsection{Manuscript 02: Executable artifacts and validation}
\label{cdc:pt:sec:implementation}

\subsubsection{What is implemented}
The supplementary files implement the core bounded polynomial theorem \wtag{} (shipped as \path{code/02-poly-trajectories-*.py} and \path{data/02-poly-trajectories-*.json} under the delivered names below):
\begin{description}[style=nextline]
\item[\code{certificates.py}]
Exact integer forward differences, binary-search generation of the
canonical sign tower, terminal padding, and an independent local
endpoint checker. A brute-force sign-run routine is used only in small
regression tests, never by the main generator or checker.
\item[\code{quartic_compiler.py}]
A reference compiler for the fixed-degree sign-table conditions and the
optional nonnegativity requirement. It exports every quadratic residual,
all variable names, and a sample complete assignment. Its symbolic
structure depends only on the nominal degree and the chosen final
assertion, not on coefficient values or the horizon.
\item[\code{run_tests.py} and \code{test_report.json}]
A deterministic regression suite and its actual output, with seed
$20260930$. The output includes operation counts, finite uniqueness tests,
and compiler structure fingerprints.
\item[\code{quartic_example.json}]
The exact degree-two quartic template in sum-of-squares form, together with
a satisfying assignment for $p(t)=(t-2)^2$, $T=6$.
\item[\code{verify_export.py}]
A separate sparse-polynomial evaluator that reads the exported JSON,
checks natural witness domains and quadratic residual degrees, and
calculates every residual without importing the compiler.
\end{description}
The higher-level Boolean-guard, arbitrary trajectory, and quasi-polynomial
front ends are proved mathematically but are not packaged as automatic
source-language compilers in this release. This distinction matters:
working core code is not the same as a complete program-analysis tool.

\subsubsection{Tests actually performed}
The included run produced the following results, with all assertions
passing:
\begin{center}
\begin{tabular}{lr}
\toprule
Check & Count\\
\midrule
Random integer polynomials, nominal degrees $0$--$8$ & $1{,}620$\\
Additional explicit edge cases & $8$\\
Degree-two coefficient triples in $\{-1,0,1\}^3$ & $27$\\
Candidate sign towers exhaustively checked on $[0,2]$ & $59{,}049$\\
Incorrect individual tag mutations rejected & $68$\\
Compiled polynomial instances, degrees $0$--$4$ & $50$\\
Negative positivity instances among those compiled & $33$\\
Nonlinear closed-form state comparisons & $2{,}500$\\
\bottomrule
\end{tabular}
\end{center}
For each of the $27$ small coefficient triples, all possible sign strings
at the first two levels and all three constant top tags were considered.
The checker accepted exactly one tower, equal to the generated tower.
This tests uniqueness in a finite domain; the general uniqueness proof is
Theorem~\ref{cdc:pt:thm:tower}.

For each compiled instance, the checker confirmed that every exported
residual has degree at most two, every sample existential coordinate is
natural, and the squared-sum equation vanishes exactly on the positive
instances among the tested assignments. Ten differently instantiated
examples at each nominal degree produced identical symbolic polynomial
structures. False instances fail the required positivity residuals even
though their full sign tables remain valid.

\subsubsection{Concrete compiler size}
Table~\ref{cdc:pt:tab:compiler} reports the unoptimized reference implementation.
It is not a lower bound or a record claim.

\begin{table}[H]
\centering
\begin{tabular}{rrrrr}
\toprule
$d$ & Primary coordinates & Auxiliary coordinates & Total witnesses & Residuals\\
\midrule
0 & 1 & 24 & 25 & 42\\
1 & 6 & 202 & 208 & 312\\
2 & 15 & 834 & 849 & 1218\\
3 & 28 & 2192 & 2220 & 3132\\
4 & 45 & 4540 & 4585 & 6414\\
\bottomrule
\end{tabular}
\caption{Actual variable and residual counts for the supplied compiler.
The final equation is the sum of the squares of all residuals.}
\label{cdc:pt:tab:compiler}
\end{table}

These constants are intentionally conservative. Repeated comparisons,
zero-test slack variables, and straightforward Boolean wiring create many
auxiliaries. Eliminating a redundant gate is safe only when it preserves
both logical equivalence and uniqueness of the remaining witness.

\subsubsection{Validation limits}
The numerical tests do not prove soundness on all integer assignments,
completeness for arbitrary horizons, or historical novelty. Those are
separate questions: soundness and completeness are addressed by the
mathematical proofs, and novelty requires further literature comparison.
No Lean or Coq formalization of the new construction was compiled here.
The ProveIt repository was inspected as a source, not independently rebuilt.

The independent JSON checker certifies an exported \emph{assignment}'s
arithmetic validity. It is not a solver and does not verify that the
compiler correctly represents the intended relation for all inputs.
That higher-level correspondence is exactly the theorem a future
proof-assistant implementation should establish.


% ---- manuscript 03, lines 1350-1458 ----
\subsection{Manuscript 03: A hand-checkable example and executable validation}\label{cdc:mem:sec:validation}
\subsubsection{Four events}
Consider the zero-initialized memory log
\begin{center}
\begin{tabular}{@{}rrrrl@{}}
\toprule
$t$ & $a_t$ & $w_t$ & $v_t$ & Interpretation \\
\midrule
0 & 7 & 0 & 0 & Read an untouched address.\\
1 & 2 & 1 & 9 & Write 9 at address 2.\\
2 & 7 & 1 & 4 & Write 4 at address 7.\\
3 & 2 & 0 & 9 & Read back 9 at address 2.\\
\bottomrule
\end{tabular}
\end{center}
Here $N=4$ and $B=5$. The chronological keys are $40,16,42,18$. The sorted
records $(\alpha,\kappa,w,v)$ are
\[
 (3,16,1,9),\quad(3,18,0,9),\quad(8,40,0,0),\quad(8,42,1,4).
\]
The adjacent address differences are $0,5,0$. Their zero-test pairs $(z,h)$ are
$(1,0),(0,4),(1,0)$, and their previous-value coordinates are $9,0,0$.
The first sorted record is a write. The next two records are reads returning
9 and 0 respectively; the final write can store any natural value. All memory
checks pass.

There are $C_4=6$ comparators, $W_4=77$ auxiliaries, and $R_4=85$ residuals,
besides the 12 free event parameters. The supplied export contains the actual
expanded polynomial: degree four, with 587 distinct nonzero monomials.
Changing the final read value from 9 to 8 makes the appropriate read residual
nonzero. Because the preceding circuit has a unique assignment, no alternate
routing or slack assignment can repair that invalid log.

\subsubsection{What the implementation does}
The file \repo{code/03-linear-memory-canonical_memory.py} is a dependency-free Python
implementation of the memory theorem, using exact integers and sparse
polynomials. Its representation stores each monomial as a sorted tuple of
variable indices. It rejects residuals of degree greater than two. The compiler
creates symbolic event coordinates, so its polynomial structure is determined
by $L$, not by the particular numerical log used to generate the witness.

The JSON export includes the natural-domain convention, every parameter and
witness name, the generated assignment, and every quadratic residual. With
\texttt{--expanded}, it also exports the actual sum-of-squares polynomial.
Canonical values are serialized as decimal strings to avoid accidental loss
of precision in consumers that treat JSON numbers as floating point.
Polynomial coefficients in this backend are small ordinary JSON integers.
The independent semantic checker uses a dictionary as memory and compares
reads with its current contents.

Both memory backends are implemented in this archive;
\repo{code/03-linear-memory-linear_memory.py} provides the large-base alternative. The RAM, graph,
counting, and probability results are mathematical constructions proved in the
preceding sections; a complete machine-language frontend is not supplied.
The probability test checks the stated weight formula, not an implemented
probabilistic RAM compiler. No new Lean or Rocq proof artifact is claimed.

\subsubsection{Executed checks}
The complete test driver is \repo{code/03-linear-memory-test_certificates.py}; its recorded
results are in \repo{data/03-linear-memory-test_results.json}. The following tests were
actually executed in this environment:
\begin{center}
\begin{tabularx}{\textwidth}{@{}p{0.31\textwidth}X@{}}
\toprule
Test family & Result \\
\midrule
Comparison and zero tests & Exhaustive candidate witnesses for 81 ordered
input pairs and 21 zero-test inputs; exactly the stated witnesses survived.\\
Sorting networks & 105,883 cases: every permutation at widths 1, 2, 4, and 8,
and all binary inputs at width 16.\\
Small memory logs & All 585 logs of lengths 0 through 3 over addresses,
write flags, and values in $\{0,1\}$: 259 valid and 326 invalid, in exact
agreement with the independent semantic checker.\\
Large integers & 150 generated valid logs, with addresses up to 160 bits and
values up to 200 bits; 143 deliberately corrupted read logs were rejected.\\
Auxiliary corruption & All 273 single-coordinate increments of the auxiliary
witness for a six-event valid log were rejected.\\
Expansion & The four-event quartic has 587 monomials and degree four;
20 additional assignments confirmed equality with the residual sum of squares.\\
Path weights & For each horizon 1 through 12, the stop-on-first-one weight sum
was exactly $2^T-1$.\\
\bottomrule
\end{tabularx}
\end{center}
The six-event test has $N=8$, 24 comparators, 273 witnesses, and 287 residuals,
in agreement with the exact formulas. All recorded tests passed.

Finite testing is not a proof that arbitrary natural assignments have no
unexpected roots. That conclusion rests on the forcing argument in
Theorem~\ref{cdc:mem:thm:memory}. The tests instead check correspondence between the
written compiler and the equations, exercise edge cases, and make regressions
observable. The mutation test in particular is not presented as an exhaustive
search over all possible witness tuples.

\subsubsection{Reproduction}
From the archive's root, run \wtag{} (the delivered layout; see the README)
\begin{lstlisting}
python3 code/test_certificates.py
python3 code/canonical_memory.py artifacts/example/events.json \
    artifacts/rebuilt_example --expanded
pdflatex -interaction=nonstopmode -halt-on-error article.tex
pdflatex -interaction=nonstopmode -halt-on-error article.tex
\end{lstlisting}
The test script writes its result file relative to its own location. The
\texttt{Makefile} supplies equivalent \texttt{test}, \texttt{example}, and
\texttt{pdf} targets. The Python files use only the standard library; Python
3.10 or later is required for their type-annotation syntax. The article uses
standard pdf\LaTeX{} packages. The included PDF is already compiled. \wtag{} (The PDF is not shipped in ProveIt, and the shipped Makefile, \path{code/03-linear-memory-Makefile}, names delivered paths; see the README.)


% ---- manuscript 04, lines 755-797 ----
\subsection{Manuscript 04: Executable artifacts and exact validation}
\subsubsection{Contents and mathematical interpretation}
The package contains \texttt{code/diophantine\_compiler.py} \wtag{} (shipped as \path{code/04-witness-faithful-diophantine_compiler.py}, with the test driver \path{code/04-witness-faithful-run_tests.py}, the results \path{data/04-witness-faithful-results.json} and the four examples \path{data/04-witness-faithful-*.json}), an exact sparse-integer-polynomial implementation; \texttt{tests/run\_tests.py}; a generated \texttt{tests/results.json}; and four JSON examples. Each example lists all natural variables, all polynomial residuals as integer coefficient/monomial data, and a complete witness assignment. The JSON convention is
\[
 \texttt{combination = sum of squared residuals}.
\]
Thus the residual list defines one explicit ordinary polynomial without printing an unnecessarily expanded quartic. No oracle, numerical tolerance, floating-point root test, or general Diophantine solver is used. The separate \texttt{code/verify\_certificate.py} verifier does not import the compiler; it independently parses and evaluates the exported JSON residuals and checks the supplied witness and degree bound.

From the package directory \wtag{} (in the delivered layout; see the README), run
\begin{verbatim}
python tests/run_tests.py
\end{verbatim}
Python 3.10 or later is sufficient; only its standard library is used. To rebuild this article, run \texttt{pdflatex article.tex} twice, or \texttt{latexmk -pdf article.tex}. The bibliography is included in the source, so a separate bibliography tool is unnecessary.

\subsubsection{What was checked}
The normal-form check enumerates every symmetric irreflexive independence relation on alphabets of sizes one through four and all words of lengths zero through five. It constructs commutation classes independently using union-find on adjacent allowed swaps, then compares the class minima with the monitor. This checks $75$ graphs, $90{,}404$ words, and $39{,}093$ classes. The lower-bound family is checked through $k=7$, including all $254$ prefixes in those instances and their distinguishing one-letter suffixes.

For Petri nets, $24$ seeded small nets with three places and three transitions are tested on $2{,}904$ candidate words. The test evaluates $1{,}417$ enabled-run certificates, checks that zero energy coincides with normality, verifies swap closure, and checks exact variable and residual counts. A separate smallest example exhausts $729$ natural assignments in a finite test box and finds one zero.

For FRACTRAN, the tests enumerate $132$ one- or two-fraction programs whose positive reduced numerators and denominators are at most four, use inputs $1$ through $12$, and horizons zero through three. Together with explicit normalization and denominator-one cases, they evaluate $4{,}914$ valid certificates and check their sizes and degree bounds. These checks test the constructive direction and bookkeeping; they do not exhaust all possible false assignments to the larger systems.

For SKI, $1{,}997$ application codes are inverted and rebuilt, with both children checked to be smaller. The root rules and several context paths are exercised in $1{,}875$ valid certificates. The proof of Theorem~\ref{cdc:wf:thm:ski}, rather than this finite test, establishes completeness and absence of arbitrarily large spurious witnesses.

\begin{center}
\small
\begin{tabular}{lrrrr}
\toprule
Example & Variables & Residuals & Largest coordinate & Rejected mutations\\
\midrule
Petri fork/join &56&67&1&69\\
FRACTRAN first halt &103&120&125&161\\
SKI root schedule &13&11&635&26\\
SKI contextual schedule &8&8&91&16\\
\bottomrule
\end{tabular}
\end{center}
Every listed witness has exactly zero sum-of-squares value. The mutation column counts all permissible one-coordinate increments and decrements by one, each rejected by a nonzero residual. These local mutation checks are useful for catching disconnected variables but are not a proof of global uniqueness.

\subsubsection{Validation status and reproducibility limitations}
The reported suite completed successfully in the working environment. Its mathematical checks are deterministic apart from the seeded net generator, whose seed is fixed in the source; the elapsed-time field in the JSON report is not a mathematical invariant. The implementation mirrors the displayed residual systems and does not implement a fixed-arity MRDP elimination.

No new Lean or Rocq/Coq theorem was compiled as part of this package, and the repository's existing build was not rerun. The article supplies conventional proofs, executable constructions, and finite exact checks. A proof-assistant version of the new witness-level theorems remains a proposed next artifact, not an accomplished part of the present result.


% ---- manuscript 06, lines 1633-1753 ----
\subsection{Manuscript 06: Executable artifacts and verification}
\label{cdc:cs:sec:verification}

The accompanying implementation uses Python's exact integers and the standard
library. It is an executable reference model of the displayed constructions,
not a verified extraction from a proof assistant. Its sparse polynomials store
integer coefficients and monomials explicitly. A certificate file stores the
quadratic residual list and declares the represented polynomial to be their
sum of squares. This gives an exact ordinary polynomial without forcing an
unhelpful expansion of every square.

\subsubsection{What the files implement}

The module \kw{code/compiler.py} \wtag{} (shipped as \path{code/06-counter-schedules-compiler.py}) defines the expression syntax, transition
semantics, resource product and powers, the globally sound commutation
relation, the Boolean signatures, and both versions of the compiler.
\kw{Builder.construct} computes the deterministic auxiliary assignment from
free inputs. It does not make a rejected instance satisfiable: the returned
assignment has nonzero sum-of-squares value when endpoint or normality
constraints fail. Witness-construction recipes involving a minimum or a case
test are \emph{not} extra symbols in the polynomial. The exported residuals
are the sole algebraic constraints, and the proofs above establish that they
force the recipe values on every root.

The module \kw{code/substrates.py} \wtag{} (shipped as \path{code/06-counter-schedules-substrates.py}) implements coefficientwise sign splitting,
nonnegative arithmetic-circuit lifting, and the separated-schedule frontend.
It compiles a polynomial directly to one counter at its original degree,
compiles its quadratic circuit system to several counters at depth two, and
compiles the sum-of-squares quartic to one counter at depth four. Its current
coefficient encoding is unary, as specified in the proof of
\cref{cdc:cs:thm:degree_depth}.

The bounded-length construction of \cref{cdc:cs:thm:tracebijection} is proved in the
article but is not a separate polynomial-export frontend in the supplied
prototype. The tests do independently enumerate bounded trace classes to
check the normality component. No claim of complete implementation coverage
should be inferred from the scope of the mathematical results.

\subsubsection{Finite checks and their limits}

The main suite exhaustively checks small resource pairs and powers, every
independence graph on a three-letter alphabet, and every word of length at
most six over that alphabet. For the trace tests, connected components under
actual adjacent independent swaps provide an independent source of true
lexicographic representatives. The suite also checks signature composition,
power saturation, and 160 seeded expression instances. Both resource-only
and canonical compiled systems are tested on valid and invalid endpoints.
A focused catalyst test catches the zero-power error discussed above.

Mutation tests change one auxiliary coordinate of a valid certificate and
verify rejection. These are useful implementation checks, but they are not a
proof that no coordinated multi-variable mutation could satisfy the system.
The structural uniqueness proof in \cref{cdc:cs:thm:compiler} supplies that theorem.
Similarly, random expression tests do not establish cap-at-two for all
expressions; \cref{cdc:cs:thm:cap} does.

\begin{table}[htbp]
\centering
\small
\begin{tabular}{@{}p{0.77\textwidth}r@{}}
\toprule
Verification group & Checks passed\\
\midrule
Exact resource product, powers, and global commutation & 8,140\\
Trace cross-sections, memory signatures, composition, saturation & 39,032\\
Seeded compressed evaluation and cap-at-two instances & 480\\
Compiler acceptance, degree checks, invalid inputs, zero repetition & 1,180\\
Valid mutation baselines and rejected one-coordinate mutations & 569\\
A certificate for an execution with a 221-digit length & 1\\
Polynomial/circuit/separated-substrate frontend checks & 320\\
\midrule
\textbf{Total} & \textbf{49,722}\\
\bottomrule
\end{tabular}
\caption{Executed finite checks. The main suite contains 49,402 checks; the
substrate frontend adds 320. A check is a tested identity or predicate, not
an independent mathematical theorem.}
\label{cdc:cs:tab:tests}
\end{table}

The substrate test uses
\[
P(x,y,z)=xyz+2z-x^2-3
\]
on the grid $\{0,1,2,3\}^3$. Its nonnegative circuit introduces six uniquely
forced gate parameters and seven quadratic equations. The corresponding
multi-counter scheme has seven counters and depth at most two. Summing the
seven residual squares gives a quartic and hence a one-counter scheme of
depth at most four. The test compares all three presentations with direct
polynomial evaluation and checks rejection after changing a gate value.

For the multiplication schedule in \cref{cdc:cs:ex:multiplication}, the resource-only
prototype allocates 79 witnesses and 86 equations. Its canonical version
allocates 309 witnesses and 547 equations, even though the independence
relation is empty and could be optimized away. These deliberately transparent
counts are not claimed to be minimal. The large-number test uses the same
canonical compiled expression; its longest witness has 221 decimal digits.

\subsubsection{Reproduction}

From the package directory \wtag{} (in the delivered layout; see the README), run:
\begin{lstlisting}
python code/test_compiler.py
python code/substrates.py
python code/compiler.py --export examples/multiplication_certificate.json
\end{lstlisting}
The first two commands regenerate the reports in \kw{verification/} \wtag{} (shipped as \path{data/06-counter-schedules-*.json}). The
random seed in the main suite is $20260930$. The compiler has no network
requirements or third-party Python dependencies. Timing fields are
machine-dependent; the reported counts and mathematical comparisons are the
reproducible content.

Rebuild the article with:
\begin{lstlisting}
latexmk -pdf -interaction=nonstopmode -halt-on-error article.tex
\end{lstlisting}
Two or more direct \kw{pdflatex} passes can be used instead. The bibliography
is contained in this source, so no external bibliography database is needed.
The PDF uses standard installed TeX fonts; font files are not included in the
package.


% ---- manuscript 07, lines 886-949 ----
\subsection{Manuscript 07: Executable artifacts and exact validation}\label{cdc:qc:sec:validation}
\subsubsection{Contents and reproduction}
The package includes \texttt{code/queue\_certificates.py} \wtag{} (shipped, like every file named here, with the prefix \path{07-queue-causality-} in \path{code/} or \path{data/}), with three cyclic-tag modes and the general deletion-tag compiler; \texttt{code/compile\_tag.py}, a JSON-specification front end; and \texttt{code/verify\_certificate.py}, an independent evaluator that does not import the compiler. All use only the Python standard library. Python 3.10 or later is sufficient.

From the package directory \wtag{} (in the delivered layout; see the README), run
\begin{verbatim}
python tests/run_tests.py
python code/verify_certificate.py examples/stream_example.json
python code/verify_certificate.py examples/ternary_tag_example.json
\end{verbatim}
To construct the small cyclic-tag example without simulating it, run
\begin{verbatim}
python code/queue_certificates.py --appendants 10 "" \
  --initial 1 --time 3 --mode stream --no-witness \
  --output examples/symbolic_only.json
\end{verbatim}
The empty quoted argument is the second appendant. The general tag front end reads arrays of integer digits, so it is not limited to single-character decimal alphabets:
\begin{verbatim}
python code/compile_tag.py --spec examples/ternary_tag_spec.json \
  --output examples/ternary_tag_example.json
\end{verbatim}
Compile the article with two runs of \texttt{pdflatex article.tex}. Its bibliography is included in this source file.

\subsubsection{The certificate format}
Each positive example stores the variable list, topologically ordered arithmetic nodes, residual-root indices, metadata, a conservative degree upper bound, and a complete natural witness. Constants and witness values are serialized as decimal strings to avoid fixed-width integer limits in other software. The convention \texttt{combination = sum\_of\_squares} defines one polynomial from the residual list.

The independent verifier validates the node operations and references, evaluates exact integers, checks the degree bound, and evaluates every residual. It has no floating-point tolerance and calls no Diophantine solver. It proves computationally that the supplied witness is a zero of the exported polynomial. Semantic equivalence and uniqueness rely on the mathematical proofs, not on that single evaluation.

The circuit's degree analysis is deliberately conservative: it may not detect cancellation or multiplication by zero. Thus \texttt{stream\_example.json} reports degree at most six, while the simplified explicit polynomial \eqref{cdc:qc:eq:small-example} visibly has degree four. This is not a discrepancy in the represented equation.

\subsubsection{What was checked}
The reproducible results are stored in \texttt{tests/results.json}. The principal counts are:
\begin{center}
\small
\begin{tabularx}{\textwidth}{@{}Xr@{}}
\toprule
Check & Cases\\
\midrule
Cyclic-tag candidate consumed bitstrings & $92{,}610$\\
Raw word-balance solutions rejected by causality & $268$\\
Exact halting instances in that enumeration & $370$\\
Positive cyclic-tag certificates checked in two modes & $740$\\
First-failure resource-product checks & $5{,}916$\\
General deletion-tag candidate consumed words & $194{,}940$\\
No-slack cyclic-tag candidate bitstrings & $840$\\
Natural witness assignments exhausted for quartics & $79{,}631$\\
Word-monoid concatenation pairs & $3{,}969$\\
Integer table-interpolation evaluations & $423$\\
General integer-polynomial tag compiler instances & $120$\\
Candidate words evaluated in those compiled polynomials & $4{,}200$\\
\bottomrule
\end{tabularx}
\end{center}
The cyclic-tag enumeration uses all 49 period-two programs whose appendants have length at most two, all 15 inputs of length at most three, horizons one through six, and every candidate consumed bitstring. It compares the word-and-guard formulas with a direct string simulator. Every positive instance is compiled in stream and quartic modes; single-coordinate witness increments are rejected.

The resource tests include all scalar three-event histories with initial resource at most three and consumptions and productions at most two, followed by 3,000 seeded multi-resource cases. They explicitly allow the formal history to become negative after failure.

The general-tag checks compare the formulas with direct two-symbol deletion, including nonempty halted tails, over exhaustive small binary instances and seeded ternary instances. Separate compiler checks use alphabets of sizes two through four, deletion numbers one and two, and horizons one and two. They evaluate the cleared \emph{integer-polynomial circuits}, not merely the interpolation tables. They also verify that disabling witness attachment changes no polynomial nodes or residuals.

The quartic tests exhaust actual finite boxes of natural assignments, not only generated semantic witnesses. Nevertheless no finite box proves global uniqueness. That is why Theorem~\ref{cdc:qc:thm:quartic} proves that every zero propagates a valid word code and cannot pass through an empty queue. The lower bound on polynomial memory dimension is likewise a mathematical proof, not an inference from the finite word-monoid test.

\subsubsection{Limitations of the validation}
The tests do not certify every implementation behavior for arbitrarily large integers, replace formal proof of the compiler, or establish historical novelty. They do provide independently reproducible checks of the stated formulas, exact resource counts in tested instances, symbolic table clearing, and several easily missed edge cases. The full fixed-read network compiler and a Lean development are proposed extensions, not unreported completed artifacts.


\section{Formalization targets}\label{cdc:sec:formal}
\wtag{} No theorem of this report is formalized in Lean or Rocq, and no manuscript ships a proof-assistant file; the subsections below are the manuscripts' proposals, and every module name in them is proposed, not asserted to exist. The report sits beside ProveIt's Lean project \path{Computability/HilbertTenthProblem}, which proves MRDP in both directions. The manuscripts use its existence interfaces \path{Diophantine.boundedForall_dioph}, \path{Diophantine.exactIter_dioph} and \path{Diophantine.existsExactIter_dioph} (\path{Lean/Diophantine/Common/DiophantineTrace.lean}) and \path{Diophantine.mrdp} as context only: those theorems assert that representations exist and say nothing about how many witnesses a representation has. The project also formalizes the single-fold unary exponential representation of Jones and Matiyasevich (\cref{cdc:bd:rem:sfu}), and the queue content and length-marker identities of its tag-system route (\cref{cdc:qc:rem:tagnotes}). Placement beside a formal development confers no formal status on anything here.

% ---- manuscript 05, lines 878-918 ----
\subsection{Manuscript 05: A formalization route that preserves the stronger specification}
\label{cdc:sec:formalization}
The new proofs are not currently Lean-checked. A useful formalization should preserve the fibre information rather than immediately discarding it through an existential predicate. The repository's existing \code{Dioph} interfaces can remain unchanged; the stronger results can imply them as corollaries.

\subsubsection{Separate semantic and arithmetic layers}
The semantic layer should define a finite labelled family of partial maps, exact finite runs, a proved sound independence relation, and its generated word equivalence. The arithmetic layer should define a fixed finite set of variables and residual polynomials. A decoder reads the selected word from a natural zero. A witness constructor produces all variables from a canonical run. The central obligations are
\[
\operatorname{decode}(\operatorname{encode}(w))=w,
\qquad
\operatorname{encode}(\operatorname{decode}(z))=z.
\]
The second identity is what excludes free arithmetic witnesses. Proving only the first would establish completeness but not injectivity of the representation.

An appropriate result object is an equivalence between a finite execution-class type and a subtype of natural assignments satisfying the polynomial equation. Cardinality and weighted-enumerator statements then follow from that equivalence rather than from separate ad hoc proofs. The quotient can be replaced by the subtype of normal runs, using the unique-normal-form theorem to build an equivalence with classes.

\subsubsection{Suggested proof decomposition}
A practical decomposition has six modules, whose names below are proposed rather than asserted to exist.

\paragraph{\code{PetriDemand}.}
Formalize \cref{cdc:thm:commute}, including both implications and the empty-coordinate case. Keep the distinction between integer vector updates and natural state domains explicit. Prove equality of partial functions, not merely equality of results when both are defined.

\paragraph{\code{IntervalGuardCommutation}.}
Represent upper endpoints by a type with an infinity element. Prove the composed-box formula, canonicalize empty boxes, and derive \cref{cdc:thm:interval}. Equality of nonempty boxes must include unbounded-coordinate cases.

\paragraph{\code{TraceNormalForm}.}
Prove the forbidden-factor characterization and the flag invariant. Establish uniqueness of the lexicographically least representative using finite word classes. Prove that the recognizer's natural recurrence forces Boolean values by induction, rather than adding unnecessary Boolean polynomial constraints.

\paragraph{\code{UniqueAffineExtension}.}
Formalize unique slack, the affine AND gadget, and deterministic circuit extension. A reusable combinator should compose witness equivalences over fixed inputs. Sum-of-squares scalarization should be stated over an ordered ring, then specialized to integer polynomials evaluated at natural assignments.

\paragraph{\code{CanonicalTraceCertificate}.}
Define the variable index type as a tagged finite union of the arrays in the construction. Prove decoding, existence, and coordinatewise uniqueness. This layer also carries the exact variable and residual counts, including $T=0$ and $d=0$.

\paragraph{\code{TracePolytope}.}
Reuse the affine system over the reals, prove each coordinate bound, and package boundedness and convexity. Relate the natural assignment subtype to lattice points without changing domains by an uncontrolled integer encoding.

\subsubsection{Integration with the current repository}
The existing \code{exactIter_dioph} theorem can provide the unbounded existence-level consequence once a substrate's one-step relation is Diophantine \cite{repo-trace}. It should not be renamed or reinterpreted as a multiplicity theorem. Instead, add a clearly separate bounded compiler interface carrying unique extension, and document precisely what is forgotten when projecting to \code{Dioph}.

The most valuable first milestone is the raw affine certificate plus unique extension; it is small, reusable, and independent of trace-monoid infrastructure. The second is the exact flag recognizer, and the third is the combined bijection. Only after those should cardinality and the full real-polytope statements be layered on. No project axiom or placeholder proof is needed in the mathematical design; an actual implementation must still pass the repository's normal kernel and transitive-axiom checks.


% ---- manuscript 01, lines 802-812 ----
\subsection{Manuscript 01: A proposed Lean integration boundary}
A useful formalization should keep the new proof obligations separate from existing ordinary-Diophantine closure. One possible layering is:
\begin{enumerate}[label=\arabic*.]
\item Define box-guarded partial translations, word execution, and equality of partial maps. Prove Proposition~\ref{cdc:ct:prop:boxes} and the natural-state invariant without importing the MRDP machinery.
\item Prove the envelope theorem using counts of positive and negative displacements in prefixes. State the theorem over arbitrary finite action and coordinate types; prove the Petri formula by simplification.
\item Introduce a representation structure containing a polynomial, a semantic decoder, and an equivalence between its natural zero set and the specified semantic type. Ordinary Diophantineness is then a projection, not the entire contract.
\item Formalize conditional slack canonicalization and the support gadget. Assemble the accelerator polynomial with finite indexed residual families and a sum-of-squares lemma.
\item Develop the dependence-poset normal form and its padded layer characterization, then prove the full-root bijection of the history compiler. Only after that should an adapter export ordinary reachability through the existing Diophantine interfaces.
\end{enumerate}
The key extra structure is the \emph{equivalence of fibres}. A theorem that a decoder is surjective onto semantic histories proves completeness but not uniqueness. Injectivity of that decoder on the complete zero set must also be retained. No uncompiled Lean skeleton is presented as a completed formalization in this package.


% ---- manuscript 02, lines 1500-1525 ----
\subsection{Manuscript 02: Integration with ProveIt and further research}
\label{cdc:pt:sec:agenda}

\subsubsection{An interface worth adding}
A useful extension of the repository would distinguish three levels of
representation explicitly. An ordinary Diophantine interface stores an
existence equivalence. A finite-fold interface additionally proves finite
fibers. A single-fold interface additionally proves subsingleton fibers.
Conversions should retain these properties in their types or theorem
statements rather than discard them in an undifferentiated existence
predicate.

For the present construction, the clean decomposition is a theorem about
unique sign tables, a theorem about canonical circuit evaluation, and a
composition theorem relating the compiled quartic to the checking formula.
The ordinary existence result can then be exported to the current
Diophantine interfaces without making the stronger witness theorem depend
on a noncanonical general MRDP backend.

The key invariants for a formal implementation are: ordered cuts cover the
entire integer domain; empty slots form a terminal suffix; their tags are
fixed; adjacent active tags differ; every parent point is covered by a
child monotonicity interval; all internal arithmetic computations have a
unique normal form. These invariants are small enough to isolate into
reusable modules, but no claim of completed formalization is made here.


% ---- manuscript 03, lines 1459-1530 ----
\subsection{Manuscript 03: A proposed ProveIt formalization architecture}\label{cdc:mem:sec:formalization}
The repository already exposes ordinary Diophantine closure interfaces, including
\repo{Diophantine.boundedForall_dioph}, \repo{Diophantine.exactIter_dioph}, and
\repo{Diophantine.existsExactIter_dioph}, according to the inspected status
register \cite{repo-status}. These are useful downstream interfaces, but their
existential statements should not be treated as uniqueness theorems.
A new compiler should retain the stronger specification until an application
explicitly chooses to forget it.

A natural representation is a finite list of polynomials over $\Z$, evaluated
on natural-number assignments cast into $\Z$. A semantic certificate object
should record a decoder from satisfying assignments to source objects and a
canonical encoder in the reverse direction, with both inverse laws on their
specified domains. Counting is then a consequence of an equivalence of types
or finite sets, rather than a separate combinatorial assertion.

The following module names are \emph{proposals}, not files claimed to exist in
the inspected repository:
\begin{center}
\begin{tabularx}{\textwidth}{@{}p{0.31\textwidth}X@{}}
\toprule
Proposed module & Main proof obligation\\
\midrule
\texttt{CanonicalQuadratic} & Natural-domain residual semantics, sum-of-squares
combination, fresh-variable composition, and fiber equivalences.\\
\texttt{CanonicalCompare} & The two-equation signed-gap theorem, zero tests,
unique bound slacks, multiplexing, and exact gadget counts.\\
\texttt{CanonicalSorting} & Record-preserving comparators, fixed network
semantics, the bitonic generator, and its closed count formula.\\
\texttt{ExactProductFingerprint} & Base-expansion coefficient recovery,
multiset equality from product identity, and injective record packing.\\
\texttt{CanonicalMemory} & Both backend theorems, chronological semantics,
exact counts, and distinct height bounds.\\
\texttt{CanonicalRAM} & The request-before-response interface, temporal
soundness induction, normalization of inactive instructions, and halted padding.\\
\texttt{CanonicalGraphRuntime} & Immutable heap unfolding, spine evaluation,
microstep accounting, and the explicitly stated strategy correspondence.\\
\texttt{CertificateAudit} & Emitted-polynomial reflection, checked examples,
domain conventions, and transitive assumption reports.\\
\bottomrule
\end{tabularx}
\end{center}

The linear backend offers an especially short conceptual trusted boundary.
Its nonlocal step is the multiset lemma, reducible to elementary base expansion
and polynomial cancellation; it requires no asymptotically optimal sorting
network formalization. The low-height backend remains valuable when later
applications need tight bit bounds.

An untrusted compiler may emit a residual syntax tree, but the kernel should
check that the emitted object is the polynomial claimed by the compiler
correctness theorem. Testing the value at one supplied witness checks only
that instance. The complete theorem must quantify over every natural assignment
to prove that there are no spurious roots, and over every source execution to
prove that none are lost. A degree or variable-count audit should inspect the
actual emitted object, not just a separately maintained formula.

A sensible order of work is to formalize the common residual semantics, then
the product and comparison lemmas independently, then the two memory backends,
and finally one small RAM instruction language. The combinator frontend should
come after the precise runtime strategy and heap invariants have been fixed.
Existing ProveIt MRDP interfaces can then export an ordinary Diophantine
reachability statement, with an explicit note when that export forgets the
canonical-fiber contract.

The inspected status register lists further work on the 1984 paper's specific
NP characterizations (53) and (54) \cite{repo-status}. The bounded-counting
corollary here is not a formal proof of those numbered statements or a
replacement for their required bounded-exponentiation theory. It should be
integrated as a distinct result unless an exact implication is separately
proved.


% ---- manuscript 04, lines 671-683 ----
\subsection{Manuscript 04: How these results extend the inspected ProveIt development}
The mathematical extension point is a \emph{witness-level} interface. An ordinary Diophantine theorem records an existential equivalence. A stronger reusable interface would record a semantic object, a decoder from every zero, a canonical encoder, and proofs that the encoder and decoder are inverse. Section~\ref{cdc:wf:sec:interfaces} shows how circuit compilation can preserve that interface, and the three substrate constructions instantiate it in substantially different ways.

The repository snapshot used for comparison is
\begin{center}\small\path{e8bb0931d67f80d9fce87a8cddb0f661ff19f956}.\end{center}
The inspection was focused on its Hilbert-tenth and combinatory-logic branches; it was not a full audit of every repository file. Its existing exact-iteration development already proves the existential closure we need \cite{repo-trace}. The present paper should therefore be integrated as a refinement tracking fibers, not as a duplicate theorem that finite reachability is Diophantine.

There are three natural mathematical modules. The first is an injective-state-coding and unique-extension layer, including canonical slacks and circuit wires. The second is a semantic quotient layer: an independence relation, enabled-swap proof, normal-form monitor, and class-to-zero bijection. The third is a resource layer measuring variables, residuals, degree, coefficient bit lengths, and witness heights separately. A verified composition theorem should carry all three layers rather than only a proposition about nonempty solution sets.

The numerical bounds in this article do not compete with the repository's universal polynomial pairs or its optimized straight-line universal certificates \cite{repo-h10}. Here the horizon and often the schedule are external compilation data, whereas a universal equation encodes unbounded computation with fixed arity. Comparing these counts without recording that distinction would be misleading.

No new repository proof-assistant build is reported in this package. The executable artifacts are standard-library Python implementations of the explicitly displayed bounded compilers. A separate source-audit note records the inspected paths and the claims relied upon, without treating repository documentation as an independent verification of the new results.


% ---- manuscript 06, lines 1956-2000 ----
\subsection{Manuscript 06: Proof dependencies and precise completion status}
\label{cdc:cs:app:status}

\begin{longtable}{@{}>{\raggedright\arraybackslash}p{0.29\textwidth}>{\raggedright\arraybackslash}p{0.40\textwidth}>{\raggedright\arraybackslash}p{0.24\textwidth}@{}}
\toprule
Result & Essential ingredients & Status in this package\\
\midrule
\endfirsthead
\toprule
Result & Essential ingredients & Status in this package\\
\midrule
\endhead
Exact resource summaries & Partial translations, coordinatewise minimum,
word induction & Full proof; implemented and tested\\[3pt]
Global swap criterion & Exact two-letter resource products & Full proof;
implemented and tested\\[3pt]
Compact normality signature & Classical forbidden factors; Boolean affine
coordinate maps & Full proof; implemented and tested\\[3pt]
Cap-at-two saturation & Signature product; $\alpha_a\beta_a=0$ & Full proof;
implemented and tested\\[3pt]
Compressed quartic compiler & Unique local gates, syntax induction,
sum of squares & Full proof; implemented and tested\\[3pt]
Bounded trace-class bijection & One-hot words, deterministic memory,
trace-invariant resource action & Full proof; separate export frontend
not implemented\\[3pt]
Flat-scheme decidability & Affine positive-power formula; Presburger
case splitting & Full reduction proof; decision frontend not implemented\\[3pt]
Degree--depth correspondence & Length polynomials; signed monomial splitting
& Full proof; frontend tested\\[3pt]
Depth-two multi-counter universality & MRDP; unique arithmetic-circuit lift
& Full reduction proof using classical MRDP; frontend tested on a sample\\[3pt]
Finite-fold equivalence & Unique quartic lift; exact separated length balance
& Full equivalence proof; no solution of the general conjectures\\[3pt]
Formal verification & A proof-assistant implementation of all new definitions
and compiler semantics & Proposed; not supplied or claimed\\
\bottomrule
\end{longtable}

The only deep external theorem used for the universality conclusions is MRDP.
The main compressed compiler, its cap-at-two theorem, and the bounded
trace-counting construction do not depend on it. The quadratic one-counter
decidability observation cites a separate classical decision theorem; it is
not used to prove the flat/nested multi-counter boundary.

\clearpage

% ---- manuscript 07, lines 816-839 ----
\subsection{Manuscript 07: Integration with ProveIt and the formal trust boundary}
\subsubsection{What was inspected}
The repository context was pinned at commit
\begin{center}\small\path{725d2ebb6909fe11a13a92354c0f47367a3cbbf5}\end{center}
The inspected Hilbert-tenth-problem README and MRDP guide describe \texttt{Diophantine.mrdp}, \texttt{Diophantine.mrdp\_iff}, and the finite-polynomial and trace interfaces \cite{repo-mrdp,repo-h10}. The guide reports an axiom-audited development and distinguishes the current forward proof from the older shared trace route. Those are repository reports, not build results independently reproduced for this package.

No theorem about the queues in this article is inferred from an unrelated research claim elsewhere in the repository. The explicit queue and word-memory proofs are independent mathematical arguments. Existing MRDP infrastructure is relevant to eventual fixed-arity representability, not needed for the bounded constructions themselves.

\subsubsection{A proposed formalization sequence}
A manageable formalization should separate word semantics from polynomial syntax. First define the cyclic-tag step as a partial function on finite Boolean lists with a phase parameter, stopping on the empty list. Prove the radix concatenation identity and fixed-length injectivity, then prove the suffix reconstruction theorem at the word level. These proofs need no Diophantine machinery.

Next formalize Theorem~\ref{cdc:qc:thm:first-failure} over integer-valued formal resource histories. The invariant before the first disabled event is nonnegativity. This is the crucial point where natural-number truncation must not enter the formal recurrence. The first disabled event supplies a specific zero factor in the finite product.

Only then define the representing polynomial in
\[
 \texttt{MvPolynomial (Fin (T+1))}\ \Z.
\]
The first $T$ coordinates are consumed symbols and the last is the guard slack. Prove that evaluating each symbolic residual gives its semantic integer expression; combine this with the sum-of-squares zero criterion. Soundness, completeness, and uniqueness should be separate theorem endpoints, followed by degree and variable counts. The quartic compiler should have its own endpoint because its preservation-of-valid-scale argument is different.

Finally, a reified arithmetic-DAG checker can connect exported data to these polynomial definitions. JSON parsing and Python generation should remain outside the trusted boundary. The checker must prove that accepted nodes define the intended polynomial and that the displayed witness evaluates to zero. Checking one witness alone must not be presented as a proof that no other witness exists.

\subsubsection{What is shipped now}
The package supplies human-readable proofs and tested Python compilers, including integer-only finite-table interpolation for deletion-tag systems. It supplies neither a Lean proof file nor an assertion that the new theorems have passed a kernel audit. The file \texttt{notes/lean\_integration.md} \wtag{} (shipped as \path{07-queue-causality-lean_integration.md}) records proposed modules and obligations, so a later formalization can begin without confusing an interface plan with completed code.


% frag_questions.tex -- merged "Research questions" section of the
% canonical-diophantine-certificates report. Sources: 01 sec. 12
% (813-855), 02 secs 12.2-12.3 (1526-1626), 03 sec. 13 (1531-1627),
% 04 sec. 11 (684-746), 05 sec. 13 (919-982), 06 sec. 12 (1754-1898),
% 07 sec. 13 (840-885).
\section{Research questions}\label{cdc:sec:questions}
\label{cdc:ct:sec:questions}\label{cdc:mem:sec:questions}\label{cdc:sec:research}\label{cdc:cs:sec:questions}\label{cdc:qc:sec:questions}

\wtag{} This section prints every research question of the seven
manuscripts: ten from manuscript~01 (its Section~12), ten from
manuscript~02 (its Section~12.2), ten from manuscript~03 (its
Section~13), ten from manuscript~04 (its Section~11), twelve from
manuscript~05 (its Section~13), twelve from manuscript~06 (its
Section~12), and ten from manuscript~07 (its Section~13): seventy-four
questions in all. Manuscripts~03, 04 and~06 posed their questions under
subsection headings; those headings are the titles here, and 06's
question numbers are kept in the titles. Where several manuscripts ask
the same question, the questions are merged into one \texttt{question}
environment whose title names every source; inside it each source's
question text is printed verbatim under \emph{Source~NN}, and each
source's follow-up paragraphs are printed verbatim after the environment
under the same heading. Seven merges reduce the seventy-four questions to
fifty environments. Listing a question here does not claim that it is
answered, by this report or elsewhere. Two questions are partly answered
inside this report, and a \wtag{} note at each says how far:
manuscript~01's question on sharp degree, arity and domain comparisons
(\cref{cdc:q:degree}), by the degree-two certificate of Part~II
(manuscript~05); and manuscript~06's Question~9 on new storage monoids
(\cref{cdc:q:storage}), for bounded runs of fixed-read FIFO networks, by
Part~VI (manuscript~07). Manuscript~02's closing paragraph on what would
constitute a stronger result is printed as the last subsection.

\subsection{How the manuscripts frame their questions}

\emph{Source 01.} The following questions separate tractable extensions
of the present construction from the genuinely global compression
problem. They are proposed investigations, not claims of solutions.

\emph{Source 03.} The following questions separate tractable
implementation work, quantitative mathematics, and the much stronger
fixed-variable frontier. They are proposals from this investigation;
calling them questions does not assert that each is unanswered in all
existing literature.

\emph{Source 04.} The following questions range from bounded
construction problems to genuinely stronger fixed-arity goals. They are
proposed research targets, not all asserted to be previously named open
problems. Each is stated so that a success criterion and a likely
failure mode are visible.

\emph{Source 05.} The following agenda separates finite, testable
development tasks from the genuinely stronger uniform arithmetic
questions. None of the questions is claimed to have a new answer merely
because it is listed here.

\emph{Source 06.} The following questions are proposed continuations,
not a claim that every one is a newly identified open problem in the
literature. They are organized to separate tasks that build directly on
the proved compiler from those that confront the finite-fold frontier.

\emph{Source 07.} The questions in this section are unresolved by this
article. Except for the explicitly identified finite-fold topic, they are
proposed construction and optimization problems, not assertions that a
literature-wide open-problem status has been established.

\emph{Source 07, closing paragraph of its question section.} The most
immediate projects are the first-failure and FIFO formalizations, a
complete fixed-read network compiler, and optimization of the quartic
variable count. They have precise interfaces and can be validated
independently. The fixed-arity problem is a separate, more demanding line
of investigation, for which the present explicit unique traces identify
exactly what information a successful compression must preserve.

\wtag{} Manuscript~02 printed its ten questions under the heading
``Research questions with concrete targets'' without an introductory
paragraph.

% ---------------------------------------------------------------------
\subsection{Verified compilers and fibre-preserving interfaces}

\begin{question}[A verified single-fold compiler interface; sources 01, 02, 03, 04, 05, 06, 07]\label{cdc:q:verified}
\emph{Source 01 (Formalizing the exact resource envelope).}
Can ProveIt obtain a reusable finite-type theorem for box-guarded translations whose specialization automatically produces Petri and finite-control counter adapters? A useful first milestone is a kernel-checked equivalence between all multiset serializations and~\eqref{cdc:ct:eq:envelope-lower}--\eqref{cdc:ct:eq:envelope-upper}, including empty support, finite upper bounds, and empty composition domains. The elementary finite combinatorics make this a focused formalization target.

\emph{Source 02 (A verified single-fold compiler).}
Formalize the total sign gadget, signed normalization, and the endpoint
propagation theorem, then extract a compiler whose output includes a
proof of both equivalence and unique extension. A useful first milestone
is the complete degree-two template, including zero polynomials and
$T=0$, rather than a generic compiler with unproved boundary cases.

\emph{Source 03 (What is the right verified compiler interface?).}
Build an actual Lean polynomial emitter whose contract contains soundness,
completeness, inverse decoding, degree, sparsity, and height bounds. The two
backends provide a test of whether the abstraction separates semantics from
resource properties cleanly. A useful first milestone is formal verification
of both four-event exports in this archive, followed by the symbolic theorem
for arbitrary $L$; finite examples alone are not the endpoint.

\emph{Source 04 (Canonicity-preserving arithmetic library design).}
Develop a library of polynomial gadgets certified not only by existential equivalences but by explicit solution-set bijections, together with exact cost and height annotations.

\emph{Source 05 (Kernel-checked fibre preservation).}
Can the full equivalence of \cref{cdc:thm:main}, including every slack coordinate, be formalized as a reusable compiler theorem in ProveIt? A concrete first target is to reproduce the raw, quartic, and quadratic variable counts with a finite-indexed polynomial definition and an executable witness constructor.

\emph{Source 06, Question 10 (A small verified compiler and a multiplicity audit).}
Formalize the resource pair action, the minimum and
selector gadgets, the Boolean signature product, and the structural
single-fold theorem as a compact independent layer. The decisive endpoint
should state a bijection between the specified semantic witnesses and
polynomial roots, not merely an implication from a generated certificate to
acceptance. A subsequent integration with ProveIt's existing Diophantine
interfaces could preserve that stronger contract. The audit should explicitly
check natural versus integer subtraction, zero-length traces, inactive
branches, and every projection that turns a parameter into a witness.

\emph{Source 07 (Formal arithmetic-DAG compilation).}
Implement the proposed ProveIt modules for first-failure products, word reconstruction, and certificate evaluation. The first deliverable should be the full equivalence and uniqueness theorem for the cyclic-tag stream compiler, not just a theorem that a supplied example evaluates to zero. Degree-bound reflection can follow independently.
\end{question}

\emph{Source 04.} Priority selection, canonical remainder, fixed-pattern matching, context reconstruction, and unique-extension quadratization are concrete initial modules. The important regression tests are inactive-variable freedom, integer-versus-natural subtraction, duplicate state encodings, and mismatched semantics between compiler stages. Formalizing the Petri-net theorem is an especially contained first milestone because its state and monitor equations are only quadratic.

\emph{Source 05.} The decisive acceptance test is the identity $\operatorname{encode}\circ\operatorname{decode}=\mathrm{id}$ on the full solution subtype. A proof only of solvability equivalence is not sufficient for the counting applications.

\begin{question}[Which normal forms preserve canonical fibers?; source 03]
Investigate small coefficient sets, bounded variable occurrence, and more
restrictive quadratic gate types. Introducing a uniquely defined copy of a
coordinate preserves parsimony, but replacing natural coordinates by ordinary
integer sum-of-squares representations generally does not. Identify exactly
which normal-form transformations preserve existence only, finite fibers, or
singleton fibers. This could become a reusable theorem library independent of
memory checking.
\end{question}

\begin{question}[Substrate interfaces with explicit witness contracts; source 02]
Build counter-machine and FRACTRAN libraries whose one-step arithmetic
certificates are single-fold and whose composition interfaces state exactly
which property survives. Use the polynomial accelerator only when a
checked trajectory certificate is available; otherwise fall back to an
ordinary existence-preserving backend. The resulting mixed system would
make the boundary between established local uniqueness and unresolved
unbounded compression visible to users.
\end{question}

% ---------------------------------------------------------------------
\subsection{Commutation, independence and canonical representatives}

\begin{question}[State-dependent independence and commutation beyond a fixed relation; sources 01, 04, 05, 06, 07]\label{cdc:q:dynamic}
\emph{Source 01 (Dynamic rather than global independence).}
Can state-dependent commuting diamonds support a canonical arithmetic trace quotient? Global partial-map equality gives a congruence for free. A rule that swaps actions only when both orders happen to work at the current state may lose a simple global normal form. A successful extension needs explicit coherence conditions and a proof of canonicality, not merely local diamond checks. Small counterexamples should be sought before designing the polynomial layer.

\emph{Source 04 (Canonical certificates for dynamically independent actions).}
Can the quadratic monitor be generalized to a sound state-dependent independence relation while retaining one solution per chosen execution-equivalence class and a polynomial-size bounded certificate?

\emph{Source 04 (Beyond disjoint supports).}
Find effectively checkable families of transition pairs with overlapping supports that satisfy the enabled-swap property for all reachable markings, and derive more aggressive trace quotients without losing witness-faithfulness.

\emph{Source 05 (State-sensitive commutation).}
Can the static relation be replaced by a state-dependent diamond test while retaining a unique polynomial-size normal-form verifier? What precise consistency conditions replace the trace-monoid forbidden-factor language?

\emph{Source 06, Question 5 (Invariant-sensitive and state-dependent swaps).}
Replace equality of partial actions on all markings by
equality on a prescribed invariant, such as exactly one control token. This
could admit more commutations than \eqref{cdc:cs:eq:commutation}. Which classes of
invariants permit an effective exact test, and when do the resulting swaps
support a sound canonical cross-section? If the permitted swap depends on
the current marking, a fixed trace monoid may no longer be the right object;
confluence and uniqueness then require new proofs, not just a larger pair
relation.

\emph{Source 07 (Concurrent traces rather than schedules).}
For communicating FIFO processes, can solutions represent equivalence classes of independent event schedules while preserving channel order? A swap can change which symbols an event observes, so a static independence relation is not automatically sufficient. A semantic commutation criterion should be proved before importing schedule-normalization techniques.
\end{question}

\emph{Source 04 (Canonical certificates for dynamically independent actions).} The current proof relies on a fixed trace congruence. If independence changes with state, a swap can change later independence tests and equivalence need not be described by a fixed graph. A meaningful first target is a finite-state system with a certified local diamond property, together with a canonicalization theorem that explicitly proves class preservation. Simply evaluating a different $I$ at every row of \eqref{cdc:wf:eq:flagpoly} is not justified by Section~\ref{cdc:wf:sec:monitor}.

\emph{Source 04 (Beyond disjoint supports).} For example, shared read-only resources can permit commutation beyond literal support disjointness. The right hypothesis is preservation of both enabling and outcome, not merely commutativity of the net incidence vectors. A certificate for an invariant region of markings may justify additional swaps. Any added pair must preserve the complete endpoint language used in Theorem~\ref{cdc:wf:thm:petri}.

\emph{Source 05.} A local diamond at one marking is insufficient to justify arbitrary sequences of swaps. A satisfactory answer must show invariance of enabledness across the whole generated equivalence and uniqueness of the chosen representative, not just local endpoint equality.

\wtag{} The hypothesis every part of this question asks to relax is the
one fixed in \cref{cdc:thm:main}: a single symmetric irreflexive relation
$I\subseteq I_{\max}$, where $I_{\max}$ of \eqref{cdc:eq:imax} contains
exactly the pairs whose two-step partial actions agree at every marking.
Under that hypothesis a swap is sound in every context; the six texts
above ask what replaces it when commutation is allowed to depend on the
current state, on a prescribed invariant, or, for FIFO channels, on the
symbols an event observes.

\begin{question}[Structure of independence graphs; source 05]
Which graph parameters permit smaller flag systems or tractable exact class counting? Candidates include neighbourhood inclusion, modular decomposition, and widths of a graph encoding the residual interactions.
\end{question}
A useful project would report both coefficient sparsity and actual enumeration behaviour on a library of graphs. \Cref{cdc:thm:graph} rules out assuming that Petri-net semantic commutation graphs belong to a narrow graph class in general.

\begin{question}[Minimal signatures and algebraic structure; source 06, Question 1]
For a given independence graph and alphabet order, what
is the smallest exact block signature for normality? Our upper bound of
$2\cdot6^q$ possible records is often redundant, especially when the scalar
acceptance bit is zero. Determine the reachable signature submonoid, compare
it with the syntactic monoid of the normal-form language, and characterize
when every reachable signature is already idempotent. The general exponent
two is sharp, but its dependence on graph structure can be finer. A useful
outcome would be a smaller circuit representation together with a proof that
it preserves the $S^2=S^3$ law.
\end{question}

\begin{question}[Compressed normalization rather than testing; source 06, Question 2]
Given a parametrized repetition expression, can one
produce a comparably small expression for its lexicographic normal word,
uniformly in the parameters? The present result only tests whether the
original expansion is normal. In the fully commuting case, sorting is
controlled by letter counts, but general dependencies retain nontrivial
ordering information. The key issues are the output representation, whether
new parameter arithmetic is allowed, and whether polynomial-size output can
always exist. A useful first target is bounded dependence width or a fixed
small alphabet.
\end{question}

\begin{question}[Trace-injective parameterizations; source 06, Question 3]
Which repetition templates have at most one parameter
tuple per trace, or finitely many tuples per trace? The example
$(a^{k_1})^{k_2}$ shows why empty powers and factorization can destroy
injectivity. Can redundant parameters be removed in a useful decidable
fragment while retaining a compact expression and a witness-preserving
compiler? A classification for flat schemes would connect the Presburger
execution relation with trace equality, but execution equality alone is not
a sufficient substitute for the latter.
\end{question}

\begin{question}[Unique acceleration of repeated causal patterns; source 01]
The accelerator compresses one commuting multiset, and Proposition~\ref{cdc:ct:prop:word-power} already handles a prescribed fixed word repeated an arbitrary number of times. Can a general causal trace be assigned a unique compressed segmentation into such patterns? The unresolved step is to infer the decomposition without multiple period or boundary choices representing the same trace. A useful restricted target is a class with a canonical primitive-period factorization, together with a proof that the arithmetic witnesses preserve that factorization uniquely.
\end{question}

\begin{question}[Canonicality inside restricted schedule families; source 06, Question 4]
Find effective syntactic conditions ensuring that a
schedule family contains the normal representative of every trace it
contains. Under such a trace-closure condition, the canonical compiler would
preserve existential reachability within the family, avoiding the failure in
\cref{cdc:cs:rem:filter}. One may alternatively seek a certified expansion of the
family that is trace-closed but has controlled size. The cap-at-two theorem
makes testing a proposed representative cheap; it does not by itself
construct a family containing that representative.
\end{question}

% ---------------------------------------------------------------------
\subsection{Counting, weights and relaxations}

\begin{question}[Counting trace classes, histories and certificates; sources 01, 03, 04, 05, 06]\label{cdc:q:counting}
\emph{Source 01 (Counting and generating functions).}
For a finite invariant marking region, can the canonical-layer transition system yield useful exact rational generating functions for trace counts, weighted by event number or cost? A finite-state construction is plausible, but extracting small formulas and comparing different valid independence relations remain worthwhile tasks. The quartic roots give a direct arithmetic interpretation of those counts and a way to test candidate recurrences on finite instances.

\emph{Source 03 (Can concurrency retain meaningful counting semantics?).}
A sequential log gives one total order of accesses. For a concurrent execution,
different legal linearizations may describe the same partially ordered history.
An existential timestamp certificate would count linearizations, not necessarily
semantic histories. Develop canonical linearization conventions or explicitly
weighted quotient semantics, distinguishing sequential consistency from weaker
memory models. The local memory checker is a backend, not a solution to those
semantic choices.

\emph{Source 04 (Counting directly from structured residual systems).}
For which structural restrictions on a net or independence graph can the quartic certificates be counted or sampled efficiently without expanding the entire reachable-state automaton?

\emph{Source 05 (Counting algorithms under explicit parameters).}
Which parameters yield fixed-parameter or output-sensitive algorithms for counting the lattice points of these particular trace polytopes? Parameters should be stated precisely: fixed horizon, small active transition set, bounded residual-graph width, or a restricted control architecture give different problems.

\emph{Source 06, Question 11 (Counting beyond bounded words).}
Extend the bounded trace-class bijection to compressed
families with a known trace-injective parameterization, and then to weighted
trace counting. A weight assigned to an execution must be invariant under
the chosen commutations before it can be interpreted as a trace weight.
The main obstruction for unbounded families is not the local quartic
encoding; it is controlling the number of parameterizations of a behavior
and the witnesses introduced by any unbounded coding mechanism.
\end{question}

\emph{Source 04.} Potential parameters include the interaction-graph treewidth, the size of conflict sets, and the number of shared resource places. Theorem~\ref{cdc:wf:thm:counting} rules out a general polynomial-time exact algorithm under standard complexity expectations, but does not rule out parameterized algorithms. The natural comparison point is current trace counting and sampling for regular languages \cite{decolnet2026}; input representation and trace closure must remain explicit.

\emph{Source 05.} An empirical benchmark should compare direct canonical-word generation, lattice-point methods, and dynamic programming. The $\sharpP$-completeness result gives a worst-case boundary, not a prohibition on useful structured algorithms.

\begin{question}[Weight-preserving quotients; source 05]
For which state-dependent costs, probabilities, or resource observations does a commutation quotient preserve the intended observable? Can the necessary observation data be augmented without destroying a useful independence relation?
\end{question}
Label-additive weights are fully handled by \cref{cdc:cor:weighted}. Peak space, scheduler probabilities, and costs depending on the current marking require additional analysis. The correct target may be a weighted sum over linear extensions rather than one unit per class.

\begin{question}[What extends to weighted, probabilistic, or algebraic computation?; source 03]
The fair-bit formula handles exact positive path weights with canonical stopped
padding. Extend the construction to bounded-denominator rational transition
probabilities and track a uniquely defined common denominator. For quantum
circuits, path counting alone is insufficient because amplitudes can cancel;
one must specify exact algebraic coefficient representations and certify the
amplitude arithmetic and measurement rule. The distinction between counting
paths and summing amplitudes is fundamental.
\end{question}

\begin{question}[Nondeterminism and path multiplicities; source 02]
For a finite family of polynomial phases, can a certificate have one
solution per distinct execution, rather than one solution per endpoint?
Canonical encodings of branch sequences and phase boundaries are needed,
especially when two phase descriptions overlap. A first tractable case is
a fixed acyclic control graph with locally deterministic phases and
explicitly labeled nondeterministic edges.
\end{question}

\begin{question}[Integral or ideal relaxations; source 05]
For which subclasses of nets and independence relations is the canonical trace polytope integral? Can one characterize subclasses admitting a polynomial-size ideal extension, meaning an extension whose real projection is the convex hull of the relevant discrete encodings?
\end{question}
\Cref{cdc:ex:gap} supplies a small explicit obstruction to unrestricted integrality. Total unimodularity of a suitable residual matrix would be a sufficient condition, but the canonicalization constraints need not preserve it. The general counting hardness theorem also prevents treating a compact formulation as a general counting algorithm.

% ---------------------------------------------------------------------
\subsection{Size, degree and height of certificates}

\begin{question}[Sharp degree, arity, domain and variable-count comparisons; sources 01, 07]\label{cdc:q:degree}
\emph{Source 01 (Sharp degree, arity, and domain comparisons).}
How much of the quartic construction is essential? Totals, forced-zero gains, and some blocker products can be eliminated in special cases. A systematic optimization pass should preserve the complete-root bijection, rather than only satisfiability. Extending the result from natural witnesses to unrestricted integer witnesses is also not automatic: a four-squares conversion can introduce many witnesses for one natural value. A separate canonical integer-domain encoding is needed for a parsimonious statement over $\Z$.

\emph{Source 07 (Nearly trace-minimal quartics).}
Can a structural cyclic-tag compiler retain the actual $T$ consumed bits, use only $T+O(1)$ total natural variables, have degree at most four, and have polynomial, preferably linear, arithmetic-circuit size? The present stream construction achieves the variable count with growing degree; the local construction achieves quartic degree with $3T-2$ coordinates.
\end{question}

\emph{Source 07.} A useful first milestone is to improve the coefficient of $T$ in the quartic construction while retaining complete witness uniqueness. Any lower-bound attempt must exclude the trivial compiler that first simulates the bounded run and outputs a constant.

\wtag{} Manuscript~05's convex quadratic certificate
(\cref{cdc:thm:main}, degree two, Part~II) answers part of manuscript~01's
question: for length-$T$ trace classes of Petri nets and of
interval-guarded translations (\cref{cdc:cor:guarded}), degree four is not
needed. The height-bounded quartic of \cref{cdc:ct:thm:history} and the
integer-domain part remain open, and so does manuscript~07's variable-count
question for cyclic-tag compilers, which the degree-two certificate does not
address.

\begin{question}[Smaller canonical predecessor circuits; source 01]
Can the $(H-1)d$ blocker witnesses be reduced by sharing common subproducts among the dependence neighborhoods $D(a)$? The objective is a smaller \emph{functional} Boolean circuit, so every introduced gate remains uniquely determined. Graph classes with repeated neighborhoods, bounded treewidth, or a small neighborhood basis offer concrete cases. The benchmark should report witness count, residual count, and expanded polynomial size separately.
\end{question}

\begin{question}[Smaller uniquely extended affine systems; source 05]
Can the five local auxiliary variables per time--label pair in \eqref{cdc:eq:q1}--\eqref{cdc:eq:q5} be reduced uniformly while retaining a convex quadratic sum of affine squares, bounded real relaxation, and unique natural extension?
\end{question}
The displayed constants are not optimality claims. Eliminating variables symbolically, sharing complementary quantities, or exploiting implication relations among flags may help. The automaton lower bound concerns Boolean memory and does not obstruct such arithmetic improvements by itself.

\begin{question}[Witness and coefficient optimization; source 06, Question 8]
How far can the explicit variable bounds be reduced
while preserving unique lifting and manageable coefficient heights? Potential
improvements include constant propagation, shared exponent selectors,
common-subexpression elimination, suppression of determined bit equations,
and direct compilation of the square-saturation formula rather than a full
signature square. The relevant objective is a tuple of costs: witness count,
residual count, degree, sparsity, coefficient bit length, and verification
depth. Optimizing only an expanded monomial count can make the actual
certificate worse. Comparisons with universal-equation operation counts must
use the same charging conventions and witness domains.
\end{question}

\begin{question}[Reduce the cubic arithmetic overhead; source 02]
Can the full certificate be reduced from $O(d^3)$ to $O(d^2)$ natural
coordinates while retaining a quartic equation and a polynomial-time
witness generator? The geometric coordinates are already quadratic.
A promising target is a canonical merge of neighboring sign partitions,
avoiding all parent--child pairs, together with shared multipoint
polynomial evaluation. Both the merge and its arithmetic intermediates
must have unique encodings; an arbitrary permutation witness would not do.
\end{question}

\begin{question}[Quantitative lower bounds; source 02]
What lower bounds can be proved for canonical sign-table certificates
under a specified gate model? A bound on our particular interval layout
would not be a lower bound on all Diophantine representations. Separating
geometric information, arithmetic circuit size, and bit height is necessary
before any optimality statement is meaningful.
\end{question}

\begin{question}[Certificate height and practical checking; source 02]
Derive tighter coefficient-sensitive bounds for every intermediate wire,
not only asymptotic bit bounds. These can guide bounded exact search,
serialization, and independently checked resource estimates. Compare the
canonical root cutoff with a minimal stabilization point; a smaller cutoff
is useful only when its certification does not cost more than the saved
sign-table work.
\end{question}

\begin{question}[Can the size--height tradeoff be improved substantially?; source 03]
The linear backend has $O(L)$ scalar witnesses but
$O(L^2(b+\log L))$ maximum bit length, whereas the comparison backend has near-linear
scalar size and essentially input-sized witnesses. Is there a canonical
$O(L)$-coordinate memory certificate with $O(L(b+\log L))$, or even
$O(b+\log L)$, bits per coordinate? A first concrete target is an exact
multiset certificate that avoids a base of size $D^{4L}$ while introducing no
choice of primes, hash challenges, or noncanonical factorization witnesses.
The comparison lower bound does not settle this question.
\end{question}

\begin{question}[What is the best explicit scalar constant?; source 03]
The count $24L$ deliberately uses a simple length-$L$ power chain. Binary
powering reduces the power stage to $O(\log L)$ coordinates. Eliminating
redundant output fields, optimizing record packing, and sharing certified
subexpressions may improve the leading constant further. Every proposed
optimization must preserve injectivity of the packed record and uniqueness of
all discarded or shared witnesses. Publish the exact generated residuals and
count proofs, not only an operation-count estimate.
\end{question}

\begin{question}[Can one minimize total witness bit volume?; source 03]
Maximum coordinate height is not the only meaningful metric. Balanced product
trees can reduce the sum of intermediate bit lengths while retaining linear
many product nodes. A systematic comparison of sequential products, balanced
products, carry-free block packing, and sorting networks could yield an explicit
Pareto frontier for scalar count, total certificate bits, arithmetic work, and
parallel depth. The current implementation uses sequential products for a
transparent exact count; it is not claimed optimal in those other metrics.
\end{question}

\begin{question}[Can canonical routing be made practical at optimal comparison order?; source 03]
The AKS theorem supplies the asymptotic bound but not a small practical backend
here. A routing-based memory checker may have fewer constraints, yet often
admits multiple internal routes for one permutation. Can those routes be
canonically normalized with low overhead and a short proof? A deterministic
routing algorithm is not enough unless its own computation is certified without
reintroducing the original overhead or unbounded auxiliary choices.
\end{question}

\begin{question}[Locality with polynomial bit-height; source 04]
Can a canonical DAG or balanced-tree coding retain context-local certificate size and also provide polynomial witness bit lengths in the size of a succinct input representation and the performed updates?
\end{question}
The pairing code pays for repeated constructor nesting by repeated squaring. A DAG can share repeated subterms, but arbitrary node numbering introduces enormous witness multiplicity. One needs canonical numbering or a quotient theorem for node-renaming equivalence, plus unique proofs of pointer validity and acyclicity. A first useful milestone is a linear-height representation of purely applicative terms with a canonical topological order.

\begin{question}[Sharp arithmetic-memory tradeoffs; source 04]
Can the $r$ flags in the monitor be replaced by fewer natural variables without hiding exponential bit-height, high degree, or a costly digit-extraction gadget?
\end{question}
Theorem~\ref{cdc:wf:thm:memory} gives a streaming bit lower bound, not a variable lower bound. Packing all flags into one integer is easy at the data level but makes bit access a new arithmetic obligation. A useful theorem would state a tradeoff among witness arity, witness height, residual degree, and circuit description size, with an explicit computational model rather than an informal ``compression'' claim.

\begin{question}[A linear-size zero-slack guard; source 07]
Can the no-slack construction be made linear-size in $T$ for a fixed program, while retaining degree $O(T)$ and precisely the consumed-symbol coordinates? The bounded-root guard uses quadratic size in the relevant emptying regime. The question concerns circuit structure, not existence of some polynomial with the correct finite zero set.
\end{question}

\begin{question}[Smaller uniquely determined guard values; source 07]
Can a guard with similarly small circuitry replace the product of all queue lengths by a substantially smaller canonical natural witness? A useful objective is a bound linear in the trace length plus the logarithm of the maximum queue length, rather than their product, without introducing many separate slacks.
\end{question}

% ---------------------------------------------------------------------
\subsection{Richer actions, storage and trajectories}

\begin{question}[Beyond translations; source 01]
Which reset, affine, stack, or indirect-addressing primitives admit exact fixed-arity unique-witness acceleration? The present extremal-prefix argument works because displacement sums are independent of the current value. Reset and multiplication destroy that property. One concrete target is diagonal affine actions on invariant regions where products and guards have a controlled closed form; an answer must separately account for witnesses used to represent any exponentiation.
\end{question}

\begin{question}[Reset, copy, and richer guarded updates; source 05]
Can \cref{cdc:thm:interval} be extended to partial affine maps with reset or copy operations, with a similarly explicit exact commutation criterion?
\end{question}
Unlike translations, different affine expressions can agree on a lower-dimensional guarded domain without having equal coefficient matrices. Any exact criterion must account for those domain equalities as well as for empty compositions. A bounded compiler can still use finite Boolean verification, but that does not supply the desired symbolic all-state criterion.

\begin{question}[New storage monoids and genuine exponential growth; source 06, Question 9]\label{cdc:q:storage}
Identify further action algebras satisfying the local
unique-summary contract of \cref{cdc:cs:prop:abstract}. Promising structured cases
include finite-state products, selected queue restrictions, and stack
fragments with bounded unmatched contexts. The essential test is whether
arbitrary powers retain an exact finite summary with controlled polynomial
witnesses. Affine multiplication of a runtime counter and unbounded string
coding should be treated as separate exponential-growth problems, not added
as innocuous gates.
\end{question}

\wtag{} The queue case is partly answered by Part~VI (manuscript~07):
\cref{cdc:qc:thm:network-certificate} gives a witness-faithful
polynomial whose natural zeros are in bijection with the exact $T$-event
executions of a fixed-read FIFO network program. That is a bounded-run
certificate for a fixed schedule; it does not supply the exact finite
summary of arbitrary powers that the contract of
\cref{cdc:cs:prop:abstract} asks for. Networks whose read counts depend on
data or on emptiness tests stay open (\cref{cdc:q:readcounts}).

\begin{question}[Data-dependent read counts and zero tests; source 07]\label{cdc:q:readcounts}
Extend Theorem~\ref{cdc:qc:thm:network-certificate} from fixed phase-dependent read counts to finite-control networks whose consumption sizes depend on data or on emptiness tests. Specify exact witness counts and ensure that local branch selection and global availability do not admit circular justifications. The fixed-read theorem gives a controlled starting point rather than a completed compiler for every actor language.
\end{question}

\begin{question}[Rational and piecewise-polynomial memory; source 07]
Which extensions of the scalar insertion architecture invalidate Theorem~\ref{cdc:qc:thm:one-dimensional}? Study rational insertion maps with poles excluded from the code image, or a bounded number of polynomial branches. The corresponding centralizer restrictions and the natural-number image condition both matter. A classification should separate genuine scalar memory from hidden finite-state coordinates.
\end{question}

\begin{question}[Automatic trajectory discovery; source 02]
Develop a front end recognizing polynomially iterable updates and
producing the trajectory identities as checked certificates. For affine
maps, a nilpotence witness or an identity of the form
$M^h(M^p-I)^r=0$ is an immediate target. For nonlinear maps, triangular
coordinate changes and verified discrete antidifferentiation could enlarge
the recognized class. The synthesis procedure may be incomplete while its
accepted certificates remain sound.
\end{question}

\begin{question}[Variable periods and variable degrees; source 02]
What uniformity survives when the degree bound or residue period is part
of the numerical input rather than fixed externally? A family of
polynomials with growing arity does not answer this question. Any proposed
single fixed polynomial must account for the encoding and checking of
variable-length coefficient data without silently using an ordinary MRDP
representation that loses single-foldness.
\end{question}

\begin{question}[Nonlinear guards on fixed-residue trajectories; source 02]
Implement the residue-class front end of Theorem~\ref{cdc:pt:thm:quasi}, with
canonical Euclidean division and empty-class normalization. Extend its
nontermination theorem using one canonical root bound per residue class.
This is a concrete development target within the proved framework,
not a claim about arbitrary exponential recurrences.
\end{question}

\begin{question}[Parameter sharing, depth, and resource tradeoffs; source 06, Question 6]
Determine which restrictions on global parameter reuse
restore decidability at nesting depth two. The universality proof reuses
numerical variables across several monomial blocks and circuit equations.
Restrictions such as one syntactic occurrence per parameter, bounded
occurrence count, or an acyclic parameter-incidence pattern may separate
tractable subclasses. Independently, determine quantitative tradeoffs between
the number of resources and repetition depth. One counter at depth four and
finitely many counters at depth two suffice; exact small-resource bounds
would require an explicit, audited construction rather than a qualitative
MRDP invocation.
\end{question}

\begin{question}[The separated cubic interface; source 06, Question 7]
Use \cref{cdc:cs:thm:degree_depth} as an exact interface for
studying separated one-counter schemes of depth three. Every theorem about
this fragment translates to a theorem about natural solutions of one cubic
polynomial, and conversely. Useful restricted targets include bounded numbers
of parameters, sparse support, prescribed sign patterns, or limited parameter
reuse. The correspondence prevents one from mistaking an unproved general
cubic decision assertion for a routine consequence of counter simplicity.
\end{question}

\begin{question}[Unbounded but structurally restricted finite-foldness; source 01]
What larger classes of fixed-skeleton counter programs admit endpoint-level finite-fold certificates? Corollaries~\ref{cdc:ct:cor:rank} and~\ref{cdc:ct:cor:progress} give simple sufficient conditions. Potential extensions use guard-induced count bounds even when the displacement matrix has a nonnegative kernel, or unique solutions of a lower-triangular count system when full column rank fails. The target should be a checkable structural criterion with a proved fibre bound.
\end{question}

% ---------------------------------------------------------------------
\subsection{Rewriting, graph evaluation and comparisons across substrates}

\begin{question}[Bridges between combinatory or graph evaluation and certificates; sources 01, 03, 04]\label{cdc:q:ski}
\emph{Source 01 (A fully specified SKI/Iota bridge).}
Can the local quadratic application-code equations be assembled into a parsimonious representation of a chosen bounded evaluation semantics? The immediate milestones are a unique-witness term-validity checker under a stated size bound, a canonical redex-position encoding, and an evaluation-strategy-specific reflection theorem. The repository's distinction between scoped contextual lambda reduction and weak call-by-value semantics must remain explicit.

\emph{Source 03 (How should graph strategies connect to the existing calculi?).}
Establish the exact relation between the immutable spine evaluator and the
repository's contextual combinator reductions, then between the particular
lambda semantics already present in its separate developments. Which source
observables are preserved and reflected? Which administrative steps should be
counted, erased, or quotiented? For probabilistic or nondeterministic evaluators,
a quotient on histories can change path multiplicity and therefore requires a
measure-preserving argument, not just a reachability simulation.

\emph{Source 04 (Existential SKI schedules without path duplication).}
For horizon $T$ and a maximum context depth $D$, construct a certificate with $\poly(T,D)$ local variables whose solutions are in bijection with valid scheduled SKI reductions, while the redex paths themselves are chosen existentially.
\end{question}

\emph{Source 04.} The fixed-schedule theorem is not this result. A direct disjunction over all paths of length at most $D$ has exponentially many branches. A prospective construction should choose one direction bit per active depth, enforce one canonical stopping point, and force all unused suffix variables to zero. It must also guarantee that a shorter path does not acquire many padded encodings. An even stronger version would choose a canonical reduction strategy, rather than count arbitrary schedules.

\begin{question}[Faithful graph-rewriting front ends; source 05]
Can bounded term-graph or interaction-net histories be compiled with a canonical node-allocation discipline and a proved commutation quotient that respects shared substructure?
\end{question}
The important distinction is between counting graph-rewrite histories, tree interpretations, and observable normal forms. Allocation symmetries and accidental free names can destroy parsimony even when the underlying reduction relation is well understood.

\begin{question}[Cost-preserving compilation between substrates; source 04]
Build translations among counter machines, FRACTRAN, and a deterministic combinatory evaluator that preserve a canonical execution object and give explicit arithmetic-certificate overheads.
\end{question}
Forward simulation is not enough. The target must not introduce multiple admissible administrative schedules for one source step, unless those schedules are canonically quotiented. A useful statement would identify a unique target macrotrace for each source step, prove reflection at designated boundaries, and bound variables, time expansion, and bit-height. This would make cross-substrate comparisons quantitative and semantically trustworthy.

\begin{question}[Spatial certificate transport with explicit clocks; source 07]
Build a formally checked bridge from a concrete cyclic-tag program to its glider or Rule~110 realization, recording exact macro-step boundaries, periodic background conventions, and observable halting events. Then compare a transported queue certificate with a direct space-time polynomial, including bit costs as well as variable counts.
\end{question}

\begin{question}[Quantitative universal compilation; source 07]
Compare universal simulation routes by the final arithmetic resource vector: number of consumed symbols, polynomial degree, circuit size, coefficient bit size, and witness bit height. Simulation time alone does not determine that vector. A useful experiment is to compile the same family of small algorithms through both deletion-tag and cyclic-tag routes, with formal clock correspondence.
\end{question}

\wtag{} This question compares horizon-indexed compilation routes by
their resource vectors; it asks nothing about fixed arity or fibres, so it
is printed here and not merged into \cref{cdc:q:compression}.

\begin{question}[Representation-size versus operation-count optimization; source 05]
Can ProveIt's compact universal-equation and straight-line-certificate work be given a semantic interface that records which optimizations preserve arithmetic fibres, and which preserve only solvability?
\end{question}
Operation counts for fixed universal certificates and variable counts for a time-unrolled compiler measure different resources. A useful comparative study should report both without conflating them, and should test whether elimination or relation-combining steps introduce multiple witnesses.

\wtag{} This question stays a question. It concerns ProveIt's
operation-counted straight-line certificates (the Hilbert's-tenth-problem
project's 75-operation universal certificate \cite{repo-h10}), which
measure a different resource from the witness counts of this report, and
no work in that project is opened here.

% ---------------------------------------------------------------------
\subsection{The fixed-arity frontier}

\begin{question}[Fixed-arity compression with finite or single fibres; sources 01, 02, 03, 04, 06, 07]\label{cdc:q:compression}
\emph{Source 01 (The universal compression problem).}
Can a fixed-arity ordinary polynomial represent first-halt computations of a universal interpreter with finite, or single, witness fibres? Theorem~\ref{cdc:ct:thm:universal} identifies this as the full finite-fold or single-fold MRDP principle, not an omitted routine assembly lemma. A useful contribution would isolate a genuinely canonical unbounded coding primitive and prove that every helper variable has controlled multiplicity. Merely choosing one semantic run or its least numerical code does not discharge that obligation.

\emph{Source 02 (The first genuinely exponential extension).}
Identify a meaningful subclass of exponential-polynomial trajectories
that admits finite-fold or single-fold ordinary polynomial certificates
without assuming such a representation for $2^n$. Any extension including
the doubling graph already has the strength established in
Proposition~\ref{cdc:pt:prop:exponential-boundary}. A conditional theorem that
exposes this dependency is useful; hiding it inside an exponentiation
predicate is not.

\emph{Source 03 (Can one make genuine progress on fixed-variable compression?).}
Theorem~\ref{cdc:mem:thm:compression} identifies the strongest frontier precisely. A
single-fold polynomial for one universal halting relation would settle the
general single-fold question; a finite-fold one would settle its finite-fold
counterpart. Investigate restricted trace families whose coding and all decoding
relations have independently proved finite fibers. Any attempted extension to
universal traces must explicitly discharge the new uniqueness or finiteness
obligations. The fact that the source machine is deterministic is not such a
proof, and no total computable global root bound can be supplied by
Proposition~\ref{cdc:mem:prop:no-bound}.

\emph{Source 04 (The universal fixed-arity challenge).}
Can one preserve the first-halting uniqueness of Corollary~\ref{cdc:wf:cor:halt} while eliminating the horizon-dependent number of variables for a universal interpreter?

\emph{Source 06, Question 12 (Finite-foldness in its minimal substrate).}
Study finite or singleton successful-parameter fibers
directly in the separated normal forms of \cref{cdc:cs:thm:finitefold}. The theorem
shows that an unrestricted positive result would have the full force of the
corresponding Diophantine conjecture. More focused questions concern special
computably enumerable relations or special polynomial support patterns whose
schedule fibers can be controlled constructively. An answer must constrain
the \emph{parameter tuples}; uniqueness of the auxiliary arithmetic
certificate is already supplied here and is not the missing step.

\emph{Source 07 (Fixed-arity compression with controlled fibers).}
Can one compress these canonical exact-halting certificates into a fixed number of Diophantine coordinates while preserving finite or unique fibers? An approach must identify and control every auxiliary introduced by word coding and bounded verification. This is the link to the published finite-fold research program \cite{cco-six}, not a consequence already contained in the present theorems.
\end{question}

\emph{Source 04.} Theorem~\ref{cdc:wf:thm:universal} explains the strength of a positive answer. This is the genuinely global compression problem, rather than a consequence of the existing deterministic gadgets. Any proposal should identify the exact number-theoretic lemma controlling the fibers of trace encoding, bounded universal elimination, or exponentiation replacement. The witness-height obstruction also shows that no computably uniform bound on positive witnesses can accompany a universal solution.

\wtag{} The source theorems cited in this question are merged in Part~IX
as \cref{cdc:bd:thm:universal} (the universal-halting equivalence) and
\cref{cdc:bd:prop:nobound} (no computable witness-height bound).

\begin{question}[Restricted unbounded classes with fixed-arity finite-fold or single-fold compression; sources 04, 05]\label{cdc:q:restricted}
\emph{Source 04 (Finite-fold compression for restricted loop patterns).}
Identify nontrivial unbounded program classes whose canonical bounded tableaux can be compressed into one fixed-arity finite-fold polynomial, with an explicit proof controlling every auxiliary fiber.

\emph{Source 05 (Restricted fixed-dimensional compression).}
Which explicitly described nonuniversal families admit fixed-dimensional single-fold arithmetic representations of their histories? Can a theorem state a sufficient structural hypothesis on the transition system together with the exact arithmetic primitives needed for compression?
\end{question}

\emph{Source 04.} A tractable starting point is flat control flow with a fixed number of loops and affine or multiplicative updates. The challenge is not merely to summarize a loop existentially, but to show that loop counts, intermediate summaries, and all number-theoretic auxiliaries admit finitely many choices. One should state precisely which restrictions prevent the construction from immediately covering a universal halting relation.

\emph{Source 05.} A computable time bound alone is not a proof that its graph, or the associated sequence coding, has a single-fold polynomial representation. A successful restricted result must audit those arithmetic fibres. \Cref{cdc:thm:boundary} provides a sharp warning whenever the family contains canonical universal first-halting histories.

\begin{question}[Finite-fold arithmetic primitives; source 05]
Can one establish finite-fold or single-fold polynomial representations for the sequence or growth relations actually used in a uniform history compressor, rather than for an unspecified black-box decoder?
\end{question}
This connects directly to the arithmetic questions in \cite{mat2010,cco2024}. The computational front end developed here isolates what such a result would need to preserve. It neither supplies nor assumes the missing number-theoretic theorem.

% ---------------------------------------------------------------------
\subsection{What would constitute a stronger result (manuscript 02)}

A successful optimization of the present compiler, or its formal
verification, would strengthen a restricted theorem already proved here.
A genuine extension to universal first-halting computations, or to a
single-fold polynomial graph of $2^n$, would cross a different boundary and
have general representability consequences. Keeping those two kinds of
progress separate prevents a modest but useful construction from being
mistaken for a resolution of the finite-fold program.

\section{Conclusions of manuscripts 05, 01, 04 and 06}\label{cdc:sec:conclusions}
\srcnote{\wtag{} The conclusions of manuscripts~02, 03 and~07 close Parts~IV, V and~VI.}

% ---- manuscript 05, lines 983-989 ----
\subsection{Manuscript 05: Conclusion}
The existence of a Diophantine representation and the preservation of a computation's witnesses are different mathematical requirements. For bounded executions, a fully explicit solution to the stronger requirement is available: compute exact partial-map commutations, select lexicographic representatives with a small Boolean state, and express the resulting verifier by affine equations with unique natural slacks. The final polynomial is convex quadratic, while its natural zeros still carry the full discrete computational content.

The resulting theorem package does more than certify reachability. It preserves class counts and label-additive weights, gives a bounded rational polytope with explicit coordinate bounds, handles interval-guarded counter programs, and explains both an automaton-size obstruction and a counting-complexity obstruction. The elementary ballot family and the exhaustive checks provide independent reference cases for an eventual formal implementation.

At the same time, a growing family of bounded certificates is not one fixed-dimensional universal equation. The exact equivalence in \cref{cdc:thm:boundary} identifies where a proposed uniform finite-fold compressor reaches the full strength of the classical finite-fold problem. This separation makes the construction useful without overstating it: the bounded theorem is proved and executable, and the stronger arithmetic frontier is stated precisely enough to guide further research.


% ---- manuscript 01, lines 856-862 ----
\subsection{Manuscript 01: Conclusion}
For box-guarded translations, the exact resources needed for all serializations of a repeated block admit a surprisingly compact description: total coordinatewise losses plus the largest retained lower margin, and dually the smallest upper margin minus total gains. Global commutation turns this universal condition into exact existential acceleration. The resulting ordinary quartic has a unique complete natural witness and constant arity with respect to the repetition counts.

Combining that resource calculation with canonical causal layers yields an explicit quartic whose roots correspond one-for-one to executable trace classes of bounded height. This preserves a stronger semantic object than reachability alone, and it exposes the otherwise easy-to-miss roles of inactive slacks, padding, and auxiliary circuit variables. Fixed block programs provide additional unbounded-repetition cases, with rank and progress criteria for single-fold or finite-fold endpoint fibres.

The general unbounded problem remains sharply separated. Infinitely many real histories cannot all be encoded by finitely many roots at one endpoint pair, and unique semantic execution does not by itself canonicalize arithmetic sequence codes. A universal fixed-arity single-fold halting polynomial would settle the corresponding full MRDP principle. The present constructions provide explicit, testable components on the tractable side of that boundary, together with a proof-oriented route for extending ProveIt.


% ---- manuscript 04, lines 747-753 ----
\subsection{Manuscript 04: Conclusion}
The informative unit of Diophantine compilation is not merely a yes/no equivalence. It is a controlled fiber: the collection of arithmetic witnesses attached to a computational object. Canonical slacks handle guards, prefix circuits handle priority, symbolic flags remove redundant schedules, and local constructor equations handle rewriting without inspecting untouched subterms.

The resulting theorems give explicit quartic certificates with exact variable and residual counts for three different substrates. For concurrent Petri executions, the zeros represent trace classes rather than their many linearizations. For FRACTRAN, every bounded deterministic run has one arithmetic witness, including its first-applicable-fraction proof. For fixed SKI schedules, a certificate follows the changed context rather than the full tree. These guarantees are compatible with substantial complexity and with large witness heights.

The unbounded question remains sharply separated: fixed arity with finite or unique fibers is not supplied by finite-horizon determinism. The practical research contribution is therefore both constructive and diagnostic---a set of proved, executable witness-faithful compilers, and a precise statement of what additional compression would have to accomplish.


% ---- manuscript 06, lines 1899-1919 ----
\subsection{Manuscript 06: Conclusion}

There is a substantial middle ground between a black-box MRDP invocation and
an explicitly expanded computation trace. For token systems, exact resource
summaries collapse nested repetition into syntax-sized arithmetic. For
lexicographic trace normality, a compact Boolean signature collapses every
power to the three cases zero, one, and at least two. These two compressions
are compatible with unique quadratic witnesses and hence with a single
quartic polynomial.

The resulting guarantees are precise. A fixed schedule with fixed parameters
has one auxiliary certificate. All words of a fixed length admit a separate
polynomial whose roots count executable trace classes exactly. Neither fact
implies that arbitrary existential schedule parameters have finite fibers.
Indeed, separated increment/decrement schedules already contain the full
Diophantine finite-fold question, while their flat fragment remains
Presburger-definable. This combination of constructive positive results and
sharp boundary statements provides a concrete extension of ProveIt's
Diophantine program without attributing to the compiler a solution of problems
it does not solve.


\appendix
\crefalias{section}{appendix}
\section*{Appendices}
\addcontentsline{toc}{section}{Appendices}
\srcnote{\wtag{} The appendices of manuscripts 05, 01, 02, 03, 04 and~06, in that order, followed by this report's provenance record. Manuscript~04's first appendix and manuscript~06's second appendix are printed in \cref{cdc:sec:validation} and \cref{cdc:sec:formal}; manuscript~07 has no appendix.}

% ---- manuscript 05, lines 990-1055 ----
% (source ppendix removed)
\section{Manuscript 05: A minimal end-to-end example}
\label{cdc:app:usage}
The following program can be run from the package directory. Python indices start at zero, whereas the mathematical exposition numbers labels and places from one. The three ordered labels remain $a<b<c$.
\begin{lstlisting}[language=Python]
import sys
sys.path.insert(0, "code")
from trace_polytope import Net, compile_certificate

net = Net(pre=((0,), (0,), (1,)),
          post=((1,), (0,), (0,)),
          names=("a", "b", "c"))
cert = compile_certificate(net, (1,), (1,), 3)

# b,c,a is the canonical representative of c,a,b.
word = (1, 2, 0)
z = cert.witness(word)
assert cert.zero(z)
assert cert.energy(z) == 0
assert cert.decode(z) == word
cert.export("certificate.json")
print(len(cert.variables), len(cert.residuals))  # 73, 59

# Make b a zero-test instead of an unconditional no-op.
# The compiler recomputes the guarded commutation relation.
guarded = compile_certificate(
    net, (0,), (0,), 3,
    upper=((None,), (0,), (None,)))
z2 = guarded.witness((0, 2, 1))  # increment, decrement, zero-test
assert guarded.zero(z2)
\end{lstlisting}

The exported equation is exact even without monomial expansion. To expand a small instance, use \code{cert.expanded_polynomial()}; for large instances, retaining the sum-of-squares circuit is preferable. Expansion is not needed to evaluate the equation or to verify a proposed zero.

\section{Manuscript 05: Edge cases and specification checklist}
\label{cdc:app:edges}
\paragraph{Zero horizon.}
There are no selectors, remainders, or step-local auxiliaries. The endpoint equations require $M=N$, and the flags are forced to zero. There is exactly one empty execution class when the endpoints agree and no solution otherwise.

\paragraph{No places.}
Markings have the unique empty tuple as value. The net has only labelled identity transitions. The selector and normal-form constraints still distinguish the intended words and their classes. The full-box test uses this legitimate degenerate case to obtain a genuinely exhaustive small search.

\paragraph{Repeated effects and repeated circuit inputs.}
Two labels with equal transition data remain distinct actions and can commute. Circuit gates that mention one predecessor twice must read one token once; otherwise a supposedly $1$-safe construction would demand two tokens and fail.

\paragraph{Empty interval compositions.}
Both orders may have empty domains. They then agree as partial maps, even when their numerical composed-box endpoints differ. The guarded commutation routine canonicalizes an empty box before comparing it. Empty \emph{one-step} domains are excluded in the stated input convention.

\paragraph{Inactive guards and branches.}
An inactive upper guard has the bound $B_p$, not an unconstrained witness. An unused branch of a logical disjunction must not leave arbitrary variables. Every variable listed in the exported certificate has a unique semantic value once the canonical word is fixed.

\paragraph{The domain is part of the theorem.}
The identity $\sum_a e_a=1$ forces a one-hot selector over $\N$, but not over $\Z$ or $\R$. Real solutions are deliberately retained only as a relaxation. Natural nonnegativity must not be erased when transferring the polynomial to another system.

\paragraph{Representation versus decision.}
The program constructs and evaluates equations. It does not enumerate all natural zeros of a large certificate. The explicit height theorem makes the set finite, but the associated search and count can still be computationally hard.

\section{Manuscript 05: Source and verification provenance}
The repository sources inspected for the baseline were its Hilbert-tenth-problem README, formalization status register, and the actual \code{DiophantineTrace.lean} interface. The latter was retrieved at the pinned commit listed in \cref{cdc:sec:intro}; its Git blob identifier was
\begin{center}
\small\code{08e5796b22fc85ef0cd97255dba8b7527960160b}.
\end{center}
Repository citations below use immutable commit URLs. The literature was checked through the authors' papers, publisher pages, or the ACL Anthology. The bibliography is selective: it documents the reused conceptual background and the finite-fold boundary, not a claim to an exhaustive priority survey.

The package's \code{README.md} describes its contents and limitations. Its \code{SHA256SUMS.txt} records the distributed files, including the article source, compiled PDF, code, and test receipts. The verification receipt reports the precise finite scopes; it should not be read as a formal proof certificate. \wtag{} In ProveIt the delivered \code{README.md} is replaced by this report's README, the checksum ledger was verified (19 of 19 entries) and retired, and the PDF is not shipped (\cref{cdc:sec:provenance}).


% ---- manuscript 01, lines 863-870 ----
% (source ppendix removed)
\section{Manuscript 01: Reproducibility and source-audit notes}\label{cdc:ct:app:repro}
The package contains the self-contained \code{article.tex}, its compiled PDF, the reference compiler and test runner, machine-readable verification output, and two expanded polynomial examples. To rebuild the document from a clean directory, run \code{pdflatex article.tex} three times to stabilize the contents and cross-references. The bibliography is included directly; no external bibliography database or source checkout is required.

The JSON polynomial format lists each monomial as an integer coefficient and a list of variable names, with repetition recording powers. The declared single polynomial is the sum of the squares of the listed residuals. Exported examples also contain its expanded form. Parameters and witnesses are separately listed, and every assignment entry is an exact integer. These conventions make the files inspectable without a computer algebra system.

The repository audit was limited to relevant live files at the pinned commit: the Hilbert-tenth README, its formalization status register, the actual exact-iteration interface, and the combinatory-logic README. A full-repository novelty audit was not performed. The code and mathematical proofs in this report were not committed to, or used to modify, the user's repository. Literature was checked through the primary sources below; no conclusion of world-first priority is drawn from the search. \wtag{} In ProveIt the package's \code{article.tex} and PDF are not shipped separately (they survive in the arrival commit, \cref{cdc:sec:provenance}); its programs and data are now committed as files of this report under the names listed in the README, and no other repository file was changed.


% ---- manuscript 02, lines 1647-1728 ----
% (source ppendix removed)
\section{Manuscript 02: The finite checking formula in one place}
\label{cdc:pt:app:formula}
For readers implementing a separate compiler, the following is a complete
specification of the primary checking relation. The parameters are
$d$ fixed, $\boldsymbol c\in\Z^{d+1}$, and $T\in\N$. Compute the
coefficient tower by \eqref{cdc:pt:eq:coefficient-difference}. At level $k$, use
$L_k=2(d-k)+1$ slots, fixed cuts $b_{k,0}=0$,
$b_{k,L_k}=T+1$, and natural primary coordinates for all remaining cuts
and all tags.

Impose ordered cuts, tags in $\{0,1,2\}$, terminal zero-tag padding, and
different tags in adjacent active slots. Impose
$\tau_{d,0}=1+\sgn q_d(0)$. For every $k<d$, parent slot $j$, and child
slot $\ell$, let
\[
 a=b_{k,j},\quad b=b_{k,j+1},\quad
 c=b_{k+1,\ell},\quad e=b_{k+1,\ell+1},
\]
\[
 u=\max(a,c),\qquad v=\min(b-1,e),\qquad
 E=(a<b)\wedge(c<e)\wedge(u\leq v).
\]
Require both implications
\[
 E\Longrightarrow 1+\sgn q_k(u)=\tau_{k,j},
 \qquad
 E\Longrightarrow 1+\sgn q_k(v)=\tau_{k,j}.
\]
For a positivity certificate, also require
$(b_{0,j}<b_{0,j+1})\Longrightarrow\tau_{0,j}\neq0$ for every $j$.
Omitting only these last assertions gives the unique sign-table
certificate for \emph{every} polynomial, including those negative
somewhere in the interval.

Every comparison is evaluated, even when $E$ is false. To preserve the
cubic bound, evaluate and cache $q_k$ at all candidate expressions
$a,c,b-1,e$ before selecting their sign tags. Use canonical sign bits for
the selections and the intersection comparison. The signed values
$b-1$ are ordinary integer expressions; truncated natural subtraction
must not be substituted into the checking formula.

\section{Manuscript 02: Reproduction and the exported polynomial format}
\label{cdc:pt:app:reproduction}
The code requires Python 3.10 or later and uses only its standard library.
Run the following commands from the extracted directory \wtag{} (the delivered layout; see the README for a copy of the shipped files in that layout):
\begin{lstlisting}
python run_tests.py
python certificates.py --coefficients 8 -6 1 --horizon 6
python quartic_compiler.py --coefficients 4 -4 1 --horizon 6 \
  --output quartic_example.json
python verify_export.py quartic_example.json
\end{lstlisting}
The coefficient order is ascending: \code{8 -6 1} means $8-6t+t^2$.
The optional compiler switch \code{--sign-table-only} omits the final
positivity requirements while retaining all sign-table conditions.

An exported residual consists of terms with an integer coefficient and a
list of variable names. For example,
\begin{lstlisting}
{"coefficient": -2, "variables": ["b0_1", "w17"]}
\end{lstlisting}
means $-2b_{0,1}w_{17}$. The empty variable list denotes a constant.
Repeated names denote powers; each residual has total degree at most two.
The exact polynomial is the sum of the squares of the listed residuals.
Free coefficient parameters are signed, the free horizon is natural, and
all primary and auxiliary witness coordinates are natural.

The file includes a sample assignment but is not specialized to it.
Changing an input coefficient in that assignment does not change the
polynomial; it normally requires regenerating the unique witness. The
structure fingerprints in the test report confirm equality of the
exported symbolic template across the tested parameter assignments at
fixed nominal degree.

Compile this article with
\begin{lstlisting}
pdflatex -interaction=nonstopmode -halt-on-error article.tex
pdflatex -interaction=nonstopmode -halt-on-error article.tex
\end{lstlisting}
No external bibliography processor or shell escape is required.


% ---- manuscript 03, lines 1647-1715 ----
% (source ppendix removed)
\section{Manuscript 03: Additional validation of the linear backend}\label{cdc:mem:app:linear-tests}
The independent driver \repo{code/03-linear-memory-test_linear_memory.py} runs both algebraic
and operational checks. It verifies collision-free product fingerprints for
451 bounded multisets, including zeros and repetitions. It repeats all 585
small logs, checks six non-Boolean write flags, and tests 40 valid larger-integer
logs together with 37 corrupted reads. All agree with the independent
chronological memory semantics. A further adversarial test replaces one sorted
write record by a duplicate and recomputes every dependent output-side
coordinate. All local bounds, packing, ordering, and memory checks still pass;
exactly the final product-equality residual rejects the forged multiset.

For the four-event example, the linear backend emits exactly 96 witnesses and
89 residuals. Its expanded quartic has 499 nonzero monomials. The largest
witness has 471 bits, compared with six bits for the bitonic backend on the
same log; here $D=164$. All 96 single-coordinate increments of the witness are
rejected. Twenty additional assignments verify the expanded-polynomial
identity, and compiling two different logs of the same length produces
identical symbolic variable lists and residual polynomials.

These checks are recorded in \repo{data/03-linear-memory-linear_test_results.json}; the
complete polynomial and assignment are in \repo{data/03-linear-memory-linear_example-*}.
\par\medskip
\noindent\begin{minipage}{\textwidth}
Run
\begin{lstlisting}
python3 code/test_linear_memory.py
python3 code/linear_memory.py artifacts/example/events.json \
    artifacts/rebuilt_linear_example --expanded
\end{lstlisting}
As with the comparison backend, the proof of uniqueness for arbitrary inputs
is mathematical, not inferred from the bounded test corpus.
\end{minipage}\par

\section{Manuscript 03: Claim and dependency ledger}
\begin{longtable}{@{}>{\raggedright\arraybackslash}p{0.28\textwidth}>{\raggedright\arraybackslash}p{0.38\textwidth}>{\raggedright\arraybackslash}p{0.25\textwidth}@{}}
\toprule
Claim & Dependencies & Status in this archive\\
\midrule
\endfirsthead
\toprule
Claim & Dependencies & Status in this archive\\
\midrule
\endhead
Canonical comparison & Integer algebra and nonnegative witnesses & Proved;
implemented; tested.\\
Low-height memory & Comparison, bitonic sorting, local memory criterion &
Proved; implemented; tested.\\
Optimal network size & External AKS sorting-network theorem & Proved conditional
on that established theorem; not implemented.\\
Carry-free product lemma & Base expansion and cancellation in $\Z[X]$ &
Proved; finite cases tested.\\
Linear-size memory & Product lemma, radix packing, local memory criterion &
Proved; implemented; tested.\\
Parsimonious RAM & Canonical local instruction interface and either memory
backend & Proved; complete frontend not implemented.\\
Immutable graph evaluator & Acyclic unfolding and normalized RAM implementation &
Mathematical construction and proof; no frontend implementation.\\
Counting and weights & RAM bijection and elementary counting/probability &
Proved; weight example tested.\\
Universal compression equivalence & Universal computation and specialization &
Proved elementary equivalence; neither conjecture solved.\\
Lean certification of this work & Proposed modules in Section~\ref{cdc:mem:sec:formalization}
& Not performed.\\
Literature-wide novelty & Additional independent specialist review & Not
established.\\
\bottomrule
\end{longtable}


% ---- manuscript 04, lines 754-754 ----
% (source ppendix removed)

% ---- manuscript 04, lines 798-818 ----
\section{Manuscript 04: Notation and resource-count reference}
\begin{longtable}{p{.19\linewidth}p{.73\linewidth}}
\toprule
Symbol & Meaning \\
\midrule
\endhead
$\N,\Z$ & Natural witness domain including zero; integer polynomial coefficient and evaluation domain.\\
$E,V$ & Number of displayed residuals; number of natural witness variables. Neither is claimed minimal.\\
$I$ & Fixed symmetric irreflexive independence relation on the labeled action alphabet.\\
$\NF_I(w)$ & Lexicographically least word in the class generated by independent adjacent swaps.\\
$p,r,T$ & Petri places, labeled actions, and externally fixed horizon; in FRACTRAN, $r$ is the number of ordered reduced fractions.\\
$U,V$ (net) & Consumption and production matrices. The production matrix $V$ is unrelated to the scalar variable count when its indices are displayed.\\
$m,x,e,f$ & Petri marking, action-selector, enabling-slack, and normal-form-flag variables.\\
$q,\rho,s,u,z$ & Canonical quotient, remainder, bound slack, positive-remainder slack, and divisibility bit.\\
$h,x$ (priority) & No-earlier-applicable-index bit and first-applicable selector.\\
$\App(\ell,r)$ & Natural code for binary application; leaves $S,K,I$ have codes $0,1,2$.\\
$\sigma,H,n_I,n_K,n_S$ & Fixed SKI reduction schedule, its total context depth, and counts of each root rule.\\
$P_T$ vs. $P(T,\cdot)$ & A horizon-indexed polynomial family may have growing arity; one polynomial with horizon as an input has fixed arity.\\
\bottomrule
\end{longtable}


% ---- manuscript 06, lines 1920-1955 ----
% (source ppendix removed)
\section{Manuscript 06: A completely displayed local power polynomial}
\label{cdc:cs:app:local}

To make the ordinary-polynomial nature of the encoding visible without any
compiler machinery, consider one resource-summary pair $(u,v)$, an exponent
$k$, and proposed powered coordinates $(U,V)$. Introduce
$c,\alpha,\beta,e,\ell\in\N$ and set
\begin{align}
\label{cdc:cs:eq:localpoly}
\mathcal P={}&(c+\alpha-u)^2+(c+\beta-v)^2+(\alpha\beta)^2
 +(e(e-1))^2\\
&+(k-e-\ell)^2+((1-e)\ell)^2
 +(U-ec-k\alpha)^2+(V-ec-k\beta)^2.
\notag
\end{align}
This is a polynomial of total degree at most four in
$u,v,k,U,V,c,\alpha,\beta,e,\ell$. For fixed $u,v,k$, its natural solutions
in all seven remaining coordinates $(U,V,c,\alpha,\beta,e,\ell)$ form a
singleton. The first three residuals force
$c=\min(u,v)$, $\alpha=u-c$, $\beta=v-c$. The next three force
$e=\mathbf1_{k>0}$ and $\ell=k-e$. The final two give exactly the powered
summary. Thus the uniqueness claim includes both the result and its internal
witnesses.

For $u=v=1$ and $k=0$, the forced values are
\[
c=1,\quad\alpha=\beta=e=\ell=0,\quad U=V=0.
\]
The internal minimum remains one, but the powered action is the identity.
For the same pair and any $k>0$, the forced values are $e=1$, $\ell=k-1$,
$U=V=1$. This small example shows simultaneously why inactive internal
witnesses must remain fixed and why their values must not be confused with
the resources required by a zero power.

\clearpage

\section{Provenance of this report}\label{cdc:sec:provenance}
\wtag{} This report was built from seven manuscripts. Six arrived in batch~60 of ProveIt's incoming reports (arrival commits \texttt{725d2ebb6} and \texttt{c1f56f842}) and were placed by commit \texttt{7498484af}; the seventh arrived alone in batch~61 (\texttt{b998f70c6}) and was placed by \texttt{3bf66a8bc} as the report's seventh member, before the report was written. Each manuscript's source and delivery README survive only in its arrival commit, except manuscript~05's source, which is the base of this article; no manuscript PDF is shipped.

\begingroup\footnotesize\setlength{\tabcolsep}{3pt}
\begin{longtable}{@{}p{0.035\textwidth}>{\raggedright\arraybackslash}p{0.25\textwidth}p{0.1\textwidth}p{0.1\textwidth}p{0.1\textwidth}>{\raggedright\arraybackslash}p{0.3\textwidth}@{}}
\toprule
No. & Archive (batch, manuscript) & Arrival & Pin & Placed & Contribution (Parts)\\
\midrule
\endhead
01 & \path{Diophantine_Causal_Computation.zip} (60, 01) & \texttt{725d2ebb6} & \texttt{e8bb0931d} & \texttt{7498484af} & box test and resource envelope (I); causal-height quartic (II); accelerator, block programs, word powers, examples (III); substrate adapters (VII, VIII); boundary (IX)\\
02 & \path{canonical_diophantine_certificates.zip} (60, 02) & \texttt{725d2ebb6} & \texttt{e8bb0931d} & \texttt{7498484af} & signed coordinates and total circuits (conventions); polynomial trajectories (IV); FRACTRAN step (VII); substrates (VIII); exponentiation boundary (IX)\\
03 & \path{ProveIt_Linear_Size_Diophantine_Certificates.zip} (60, 03) & \texttt{725d2ebb6} & \texttt{e8bb0931d} & \texttt{7498484af} & memory logs, RAM, graph evaluation, counting (V); boundary (IX)\\
04 & \path{Witness_Faithful_Diophantine_Compilation.zip} (60, 04) & \texttt{c1f56f842} & \texttt{e8bb0931d} & \texttt{7498484af} & monitor, Petri quartic, counting (II, second route); FRACTRAN and SKI (VII); substrates (VIII); boundary (IX)\\
05 & \path{canonical_trace_polytopes.zip} (60, 05), base & \texttt{c1f56f842} & \texttt{e8bb0931d} & \texttt{7498484af} & introduction; commutation (I); trace polytopes, costs, counting (II); rewriting (VII); substrates (VIII); boundary (IX)\\
06 & \path{diophantine_counter_programs.zip} (60, 06) & \texttt{c1f56f842} & \texttt{e8bb0931d} & \texttt{7498484af} & resource summaries (I); signatures and trace bijection (II, second route); saturation, compiler, nesting boundary (III); substrates (VIII); finite-fold substrate equivalence (IX)\\
07 & \path{Causal_Diophantine_Queue_Compilation.zip} (61, 01) & \texttt{b998f70c6} & \texttt{725d2ebb6} & \texttt{3bf66a8bc} & first-failure product (I); queues, tag systems, word memory, FIFO networks (VI); boundary (IX)\\
\bottomrule
\end{longtable}
\endgroup

\paragraph{Where the merge had to choose.} \wtag{}
\begin{itemize}
\item \emph{Base.} Manuscript~05, whose trace-class theorem needs the weakest hypothesis (any sound $I\subseteq I_{\max}$, with interval guards) and gives the strongest conclusion (a convex quadratic whose real zero set is a rational polytope with exactly the natural zeros as lattice points). Its text supplies the title page layout, the preamble, the introduction and the first place of every shared result.
\item \emph{Printed once.} The universal-halting equivalence (manuscripts 05, 01, 03, 04) is \cref{cdc:bd:thm:universal}; the absence of a computable witness-height bound (02, 01, 03, 04, 05) is \cref{cdc:bd:prop:nobound}. The labels of the four equivalence theorems and the five height statements (four propositions and manuscript~04's theorem) that they replace point to the merged statements, and each manuscript's proof is printed in its own section after a note mapping its letters.
\item \emph{Second routes.} Manuscripts~04's and~06's trace-class bijections (\cref{cdc:wf:thm:petri,cdc:cs:thm:tracebijection}), manuscript~01's box test (\cref{cdc:ct:prop:boxes}) and manuscript~06's commutation criterion (\cref{cdc:cs:thm:commutation}) are kept with their own statements, counts and proofs, marked in their titles; manuscript~04's monitor, lower-bound family and counting theorem, and the Parikh refinements of manuscripts~01, 04 and~06, likewise point to manuscript~05's versions. Manuscript~01's height-bounded compiler is a different theorem and is kept whole.
\item \emph{Moved text.} Manuscript~05's interval-guard theorem and guarded-program corollary were moved from its section on other substrates to Parts~I and~II, and its rewriting subsection to Part~VII; the substrate surveys of manuscripts~01, 02 and~05 are split between Parts~VII, VIII and~IX; manuscript~02's first two introductory subsections open Part~IV and its Section~2 stands in the conventions; manuscript~07's word-code subsection opens Part~VI and its first-failure section closes Part~I. Every move is announced by a note at both ends.
\item \emph{Back matter.} Implementation and validation, formalization proposals and conclusions are collected by manuscript (the conclusions of the single-manuscript Parts~IV--VI close those Parts). The research questions of all seven manuscripts (74) are merged into 50 questions; each merged question names its sources and prints each source's text (\cref{cdc:sec:questions}).
\item \emph{Bibliography.} One list; duplicates merged, each item naming the manuscripts that cite it. Manuscript~07 cites a different Cantone--Cuzziol--Omodeo article (\cite{cco-six}) from the one cited by manuscripts~02, 04, 05 and~06 (\cite{cco2024}); they are kept apart. Manuscript~07's two unpublished ``companion'' items are manuscripts~04 and~06 themselves; its sentence citing them now names them.
\item \emph{Corrections.} Manuscript~06 cited Jones--Matiyasevich for ProveIt's documentation of MRDP and its interfaces; it now cites the repository. Manuscript~04's hard-coded section numbers (``Section~2'', ``Section~3'', ``Section~7'') now point to its sections here. Sentences in which manuscripts~01 and~05 describe their own delivered packages carry notes on what is and is not shipped.
\item \emph{Additions.} The conventions (\cref{cdc:conv:sec}), the Part openings, the remarks on the formalized exponential counterpart (\cref{cdc:bd:rem:sfu}) and on the project's tag-system notes (\cref{cdc:qc:rem:tagnotes}), and every other paragraph marked \wtag{}.
\end{itemize}

\paragraph{Labels.} \wtag{} Every label carries the prefix \texttt{cdc:}. Manuscript~05's labels keep their names after it (\texttt{cdc:thm:main}); the other manuscripts' labels carry a second prefix: 01 \texttt{cdc:ct:}, 02 \texttt{cdc:pt:}, 03 \texttt{cdc:mem:}, 04 \texttt{cdc:wf:}, 06 \texttt{cdc:cs:}, 07 \texttt{cdc:qc:}. Labels written in the merge use \texttt{cdc:conv:}, \texttt{cdc:bd:}, \texttt{cdc:part:}, \texttt{cdc:sec:}, \texttt{cdc:q:}, or a manuscript's prefix for a section it did not label (\texttt{cdc:wf:sec:interfaces}).

\clearpage
\phantomsection
\addcontentsline{toc}{section}{References}
\begin{thebibliography}{99}
\small
\raggedright
\item[] \wtag{} One bibliography for the seven manuscripts. An item cited by several manuscripts is printed once; the manuscripts citing it are listed at its end, and a detail given by only one of them is kept and attributed. Repository items are pinned at \texttt{e8bb0931d67f80d9fce87a8cddb0f661ff19f956} (manuscripts 01--06) or \texttt{725d2ebb6909fe11a13a92354c0f47367a3cbbf5} (manuscript~07); the cited files of \path{Computability/HilbertTenthProblem} are identical at the two commits. Manuscript~07's two ``companion'' items were manuscripts~04 and~06 of this report and are not repeated.

\bibitem{repo-h10}
V. Reshetnikov (and contributors).
\emph{ProveIt: Hilbert's tenth problem} --- project README (manuscripts 01, 03, 04, 07), with its corrected Jones articles and Diophantine developments (04), and, as one item, README and formalization status register (02).
\path{Computability/HilbertTenthProblem/README.md}, inspected September 30, 2026.
\url{https://github.com/VladimirReshetnikov/ProveIt/blob/e8bb0931d67f80d9fce87a8cddb0f661ff19f956/Computability/HilbertTenthProblem/README.md};
at manuscript~07's pin \url{https://github.com/VladimirReshetnikov/ProveIt/blob/725d2ebb6909fe11a13a92354c0f47367a3cbbf5/Computability/HilbertTenthProblem/README.md};
the development's tree as cited by manuscript~02:
\url{https://github.com/VladimirReshetnikov/ProveIt/tree/e8bb0931d67f80d9fce87a8cddb0f661ff19f956/Computability/HilbertTenthProblem}.
Manuscript~02: the repository supplies corrected editions and source developments; inspection here does not constitute an independent build audit.
Cited by manuscripts 01, 02, 03, 04, 07.

\bibitem{repo-status}
V. Reshetnikov (and contributors).
\emph{ProveIt: Hilbert's tenth problem, formalization status}.
\path{Computability/HilbertTenthProblem/Lean/STATUS.md}, inspected September 30, 2026.
\url{https://github.com/VladimirReshetnikov/ProveIt/blob/e8bb0931d67f80d9fce87a8cddb0f661ff19f956/Computability/HilbertTenthProblem/Lean/STATUS.md}.
Cited by manuscripts 03, 05.

\bibitem{repo-trace}
V. Reshetnikov (and contributors).
\emph{ProveIt: Diophantine bounded universals and finite traces}.
\path{Computability/HilbertTenthProblem/Lean/Diophantine/Common/DiophantineTrace.lean}, inspected September 30, 2026.
\url{https://github.com/VladimirReshetnikov/ProveIt/blob/e8bb0931d67f80d9fce87a8cddb0f661ff19f956/Computability/HilbertTenthProblem/Lean/Diophantine/Common/DiophantineTrace.lean}.
Cited by manuscripts 01, 02, 04, 05.

\bibitem{repo-cl}
V. Reshetnikov.
\emph{ProveIt: Combinatory logic --- SK, SKI, and Iota universality}.
\path{Computability/CombinatoryLogic/README.md}, inspected September 30, 2026.
\url{https://github.com/VladimirReshetnikov/ProveIt/blob/e8bb0931d67f80d9fce87a8cddb0f661ff19f956/Computability/CombinatoryLogic/README.md}.
Cited by manuscripts 01, 03, 04.

\bibitem{repo-mrdp}
V. Reshetnikov, ProveIt repository.
\emph{MRDP: recursively enumerable if and only if Diophantine}.
\path{Computability/HilbertTenthProblem/Lean/MRDP.md}, inspected at commit \texttt{725d2ebb6909fe11a13a92354c0f47367a3cbbf5}, September 30, 2026.
\url{https://github.com/VladimirReshetnikov/ProveIt/blob/725d2ebb6909fe11a13a92354c0f47367a3cbbf5/Computability/HilbertTenthProblem/Lean/MRDP.md}.
Cited by manuscript 07.

\bibitem{repo}
V. Reshetnikov.
\emph{ProveIt}, public research repository, inspected September 30, 2026.
Repository tree at commit \texttt{e8bb0931d67f}; relevant documents: \path{Computability/HilbertTenthProblem/README.md} and \path{Computability/HilbertTenthProblem/Lean/STATUS.md}.
\href{https://github.com/VladimirReshetnikov/ProveIt/tree/e8bb0931d67f80d9fce87a8cddb0f661ff19f956/Computability/HilbertTenthProblem}{Pinned Diophantine development}.
The repository descriptions were inspected; its proofs were not rebuilt for manuscript~06.
Cited by manuscript 06.

\bibitem{mat1970}
Yu. V. Matiyasevich.
Enumerable sets are Diophantine.
\emph{Soviet Mathematics Doklady} \textbf{11} (1970), 354--358.
(Manuscript~04 gives the pages as 354--357.)
Cited by manuscripts 04, 06.

\bibitem{jm1984}
J. P. Jones and Yu. V. Matiyasevich.
Register machine proof of the theorem on exponential Diophantine representation of enumerable sets.
\emph{Journal of Symbolic Logic} \textbf{49}(3) (1984), 818--829.
\href{https://doi.org/10.2307/2274135}{doi:10.2307/2274135}.
Bibliographic details, a corrected edition and its formalization status are catalogued in \cite{repo-h10} (manuscripts 01 and 06).
Cited by manuscripts 01, 03, 04, 05, 06.

\bibitem{mat2010}
Yu. V. Matiyasevich.
Towards finite-fold Diophantine representations.
\emph{Journal of Mathematical Sciences} \textbf{171} (2010), 745--752.
\href{https://doi.org/10.1007/s10958-010-0179-4}{doi:10.1007/s10958-010-0179-4}.
Manuscript~03: this reference documents the finite-fold/single-fold frontier; its universal-substrate equivalence is proved directly (\cref{cdc:bd:thm:universal}).
Cited by manuscripts 01, 02, 03, 05, 06.

\bibitem{cco2024}
D. Cantone, L. Cuzziol, and E. G. Omodeo.
On Diophantine singlefold specifications.
\emph{Le Matematiche} \textbf{79}, no.~2 (2024), 585--620. Published December 28, 2024.
\href{https://doi.org/10.4418/2024.79.2.18}{doi:10.4418/2024.79.2.18}.
Journal page: \url{https://lematematiche.dmi.unict.it/index.php/lematematiche/article/view/2703}.
Author-deposited text: \url{https://arts.units.it/bitstream/11368/3101478/1/CCO24.pdf}.
Cited by manuscripts 02, 04, 05, 06.

\bibitem{cco-six}
D. Cantone, L. Cuzziol, and E. G. Omodeo.
\emph{Six equations in search of a finite-fold-ness proof}.
2023; version 3, September 30, 2024.
\url{https://arxiv.org/abs/2303.02208v3}.
A different article from \cite{cco2024}.
Cited by manuscript 07.

\bibitem{ccfo2019}
D. Cantone, A. Casagrande, F. Fabris, and E. Omodeo.
Does every recursively enumerable set admit a finite-fold Diophantine representation?
\emph{CEUR Workshop Proceedings} \textbf{2396}, paper 11 (2019).
\url{https://ceur-ws.org/Vol-2396/paper11.pdf}.
Cited by manuscript 01.

\bibitem{lf2019}
D. Larchey-Wendling and Y. Forster.
Hilbert's tenth problem in Coq.
In \emph{FSCD 2019}, LIPIcs \textbf{131} (2019), 27:1--27:20.
\href{https://doi.org/10.4230/LIPIcs.FSCD.2019.27}{doi:10.4230/LIPIcs.FSCD.2019.27}.
Cited by manuscripts 04, 07.

\bibitem{lf2022}
D. Larchey-Wendling and Y. Forster.
Hilbert's tenth problem in Coq (extended version).
\emph{Logical Methods in Computer Science} \textbf{18}(1), article 35 (2022).
\href{https://doi.org/10.46298/lmcs-18(1:35)2022}{doi:10.46298/lmcs-18(1:35)2022}.
Author version: \url{https://arxiv.org/abs/2003.04604} (version 5: \url{https://arxiv.org/abs/2003.04604v5}).
Cited by manuscripts 01, 02.

\bibitem{carneiro}
M. Carneiro.
\emph{A Lean formalization of Matiyasevi\v{c}'s Theorem}.
arXiv:1802.01795, 2018.
\url{https://arxiv.org/abs/1802.01795}.
Cited by manuscripts 03, 07.

\bibitem{cartier-foata}
P. Cartier and D. Foata.
\emph{Probl\`emes combinatoires de commutation et r\'earrangements}.
Lecture Notes in Mathematics \textbf{85}, Springer, 1969.
Cited in manuscript~01 for the classical partially commutative monoid and normal-form framework; see also \cite{krattenthaler}.
Cited by manuscript 01.

\bibitem{krattenthaler}
C. Krattenthaler.
\emph{The theory of heaps and the Cartier--Foata monoid}, author-hosted exposition.
\url{https://www.mat.univie.ac.at/~kratt/artikel/heaps.pdf}.
Cited by manuscript 01.

\bibitem{barthelemy}
F. Barth\'elemy.
Using Mazurkiewicz trace languages for partition-based morphology.
In \emph{Proceedings of the 45th Annual Meeting of the Association of Computational Linguistics} (2007), 928--935.
See the lexicographic normal-form characterization in Section 2.2, p.~930.
\url{https://aclanthology.org/P07-1117/}.
Cited by manuscript 05.

\bibitem{anisimov-knuth}
A. V. Anisimov and D. E. Knuth.
\emph{Inhomogeneous sorting}.
International Journal of Computer \& Information Sciences \textbf{8} (1979), 255--260.
\href{https://doi.org/10.1007/BF00993053}{doi:10.1007/BF00993053}.
Cited by manuscript 06.

\bibitem{lohrey2024}
M. Lohrey, F. Stober, and A. Wei\ss.
The power word problem in graph products.
\emph{Theory of Computing Systems} \textbf{68} (2024), 403--464.
\href{https://doi.org/10.1007/s00224-024-10173-z}{doi:10.1007/s00224-024-10173-z}.
Cited by manuscript 04.

\bibitem{decolnet2026}
A. de Colnet, K. S. Meel, and U. Mathur.
Counting and sampling traces in regular languages.
\emph{Proceedings of the ACM on Programming Languages}, POPL 2026.
\href{https://doi.org/10.1145/3776723}{doi:10.1145/3776723}.
Preprint \href{https://arxiv.org/abs/2512.00314}{arXiv:2512.00314}.
Cited by manuscript 04.

\bibitem{valiant}
L. G. Valiant.
The complexity of enumeration and reliability problems.
\emph{SIAM Journal on Computing} \textbf{8}(3) (1979), 410--421.
\url{https://doi.org/10.1137/0208032}.
Cited by manuscript 05.

\bibitem{zetzsche}
G. Zetzsche.
\emph{The Emptiness Problem for Valence Automata over Graph Monoids}.
arXiv:1710.07528 (2017).
\url{https://arxiv.org/abs/1710.07528}.
Cited by manuscript 06.

\bibitem{grunewald-segal}
F. Grunewald and D. Segal.
\emph{On the integer solutions of quadratic equations}.
Journal f\"ur die reine und angewandte Mathematik \textbf{569} (2004), 13--45.
\href{https://doi.org/10.1515/crll.2004.023}{doi:10.1515/crll.2004.023}.
Cited by manuscript 06.

\bibitem{hfg2024}
M. Hark, F. Frohn, and J. Giesl.
\emph{Termination of Triangular Polynomial Loops}.
arXiv:1910.11588, version 7, February 13, 2024.
\url{https://arxiv.org/abs/1910.11588v7}.
Cited by manuscript 02.

\bibitem{lmg2024}
N. Lommen, F. Meyer, and J. Giesl.
\emph{Automatic Complexity Analysis of Integer Programs via Triangular Weakly Non-Linear Loops}.
arXiv:2205.08869, version 3, November 15, 2024.
\url{https://arxiv.org/abs/2205.08869v3}.
Cited by manuscript 02.

\bibitem{snark}
E. Ben-Sasson, A. Chiesa, D. Genkin, E. Tromer, and M. Virza.
SNARKs for C: Verifying Program Executions Succinctly and in Zero Knowledge.
In \emph{Advances in Cryptology --- CRYPTO 2013}, Part II, LNCS 8043, Springer, 2013, 90--108.
\url{https://doi.org/10.1007/978-3-642-40084-1_6}.
Extended version, Cryptology ePrint Archive 2013/507: \url{https://eprint.iacr.org/2013/507}.
The extended version's memory-consistency discussion is the relevant prior art.
Cited by manuscript 03.

\bibitem{batcher}
K. E. Batcher.
Sorting networks and their applications.
In \emph{Proceedings of the April 30--May 2, 1968, Spring Joint Computer Conference}, AFIPS \textbf{32}, 307--314.
\url{https://doi.org/10.1145/1468075.1468121}.
Cited by manuscript 03.

\bibitem{aks}
M. Ajtai, J. Koml\'os, and E. Szemer\'edi.
An $O(n\log n)$ sorting network.
In \emph{Proceedings of the Fifteenth Annual ACM Symposium on Theory of Computing}, 1983, 1--9.
\url{https://doi.org/10.1145/800061.808726}.
Cited by manuscript 03.

\bibitem{conway1987}
J. H. Conway.
FRACTRAN: A simple universal programming language for arithmetic.
In \emph{Open Problems in Communication and Computation}, Springer, New York, 1987, pp.~4--26.
Cited by manuscript 04.

\bibitem{cook}
M. Cook.
\emph{Universality in Elementary Cellular Automata}.
Complex Systems \textbf{15}(1) (2004), 1--40.
\href{https://doi.org/10.25088/ComplexSystems.15.1.1}{doi:10.25088/ComplexSystems.15.1.1}.
The cyclic-tag definition and simulation are in Sections 2.1--2.2.
Cited by manuscript 07.

\bibitem{woods-neary}
D. Woods and T. Neary.
\emph{On the time complexity of 2-tag systems and small universal Turing machines}.
Proceedings of FOCS 2006.
\href{https://doi.org/10.1109/FOCS.2006.58}{doi:10.1109/FOCS.2006.58};
\url{https://arxiv.org/abs/cs/0612089}.
Cited by manuscript 07.

\bibitem{atela-hu}
P. Atela and J. Hu.
\emph{Commuting polynomials and polynomials with same Julia set}.
1995. \url{https://arxiv.org/abs/math/9504210}.
Cited for the broader centralizer context, not as a dependency of the elementary proof in manuscript~07.
Cited by manuscript 07.
\end{thebibliography}
\end{document}
